\documentclass[10pt,a4paper]{article}
\usepackage[utf8]{inputenc}
\usepackage[T1]{fontenc}
\usepackage{lmodern}
\usepackage[left=18mm,right=18mm,top=20mm,bottom=19mm,headheight=15pt]{geometry}
\usepackage{xcolor}
\definecolor{ink}{HTML}{193B4C}\definecolor{teal}{HTML}{0C757D}\definecolor{muted}{HTML}{606D76}
\usepackage[most]{tcolorbox}
\usepackage{hyperref}
\hypersetup{colorlinks=true,linkcolor=teal,urlcolor=teal}
\usepackage{fancyhdr,enumitem,microtype,needspace,titlesec}
\setlength{\parindent}{0pt}
\setlength{\parskip}{4pt}
\setlist[enumerate]{topsep=3pt,itemsep=2pt}
\titleformat{\section}{\LARGE\bfseries\color{ink}}{\thesection}{0.7em}{}
\titleformat{\subsection}{\large\bfseries\color{teal}}{\thesubsection}{0.7em}{}
\pagestyle{fancy}\fancyhf{}
\lhead{\footnotesize Microsoft Office | Informatika attestatsiyasi}
\rhead{\footnotesize 2026}
\cfoot{\thepage}
\renewcommand{\headrulewidth}{0.4pt}
\emergencystretch=3em\sloppy\hbadness=10000\hfuzz=2000pt
\begin{document}
\begin{titlepage}
\vspace*{33mm}
{\color{teal}\rule{\linewidth}{2pt}}\par
\vspace{17mm}
{\Large INFORMATIKA ATTESTATSIYASI \quad 2026}\par
\vspace{13mm}
{\fontsize{31}{38}\selectfont\bfseries\color{ink}MICROSOFT OFFICE}\par
\vspace{5mm}
{\fontsize{21}{26}\selectfont WORD \quad EXCEL \quad POWERPOINT}\par
\vspace{15mm}
{\Large Savollar va har biri ostidagi izohli javoblar}\par
\vspace{20mm}
{\Huge\bfseries\color{teal}389}\par
{\large 381 published \quad 8 review}\par
\vfill
{\color{muted}2026-yil attestatsiya spetsifikatsiyasi bo'yicha 6--10-savollar yo'nalishi.}
\par\vspace{8mm}{\color{teal}\rule{\linewidth}{2pt}}
\end{titlepage}
\section*{Foydalanish tartibi}
\addcontentsline{toc}{section}{Foydalanish tartibi}
Bu kitobda Microsoft Word, Excel va PowerPoint bo'yicha 389 ta amaldagi savol mavjud.
Har bir savoldan so'ng \textbf{To'g'ri javob} va \textbf{Izoh} bloki berilgan.
Y1 - variant tanlash; Y2 - moslashtirish; Y3 - amallarni ketma-ket joylashtirish.
\par\textbf{Manba:} attestatsiya100\% Supabase bazasi: public.questions, question\_options, question\_keys va lessons. 
\par\textbf{Holat:} published (381) va review (8). Review holatidagi savollar yakuniy ekspert tasdig'i emas.
Arxiv (148) va draft (7) savollar kiritilmagan. 3 ta published savolda tasvirga havola mavjud.
\par\textbf{Eslatma:} ruscha menyu nomlari PDF kompilyatsiyasi mosligi uchun lotin transliteratsiyasi bilan berilgan.
Asl yozilish va baza ID-larini manba LaTeX izohlari orqali tekshirish mumkin.
\tableofcontents\clearpage
\section{Microsoft Word}
\subsection{MS Word: hujjat formatlash va tuzilma}

% DB 82285d56-0758-4d59-ad82-d23eee834721 status=published
\needspace{6\baselineskip}
\textbf{\color{ink}1.} {ootnotesizecolor{muted}Y1 | published | 82285d56}
\par
MS Word atamalarini ta'riflari bilan moslang. \newline  \newline 1. Justify \newline 2. Border \newline 3. Jadval \newline 4. Ro'yxat \newline  \newline a. Matnni o'qishni osonlashtiruvchi maket \newline b. Bir-biriga bog'liq elementlar ketma-ketligi \newline c. Matnni sahifaning ikki cheti bo'ylab tekislash \newline d. Tasvir yoki matn bloki atrofidagi hoshiya chizig'i \newline e. Ro'yxat yaratishda qo'llanadigan dizayn
\begin{enumerate}[label=\Alph*),leftmargin=9mm]
\item 1-c, 2-d, 3-a, 4-b
\item 1-c, 2-d, 3-b, 4-a
\item 1-c, 2-d, 3-a, 4-e
\item 1-c, 2-e, 3-a, 4-b
\end{enumerate}
\begin{tcolorbox}[enhanced,breakable,colback=teal!4,colframe=teal!55!black,boxrule=0.45pt,arc=1.2mm,left=2.8mm,right=2.8mm,top=2mm,bottom=2mm]
\textbf{\color{teal}To'g'ri javob:} \textbf{A) 1-c, 2-d, 3-a, 4-b}
\par\smallskip
\textbf{Izoh:} Justify - ikki chetga tekislash; Border - hoshiya; jadval axborotni o'qishni osonlashtiradi; ro'yxat - bog'liq elementlar ketma-ketligi. [Yutuqli\_test, test kodi 12; Telegram xabarlar 6446 va 6556.]
\end{tcolorbox}
\vspace{2mm}\hrule height .2pt\vspace{2mm}

% DB 05dc9f05-e9db-4f64-a622-6662ca2b142b status=published
\needspace{6\baselineskip}
\textbf{\color{ink}2.} {ootnotesizecolor{muted}Y1 | published | 05dc9f05}
\par
Word hujjatining har bir sahifasida bir xil maktab nomi yuqori qismda ko'rinishi kerak. Qaysi elementdan foydalanish ma'qul?
\begin{enumerate}[label=\Alph*),leftmargin=9mm]
\item Header
\item Footnote
\item Text Highlight
\item Page Break
\end{enumerate}
\begin{tcolorbox}[enhanced,breakable,colback=teal!4,colframe=teal!55!black,boxrule=0.45pt,arc=1.2mm,left=2.8mm,right=2.8mm,top=2mm,bottom=2mm]
\textbf{\color{teal}To'g'ri javob:} \textbf{A) Header}
\par\smallskip
\textbf{Izoh:} Header har bir sahifaning yuqori qismida takrorlanuvchi ma'lumotni ko'rsatadi.
\end{tcolorbox}
\vspace{2mm}\hrule height .2pt\vspace{2mm}

% DB 0832b113-2721-46ae-845c-e611989de704 status=published
\needspace{6\baselineskip}
\textbf{\color{ink}3.} {ootnotesizecolor{muted}Y1 | published | 0832b113}
\par
Yangi bobni Word hujjatida aniq keyingi sahifadan boshlash kerak. Eng to'g'ri amal qaysi?
\begin{enumerate}[label=\Alph*),leftmargin=9mm]
\item Page Break kiritish
\item Bir necha marta Enter bosish
\item Shriftni kattalashtirish
\item Matn rangini oq qilish
\end{enumerate}
\begin{tcolorbox}[enhanced,breakable,colback=teal!4,colframe=teal!55!black,boxrule=0.45pt,arc=1.2mm,left=2.8mm,right=2.8mm,top=2mm,bottom=2mm]
\textbf{\color{teal}To'g'ri javob:} \textbf{A) Page Break kiritish}
\par\smallskip
\textbf{Izoh:} Page Break keyingi mazmunni yangi sahifadan barqaror boshlaydi.
\end{tcolorbox}
\vspace{2mm}\hrule height .2pt\vspace{2mm}

% DB 15b264e9-0ec9-4472-96b3-3e87a061fbf2 status=published
\needspace{6\baselineskip}
\textbf{\color{ink}4.} {ootnotesizecolor{muted}Y1 | published | 15b264e9}
\par
Word hujjatida barcha \texttt{2025} yozuvlarini bir amal bilan \texttt{2026} ga almashtirish kerak. Qaysi vosita eng mos?
\begin{enumerate}[label=\Alph*),leftmargin=9mm]
\item Find and Replace
\item Zoom
\item Print Preview
\item WordArt
\end{enumerate}
\begin{tcolorbox}[enhanced,breakable,colback=teal!4,colframe=teal!55!black,boxrule=0.45pt,arc=1.2mm,left=2.8mm,right=2.8mm,top=2mm,bottom=2mm]
\textbf{\color{teal}To'g'ri javob:} \textbf{A) Find and Replace}
\par\smallskip
\textbf{Izoh:} Find and Replace takrorlanuvchi matnni ommaviy almashtirish uchun ishlatiladi.
\end{tcolorbox}
\vspace{2mm}\hrule height .2pt\vspace{2mm}

% DB 5e0c753a-095a-42dc-bcba-fee64094bd2c status=published
\needspace{6\baselineskip}
\textbf{\color{ink}5.} {ootnotesizecolor{muted}Y1 | published | 5e0c753a}
\par
Word hujjatida matnning ikki cheti ham tekis ko'rinishi kerak. Qaysi tekislash tanlanadi?
\begin{enumerate}[label=\Alph*),leftmargin=9mm]
\item Justify
\item Center
\item Left
\item Right
\end{enumerate}
\begin{tcolorbox}[enhanced,breakable,colback=teal!4,colframe=teal!55!black,boxrule=0.45pt,arc=1.2mm,left=2.8mm,right=2.8mm,top=2mm,bottom=2mm]
\textbf{\color{teal}To'g'ri javob:} \textbf{A) Justify}
\par\smallskip
\textbf{Izoh:} Justify matnni chap va o'ng chegaralar bo'yicha tekislaydi.
\end{tcolorbox}
\vspace{2mm}\hrule height .2pt\vspace{2mm}

% DB c1da0e32-3a7a-4db3-bc74-affdc080b80d status=published
\needspace{6\baselineskip}
\textbf{\color{ink}6.} {ootnotesizecolor{muted}Y1 | published | c1da0e32}
\par
Word hujjatida asosiy sarlavha va uning ostidagi bo'lim sarlavhasi Navigation Pane va avtomatik mundarijada ierarxiya sifatida chiqishi kerak. Qaysi uslub jufti mos?
\begin{enumerate}[label=\Alph*),leftmargin=9mm]
\item Asosiy sarlavhaga Title, bo'lim sarlavhasiga Subtitle uslubini berish
\item Asosiy sarlavhaga Heading 1, bo'lim sarlavhasiga Heading 2 uslubini berish
\item Asosiy sarlavhaga Heading 2, bo'lim sarlavhasiga esa Heading 1 berish
\item Asosiy sarlavhani Bold, bo'lim sarlavhasini Italic qilib formatlash
\end{enumerate}
\begin{tcolorbox}[enhanced,breakable,colback=teal!4,colframe=teal!55!black,boxrule=0.45pt,arc=1.2mm,left=2.8mm,right=2.8mm,top=2mm,bottom=2mm]
\textbf{\color{teal}To'g'ri javob:} \textbf{B) Asosiy sarlavhaga Heading 1, bo'lim sarlavhasiga Heading 2 uslubini berish}
\par\smallskip
\textbf{Izoh:} Word'dagi Heading 1 va Heading 2 uslublari matnning tuzilmaviy darajasini belgilaydi. Oddiy qalin yoki kursiv format tuzilmaviy darajani o'zi belgilamaydi.
\end{tcolorbox}
\vspace{2mm}\hrule height .2pt\vspace{2mm}

% DB f2e2ff30-d6cb-4669-83e7-e9f4d65bea01 status=published
\needspace{6\baselineskip}
\textbf{\color{ink}7.} {ootnotesizecolor{muted}Y1 | published | f2e2ff30}
\par
MS Word jadvalida birinchi qator (\texttt{Familiya}, \texttt{Ism}, \texttt{Sinf}) qalin va markazga tekislangan. Shuningdek, \texttt{Karimova}, \texttt{Muhammadiyeva}, \texttt{Doniyor}, \texttt{Layli} yozilgan kataklar och to'q sariq rangga bo'yalgan. Qaysi amallar bajarilgan?
\begin{enumerate}[label=\Alph*),leftmargin=9mm]
\item Birinchi qatorga Ctrl+B va Ctrl+J; so'zlarni Shift+ikki marta bosish bilan ajratish; shading
\item Birinchi qatorga Ctrl+G va Ctrl+E; so'zlarni Ctrl+ikki marta bosish bilan ajratish; text highlight
\item Jadvalning 1-qatorini tanlab Ctrl+B va Ctrl+E ni qo'llash; kerakli kataklarni alohida tanlab, Shading orqali ularning fonini bo'yash.
\item Birinchi qatorga Ctrl+A va Ctrl+E; so'zlarni Alt+ikki marta bosish bilan ajratish; text highlight
\end{enumerate}
\par\smallskip\textit{Rasmli savol: }\href{https://github.com/sarvar9417/attestatsiya/blob/main/public/kuzgi-demo/demo_q19_wordtable.svg}{Rasmni onlayn ochish}
\begin{tcolorbox}[enhanced,breakable,colback=teal!4,colframe=teal!55!black,boxrule=0.45pt,arc=1.2mm,left=2.8mm,right=2.8mm,top=2mm,bottom=2mm]
\textbf{\color{teal}To'g'ri javob:} \textbf{C) Jadvalning 1-qatorini tanlab Ctrl+B va Ctrl+E ni qo'llash; kerakli kataklarni alohida tanlab, Shading orqali ularning fonini bo'yash.}
\par\smallskip
\textbf{Izoh:} Ctrl+B tanlangan sarlavha qatorini qalin qiladi; Ctrl+E markazga tekislaydi. Kerakli jadval kataklari tanlanib, Shading vositasi bilan ularning orqa foni ranglanadi. Word'da Ctrl+ikki marta bosish katakni ishonchli tanlash usuli emas, shu sabab manba variantidagi amaliy ketma-ketlik aniqlashtirildi.
\end{tcolorbox}
\vspace{2mm}\hrule height .2pt\vspace{2mm}
\section{Microsoft Excel}
\subsection{MS Excel: formulalar va funksiyalar}

% DB 8b72860b-ccb8-4995-870f-75f30e6c24eb status=published
\needspace{6\baselineskip}
\textbf{\color{ink}8.} {ootnotesizecolor{muted}Y1 | published | 8b72860b}
\par
Excelda A1:A8 diapazonidagi sonlar yig'indisi 400. A9 katakka \texttt{=AVERAGE(A1:A8)-2} formulasi yozilgan. A9 qiymatini toping.
\begin{enumerate}[label=\Alph*),leftmargin=9mm]
\item 24
\item 48
\item 398
\item 400
\end{enumerate}
\begin{tcolorbox}[enhanced,breakable,colback=teal!4,colframe=teal!55!black,boxrule=0.45pt,arc=1.2mm,left=2.8mm,right=2.8mm,top=2mm,bottom=2mm]
\textbf{\color{teal}To'g'ri javob:} \textbf{B) 48}
\par\smallskip
\textbf{Izoh:} Sakkizta sonning o'rtachasi 400/8=50; 50-2=48.
\end{tcolorbox}
\vspace{2mm}\hrule height .2pt\vspace{2mm}

% DB 329e48e4-a68c-46ea-a977-74bef7aa326a status=published
\needspace{6\baselineskip}
\textbf{\color{ink}9.} {ootnotesizecolor{muted}Y1 | published | 329e48e4}
\par
Excelda A1=1 va B1=2 kataklari belgilanib, to'ldirish markeri BA1 gacha tortildi. Ketma-ketlik har ustunda 1 ga oshsa, 33 soni qaysi katakda hosil bo'ladi?
\begin{enumerate}[label=\Alph*),leftmargin=9mm]
\item BA1
\item AE1
\item AG1
\item AH1
\end{enumerate}
\begin{tcolorbox}[enhanced,breakable,colback=teal!4,colframe=teal!55!black,boxrule=0.45pt,arc=1.2mm,left=2.8mm,right=2.8mm,top=2mm,bottom=2mm]
\textbf{\color{teal}To'g'ri javob:} \textbf{C) AG1}
\par\smallskip
\textbf{Izoh:} Ustun tartib raqami qiymatga teng: Z=26, AA=27, ..., AG=33. Demak 33 soni AG1 da.
\end{tcolorbox}
\vspace{2mm}\hrule height .2pt\vspace{2mm}

% DB 7e9d82cc-fcdf-4954-9695-a049c0ba16e4 status=published
\needspace{6\baselineskip}
\textbf{\color{ink}10.} {ootnotesizecolor{muted}Y1 | published | 7e9d82cc}
\par
Excel jadvali: A1=4, B1=-1, C1=3; A2=5, B2=4, C2=-3; A3=13, B3=9, C3=0. D4 katakda \texttt{=B1*\$C\$2+B3} formula bor. Formula D4 dan C5 ga ko'chirilsa, natija nima bo'ladi?
\begin{enumerate}[label=\Alph*),leftmargin=9mm]
\item -15
\item 15
\item -16
\item 16
\end{enumerate}
\begin{tcolorbox}[enhanced,breakable,colback=teal!4,colframe=teal!55!black,boxrule=0.45pt,arc=1.2mm,left=2.8mm,right=2.8mm,top=2mm,bottom=2mm]
\textbf{\color{teal}To'g'ri javob:} \textbf{A) -15}
\par\smallskip
\textbf{Izoh:} Bir ustun chapga va bir qator pastga ko'chirganda B1\ensuremath{\rightarrow}A2, \$C\$2 o'zgarmaydi, B3\ensuremath{\rightarrow}A4. A4 bo'sh katak bo'lib 0 sifatida olinadi: 5\ensuremath{\cdot}(-3)+0=-15.
\end{tcolorbox}
\vspace{2mm}\hrule height .2pt\vspace{2mm}

% DB 91add323-806c-4342-8793-40260bb3f33c status=published
\needspace{6\baselineskip}
\textbf{\color{ink}11.} {ootnotesizecolor{muted}Y1 | published | 91add323}
\par
Excelda A1=4, B1=3, C1=4, D1=\texttt{=A1+C1}; B2=\texttt{=B1+3}, C2=\texttt{=A1+C1/2}, D2=\texttt{=D1-2}. Diagrammada B2, C2, D2 qiymatlari 6, A2 qiymati esa 5 bo'lishi kerak. A2 uchun qaysi formula mos?
\begin{enumerate}[label=\Alph*),leftmargin=9mm]
\item \texttt{=B1+A1}
\item \texttt{=A1+C1/4}
\item \texttt{=A1+C1/2}
\item \texttt{=C1*A1-D1}
\end{enumerate}
\begin{tcolorbox}[enhanced,breakable,colback=teal!4,colframe=teal!55!black,boxrule=0.45pt,arc=1.2mm,left=2.8mm,right=2.8mm,top=2mm,bottom=2mm]
\textbf{\color{teal}To'g'ri javob:} \textbf{B) \texttt{=A1+C1/4}}
\par\smallskip
\textbf{Izoh:} \texttt{=A1+C1/4} = 4+4/4 = 5. Qolgan formulalar 7, 6 va 8 natijalarini beradi.
\end{tcolorbox}
\vspace{2mm}\hrule height .2pt\vspace{2mm}

% DB a5f4df57-6ffc-436f-ae18-363697a9a4bc status=published
\needspace{6\baselineskip}
\textbf{\color{ink}12.} {ootnotesizecolor{muted}Y1 | published | a5f4df57}
\par
MS Excel jadvali: A1=5, B1=1, C1=3; A2=16, B2=4, C2=6; A3=-7, B3=3, C3=-2. \texttt{=NOK(KOREN'(B2);ZNAK(B2);SChYoTESLI(A1:C3;"\textless{}=0"))} formulasining natijasini toping. (\texttt{NOK}=LCM, \texttt{KOREN'}=SQRT, \texttt{ZNAK}=SIGN, \texttt{SChYoTESLI}=COUNTIF.)
\begin{enumerate}[label=\Alph*),leftmargin=9mm]
\item 8
\item 4
\item 2
\item 6
\end{enumerate}
\begin{tcolorbox}[enhanced,breakable,colback=teal!4,colframe=teal!55!black,boxrule=0.45pt,arc=1.2mm,left=2.8mm,right=2.8mm,top=2mm,bottom=2mm]
\textbf{\color{teal}To'g'ri javob:} \textbf{C) 2}
\par\smallskip
\textbf{Izoh:} SQRT(4)=2, SIGN(4)=1, diapazondagi noldan katta bo'lmagan qiymatlar -7 va -2, ya'ni 2 ta. LCM(2,1,2)=2.
\end{tcolorbox}
\vspace{2mm}\hrule height .2pt\vspace{2mm}

% DB 0e99011d-ce04-41d0-8525-86418b1f042c status=published
\needspace{6\baselineskip}
\textbf{\color{ink}13.} {ootnotesizecolor{muted}Y1 | published | 0e99011d}
\par
A2=12, B2=5 bo'lsa, \texttt{=MAX(A2,B2)*2} formula natijasi qancha?
\begin{enumerate}[label=\Alph*),leftmargin=9mm]
\item 24
\item 17
\item 10
\item 7
\end{enumerate}
\begin{tcolorbox}[enhanced,breakable,colback=teal!4,colframe=teal!55!black,boxrule=0.45pt,arc=1.2mm,left=2.8mm,right=2.8mm,top=2mm,bottom=2mm]
\textbf{\color{teal}To'g'ri javob:} \textbf{A) 24}
\par\smallskip
\textbf{Izoh:} MAX(12,5)=12 va 12*2=24.
\end{tcolorbox}
\vspace{2mm}\hrule height .2pt\vspace{2mm}

% DB 11cb237e-3317-4838-b18a-d5bda5513df6 status=published
\needspace{6\baselineskip}
\textbf{\color{ink}14.} {ootnotesizecolor{muted}Y1 | published | 11cb237e}
\par
Excelda A1=6, B1=0, C1=12, A2=-2, B2=7, C2=5. E5 katakdagi \texttt{=SUM(\$A\$1:\$B\$1)+\$A2*C2} formula D4 ga ko'chirilsa, natija nechaga teng?
\begin{enumerate}[label=\Alph*),leftmargin=9mm]
\item 3
\item -4
\item -8
\item 6
\end{enumerate}
\begin{tcolorbox}[enhanced,breakable,colback=teal!4,colframe=teal!55!black,boxrule=0.45pt,arc=1.2mm,left=2.8mm,right=2.8mm,top=2mm,bottom=2mm]
\textbf{\color{teal}To'g'ri javob:} \textbf{D) 6}
\par\smallskip
\textbf{Izoh:} E5 dan D4 ga 1 ustun chap, 1 qator yuqori siljiydi. Mutlaq oraliq SUM(A1:B1)=6; \$A2\ensuremath{\rightarrow}\$A1=6, C2\ensuremath{\rightarrow}B1=0; natija 6.
\end{tcolorbox}
\vspace{2mm}\hrule height .2pt\vspace{2mm}

% DB 247161d4-ff33-4d27-b298-d486beae88ed status=published
\needspace{6\baselineskip}
\textbf{\color{ink}15.} {ootnotesizecolor{muted}Y1 | published | 247161d4}
\par
MS Excel'da B2 katakda \texttt{=A2*\$F\$1} formula bor. Formula B3 katakka nusxalansa, qaysi ko'rinishga keladi?
\begin{enumerate}[label=\Alph*),leftmargin=9mm]
\item \texttt{=A2*\$F\$1}
\item \texttt{=A3*\$F\$1}
\item \texttt{=A3*\$F\$2}
\item \texttt{=\$A\$3*F1}
\end{enumerate}
\begin{tcolorbox}[enhanced,breakable,colback=teal!4,colframe=teal!55!black,boxrule=0.45pt,arc=1.2mm,left=2.8mm,right=2.8mm,top=2mm,bottom=2mm]
\textbf{\color{teal}To'g'ri javob:} \textbf{B) \texttt{=A3*\$F\$1}}
\par\smallskip
\textbf{Izoh:} A2 nisbiy murojaat bo'lgani uchun bir qator pastga ko'chishda A3 ga o'zgaradi. \$F\$1 mutlaq murojaat bo'lgani uchun o'zgarmaydi.
\end{tcolorbox}
\vspace{2mm}\hrule height .2pt\vspace{2mm}

% DB 355c3f26-3f54-49b4-8ed4-6b302d4cdc30 status=published
\needspace{6\baselineskip}
\textbf{\color{ink}16.} {ootnotesizecolor{muted}Y1 | published | 355c3f26}
\par
A1=8, A2=3, A3=11 bo'lsa, \texttt{=MIN(A1:A3)+MAX(A1:A3)} natija qancha?
\begin{enumerate}[label=\Alph*),leftmargin=9mm]
\item 14
\item 11
\item 8
\item 22
\end{enumerate}
\begin{tcolorbox}[enhanced,breakable,colback=teal!4,colframe=teal!55!black,boxrule=0.45pt,arc=1.2mm,left=2.8mm,right=2.8mm,top=2mm,bottom=2mm]
\textbf{\color{teal}To'g'ri javob:} \textbf{A) 14}
\par\smallskip
\textbf{Izoh:} MIN=3, MAX=11; yig'indi 14.
\end{tcolorbox}
\vspace{2mm}\hrule height .2pt\vspace{2mm}

% DB 647ac0b5-29c9-4704-bd3d-e7803f7ea3ab status=published
\needspace{6\baselineskip}
\textbf{\color{ink}17.} {ootnotesizecolor{muted}Y1 | published | 647ac0b5}
\par
C2 katakda \texttt{=A2+B2} formula bor. Formula C5 gacha pastga nusxalansa, C5 dagi formula qanday bo'ladi?
\begin{enumerate}[label=\Alph*),leftmargin=9mm]
\item =A5+B5
\item =A2+B2
\item =\$A\$2+\$B\$2
\item =A5+B2
\end{enumerate}
\begin{tcolorbox}[enhanced,breakable,colback=teal!4,colframe=teal!55!black,boxrule=0.45pt,arc=1.2mm,left=2.8mm,right=2.8mm,top=2mm,bottom=2mm]
\textbf{\color{teal}To'g'ri javob:} \textbf{A) =A5+B5}
\par\smallskip
\textbf{Izoh:} Nisbiy murojaatlar formula bilan birga qatorlar bo'yicha siljiydi.
\end{tcolorbox}
\vspace{2mm}\hrule height .2pt\vspace{2mm}

% DB 670efda0-efee-4b07-ba59-0fc5a435c98b status=published
\needspace{6\baselineskip}
\textbf{\color{ink}18.} {ootnotesizecolor{muted}Y2 | published | 670efda0}
\par
Excel funksiyalarini ularning vazifalari bilan juftlang.
\textbf{Chap ustun:}\begin{enumerate}[label=\arabic*),leftmargin=9mm]
\item SUM
\item AVERAGE
\item MAX
\item IF
\end{enumerate}
\textbf{O'ng ustun:}\begin{enumerate}[label=\Alph*),leftmargin=9mm]
\item Eng katta qiymatni topish
\item Qiymatlarni qo'shish
\item Shartga qarab turli natija qaytarish
\item Arifmetik o'rtachani hisoblash
\end{enumerate}
\begin{tcolorbox}[enhanced,breakable,colback=teal!4,colframe=teal!55!black,boxrule=0.45pt,arc=1.2mm,left=2.8mm,right=2.8mm,top=2mm,bottom=2mm]
\textbf{\color{teal}To'g'ri javob:} \textbf{1-B, 2-D, 3-A, 4-C}
\par\smallskip
\textbf{Izoh:} SUM - yig'indi, AVERAGE - o'rtacha, MAX - maksimum, IF - shartli natija.
\end{tcolorbox}
\vspace{2mm}\hrule height .2pt\vspace{2mm}

% DB 6788edd1-6ae6-4d03-b7e2-986c5758e670 status=published
\needspace{6\baselineskip}
\textbf{\color{ink}19.} {ootnotesizecolor{muted}Y1 | published | 6788edd1}
\par
A1:A10 kataklaridagi sonlarning yig'indisini hisoblash uchun qaysi formula to'g'ri?
\begin{enumerate}[label=\Alph*),leftmargin=9mm]
\item =SUM(A1:A10)
\item =COUNT(A1:A10)
\item =MAX(A1:A10)
\item =IF(A1:A10)
\end{enumerate}
\begin{tcolorbox}[enhanced,breakable,colback=teal!4,colframe=teal!55!black,boxrule=0.45pt,arc=1.2mm,left=2.8mm,right=2.8mm,top=2mm,bottom=2mm]
\textbf{\color{teal}To'g'ri javob:} \textbf{A) =SUM(A1:A10)}
\par\smallskip
\textbf{Izoh:} SUM diapazondagi sonlarni qo'shadi.
\end{tcolorbox}
\vspace{2mm}\hrule height .2pt\vspace{2mm}

% DB 7af54460-d8ed-422c-b0aa-50f551c91774 status=published
\needspace{6\baselineskip}
\textbf{\color{ink}20.} {ootnotesizecolor{muted}Y3 | published | 7af54460}
\par
Excel jadvalida bir xil formulani ustun bo'ylab qo'llash jarayonini tartiblang.
\begin{enumerate}[label=\Alph*),leftmargin=9mm]
\item Formuladagi nisbiy/mutlaq murojaatlarni tekshirish
\item Kirish ma'lumotlarini kataklarga joylashtirish
\item Formulani kerakli kataklarga nusxalab natijani tekshirish
\item Birinchi natija katagiga formulani yozish
\end{enumerate}
\begin{tcolorbox}[enhanced,breakable,colback=teal!4,colframe=teal!55!black,boxrule=0.45pt,arc=1.2mm,left=2.8mm,right=2.8mm,top=2mm,bottom=2mm]
\textbf{\color{teal}To'g'ri javob:} \textbf{B - D - A - C}
\par\smallskip
\textbf{Izoh:} Ma'lumot \ensuremath{\rightarrow} formula \ensuremath{\rightarrow} reference tekshiruvi \ensuremath{\rightarrow} copy va natija nazorati.
\end{tcolorbox}
\vspace{2mm}\hrule height .2pt\vspace{2mm}

% DB 9a57353a-3a70-4078-b764-37f22308490e status=published
\needspace{6\baselineskip}
\textbf{\color{ink}21.} {ootnotesizecolor{muted}Y1 | published | 9a57353a}
\par
B2:B6 kataklaridagi qiymatlarning o'rtachasini topish uchun qaysi formula mos?
\begin{enumerate}[label=\Alph*),leftmargin=9mm]
\item =AVERAGE(B2:B6)
\item =SUM(B2:B6)
\item =MIN(B2:B6)
\item =COUNTIF(B2:B6)
\end{enumerate}
\begin{tcolorbox}[enhanced,breakable,colback=teal!4,colframe=teal!55!black,boxrule=0.45pt,arc=1.2mm,left=2.8mm,right=2.8mm,top=2mm,bottom=2mm]
\textbf{\color{teal}To'g'ri javob:} \textbf{A) =AVERAGE(B2:B6)}
\par\smallskip
\textbf{Izoh:} AVERAGE sonlarning arifmetik o'rtachasini hisoblaydi.
\end{tcolorbox}
\vspace{2mm}\hrule height .2pt\vspace{2mm}

% DB a9aa1ba0-69c0-47ff-ad23-fbca32df51fd status=published
\needspace{6\baselineskip}
\textbf{\color{ink}22.} {ootnotesizecolor{muted}Y1 | published | a9aa1ba0}
\par
D2 katakda \texttt{=C2*\$H\$1} formula bor. Uni D4 ga nusxalaganda qaysi formula hosil bo'ladi?
\begin{enumerate}[label=\Alph*),leftmargin=9mm]
\item =C4*\$H\$1
\item =C2*\$H\$1
\item =C4*\$H\$3
\item =\$C\$4*H1
\end{enumerate}
\begin{tcolorbox}[enhanced,breakable,colback=teal!4,colframe=teal!55!black,boxrule=0.45pt,arc=1.2mm,left=2.8mm,right=2.8mm,top=2mm,bottom=2mm]
\textbf{\color{teal}To'g'ri javob:} \textbf{A) =C4*\$H\$1}
\par\smallskip
\textbf{Izoh:} C2 nisbiy bo'lgani uchun C4 ga siljiydi, \$H\$1 esa mutlaq qoladi.
\end{tcolorbox}
\vspace{2mm}\hrule height .2pt\vspace{2mm}

% DB ba64490e-eee8-43ac-a395-551d1c8eb2bc status=published
\needspace{6\baselineskip}
\textbf{\color{ink}23.} {ootnotesizecolor{muted}Y1 | published | ba64490e}
\par
MS Excel. D4 katakchasida \texttt{=XABX1+B2} ko'rinishida formula mavjud. Jadvalda A1=5, B1=1, C1=3; A2=16, B2=4, C2=6; A3=-7, B3=3, C3=-2. D4 katakchasidagi formula natijasini toping.
\begin{enumerate}[label=\Alph*),leftmargin=9mm]
\item \#IMYa? (\#NAME?)
\item \#ZNACh! (\#VALUE!)
\item \#ChISLO! (\#NUM!)
\item \#SSYLKA! (\#REF!)
\end{enumerate}
\begin{tcolorbox}[enhanced,breakable,colback=teal!4,colframe=teal!55!black,boxrule=0.45pt,arc=1.2mm,left=2.8mm,right=2.8mm,top=2mm,bottom=2mm]
\textbf{\color{teal}To'g'ri javob:} \textbf{A) \#IMYa? (\#NAME?)}
\par\smallskip
\textbf{Izoh:} Formulada Excel tomonidan tanilmagan \texttt{XABX1} nomi ishlatilgani uchun \#NAME? xatosi hosil bo'ladi.
\end{tcolorbox}
\vspace{2mm}\hrule height .2pt\vspace{2mm}

% DB bd4962ac-67fa-4610-8af6-672e33b3d2c8 status=published
\needspace{6\baselineskip}
\textbf{\color{ink}24.} {ootnotesizecolor{muted}Y1 | published | bd4962ac}
\par
M2=45 bo'lsa, \texttt{=IF(M2\textgreater{}=50,"O'tdi","Yiqildi")} natijasi qaysi?
\begin{enumerate}[label=\Alph*),leftmargin=9mm]
\item Yiqildi
\item O'tdi
\item 45
\item FALSE
\end{enumerate}
\begin{tcolorbox}[enhanced,breakable,colback=teal!4,colframe=teal!55!black,boxrule=0.45pt,arc=1.2mm,left=2.8mm,right=2.8mm,top=2mm,bottom=2mm]
\textbf{\color{teal}To'g'ri javob:} \textbf{A) Yiqildi}
\par\smallskip
\textbf{Izoh:} 45 \textgreater{}= 50 sharti yolg'on, shuning uchun ikkinchi matn qaytariladi.
\end{tcolorbox}
\vspace{2mm}\hrule height .2pt\vspace{2mm}

% DB d9c447f5-1777-494d-9768-bbc6a304a269 status=published
\needspace{6\baselineskip}
\textbf{\color{ink}25.} {ootnotesizecolor{muted}Y1 | published | d9c447f5}
\par
M2 katakdagi qiymat 75. Formula \texttt{=IF(M2\textgreater{}=80,"A",IF(M2\textgreater{}=70,"B","C"))} qanday natija qaytaradi?
\begin{enumerate}[label=\Alph*),leftmargin=9mm]
\item A
\item B
\item C
\item FALSE
\end{enumerate}
\begin{tcolorbox}[enhanced,breakable,colback=teal!4,colframe=teal!55!black,boxrule=0.45pt,arc=1.2mm,left=2.8mm,right=2.8mm,top=2mm,bottom=2mm]
\textbf{\color{teal}To'g'ri javob:} \textbf{B) B}
\par\smallskip
\textbf{Izoh:} 75 soni 80 dan kichik, ammo 70 dan katta yoki teng. Shuning uchun ichki IF \texttt{B} natijasini qaytaradi.
\end{tcolorbox}
\vspace{2mm}\hrule height .2pt\vspace{2mm}

% DB e6fe4838-2d2f-40ec-8807-4633c17af713 status=published
\needspace{6\baselineskip}
\textbf{\color{ink}26.} {ootnotesizecolor{muted}Y1 | published | e6fe4838}
\par
Excelda B1 katakda 13\% soliq stavkasi, B3:B6 kataklarda esa xodimlarning maoshlari berilgan. C3 katakda soliq summasini hisoblab, formulani C6 gacha pastga ko'chirmoqchisiz. Har safar maosh tegishli qatordan, soliq stavkasi esa aynan B1 dan olinishi uchun C3 ga qaysi formula kiritiladi?
\begin{enumerate}[label=\Alph*),leftmargin=9mm]
\item =B3*B1
\item =B3*\$B1
\item =B3*\$B\$1
\item =\$B\$3*\$B\$1
\end{enumerate}
\begin{tcolorbox}[enhanced,breakable,colback=teal!4,colframe=teal!55!black,boxrule=0.45pt,arc=1.2mm,left=2.8mm,right=2.8mm,top=2mm,bottom=2mm]
\textbf{\color{teal}To'g'ri javob:} \textbf{C) =B3*\$B\$1}
\par\smallskip
\textbf{Izoh:} C3 da maoshga B3 nisbiy murojaat bo'lib, ko'chirilganda B4, B5, B6 ga o'tadi. Soliq stavkasi \$B\$1 mutlaq murojaatdir. Maosh turgan B3 katakka B3 qatnashgan formula kiritilsa aylana murojaat kelib chiqadi. \newline Manba: https://support.microsoft.com/en-us/excel/switch-between-relative-absolute-and-mixed-references
\end{tcolorbox}
\vspace{2mm}\hrule height .2pt\vspace{2mm}

% DB bcdc9879-05c7-42ea-b503-3db497fdbc3e status=published
\needspace{6\baselineskip}
\textbf{\color{ink}27.} {ootnotesizecolor{muted}Y1 | published | bcdc9879}
\par
Quyidagi Excel jadvalida C - xarid narxi, D - xarid miqdori, E - sotuv narxi, F - sotuv miqdori: \newline  \newline | Qator | Meva | C (so'm) | D (kg) | E (so'm) | F (kg) | \newline |---|---|---:|---:|---:|---:| \newline | 2 | Olma | 18 000 | 24 | 25 000 | 17 | \newline | 3 | Apelsin | 24 000 | 18 | 33 000 | 13 | \newline  \newline G2:G3 uchun shartli formatlash: qiymat 0 dan katta bo'lsa \textbf{yashil}, 0 dan kichik bo'lsa \textbf{qizil}. G2 ga \texttt{=(E2*F2)-(C2*D2)} formula kiritilib G3 ga nusxalandi. G2 va G3 kataklarning fon rangi qanday bo'ladi?
\begin{enumerate}[label=\Alph*),leftmargin=9mm]
\item yashil / yashil
\item yashil / qizil
\item qizil / qizil
\item qizil / yashil
\end{enumerate}
\begin{tcolorbox}[enhanced,breakable,colback=teal!4,colframe=teal!55!black,boxrule=0.45pt,arc=1.2mm,left=2.8mm,right=2.8mm,top=2mm,bottom=2mm]
\textbf{\color{teal}To'g'ri javob:} \textbf{C) qizil / qizil}
\par\smallskip
\textbf{Izoh:} G2 = 25 000\ensuremath{\times}17 - 18 000\ensuremath{\times}24 = -7 000; G3 = 33 000\ensuremath{\times}13 - 24 000\ensuremath{\times}18 = -3 000. Har ikkalasi manfiy, shartli formatlash ikkalasini qizil qiladi. Manbaning asl jadvalidagi raqamlar qo'lda solishtirildi. Individual human editorial approval remains pending.
\end{tcolorbox}
\vspace{2mm}\hrule height .2pt\vspace{2mm}
\subsection{MS Excel: filtr, saralash va diagramma tahlili}

% DB 07badc86-a5e7-4c7f-a594-1c0954c5f5b7 status=published
\needspace{6\baselineskip}
\textbf{\color{ink}28.} {ootnotesizecolor{muted}Y1 | published | 07badc86}
\par
Excel jadvalida A1=2, B1=x, C1=6. Formulalar: A2=\texttt{A1*(B1-C1)}, B2=\texttt{(C1-A1)/2}, C2=\texttt{(C1-A1)*(B1-C1)}. A2:C2 dan tuzilgan doiraviy diagrammada A2 va B2 sektorlarining har biri 1/4, C2 sektori 1/2 qismini egallaydi. B1 katagidagi butun sonni toping.
\begin{enumerate}[label=\Alph*),leftmargin=9mm]
\item 7
\item 5
\item 6
\item 9
\end{enumerate}
\begin{tcolorbox}[enhanced,breakable,colback=teal!4,colframe=teal!55!black,boxrule=0.45pt,arc=1.2mm,left=2.8mm,right=2.8mm,top=2mm,bottom=2mm]
\textbf{\color{teal}To'g'ri javob:} \textbf{A) 7}
\par\smallskip
\textbf{Izoh:} B2=(6-2)/2=2. Diagrammada A2 ham 2: 2(x-6)=2, bundan x=7. Tekshiruv: C2=4(x-6)=4, demak nisbat 2:2:4.
\end{tcolorbox}
\vspace{2mm}\hrule height .2pt\vspace{2mm}

% DB 2940d869-2ef9-4e91-98a8-6194539a3266 status=published
\needspace{6\baselineskip}
\textbf{\color{ink}29.} {ootnotesizecolor{muted}Y1 | published | 2940d869}
\par
Excel jadvalida faqat bali 80 va undan yuqori bo'lgan o'quvchilar satrlarini ko'rsatish kerak. Qaysi amal eng mos?
\begin{enumerate}[label=\Alph*),leftmargin=9mm]
\item Saralash (Sort)
\item Filtrlash (Filter)
\item Diagramma yaratish
\item Kataklarni birlashtirish
\end{enumerate}
\begin{tcolorbox}[enhanced,breakable,colback=teal!4,colframe=teal!55!black,boxrule=0.45pt,arc=1.2mm,left=2.8mm,right=2.8mm,top=2mm,bottom=2mm]
\textbf{\color{teal}To'g'ri javob:} \textbf{B) Filtrlash (Filter)}
\par\smallskip
\textbf{Izoh:} Filtrlash berilgan shartga mos satrlarni ajratib ko'rsatadi; saralash esa yozuvlarning tartibini o'zgartiradi.
\end{tcolorbox}
\vspace{2mm}\hrule height .2pt\vspace{2mm}

% DB 56255278-1eb7-48b7-811b-30892dca71f8 status=published
\needspace{6\baselineskip}
\textbf{\color{ink}30.} {ootnotesizecolor{muted}Y1 | published | 56255278}
\par
Sotuv jadvalida faqat \texttt{Hudud=Navoiy} bo'lgan satrlarni vaqtincha ko'rsatish kerak. Qaysi amal mos?
\begin{enumerate}[label=\Alph*),leftmargin=9mm]
\item Filter
\item Sort
\item Wrap Text
\item Conditional Formatting
\end{enumerate}
\begin{tcolorbox}[enhanced,breakable,colback=teal!4,colframe=teal!55!black,boxrule=0.45pt,arc=1.2mm,left=2.8mm,right=2.8mm,top=2mm,bottom=2mm]
\textbf{\color{teal}To'g'ri javob:} \textbf{A) Filter}
\par\smallskip
\textbf{Izoh:} Filter shartga mos yozuvlarni ko'rsatib, qolganlarini vaqtincha yashiradi.
\end{tcolorbox}
\vspace{2mm}\hrule height .2pt\vspace{2mm}

% DB 58c76288-4dbc-46bd-a47a-5f5fc73199fb status=published
\needspace{6\baselineskip}
\textbf{\color{ink}31.} {ootnotesizecolor{muted}Y1 | published | 58c76288}
\par
Jadvalni Familiya bo'yicha A dan Z gacha tartibga keltirish kerak. Qaysi amal bajariladi?
\begin{enumerate}[label=\Alph*),leftmargin=9mm]
\item Ascending Sort
\item Filter
\item Freeze Panes
\item Merge Cells
\end{enumerate}
\begin{tcolorbox}[enhanced,breakable,colback=teal!4,colframe=teal!55!black,boxrule=0.45pt,arc=1.2mm,left=2.8mm,right=2.8mm,top=2mm,bottom=2mm]
\textbf{\color{teal}To'g'ri javob:} \textbf{A) Ascending Sort}
\par\smallskip
\textbf{Izoh:} A\ensuremath{\rightarrow}Z tartibga keltirish o'sish bo'yicha saralashdir.
\end{tcolorbox}
\vspace{2mm}\hrule height .2pt\vspace{2mm}

% DB 59f44c71-11a0-48b6-b6a9-a99bfdb0aad5 status=published
\needspace{6\baselineskip}
\textbf{\color{ink}32.} {ootnotesizecolor{muted}Y1 | published | 59f44c71}
\par
Umumiy budjetning 40\%, 30\%, 20\%, 10\% qismlarga taqsimlanishini ko'rsatish uchun qaysi diagramma mos?
\begin{enumerate}[label=\Alph*),leftmargin=9mm]
\item Doiraviy (pie) diagramma
\item Chiziqli grafik
\item Faqat jadval chegarasi
\item WordArt
\end{enumerate}
\begin{tcolorbox}[enhanced,breakable,colback=teal!4,colframe=teal!55!black,boxrule=0.45pt,arc=1.2mm,left=2.8mm,right=2.8mm,top=2mm,bottom=2mm]
\textbf{\color{teal}To'g'ri javob:} \textbf{A) Doiraviy (pie) diagramma}
\par\smallskip
\textbf{Izoh:} Butunning ulushlarini ko'rsatishda doiraviy diagramma mos keladi.
\end{tcolorbox}
\vspace{2mm}\hrule height .2pt\vspace{2mm}

% DB 69195bba-c001-46f4-9846-195a740be5f1 status=published
\needspace{6\baselineskip}
\textbf{\color{ink}33.} {ootnotesizecolor{muted}Y1 | published | 69195bba}
\par
To'rtta sinfning o'rtacha ballarini bir-biri bilan solishtirish uchun qaysi diagramma turi qulay?
\begin{enumerate}[label=\Alph*),leftmargin=9mm]
\item Ustunli (column/bar) diagramma
\item Doiraviy diagramma
\item Scatter faqat koordinatasiz
\item Hech qanday diagramma emas
\end{enumerate}
\begin{tcolorbox}[enhanced,breakable,colback=teal!4,colframe=teal!55!black,boxrule=0.45pt,arc=1.2mm,left=2.8mm,right=2.8mm,top=2mm,bottom=2mm]
\textbf{\color{teal}To'g'ri javob:} \textbf{A) Ustunli (column/bar) diagramma}
\par\smallskip
\textbf{Izoh:} Kategoriyalar qiymatini taqqoslashda ustunli diagramma qulay.
\end{tcolorbox}
\vspace{2mm}\hrule height .2pt\vspace{2mm}
\section{Microsoft PowerPoint}
\subsection{MS PowerPoint: taqdimot, animatsiya va o'tish effektlari}

% DB eedc2c79-a46b-49af-9be0-7d7451ea64f8 status=published
\needspace{6\baselineskip}
\textbf{\color{ink}34.} {ootnotesizecolor{muted}Y1 | published | eedc2c79}
\par
PowerPointda o'tish effekti 4 soniya davom etishi, keyingi slaydga 2 soniyadan so'ng avtomatik o'tilishi va \texttt{Bomb} ovozi yangrashi kerak. Qaysi Timing sozlamalari mos?
\begin{enumerate}[label=\Alph*),leftmargin=9mm]
\item Sound=Bomb; Duration=2; After=4; On Mouse Click o'chiq
\item Sound=Bomb; Duration=4; After=2; On Mouse Click o'chiq
\item Sound=Bomb; Duration=3; After=4; On Mouse Click o'chiq
\item Sound=Bomb; Duration=3; After=4; On Mouse Click yoqilgan
\end{enumerate}
\begin{tcolorbox}[enhanced,breakable,colback=teal!4,colframe=teal!55!black,boxrule=0.45pt,arc=1.2mm,left=2.8mm,right=2.8mm,top=2mm,bottom=2mm]
\textbf{\color{teal}To'g'ri javob:} \textbf{B) Sound=Bomb; Duration=4; After=2; On Mouse Click o'chiq}
\par\smallskip
\textbf{Izoh:} Duration o'tish animatsiyasi davomiyligini, After esa avtomatik o'tishgacha bo'lgan vaqtni belgilaydi.
\end{tcolorbox}
\vspace{2mm}\hrule height .2pt\vspace{2mm}

% DB 0bbaadfd-f4c5-410a-a43d-1b30e7d06339 status=published
\needspace{6\baselineskip}
\textbf{\color{ink}35.} {ootnotesizecolor{muted}Y1 | published | 0bbaadfd}
\par
Taqdimotni boshidan namoyish qilish kerak. Qaysi amal eng mos?
\begin{enumerate}[label=\Alph*),leftmargin=9mm]
\item Slide Show \ensuremath{\rightarrow} From Beginning
\item Insert \ensuremath{\rightarrow} Table
\item Review \ensuremath{\rightarrow} Spelling
\item Design \ensuremath{\rightarrow} Format Background
\end{enumerate}
\begin{tcolorbox}[enhanced,breakable,colback=teal!4,colframe=teal!55!black,boxrule=0.45pt,arc=1.2mm,left=2.8mm,right=2.8mm,top=2mm,bottom=2mm]
\textbf{\color{teal}To'g'ri javob:} \textbf{A) Slide Show \ensuremath{\rightarrow} From Beginning}
\par\smallskip
\textbf{Izoh:} From Beginning namoyishni birinchi slayddan boshlaydi.
\end{tcolorbox}
\vspace{2mm}\hrule height .2pt\vspace{2mm}

% DB 12d3f839-6c13-4d22-a480-46c58c3e5501 status=published
\needspace{6\baselineskip}
\textbf{\color{ink}36.} {ootnotesizecolor{muted}Y1 | published | 12d3f839}
\par
PowerPointda 2-slayddagi shakl bosilganda taqdimotning 7-slaydiga darhol o'tish talab qilinadi. Qaysi amal bu interaktiv navigatsiyani bevosita ta'minlaydi?
\begin{enumerate}[label=\Alph*),leftmargin=9mm]
\item Shaklga Animation Path o'rnatib, 7-slayd nomini effektga yozish
\item Shaklga 7-slaydga olib boruvchi giperhavola yoki action berish
\item 7-slaydga Fade o'tishini berib, shaklni Slide Master'ga biriktirish
\item Shaklga Trigger bilan kirish animatsiyasi berib, slayd raqamini ko'rsatish
\end{enumerate}
\begin{tcolorbox}[enhanced,breakable,colback=teal!4,colframe=teal!55!black,boxrule=0.45pt,arc=1.2mm,left=2.8mm,right=2.8mm,top=2mm,bottom=2mm]
\textbf{\color{teal}To'g'ri javob:} \textbf{B) Shaklga 7-slaydga olib boruvchi giperhavola yoki action berish}
\par\smallskip
\textbf{Izoh:} Shaklning klik hodisasiga 7-slaydga olib boruvchi giperhavola yoki Action bog'lanadi. Transition slaydlar o'rtasidagi umumiy o'tish, animatsiya esa obyekt xatti-harakatini boshqaradi.
\end{tcolorbox}
\vspace{2mm}\hrule height .2pt\vspace{2mm}

% DB 3ca0bba8-bb22-4c2d-91d7-27e3e67eefd2 status=published
\needspace{6\baselineskip}
\textbf{\color{ink}37.} {ootnotesizecolor{muted}Y3 | published | 3ca0bba8}
\par
Auditoriya uchun PowerPoint taqdimotini tayyorlash jarayonidagi amallarni dastlabki rejalashtirishdan namoyish oldi tekshiruvigacha bajarilish tartibida joylashtiring.
\begin{enumerate}[label=\Alph*),leftmargin=9mm]
\item Zarur vizual, animatsiya va o'tishlarni me'yorida qo'llash
\item Maqsad va auditoriyani aniqlash
\item Slide Show rejimida tekshirib taqdimotni yakunlash
\item Slaydlar tuzilmasi va asosiy mazmunni yaratish
\end{enumerate}
\begin{tcolorbox}[enhanced,breakable,colback=teal!4,colframe=teal!55!black,boxrule=0.45pt,arc=1.2mm,left=2.8mm,right=2.8mm,top=2mm,bottom=2mm]
\textbf{\color{teal}To'g'ri javob:} \textbf{B - D - A - C}
\par\smallskip
\textbf{Izoh:} Maqsad \ensuremath{\rightarrow} struktura/mazmun \ensuremath{\rightarrow} dizayn/effect \ensuremath{\rightarrow} namoyish tekshiruvi.
\end{tcolorbox}
\vspace{2mm}\hrule height .2pt\vspace{2mm}

% DB 4bc7b59e-d6c8-4042-96da-26f9c53d8528 status=published
\needspace{6\baselineskip}
\textbf{\color{ink}38.} {ootnotesizecolor{muted}Y2 | published | 4bc7b59e}
\par
PowerPoint tushunchalarini mos vazifalari bilan juftlang.
\textbf{Chap ustun:}\begin{enumerate}[label=\arabic*),leftmargin=9mm]
\item Animation
\item Transition
\item Hyperlink/Action
\item Slide Show
\end{enumerate}
\textbf{O'ng ustun:}\begin{enumerate}[label=\Alph*),leftmargin=9mm]
\item Obyekt bosilganda boshqa slayd yoki manzilga o'tish
\item Slayd ichidagi obyektning paydo bo'lishi yoki harakatini boshqarish
\item Taqdimotni namoyish rejimida ko'rsatish
\item Slaydlar orasidagi o'tish effektini boshqarish
\end{enumerate}
\begin{tcolorbox}[enhanced,breakable,colback=teal!4,colframe=teal!55!black,boxrule=0.45pt,arc=1.2mm,left=2.8mm,right=2.8mm,top=2mm,bottom=2mm]
\textbf{\color{teal}To'g'ri javob:} \textbf{1-B, 2-D, 3-A, 4-C}
\par\smallskip
\textbf{Izoh:} Animation obyektga, Transition slaydlar oralig'iga, Hyperlink/Action navigatsiyaga, Slide Show namoyishga tegishli.
\end{tcolorbox}
\vspace{2mm}\hrule height .2pt\vspace{2mm}

% DB 8191bfcf-7a03-4ef8-8bed-58cc628834a2 status=published
\needspace{6\baselineskip}
\textbf{\color{ink}39.} {ootnotesizecolor{muted}Y1 | published | 8191bfcf}
\par
PowerPoint'da matn slayd ochilgandan keyin bosqichma-bosqich paydo bo'lishi kerak. Qaysi xususiyat qo'llanadi?
\begin{enumerate}[label=\Alph*),leftmargin=9mm]
\item Animation
\item Transition
\item Page Orientation
\item Header
\end{enumerate}
\begin{tcolorbox}[enhanced,breakable,colback=teal!4,colframe=teal!55!black,boxrule=0.45pt,arc=1.2mm,left=2.8mm,right=2.8mm,top=2mm,bottom=2mm]
\textbf{\color{teal}To'g'ri javob:} \textbf{A) Animation}
\par\smallskip
\textbf{Izoh:} Animation slayd ichidagi obyektning paydo bo'lishi yoki harakatini boshqaradi.
\end{tcolorbox}
\vspace{2mm}\hrule height .2pt\vspace{2mm}

% DB 82a944e6-892d-4079-9aa5-aea92673eb77 status=published
\needspace{6\baselineskip}
\textbf{\color{ink}40.} {ootnotesizecolor{muted}Y1 | published | 82a944e6}
\par
Slayddagi \texttt{Bosh sahifa} tugmasi bosilganda 1-slaydga qaytishi kerak. Qaysi vosita kerak?
\begin{enumerate}[label=\Alph*),leftmargin=9mm]
\item Action yoki hyperlink
\item Faqat Font Color
\item Crop
\item Slide Size
\end{enumerate}
\begin{tcolorbox}[enhanced,breakable,colback=teal!4,colframe=teal!55!black,boxrule=0.45pt,arc=1.2mm,left=2.8mm,right=2.8mm,top=2mm,bottom=2mm]
\textbf{\color{teal}To'g'ri javob:} \textbf{A) Action yoki hyperlink}
\par\smallskip
\textbf{Izoh:} Action/hyperlink obyektni ma'lum slaydga o'tish bilan bog'lash imkonini beradi.
\end{tcolorbox}
\vspace{2mm}\hrule height .2pt\vspace{2mm}

% DB fc7bc571-ca29-4bec-a971-9bc7f7a0b488 status=published
\needspace{6\baselineskip}
\textbf{\color{ink}41.} {ootnotesizecolor{muted}Y1 | published | fc7bc571}
\par
Bir slayddan keyingisiga o'tishda Fade effekti kerak. Qaysi xususiyat tanlanadi?
\begin{enumerate}[label=\Alph*),leftmargin=9mm]
\item Transition
\item Animation Pane faqat obyekt uchun
\item Spell Check
\item Crop
\end{enumerate}
\begin{tcolorbox}[enhanced,breakable,colback=teal!4,colframe=teal!55!black,boxrule=0.45pt,arc=1.2mm,left=2.8mm,right=2.8mm,top=2mm,bottom=2mm]
\textbf{\color{teal}To'g'ri javob:} \textbf{A) Transition}
\par\smallskip
\textbf{Izoh:} Transition slaydlar orasidagi o'tish effektidir.
\end{tcolorbox}
\vspace{2mm}\hrule height .2pt\vspace{2mm}
\section{Microsoft Word}
\subsection{MS Word: interfeys, hujjat yaratish va tahrirlash}

% DB 4958bc3f-9736-e34b-ce92-efd28c5db6b9 status=published
\needspace{6\baselineskip}
\textbf{\color{ink}42.} {ootnotesizecolor{muted}Y1 | published | 4958bc3f}
\par
Save As buyrug'ining muhim xususiyati qaysi?
\begin{enumerate}[label=\Alph*),leftmargin=9mm]
\item Faqat printerda chiqarish
\item Boshqa nom yoki joyda nusxa saqlash imkoniyati
\item Hujjatdan barcha rasmlarni olib tashlash
\item Barcha hujjatni avtomatik tarjima qilish
\end{enumerate}
\begin{tcolorbox}[enhanced,breakable,colback=teal!4,colframe=teal!55!black,boxrule=0.45pt,arc=1.2mm,left=2.8mm,right=2.8mm,top=2mm,bottom=2mm]
\textbf{\color{teal}To'g'ri javob:} \textbf{B) Boshqa nom yoki joyda nusxa saqlash imkoniyati}
\par\smallskip
\textbf{Izoh:} Save As saqlash nomi, manzili yoki formatini tanlaydi.
\end{tcolorbox}
\vspace{2mm}\hrule height .2pt\vspace{2mm}

% DB 4ce0ae0c-24a7-852d-16e0-a4abf712b30a status=published
\needspace{6\baselineskip}
\textbf{\color{ink}43.} {ootnotesizecolor{muted}Y1 | published | 4ce0ae0c}
\par
Word hujjatining odatiy tahrirlanuvchi kengaytmasi qaysi?
\begin{enumerate}[label=\Alph*),leftmargin=9mm]
\item .pptx
\item .docx
\item .xlsx
\item .mp3
\end{enumerate}
\begin{tcolorbox}[enhanced,breakable,colback=teal!4,colframe=teal!55!black,boxrule=0.45pt,arc=1.2mm,left=2.8mm,right=2.8mm,top=2mm,bottom=2mm]
\textbf{\color{teal}To'g'ri javob:} \textbf{B) .docx}
\par\smallskip
\textbf{Izoh:} DOCX Wordning zamonaviy hujjat formatidir.
\end{tcolorbox}
\vspace{2mm}\hrule height .2pt\vspace{2mm}

% DB 690e3e6d-3be1-f7ad-0af3-7a69a2dda28f status=published
\needspace{6\baselineskip}
\textbf{\color{ink}44.} {ootnotesizecolor{muted}Y2 | published | 690e3e6d}
\par
Quyidagi buyruqlarni vazifasi bilan moslashtiring.  1) Ctrl+C;  2) Ctrl+X;  3) Ctrl+V;  Izohlar: A) Kesib olish;  B) Nusxalash;  C) Joylashtirish;
\textbf{Chap ustun:}\begin{enumerate}[label=\arabic*),leftmargin=9mm]
\item Ctrl+C
\item Ctrl+X
\item Ctrl+V
\end{enumerate}
\textbf{O'ng ustun:}\begin{enumerate}[label=\Alph*),leftmargin=9mm]
\item Kesib olish
\item Nusxalash
\item Joylashtirish
\end{enumerate}
\begin{tcolorbox}[enhanced,breakable,colback=teal!4,colframe=teal!55!black,boxrule=0.45pt,arc=1.2mm,left=2.8mm,right=2.8mm,top=2mm,bottom=2mm]
\textbf{\color{teal}To'g'ri javob:} \textbf{1-B, 2-A, 3-C}
\par\smallskip
\textbf{Izoh:} To'g'ri moslik: Ctrl+C - Nusxalash, Ctrl+X - Kesib olish, Ctrl+V - Joylashtirish.
\end{tcolorbox}
\vspace{2mm}\hrule height .2pt\vspace{2mm}

% DB 95109ea5-cafd-5f65-ab33-6a43732726c3 status=published
\needspace{6\baselineskip}
\textbf{\color{ink}45.} {ootnotesizecolor{muted}Y1 | published | 95109ea5}
\par
MS Word bo'yicha mulohazalarning rostligini aniqlang. 1) Save joriy hujjatdagi o'zgarishlarni saqlaydi. 2) Ctrl+Left kursorni oldingi abzas boshiga o'tkazadi. 3) Ctrl+Right kursorni keyingi abzas boshiga o'tkazadi. 4) Save As faylni boshqa nom yoki formatda saqlash imkonini beradi.
\begin{enumerate}[label=\Alph*),leftmargin=9mm]
\item 1 - rost; 2 - yolg'on; 3 - yolg'on; 4 - rost
\item 1 - yolg'on; 2 - yolg'on; 3 - yolg'on; 4 - rost
\item 1 - rost; 2 - rost; 3 - yolg'on; 4 - yolg'on
\item 1 - yolg'on; 2 - yolg'on; 3 - rost; 4 - rost
\end{enumerate}
\begin{tcolorbox}[enhanced,breakable,colback=teal!4,colframe=teal!55!black,boxrule=0.45pt,arc=1.2mm,left=2.8mm,right=2.8mm,top=2mm,bottom=2mm]
\textbf{\color{teal}To'g'ri javob:} \textbf{A) 1 - rost; 2 - yolg'on; 3 - yolg'on; 4 - rost}
\par\smallskip
\textbf{Izoh:} Save o'zgarishlarni joriy faylga yozadi, Save As esa yangi nom yoki formatni tanlashga imkon beradi. Ctrl+Left va Ctrl+Right so'zlar bo'yicha harakat qiladi; abzaslar bo'yicha o'tish uchun Ctrl+Up/Ctrl+Down ishlatiladi.
\end{tcolorbox}
\vspace{2mm}\hrule height .2pt\vspace{2mm}

% DB b7ac8f64-a42c-f8fb-6751-68546b5b5167 status=published
\needspace{6\baselineskip}
\textbf{\color{ink}46.} {ootnotesizecolor{muted}Y1 | published | b7ac8f64}
\par
Worddagi Ribbon atamasi nimani bildiradi?
\begin{enumerate}[label=\Alph*),leftmargin=9mm]
\item Hujjatning faqat pastki chetini
\item Faylning doimiy xotiradagi manzilini
\item Buyruqlar guruhlangan lenta
\item Har bir sahifa orasidagi bo'shliq
\end{enumerate}
\begin{tcolorbox}[enhanced,breakable,colback=teal!4,colframe=teal!55!black,boxrule=0.45pt,arc=1.2mm,left=2.8mm,right=2.8mm,top=2mm,bottom=2mm]
\textbf{\color{teal}To'g'ri javob:} \textbf{C) Buyruqlar guruhlangan lenta}
\par\smallskip
\textbf{Izoh:} Ribbon yorliq va guruhlarga ajratilgan buyruqlar tasmasidir.
\end{tcolorbox}
\vspace{2mm}\hrule height .2pt\vspace{2mm}

% DB c07fadaf-f737-4747-9e4f-74e3a6677d53 status=published
\needspace{6\baselineskip}
\textbf{\color{ink}47.} {ootnotesizecolor{muted}Y3 | published | c07fadaf}
\par
MS Wordda yangi matnli hujjat yaratib, matn kiritish va uni DOCX formatida saqlash uchun amallarni tartiblang.
\begin{enumerate}[label=\Alph*),leftmargin=9mm]
\item Matnni kiritish
\item Word dasturini ishga tushirish
\item New buyrug'i bilan hujjat yaratish
\item Hujjatni DOCX ko'rinishida saqlash
\end{enumerate}
\begin{tcolorbox}[enhanced,breakable,colback=teal!4,colframe=teal!55!black,boxrule=0.45pt,arc=1.2mm,left=2.8mm,right=2.8mm,top=2mm,bottom=2mm]
\textbf{\color{teal}To'g'ri javob:} \textbf{B - C - A - D}
\par\smallskip
\textbf{Izoh:} Bosqichlar: Word dasturini ishga tushirish \ensuremath{\rightarrow} New buyrug'i bilan hujjat yaratish \ensuremath{\rightarrow} Matnni kiritish \ensuremath{\rightarrow} Hujjatni DOCX ko'rinishida saqlash.
\end{tcolorbox}
\vspace{2mm}\hrule height .2pt\vspace{2mm}

% DB e6089804-7928-7839-51dc-e0cdaaa1dbc8 status=published
\needspace{6\baselineskip}
\textbf{\color{ink}48.} {ootnotesizecolor{muted}Y1 | published | e6089804}
\par
Matn protsessorining asosiy vazifasi qaysi?
\begin{enumerate}[label=\Alph*),leftmargin=9mm]
\item Matnli hujjatlarni yaratish, tahrirlash va formatlash
\item Kompyuter operatsion tizimini yuklash
\item Elektr zanjirini loyihalash
\item Fayllarni antivirusdan tekshirish
\end{enumerate}
\begin{tcolorbox}[enhanced,breakable,colback=teal!4,colframe=teal!55!black,boxrule=0.45pt,arc=1.2mm,left=2.8mm,right=2.8mm,top=2mm,bottom=2mm]
\textbf{\color{teal}To'g'ri javob:} \textbf{A) Matnli hujjatlarni yaratish, tahrirlash va formatlash}
\par\smallskip
\textbf{Izoh:} Word matn protsessoridir.
\end{tcolorbox}
\vspace{2mm}\hrule height .2pt\vspace{2mm}

% DB 30fc60cb-f92b-7fbe-cf41-f3aa8d52162b status=published
\needspace{6\baselineskip}
\textbf{\color{ink}49.} {ootnotesizecolor{muted}Y2 | published | 30fc60cb}
\par
Atamalar va ularning aniq ma'nolarini moslashtiring.  1) New;  2) Open;  3) Save As;  Izohlar: A) Yangi hujjat yaratish;  B) Yangi nom yoki formatda saqlash;  C) Mavjud hujjatni ochish;
\textbf{Chap ustun:}\begin{enumerate}[label=\arabic*),leftmargin=9mm]
\item New
\item Open
\item Save As
\end{enumerate}
\textbf{O'ng ustun:}\begin{enumerate}[label=\Alph*),leftmargin=9mm]
\item Yangi hujjat yaratish
\item Yangi nom yoki formatda saqlash
\item Mavjud hujjatni ochish
\end{enumerate}
\begin{tcolorbox}[enhanced,breakable,colback=teal!4,colframe=teal!55!black,boxrule=0.45pt,arc=1.2mm,left=2.8mm,right=2.8mm,top=2mm,bottom=2mm]
\textbf{\color{teal}To'g'ri javob:} \textbf{1-A, 2-C, 3-B}
\par\smallskip
\textbf{Izoh:} To'g'ri moslik: New - Yangi hujjat yaratish, Open - Mavjud hujjatni ochish, Save As - Yangi nom yoki formatda saqlash.
\end{tcolorbox}
\vspace{2mm}\hrule height .2pt\vspace{2mm}

% DB 55d37611-4c25-c4dd-8550-983aebc9e205 status=published
\needspace{6\baselineskip}
\textbf{\color{ink}50.} {ootnotesizecolor{muted}Y1 | published | 55d37611}
\par
Open va New buyruqlarining farqi qaysi?
\begin{enumerate}[label=\Alph*),leftmargin=9mm]
\item Open matnni qalin qiladi, New chop etadi
\item Ikkalasi ham faqat PDF yaratadi
\item Open mavjud faylni ochadi, New yangi hujjat yaratadi
\item New mavjud faylni o'chiradi
\end{enumerate}
\begin{tcolorbox}[enhanced,breakable,colback=teal!4,colframe=teal!55!black,boxrule=0.45pt,arc=1.2mm,left=2.8mm,right=2.8mm,top=2mm,bottom=2mm]
\textbf{\color{teal}To'g'ri javob:} \textbf{C) Open mavjud faylni ochadi, New yangi hujjat yaratadi}
\par\smallskip
\textbf{Izoh:} Buyruqlar turli fayl amallarini bajaradi.
\end{tcolorbox}
\vspace{2mm}\hrule height .2pt\vspace{2mm}

% DB 8380f170-b219-ae74-2f20-8ddd814e20d2 status=published
\needspace{6\baselineskip}
\textbf{\color{ink}51.} {ootnotesizecolor{muted}Y1 | published | 8380f170}
\par
Hujjatdagi xato jumlani tuzatish bilan fontni o'zgartirish farqi nimada?
\begin{enumerate}[label=\Alph*),leftmargin=9mm]
\item Ikkisi ham faqat kodlash
\item Jumla tuzatish faqat Print menyusida
\item Font o'zgarishi faylni o'chiradi
\item Birinchisi tahrirlash, ikkinchisi formatlash
\end{enumerate}
\begin{tcolorbox}[enhanced,breakable,colback=teal!4,colframe=teal!55!black,boxrule=0.45pt,arc=1.2mm,left=2.8mm,right=2.8mm,top=2mm,bottom=2mm]
\textbf{\color{teal}To'g'ri javob:} \textbf{D) Birinchisi tahrirlash, ikkinchisi formatlash}
\par\smallskip
\textbf{Izoh:} Tahrir mazmunga, format ko'rinishga tegishli.
\end{tcolorbox}
\vspace{2mm}\hrule height .2pt\vspace{2mm}

% DB 993ce815-9964-137c-99f3-5ce6bd9bbef3 status=published
\needspace{6\baselineskip}
\textbf{\color{ink}52.} {ootnotesizecolor{muted}Y1 | published | 993ce815}
\par
Backspace va Delete belgilanmagan matnda qanday farqlanadi?
\begin{enumerate}[label=\Alph*),leftmargin=9mm]
\item Backspace odatda chapdagi, Delete o'ngdagi belgini o'chiradi
\item Delete faqat jadval yaratadi
\item Backspace matnni saqlaydi, Delete chop etadi
\item Ikkisi ham faqat butun sahifani o'chiradi
\end{enumerate}
\begin{tcolorbox}[enhanced,breakable,colback=teal!4,colframe=teal!55!black,boxrule=0.45pt,arc=1.2mm,left=2.8mm,right=2.8mm,top=2mm,bottom=2mm]
\textbf{\color{teal}To'g'ri javob:} \textbf{A) Backspace odatda chapdagi, Delete o'ngdagi belgini o'chiradi}
\par\smallskip
\textbf{Izoh:} Kursorning joyi muhim.
\end{tcolorbox}
\vspace{2mm}\hrule height .2pt\vspace{2mm}

% DB ab9ed652-18ae-4435-9e5e-8514acb8a8e7 status=published
\needspace{6\baselineskip}
\textbf{\color{ink}53.} {ootnotesizecolor{muted}Y1 | published | ab9ed652}
\par
Word hujjatida tanlangan jumlani bir joydan boshqa joyga ko'chirish yoki asl joyida qoldirib nusxalash kerak. Ctrl+X va Ctrl+C o'rtasidagi asosiy farq qaysi?
\begin{enumerate}[label=\Alph*),leftmargin=9mm]
\item Ctrl+X asl jumlani qoldiradi, Ctrl+C esa uni joyidan olib tashlaydi
\item Ctrl+X matnni chop etadi, Ctrl+C faqat uning formatini o'zgartiradi
\item Ctrl+X faylni saqlaydi, Ctrl+C esa barcha sahifani qayta tartiblaydi
\item Ctrl+X manba tanlovni olib tashlaydi, Ctrl+C esa qoldiradi
\end{enumerate}
\begin{tcolorbox}[enhanced,breakable,colback=teal!4,colframe=teal!55!black,boxrule=0.45pt,arc=1.2mm,left=2.8mm,right=2.8mm,top=2mm,bottom=2mm]
\textbf{\color{teal}To'g'ri javob:} \textbf{D) Ctrl+X manba tanlovni olib tashlaydi, Ctrl+C esa qoldiradi}
\par\smallskip
\textbf{Izoh:} Ctrl+X kesish bo'lib tanlangan elementni joyidan olib almashish buferiga o'tkazadi; Ctrl+C esa asl matnni saqlagan holda nusxalaydi. Ikkisidan keyin Ctrl+V joylashtiradi.
\end{tcolorbox}
\vspace{2mm}\hrule height .2pt\vspace{2mm}

% DB ba1b5ab7-43ef-406f-819f-2623ba955b77 status=published
\needspace{6\baselineskip}
\textbf{\color{ink}54.} {ootnotesizecolor{muted}Y3 | published | ba1b5ab7}
\par
MS Word hujjatida takrorlangan iborani boshqa ibora bilan almashtirishdan oldin kerakli matnni topish va tekshirish, so'ng Replace bajarish amallarini tartiblang.
\begin{enumerate}[label=\Alph*),leftmargin=9mm]
\item Hujjatni ochish
\item Replace bilan almashtirish
\item Find maydoniga ibora yozish
\item Topilgan iborani tekshirish
\end{enumerate}
\begin{tcolorbox}[enhanced,breakable,colback=teal!4,colframe=teal!55!black,boxrule=0.45pt,arc=1.2mm,left=2.8mm,right=2.8mm,top=2mm,bottom=2mm]
\textbf{\color{teal}To'g'ri javob:} \textbf{A - C - D - B}
\par\smallskip
\textbf{Izoh:} Bosqichlar: Hujjatni ochish \ensuremath{\rightarrow} Find maydoniga ibora yozish \ensuremath{\rightarrow} Topilgan iborani tekshirish \ensuremath{\rightarrow} Replace bilan almashtirish.
\end{tcolorbox}
\vspace{2mm}\hrule height .2pt\vspace{2mm}

% DB 05aa0a7b-e650-44ff-bba6-524ee54b5a01 status=published
\needspace{6\baselineskip}
\textbf{\color{ink}55.} {ootnotesizecolor{muted}Y1 | published | 05aa0a7b}
\par
Word'da foydalanuvchi harflar tugmalarini bosganda Shiftni ushlab turmay, "kompyuter" o'rniga "KOMPYUTER" yozildi. Sababini eng to'g'ri tushuntiradigan klaviatura holatini tanlang.
\begin{enumerate}[label=\Alph*),leftmargin=9mm]
\item Caps Lock o'chiq, Shift ushlab turilmagan
\item Caps Lock yoqilgan, Shift ushlab turilmagan
\item Num Lock o'chirilgan, Shift ushlab turilmagan
\item Insert rejimi yoqilgan, Shift ushlab turilmagan
\end{enumerate}
\begin{tcolorbox}[enhanced,breakable,colback=teal!4,colframe=teal!55!black,boxrule=0.45pt,arc=1.2mm,left=2.8mm,right=2.8mm,top=2mm,bottom=2mm]
\textbf{\color{teal}To'g'ri javob:} \textbf{B) Caps Lock yoqilgan, Shift ushlab turilmagan}
\par\smallskip
\textbf{Izoh:} Caps Lock yoqilganda Shift ushlab turilmasdan bosilgan lotin harf tugmalari katta harf hosil qiladi. Num Lock raqam klaviaturasiga, Insert esa kiritish/ustidan yozish rejimiga tegishli. Shart Shift bosilmaganini aniq belgilaydi.
\end{tcolorbox}
\vspace{2mm}\hrule height .2pt\vspace{2mm}

% DB 1f79b926-91e8-a628-cc22-5f4177371f83 status=published
\needspace{6\baselineskip}
\textbf{\color{ink}56.} {ootnotesizecolor{muted}Y1 | published | 1f79b926}
\par
Kompyuter o'chishidan avval hujjat yo'qolmasin: eng to'g'ri harakat?
\begin{enumerate}[label=\Alph*),leftmargin=9mm]
\item Font rangini oq qilish
\item Save bilan saqlash va saqlangan joyni tekshirish
\item Rasmga ramka qo'yish
\item Faqat matnni belgilash
\end{enumerate}
\begin{tcolorbox}[enhanced,breakable,colback=teal!4,colframe=teal!55!black,boxrule=0.45pt,arc=1.2mm,left=2.8mm,right=2.8mm,top=2mm,bottom=2mm]
\textbf{\color{teal}To'g'ri javob:} \textbf{B) Save bilan saqlash va saqlangan joyni tekshirish}
\par\smallskip
\textbf{Izoh:} Saqlash hujjatni doimiy xotiraga yozadi.
\end{tcolorbox}
\vspace{2mm}\hrule height .2pt\vspace{2mm}

% DB 3d46211a-bdb5-bf56-3738-9278d686a3c2 status=published
\needspace{6\baselineskip}
\textbf{\color{ink}57.} {ootnotesizecolor{muted}Y2 | published | 3d46211a}
\par
Amallarni izohi bilan moslang.  1) Find;  2) Replace;  3) Word Count;  Izohlar: A) Mos matnni almashtirish;  B) Matnni izlash;  C) So'zlar sonini aniqlash;
\textbf{Chap ustun:}\begin{enumerate}[label=\arabic*),leftmargin=9mm]
\item Find
\item Replace
\item Word Count
\end{enumerate}
\textbf{O'ng ustun:}\begin{enumerate}[label=\Alph*),leftmargin=9mm]
\item Mos matnni almashtirish
\item Matnni izlash
\item So'zlar sonini aniqlash
\end{enumerate}
\begin{tcolorbox}[enhanced,breakable,colback=teal!4,colframe=teal!55!black,boxrule=0.45pt,arc=1.2mm,left=2.8mm,right=2.8mm,top=2mm,bottom=2mm]
\textbf{\color{teal}To'g'ri javob:} \textbf{1-B, 2-A, 3-C}
\par\smallskip
\textbf{Izoh:} To'g'ri moslik: Find - Matnni izlash, Replace - Mos matnni almashtirish, Word Count - So'zlar sonini aniqlash.
\end{tcolorbox}
\vspace{2mm}\hrule height .2pt\vspace{2mm}

% DB 914bb59e-b66e-4907-b60d-8ceb97b793fc status=published
\needspace{6\baselineskip}
\textbf{\color{ink}58.} {ootnotesizecolor{muted}Y3 | published | 914bb59e}
\par
MS Word hujjatidagi oxirgi tahrir xato bajarilgan. Uni Ctrl+Z bilan bekor qilib, natijani tekshirish va to'g'ri nusxasini saqlash amallarini tartiblang.
\begin{enumerate}[label=\Alph*),leftmargin=9mm]
\item Natijani ko'rib chiqish
\item Ctrl+Z bilan so'nggi o'zgarishni bekor qilish
\item To'g'ri nusxani saqlash
\item Xato o'zgartirilgan hujjatni ochiq holda tekshirish
\end{enumerate}
\begin{tcolorbox}[enhanced,breakable,colback=teal!4,colframe=teal!55!black,boxrule=0.45pt,arc=1.2mm,left=2.8mm,right=2.8mm,top=2mm,bottom=2mm]
\textbf{\color{teal}To'g'ri javob:} \textbf{D - B - A - C}
\par\smallskip
\textbf{Izoh:} Bosqichlar: Xato o'zgartirilgan hujjatni ochiq holda tekshirish \ensuremath{\rightarrow} Ctrl+Z bilan so'nggi o'zgarishni bekor qilish \ensuremath{\rightarrow} Natijani ko'rib chiqish \ensuremath{\rightarrow} To'g'ri nusxani saqlash.
\end{tcolorbox}
\vspace{2mm}\hrule height .2pt\vspace{2mm}

% DB ad6917aa-ce60-995a-35f9-8a95043dc226 status=published
\needspace{6\baselineskip}
\textbf{\color{ink}59.} {ootnotesizecolor{muted}Y1 | published | ad6917aa}
\par
O'qituvchi asl faylni buzmasdan tajriba qilmoqchi. Nima qiladi?
\begin{enumerate}[label=\Alph*),leftmargin=9mm]
\item Hujjatni saqlamasdan yopadi
\item Replace Allni bevosita asl faylda bajaradi
\item Save As bilan alohida nusxa yaratadi
\item Ctrl+X bilan barcha matnni kesadi
\end{enumerate}
\begin{tcolorbox}[enhanced,breakable,colback=teal!4,colframe=teal!55!black,boxrule=0.45pt,arc=1.2mm,left=2.8mm,right=2.8mm,top=2mm,bottom=2mm]
\textbf{\color{teal}To'g'ri javob:} \textbf{C) Save As bilan alohida nusxa yaratadi}
\par\smallskip
\textbf{Izoh:} Yangi nusxa dastlabki faylni himoyalaydi.
\end{tcolorbox}
\vspace{2mm}\hrule height .2pt\vspace{2mm}

% DB d1cc1f23-1547-2bd4-709a-aba9af56bcf5 status=published
\needspace{6\baselineskip}
\textbf{\color{ink}60.} {ootnotesizecolor{muted}Y1 | published | d1cc1f23}
\par
Yangi hujjatni tahrirlashga qaytish uchun PDFdan ko'ra nima afzal?
\begin{enumerate}[label=\Alph*),leftmargin=9mm]
\item WAV saqlash
\item DOCX saqlash
\item JPG saqlash
\item MP4 saqlash
\end{enumerate}
\begin{tcolorbox}[enhanced,breakable,colback=teal!4,colframe=teal!55!black,boxrule=0.45pt,arc=1.2mm,left=2.8mm,right=2.8mm,top=2mm,bottom=2mm]
\textbf{\color{teal}To'g'ri javob:} \textbf{B) DOCX saqlash}
\par\smallskip
\textbf{Izoh:} DOCX Worddagi tahrirlash tuzilmasini saqlaydi.
\end{tcolorbox}
\vspace{2mm}\hrule height .2pt\vspace{2mm}

% DB e6143e7e-eaba-b1b9-dd0c-7d90b18a6d29 status=published
\needspace{6\baselineskip}
\textbf{\color{ink}61.} {ootnotesizecolor{muted}Y1 | published | e6143e7e}
\par
Belgilangan xatboshini bir nusxa pastga joylash uchun qaysi kombinatsiya mos?
\begin{enumerate}[label=\Alph*),leftmargin=9mm]
\item Ctrl+C, so'ng Ctrl+V
\item Ctrl+A, so'ng Delete
\item Ctrl+X, so'ng Esc
\item Ctrl+P, so'ng Enter
\end{enumerate}
\begin{tcolorbox}[enhanced,breakable,colback=teal!4,colframe=teal!55!black,boxrule=0.45pt,arc=1.2mm,left=2.8mm,right=2.8mm,top=2mm,bottom=2mm]
\textbf{\color{teal}To'g'ri javob:} \textbf{A) Ctrl+C, so'ng Ctrl+V}
\par\smallskip
\textbf{Izoh:} Copy+Paste nusxani asl joyda saqlaydi.
\end{tcolorbox}
\vspace{2mm}\hrule height .2pt\vspace{2mm}

% DB 1e5cbb05-cba5-4c6c-b79b-5aeafc976c57 status=published
\needspace{6\baselineskip}
\textbf{\color{ink}62.} {ootnotesizecolor{muted}Y3 | published | 1e5cbb05}
\par
MS Wordda asl hujjatni saqlab qolib, alohida nusxa ustida tahrir qilish va natijani tekshirish amallarini tartiblang.
\begin{enumerate}[label=\Alph*),leftmargin=9mm]
\item Save As orqali nusxa saqlash
\item Manba hujjatni ochish
\item Asl fayldan mustaqil yangi nusxani tekshirish
\item Nusxadagi ma'lumotni tahrirlash
\end{enumerate}
\begin{tcolorbox}[enhanced,breakable,colback=teal!4,colframe=teal!55!black,boxrule=0.45pt,arc=1.2mm,left=2.8mm,right=2.8mm,top=2mm,bottom=2mm]
\textbf{\color{teal}To'g'ri javob:} \textbf{B - A - D - C}
\par\smallskip
\textbf{Izoh:} Bosqichlar: Manba hujjatni ochish \ensuremath{\rightarrow} Save As orqali nusxa saqlash \ensuremath{\rightarrow} Nusxadagi ma'lumotni tahrirlash \ensuremath{\rightarrow} Asl fayldan mustaqil yangi nusxani tekshirish.
\end{tcolorbox}
\vspace{2mm}\hrule height .2pt\vspace{2mm}

% DB 3c31df50-5257-6865-4335-16233dd3e2f3 status=published
\needspace{6\baselineskip}
\textbf{\color{ink}63.} {ootnotesizecolor{muted}Y1 | published | 3c31df50}
\par
Wordda matn nusxalangach Paste bajarildi, lekin eski matn ham qoldi. Qaysi amal ishlatilgan?
\begin{enumerate}[label=\Alph*),leftmargin=9mm]
\item Print va Save
\item Delete va Undo
\item Cut va Paste
\item Copy va Paste
\end{enumerate}
\begin{tcolorbox}[enhanced,breakable,colback=teal!4,colframe=teal!55!black,boxrule=0.45pt,arc=1.2mm,left=2.8mm,right=2.8mm,top=2mm,bottom=2mm]
\textbf{\color{teal}To'g'ri javob:} \textbf{D) Copy va Paste}
\par\smallskip
\textbf{Izoh:} Copy aslni o'chirmasdan nusxa oladi.
\end{tcolorbox}
\vspace{2mm}\hrule height .2pt\vspace{2mm}

% DB 5988a816-8bb3-a08c-4512-7a9f0a0dc71f status=published
\needspace{6\baselineskip}
\textbf{\color{ink}64.} {ootnotesizecolor{muted}Y1 | published | 5988a816}
\par
Fayldan faqat bitta jumlani boshqa joyga ko'chirish kerak. Amal ketma-ketligi?
\begin{enumerate}[label=\Alph*),leftmargin=9mm]
\item Ctrl+A va Save
\item Faqat Ctrl+P
\item Esc va Alt+F4
\item Jumlani belgilash, Ctrl+X, kerakli joyda Ctrl+V
\end{enumerate}
\begin{tcolorbox}[enhanced,breakable,colback=teal!4,colframe=teal!55!black,boxrule=0.45pt,arc=1.2mm,left=2.8mm,right=2.8mm,top=2mm,bottom=2mm]
\textbf{\color{teal}To'g'ri javob:} \textbf{D) Jumlani belgilash, Ctrl+X, kerakli joyda Ctrl+V}
\par\smallskip
\textbf{Izoh:} Cut/Paste tanlangan matnni ko'chiradi.
\end{tcolorbox}
\vspace{2mm}\hrule height .2pt\vspace{2mm}

% DB 62f7faa1-2d74-32de-88ac-57d94c44e23f status=published
\needspace{6\baselineskip}
\textbf{\color{ink}65.} {ootnotesizecolor{muted}Y1 | published | 62f7faa1}
\par
Fayl nomi va Word buyruqlari bir oynada qayerlarda ko'rinadi?
\begin{enumerate}[label=\Alph*),leftmargin=9mm]
\item Nomi Title bar, buyruqlar Ribbon qismida
\item Har ikkalasi faqat footerda
\item Faqat operatsion tizim terminalida
\item Nomi Formula Bar, buyruqlar Name Boxda
\end{enumerate}
\begin{tcolorbox}[enhanced,breakable,colback=teal!4,colframe=teal!55!black,boxrule=0.45pt,arc=1.2mm,left=2.8mm,right=2.8mm,top=2mm,bottom=2mm]
\textbf{\color{teal}To'g'ri javob:} \textbf{A) Nomi Title bar, buyruqlar Ribbon qismida}
\par\smallskip
\textbf{Izoh:} Word interfeys elementlarining vazifasi farq qiladi.
\end{tcolorbox}
\vspace{2mm}\hrule height .2pt\vspace{2mm}

% DB 992ba9e6-664b-fdac-faea-c79274d1405f status=published
\needspace{6\baselineskip}
\textbf{\color{ink}66.} {ootnotesizecolor{muted}Y1 | published | 992ba9e6}
\par
Qo'lyozmada kiritish nuqtasidan oldingi xato belgi bor. Qaysi tugma?
\begin{enumerate}[label=\Alph*),leftmargin=9mm]
\item Caps Lock
\item F5
\item Backspace
\item Delete (o'ngdagi belgi uchun)
\end{enumerate}
\begin{tcolorbox}[enhanced,breakable,colback=teal!4,colframe=teal!55!black,boxrule=0.45pt,arc=1.2mm,left=2.8mm,right=2.8mm,top=2mm,bottom=2mm]
\textbf{\color{teal}To'g'ri javob:} \textbf{C) Backspace}
\par\smallskip
\textbf{Izoh:} Backspace odatda kursor chapidagi belgini o'chiradi.
\end{tcolorbox}
\vspace{2mm}\hrule height .2pt\vspace{2mm}

% DB abfd5a3f-5d2d-7012-8953-8a44c14cd75c status=published
\needspace{6\baselineskip}
\textbf{\color{ink}67.} {ootnotesizecolor{muted}Y1 | published | abfd5a3f}
\par
Tahrirchi hujjatda Track Changes yoqilgan holda 12 ta o'zgartirish kiritdi. Yakuniy PDFda tasdiqlanmagan tahrir izlari qolmasligi uchun qaysi nazorat bajarilishi kerak?
\begin{enumerate}[label=\Alph*),leftmargin=9mm]
\item Review bo'limida o'zgarishlarni ko'rib Accept/Reject qilish va sharhlarni tekshirish
\item Ctrl+B bilan butun matnni qalinlashtirish
\item Faqat Zoom qiymatini 100\% qilish
\item Hujjatni nomini o'zgartirib saqlash
\end{enumerate}
\begin{tcolorbox}[enhanced,breakable,colback=teal!4,colframe=teal!55!black,boxrule=0.45pt,arc=1.2mm,left=2.8mm,right=2.8mm,top=2mm,bottom=2mm]
\textbf{\color{teal}To'g'ri javob:} \textbf{A) Review bo'limida o'zgarishlarni ko'rib Accept/Reject qilish va sharhlarni tekshirish}
\par\smallskip
\textbf{Izoh:} Track Changes belgilari avtomatik ma'qullanmaydi; yuborishdan oldin reviziyalar va sharhlarni alohida nazorat qilish kerak.
\end{tcolorbox}
\vspace{2mm}\hrule height .2pt\vspace{2mm}

% DB fff0d759-6429-3b4c-7486-649479be67e7 status=published
\needspace{6\baselineskip}
\textbf{\color{ink}68.} {ootnotesizecolor{muted}Y2 | published | fff0d759}
\par
Ushbu vaziyatdagi vositalar va vazifalarni moslashtiring.  1) DOCX;  2) PDF;  3) TXT;  Izohlar: A) Tahrirlanadigan matn hujjati;  B) Oddiy formatlanmagan matn;  C) Maket barqaror bo'ladigan tarqatma hujjat;
\textbf{Chap ustun:}\begin{enumerate}[label=\arabic*),leftmargin=9mm]
\item DOCX
\item PDF
\item TXT
\end{enumerate}
\textbf{O'ng ustun:}\begin{enumerate}[label=\Alph*),leftmargin=9mm]
\item Tahrirlanadigan matn hujjati
\item Oddiy formatlanmagan matn
\item Maket barqaror bo'ladigan tarqatma hujjat
\end{enumerate}
\begin{tcolorbox}[enhanced,breakable,colback=teal!4,colframe=teal!55!black,boxrule=0.45pt,arc=1.2mm,left=2.8mm,right=2.8mm,top=2mm,bottom=2mm]
\textbf{\color{teal}To'g'ri javob:} \textbf{1-A, 2-C, 3-B}
\par\smallskip
\textbf{Izoh:} To'g'ri moslik: DOCX - Tahrirlanadigan matn hujjati, PDF - Maket barqaror bo'ladigan tarqatma hujjat, TXT - Oddiy formatlanmagan matn.
\end{tcolorbox}
\vspace{2mm}\hrule height .2pt\vspace{2mm}

% DB 10788df2-5052-6b44-91ee-0cdaae03ccec status=published
\needspace{6\baselineskip}
\textbf{\color{ink}69.} {ootnotesizecolor{muted}Y1 | published | 10788df2}
\par
Hamkasb DOCXga tahrirlar va izohlar yubordi. 60 betli hujjatni yakuniy PDFga aylantirishdan oldin qaysi nazorat eng to'g'ri?
\begin{enumerate}[label=\Alph*),leftmargin=9mm]
\item Barcha matnni TXTga o'tkazish va shundan PDF chiqarish
\item Track Changes tahrirlarini ko'rib Accept/Reject qilish, izohlarni tekshirish va PDF Previewni ko'rish
\item Hujjatdagi barcha izohni avtomatik o'chirish va darhol PDF qilish
\item Faqat Word Countni tekshirib asl nusxa ustiga saqlash
\end{enumerate}
\begin{tcolorbox}[enhanced,breakable,colback=teal!4,colframe=teal!55!black,boxrule=0.45pt,arc=1.2mm,left=2.8mm,right=2.8mm,top=2mm,bottom=2mm]
\textbf{\color{teal}To'g'ri javob:} \textbf{B) Track Changes tahrirlarini ko'rib Accept/Reject qilish, izohlarni tekshirish va PDF Previewni ko'rish}
\par\smallskip
\textbf{Izoh:} Mazmuni tekshirilgan tahrirlar tasdiqlanadi, so'ng PDFning sahifalashi tekshiriladi.
\end{tcolorbox}
\vspace{2mm}\hrule height .2pt\vspace{2mm}

% DB a04012c8-d664-e29c-0c7e-95f21d765f81 status=published
\needspace{6\baselineskip}
\textbf{\color{ink}70.} {ootnotesizecolor{muted}Y1 | published | a04012c8}
\par
Tasdiqlangan DOCX hujjatini saqlagan holda yangi buyruq tayyorlash va tarqatish uchun PDF berish kerak. Xavfsiz ish tartibi qaysi?
\begin{enumerate}[label=\Alph*),leftmargin=9mm]
\item Open \ensuremath{\rightarrow} Save As (yangi DOCX) \ensuremath{\rightarrow} tahrirlash \ensuremath{\rightarrow} tekshirish \ensuremath{\rightarrow} PDF eksport
\item Open \ensuremath{\rightarrow} Edit \ensuremath{\rightarrow} asl DOCX ustiga saqlash \ensuremath{\rightarrow} chop etish
\item New \ensuremath{\rightarrow} asl DOCXni o'chirish \ensuremath{\rightarrow} nusxani printerga yuborish
\item Open \ensuremath{\rightarrow} PDF eksport \ensuremath{\rightarrow} PDFda Track Changes bilan tahrirlash
\end{enumerate}
\begin{tcolorbox}[enhanced,breakable,colback=teal!4,colframe=teal!55!black,boxrule=0.45pt,arc=1.2mm,left=2.8mm,right=2.8mm,top=2mm,bottom=2mm]
\textbf{\color{teal}To'g'ri javob:} \textbf{A) Open \ensuremath{\rightarrow} Save As (yangi DOCX) \ensuremath{\rightarrow} tahrirlash \ensuremath{\rightarrow} tekshirish \ensuremath{\rightarrow} PDF eksport}
\par\smallskip
\textbf{Izoh:} Asl faylni buzmaslik uchun avval yangi tahrirlanuvchi nusxa yaratiladi; yakunda PDF chiqariladi.
\end{tcolorbox}
\vspace{2mm}\hrule height .2pt\vspace{2mm}

% DB bcaa8e2b-2cd5-4d70-93f5-afe97d6ffbd3 status=published
\needspace{6\baselineskip}
\textbf{\color{ink}71.} {ootnotesizecolor{muted}Y3 | published | bcaa8e2b}
\par
MS Wordda mavjud hujjatni ochib, matn xatolarini tuzatish, tahrir natijasini tekshirish va hujjatning yangi nomli nusxasini Save As bilan saqlash kerak. Amallarni tartiblang.
\begin{enumerate}[label=\Alph*),leftmargin=9mm]
\item Hujjat ko'rinishini tekshirish
\item Matndagi xatolarni tahrirlash
\item Mavjud hujjatni Open bilan ochish
\item Save As bilan yangi nusxani saqlash
\end{enumerate}
\begin{tcolorbox}[enhanced,breakable,colback=teal!4,colframe=teal!55!black,boxrule=0.45pt,arc=1.2mm,left=2.8mm,right=2.8mm,top=2mm,bottom=2mm]
\textbf{\color{teal}To'g'ri javob:} \textbf{C - B - A - D}
\par\smallskip
\textbf{Izoh:} Oldin mavjud hujjat ochiladi, xatolar tuzatiladi, hujjat tekshiriladi va yangi nusxasi Save As bilan saqlanadi. Tartib: 3 \ensuremath{\rightarrow} 2 \ensuremath{\rightarrow} 1 \ensuremath{\rightarrow} 4.
\end{tcolorbox}
\vspace{2mm}\hrule height .2pt\vspace{2mm}

% DB be9e796b-c61b-b9fb-2a9a-94c40c806a53 status=published
\needspace{6\baselineskip}
\textbf{\color{ink}72.} {ootnotesizecolor{muted}Y2 | published | be9e796b}
\par
Moslashtiring.  1) Save As;  2) Replace;  3) Delete;  Ta'riflar: A) Kursor o'ngidagi belgini o'chirish;  B) Boshqa nom/joyda saqlash;  C) Topilgan matnni almashtirish;
\textbf{Chap ustun:}\begin{enumerate}[label=\arabic*),leftmargin=9mm]
\item Save As
\item Replace
\item Delete
\end{enumerate}
\textbf{O'ng ustun:}\begin{enumerate}[label=\Alph*),leftmargin=9mm]
\item Kursor o'ngidagi belgini o'chirish
\item Boshqa nom/joyda saqlash
\item Topilgan matnni almashtirish
\end{enumerate}
\begin{tcolorbox}[enhanced,breakable,colback=teal!4,colframe=teal!55!black,boxrule=0.45pt,arc=1.2mm,left=2.8mm,right=2.8mm,top=2mm,bottom=2mm]
\textbf{\color{teal}To'g'ri javob:} \textbf{1-B, 2-C, 3-A}
\par\smallskip
\textbf{Izoh:} Mos keladigan ta'riflar: 1-B, 2-C, 3-A.
\end{tcolorbox}
\vspace{2mm}\hrule height .2pt\vspace{2mm}

% DB f00d476c-8164-8d5e-dfef-7769fca10caa status=published
\needspace{6\baselineskip}
\textbf{\color{ink}73.} {ootnotesizecolor{muted}Y2 | published | f00d476c}
\par
Quyidagi tushunchalarni mazmuniga qarab bog'lang.  1) Track Changes;  2) Comments;  3) Accept;  Izohlar: A) Tahrirlarni belgilab kuzatish;  B) Taklif etilgan o'zgarishni tasdiqlash;  C) Muhokama izohi qoldirish;
\textbf{Chap ustun:}\begin{enumerate}[label=\arabic*),leftmargin=9mm]
\item Track Changes
\item Comments
\item Accept
\end{enumerate}
\textbf{O'ng ustun:}\begin{enumerate}[label=\Alph*),leftmargin=9mm]
\item Tahrirlarni belgilab kuzatish
\item Taklif etilgan o'zgarishni tasdiqlash
\item Muhokama izohi qoldirish
\end{enumerate}
\begin{tcolorbox}[enhanced,breakable,colback=teal!4,colframe=teal!55!black,boxrule=0.45pt,arc=1.2mm,left=2.8mm,right=2.8mm,top=2mm,bottom=2mm]
\textbf{\color{teal}To'g'ri javob:} \textbf{1-A, 2-C, 3-B}
\par\smallskip
\textbf{Izoh:} To'g'ri moslik: Track Changes - Tahrirlarni belgilab kuzatish, Comments - Muhokama izohi qoldirish, Accept - Taklif etilgan o'zgarishni tasdiqlash.
\end{tcolorbox}
\vspace{2mm}\hrule height .2pt\vspace{2mm}

% DB f6679c29-098f-41ef-81f4-b673e8927b11 status=published
\needspace{6\baselineskip}
\textbf{\color{ink}74.} {ootnotesizecolor{muted}Y3 | published | f6679c29}
\par
MS Word hujjatida Track Changes orqali ko'rsatilgan tahrirlar mavjud. Yakuniy hujjatni PDF shaklida tarqatishdan oldin o'zgarishlarni qabul qilish, hujjatni tekshirish va eksport qilish kerak. Amallarni tartiblang.
\begin{enumerate}[label=\Alph*),leftmargin=9mm]
\item PDF formatida eksport qilish
\item Kerakli o'zgarishlarni qabul qilish
\item Hujjatni yakuniy tekshirish
\item PDFni ochib sahifalashni tekshirish
\end{enumerate}
\begin{tcolorbox}[enhanced,breakable,colback=teal!4,colframe=teal!55!black,boxrule=0.45pt,arc=1.2mm,left=2.8mm,right=2.8mm,top=2mm,bottom=2mm]
\textbf{\color{teal}To'g'ri javob:} \textbf{B - C - A - D}
\par\smallskip
\textbf{Izoh:} Yakuniy nusxa uchun avval kerakli tahrirlar qabul qilinadi, keyin hujjat tekshiriladi. Shundan so'ng PDFga eksport qilinadi va tayyor PDFdagi sahifalar tekshiriladi. Tartib: 2 \ensuremath{\rightarrow} 3 \ensuremath{\rightarrow} 1 \ensuremath{\rightarrow} 4. Eski 3 \ensuremath{\rightarrow} 2 ketma-ketligi xato edi.
\end{tcolorbox}
\vspace{2mm}\hrule height .2pt\vspace{2mm}
\subsection{MS Word: formatlash, abzas va sahifa parametrlari}

% DB 2af67e48-e177-5628-9b3e-a2b940f686a3 status=published
\needspace{6\baselineskip}
\textbf{\color{ink}75.} {ootnotesizecolor{muted}Y1 | published | 2af67e48}
\par
MS Wordda matn qatorlari orasidagi masofani o'zgartirish uchun qaysi uskuna ishlatiladi?
\begin{enumerate}[label=\Alph*),leftmargin=9mm]
\item Margins
\item Paragraph Spacing
\item Line Spacing
\item Page Breaks
\end{enumerate}
\begin{tcolorbox}[enhanced,breakable,colback=teal!4,colframe=teal!55!black,boxrule=0.45pt,arc=1.2mm,left=2.8mm,right=2.8mm,top=2mm,bottom=2mm]
\textbf{\color{teal}To'g'ri javob:} \textbf{C) Line Spacing}
\par\smallskip
\textbf{Izoh:} Qatorlar orasidagi interval \texttt{Line Spacing} orqali boshqariladi. \texttt{Paragraph Spacing} esa abzaslar oldi yoki keyingi masofani belgilaydi.
\end{tcolorbox}
\vspace{2mm}\hrule height .2pt\vspace{2mm}

% DB 2b5b3e67-3a58-b7d4-c0db-61b026f42dec status=published
\needspace{6\baselineskip}
\textbf{\color{ink}76.} {ootnotesizecolor{muted}Y1 | published | 2b5b3e67}
\par
Bold nimani o'zgartiradi?
\begin{enumerate}[label=\Alph*),leftmargin=9mm]
\item Jadval satrlari sonini
\item Fayl kengaytmasini
\item Sahifa orientatsiyasini
\item Matn harflarining qalinligini
\end{enumerate}
\begin{tcolorbox}[enhanced,breakable,colback=teal!4,colframe=teal!55!black,boxrule=0.45pt,arc=1.2mm,left=2.8mm,right=2.8mm,top=2mm,bottom=2mm]
\textbf{\color{teal}To'g'ri javob:} \textbf{D) Matn harflarining qalinligini}
\par\smallskip
\textbf{Izoh:} Bold belgilar ko'rinishidir.
\end{tcolorbox}
\vspace{2mm}\hrule height .2pt\vspace{2mm}

% DB 40c1f2db-ba9e-3b94-f1d6-e0633d68b7d5 status=published
\needspace{6\baselineskip}
\textbf{\color{ink}77.} {ootnotesizecolor{muted}Y1 | published | 40c1f2db}
\par
Heading 1 qanday maqsadga xizmat qiladi?
\begin{enumerate}[label=\Alph*),leftmargin=9mm]
\item Asosiy bo'lim sarlavhasi uslubini beradi
\item Rasmni avtomatik kesadi
\item Faylni Worddan o'chiradi
\item Faqat raqamlarni qo'shadi
\end{enumerate}
\begin{tcolorbox}[enhanced,breakable,colback=teal!4,colframe=teal!55!black,boxrule=0.45pt,arc=1.2mm,left=2.8mm,right=2.8mm,top=2mm,bottom=2mm]
\textbf{\color{teal}To'g'ri javob:} \textbf{A) Asosiy bo'lim sarlavhasi uslubini beradi}
\par\smallskip
\textbf{Izoh:} Heading uslubi tuzilma va mundarija bilan ishlaydi.
\end{tcolorbox}
\vspace{2mm}\hrule height .2pt\vspace{2mm}

% DB 80c51bae-5668-40b2-89ad-a16896ae1e3d status=published
\needspace{6\baselineskip}
\textbf{\color{ink}78.} {ootnotesizecolor{muted}Y3 | published | 80c51bae}
\par
MS Word hujjatida bo'lim sarlavhalari asosida avtomatik mundarija yaratish kerak. Matnni tayyorlash, sarlavhaga uslub qo'llash va mundarija hosil qilish amallarini tartiblang.
\begin{enumerate}[label=\Alph*),leftmargin=9mm]
\item Sarlavhani belgilash
\item Mundarijani yaratish
\item Heading 1 uslubini qo'llash
\item Matnni kiritish
\end{enumerate}
\begin{tcolorbox}[enhanced,breakable,colback=teal!4,colframe=teal!55!black,boxrule=0.45pt,arc=1.2mm,left=2.8mm,right=2.8mm,top=2mm,bottom=2mm]
\textbf{\color{teal}To'g'ri javob:} \textbf{D - A - C - B}
\par\smallskip
\textbf{Izoh:} Bosqichlar: Matnni kiritish \ensuremath{\rightarrow} Sarlavhani belgilash \ensuremath{\rightarrow} Heading 1 uslubini qo'llash \ensuremath{\rightarrow} Mundarijani yaratish.
\end{tcolorbox}
\vspace{2mm}\hrule height .2pt\vspace{2mm}

% DB b3dcb5d4-5d9e-eca1-b721-21a7d433fc25 status=published
\needspace{6\baselineskip}
\textbf{\color{ink}79.} {ootnotesizecolor{muted}Y1 | published | b3dcb5d4}
\par
Bullets buyrug'i nima uchun?
\begin{enumerate}[label=\Alph*),leftmargin=9mm]
\item Tartib raqamisiz markerli ro'yxat uchun
\item Barcha rasmlarni kesish uchun
\item Page Break qo'yish uchun
\item Hujjatni shifrlash uchun
\end{enumerate}
\begin{tcolorbox}[enhanced,breakable,colback=teal!4,colframe=teal!55!black,boxrule=0.45pt,arc=1.2mm,left=2.8mm,right=2.8mm,top=2mm,bottom=2mm]
\textbf{\color{teal}To'g'ri javob:} \textbf{A) Tartib raqamisiz markerli ro'yxat uchun}
\par\smallskip
\textbf{Izoh:} Bullets markerli ro'yxat yaratadi.
\end{tcolorbox}
\vspace{2mm}\hrule height .2pt\vspace{2mm}

% DB ec81ff88-f92f-ee3b-8b80-70bcc5d9034e status=published
\needspace{6\baselineskip}
\textbf{\color{ink}80.} {ootnotesizecolor{muted}Y1 | published | ec81ff88}
\par
Abzasning markazga tekislanishi qanday format darajasi?
\begin{enumerate}[label=\Alph*),leftmargin=9mm]
\item Faqat belgi formatlash
\item Abzas formatlash
\item Fayl nomini o'zgartirish
\item Faylni konvertatsiya qilish
\end{enumerate}
\begin{tcolorbox}[enhanced,breakable,colback=teal!4,colframe=teal!55!black,boxrule=0.45pt,arc=1.2mm,left=2.8mm,right=2.8mm,top=2mm,bottom=2mm]
\textbf{\color{teal}To'g'ri javob:} \textbf{B) Abzas formatlash}
\par\smallskip
\textbf{Izoh:} Alignment abzasga tegishli.
\end{tcolorbox}
\vspace{2mm}\hrule height .2pt\vspace{2mm}

% DB fbecbc4c-652a-3d6a-e526-10144c4509cb status=published
\needspace{6\baselineskip}
\textbf{\color{ink}81.} {ootnotesizecolor{muted}Y2 | published | fbecbc4c}
\par
Quyidagi buyruqlarni vazifasi bilan moslashtiring.  1) Bold;  2) Italic;  3) Underline;  Izohlar: A) Qalin shrift;  B) Tagiga chizilgan shrift;  C) Qiya shrift;
\textbf{Chap ustun:}\begin{enumerate}[label=\arabic*),leftmargin=9mm]
\item Bold
\item Italic
\item Underline
\end{enumerate}
\textbf{O'ng ustun:}\begin{enumerate}[label=\Alph*),leftmargin=9mm]
\item Qalin shrift
\item Tagiga chizilgan shrift
\item Qiya shrift
\end{enumerate}
\begin{tcolorbox}[enhanced,breakable,colback=teal!4,colframe=teal!55!black,boxrule=0.45pt,arc=1.2mm,left=2.8mm,right=2.8mm,top=2mm,bottom=2mm]
\textbf{\color{teal}To'g'ri javob:} \textbf{1-A, 2-C, 3-B}
\par\smallskip
\textbf{Izoh:} To'g'ri moslik: Bold - Qalin shrift, Italic - Qiya shrift, Underline - Tagiga chizilgan shrift.
\end{tcolorbox}
\vspace{2mm}\hrule height .2pt\vspace{2mm}

% DB 15599f0c-323d-b28a-ae69-2c6ad4ea9896 status=published
\needspace{6\baselineskip}
\textbf{\color{ink}82.} {ootnotesizecolor{muted}Y1 | published | 15599f0c}
\par
Enter va Shift+Enter farqi qaysi?
\begin{enumerate}[label=\Alph*),leftmargin=9mm]
\item Ikkisi ham yangi sahifa yaratadi
\item Enter faqat rasm qo'shadi
\item Enter yangi abzas, Shift+Enter abzas ichida satr uzilishi yaratadi
\item Shift+Enter barcha matnni o'chiradi
\end{enumerate}
\begin{tcolorbox}[enhanced,breakable,colback=teal!4,colframe=teal!55!black,boxrule=0.45pt,arc=1.2mm,left=2.8mm,right=2.8mm,top=2mm,bottom=2mm]
\textbf{\color{teal}To'g'ri javob:} \textbf{C) Enter yangi abzas, Shift+Enter abzas ichida satr uzilishi yaratadi}
\par\smallskip
\textbf{Izoh:} Yangi satr va yangi abzas alohida tushunchalar.
\end{tcolorbox}
\vspace{2mm}\hrule height .2pt\vspace{2mm}

% DB 54cf9bb5-dc35-5065-c47c-9c4731e8dbaf status=published
\needspace{6\baselineskip}
\textbf{\color{ink}83.} {ootnotesizecolor{muted}Y2 | published | 54cf9bb5}
\par
Atamalar va ularning aniq ma'nolarini moslashtiring.  1) Align Left;  2) Center;  3) Justify;  Izohlar: A) Ikki chetga tekislash;  B) Markazga tekislash;  C) Chapga tekislash;
\textbf{Chap ustun:}\begin{enumerate}[label=\arabic*),leftmargin=9mm]
\item Align Left
\item Center
\item Justify
\end{enumerate}
\textbf{O'ng ustun:}\begin{enumerate}[label=\Alph*),leftmargin=9mm]
\item Ikki chetga tekislash
\item Markazga tekislash
\item Chapga tekislash
\end{enumerate}
\begin{tcolorbox}[enhanced,breakable,colback=teal!4,colframe=teal!55!black,boxrule=0.45pt,arc=1.2mm,left=2.8mm,right=2.8mm,top=2mm,bottom=2mm]
\textbf{\color{teal}To'g'ri javob:} \textbf{1-C, 2-B, 3-A}
\par\smallskip
\textbf{Izoh:} To'g'ri moslik: Align Left - Chapga tekislash, Center - Markazga tekislash, Justify - Ikki chetga tekislash.
\end{tcolorbox}
\vspace{2mm}\hrule height .2pt\vspace{2mm}

% DB 8c65e6d4-20a2-fc60-673f-7aab441ee02e status=published
\needspace{6\baselineskip}
\textbf{\color{ink}84.} {ootnotesizecolor{muted}Y1 | published | 8c65e6d4}
\par
Page Break va Section Break farqi qanday?
\begin{enumerate}[label=\Alph*),leftmargin=9mm]
\item Section Break faqat matnni qalinlashtiradi
\item Page Break faylni yopadi
\item Ikkalasi ham faqat rasmga ishlaydi
\item Page Break yangi bet, Section Break bo'limni ajratadi
\end{enumerate}
\begin{tcolorbox}[enhanced,breakable,colback=teal!4,colframe=teal!55!black,boxrule=0.45pt,arc=1.2mm,left=2.8mm,right=2.8mm,top=2mm,bottom=2mm]
\textbf{\color{teal}To'g'ri javob:} \textbf{D) Page Break yangi bet, Section Break bo'limni ajratadi}
\par\smallskip
\textbf{Izoh:} Section Break alohida sahifa parametrlari uchun kerak.
\end{tcolorbox}
\vspace{2mm}\hrule height .2pt\vspace{2mm}

% DB c84cb38b-31b7-4448-8331-8e843df6093c status=published
\needspace{6\baselineskip}
\textbf{\color{ink}85.} {ootnotesizecolor{muted}Y3 | published | c84cb38b}
\par
MS Word hujjatidagi bitta abzasni o'ngga tekislab, qator oralig'ini 1,5 ga o'rnatish kerak. Abzasni belgilash, Paragraph boshqaruvini ochish, sozlash va tekshirish amallarini tartiblang.
\begin{enumerate}[label=\Alph*),leftmargin=9mm]
\item Ko'rinishni tekshirish
\item Abzasni belgilash
\item Abzasni o'ngga tekislab, qatorlar oralig'ini 1,5 ga o'rnatish
\item Home menyusida Paragraph sozlamalarini ochish
\end{enumerate}
\begin{tcolorbox}[enhanced,breakable,colback=teal!4,colframe=teal!55!black,boxrule=0.45pt,arc=1.2mm,left=2.8mm,right=2.8mm,top=2mm,bottom=2mm]
\textbf{\color{teal}To'g'ri javob:} \textbf{B - D - C - A}
\par\smallskip
\textbf{Izoh:} Avval abzas belgilanadi, Home menyusidan Paragraph sozlamalari ochiladi, kerakli o'ngga tekislash va 1,5 qator oralig'i o'rnatiladi, so'ng ko'rinish tekshiriladi. Tartib: 2 \ensuremath{\rightarrow} 4 \ensuremath{\rightarrow} 3 \ensuremath{\rightarrow} 1.
\end{tcolorbox}
\vspace{2mm}\hrule height .2pt\vspace{2mm}

% DB d7d9e4b6-031c-d613-812e-f997e1e7113f status=published
\needspace{6\baselineskip}
\textbf{\color{ink}86.} {ootnotesizecolor{muted}Y1 | published | d7d9e4b6}
\par
Margin va Indent farqi nimada?
\begin{enumerate}[label=\Alph*),leftmargin=9mm]
\item Indent hujjatni PDFga aylantiradi
\item Margin faqat rasmlarga, Indent fayl nomiga
\item Margin sahifa chegarasi, Indent abzas chekinishi
\item Ikkisi bir xil
\end{enumerate}
\begin{tcolorbox}[enhanced,breakable,colback=teal!4,colframe=teal!55!black,boxrule=0.45pt,arc=1.2mm,left=2.8mm,right=2.8mm,top=2mm,bottom=2mm]
\textbf{\color{teal}To'g'ri javob:} \textbf{C) Margin sahifa chegarasi, Indent abzas chekinishi}
\par\smallskip
\textbf{Izoh:} Birida sahifa, ikkinchisida abzas o'lchovi.
\end{tcolorbox}
\vspace{2mm}\hrule height .2pt\vspace{2mm}

% DB f0fa1002-1fa8-8a7f-7635-13c60d8bcafc status=published
\needspace{6\baselineskip}
\textbf{\color{ink}87.} {ootnotesizecolor{muted}Y1 | published | f0fa1002}
\par
Underline va Italic farqi qaysi?
\begin{enumerate}[label=\Alph*),leftmargin=9mm]
\item Underline slayd yaratadi
\item Birinchisi tagiga chizadi, ikkinchisi qiya ko'rsatadi
\item Ikkisi faqat abzas chekinishini o'zgartiradi
\item Italic har doim PDFga eksport qiladi
\end{enumerate}
\begin{tcolorbox}[enhanced,breakable,colback=teal!4,colframe=teal!55!black,boxrule=0.45pt,arc=1.2mm,left=2.8mm,right=2.8mm,top=2mm,bottom=2mm]
\textbf{\color{teal}To'g'ri javob:} \textbf{B) Birinchisi tagiga chizadi, ikkinchisi qiya ko'rsatadi}
\par\smallskip
\textbf{Izoh:} Bu belgi formatlash buyruqlari.
\end{tcolorbox}
\vspace{2mm}\hrule height .2pt\vspace{2mm}

% DB 2543db0c-41d1-a6a8-11c9-a73195dab7f3 status=published
\needspace{6\baselineskip}
\textbf{\color{ink}88.} {ootnotesizecolor{muted}Y1 | published | 2543db0c}
\par
Bitta gapni qalin qilish uchun qanday yo'l ma'qul?
\begin{enumerate}[label=\Alph*),leftmargin=9mm]
\item Butun faylni qalin qilish
\item Page Size o'zgartirish
\item Document Protect qo'yish
\item Gapni belgilab Bold qo'llash
\end{enumerate}
\begin{tcolorbox}[enhanced,breakable,colback=teal!4,colframe=teal!55!black,boxrule=0.45pt,arc=1.2mm,left=2.8mm,right=2.8mm,top=2mm,bottom=2mm]
\textbf{\color{teal}To'g'ri javob:} \textbf{D) Gapni belgilab Bold qo'llash}
\par\smallskip
\textbf{Izoh:} Belgi formati tanlangan matnga beriladi.
\end{tcolorbox}
\vspace{2mm}\hrule height .2pt\vspace{2mm}

% DB 37e31558-1f32-4956-9f54-3ec9397a188e status=published
\needspace{6\baselineskip}
\textbf{\color{ink}89.} {ootnotesizecolor{muted}Y3 | published | 37e31558}
\par
MS Word hujjatining faqat o'rtadagi jadval joylashgan bo'limi Landscape, boshqa sahifalari Portrait bo'lishi kerak. Section Break orqali ushbu bo'limni ajratish va orientatsiyasini o'zgartirish amallarini tartiblang.
\begin{enumerate}[label=\Alph*),leftmargin=9mm]
\item O'rta bo'lim orientatsiyasini Landscape qilish
\item Jadval boshlanishiga kursor qo'yish
\item Section Break qo'shish
\item Jadvaldan keyin ikkinchi Section Break qo'shish
\end{enumerate}
\begin{tcolorbox}[enhanced,breakable,colback=teal!4,colframe=teal!55!black,boxrule=0.45pt,arc=1.2mm,left=2.8mm,right=2.8mm,top=2mm,bottom=2mm]
\textbf{\color{teal}To'g'ri javob:} \textbf{B - C - D - A}
\par\smallskip
\textbf{Izoh:} Bosqichlar: Jadval boshlanishiga kursor qo'yish \ensuremath{\rightarrow} Section Break qo'shish \ensuremath{\rightarrow} Jadvaldan keyin ikkinchi Section Break qo'shish \ensuremath{\rightarrow} O'rta bo'lim orientatsiyasini Landscape qilish.
\end{tcolorbox}
\vspace{2mm}\hrule height .2pt\vspace{2mm}

% DB 793dd784-20d8-3722-efb5-fa04ddca32e1 status=published
\needspace{6\baselineskip}
\textbf{\color{ink}90.} {ootnotesizecolor{muted}Y1 | published | 793dd784}
\par
Ko'p qatorni Space bilan tekislash beqaror. Nima afzal?
\begin{enumerate}[label=\Alph*),leftmargin=9mm]
\item Tab stop yoki jadvaldan foydalanish
\item Barcha matnni o'chirish
\item Yuzlab probel qo'yish
\item Matnni rasmga aylantirish
\end{enumerate}
\begin{tcolorbox}[enhanced,breakable,colback=teal!4,colframe=teal!55!black,boxrule=0.45pt,arc=1.2mm,left=2.8mm,right=2.8mm,top=2mm,bottom=2mm]
\textbf{\color{teal}To'g'ri javob:} \textbf{A) Tab stop yoki jadvaldan foydalanish}
\par\smallskip
\textbf{Izoh:} Tabulyatsiya/strukturali joylashuv tahrirda barqaror.
\end{tcolorbox}
\vspace{2mm}\hrule height .2pt\vspace{2mm}

% DB c29f38f0-7610-5ab0-a518-25b7f83235db status=published
\needspace{6\baselineskip}
\textbf{\color{ink}91.} {ootnotesizecolor{muted}Y1 | published | c29f38f0}
\par
Bir necha qatorli reja ketma-ket bajarilishi kerak. Qanday ro'yxat?
\begin{enumerate}[label=\Alph*),leftmargin=9mm]
\item Bullets
\item Numbering
\item Highlight
\item Crop
\end{enumerate}
\begin{tcolorbox}[enhanced,breakable,colback=teal!4,colframe=teal!55!black,boxrule=0.45pt,arc=1.2mm,left=2.8mm,right=2.8mm,top=2mm,bottom=2mm]
\textbf{\color{teal}To'g'ri javob:} \textbf{B) Numbering}
\par\smallskip
\textbf{Izoh:} Rejada tartib raqami muhim.
\end{tcolorbox}
\vspace{2mm}\hrule height .2pt\vspace{2mm}

% DB c5965deb-563a-3d2c-6e39-578d7e830190 status=published
\needspace{6\baselineskip}
\textbf{\color{ink}92.} {ootnotesizecolor{muted}Y1 | published | c5965deb}
\par
A4 hujjatni ko'ndalang qilish uchun qaysi parametr?
\begin{enumerate}[label=\Alph*),leftmargin=9mm]
\item Layout \ensuremath{\rightarrow} Orientation \ensuremath{\rightarrow} Landscape
\item Insert \ensuremath{\rightarrow} Symbol
\item Review \ensuremath{\rightarrow} Spelling
\item Home \ensuremath{\rightarrow} Bold
\end{enumerate}
\begin{tcolorbox}[enhanced,breakable,colback=teal!4,colframe=teal!55!black,boxrule=0.45pt,arc=1.2mm,left=2.8mm,right=2.8mm,top=2mm,bottom=2mm]
\textbf{\color{teal}To'g'ri javob:} \textbf{A) Layout \ensuremath{\rightarrow} Orientation \ensuremath{\rightarrow} Landscape}
\par\smallskip
\textbf{Izoh:} Landscape sahifa yo'nalishidir.
\end{tcolorbox}
\vspace{2mm}\hrule height .2pt\vspace{2mm}

% DB f0777aed-8b03-84e5-6fb7-2e2b24399b4b status=published
\needspace{6\baselineskip}
\textbf{\color{ink}93.} {ootnotesizecolor{muted}Y2 | published | f0777aed}
\par
Amallarni izohi bilan moslang.  1) Enter;  2) Shift+Enter;  3) Ctrl+Enter;  Izohlar: A) Abzas ichida yangi satr;  B) Yangi abzas;  C) Yangi sahifa;
\textbf{Chap ustun:}\begin{enumerate}[label=\arabic*),leftmargin=9mm]
\item Enter
\item Shift+Enter
\item Ctrl+Enter
\end{enumerate}
\textbf{O'ng ustun:}\begin{enumerate}[label=\Alph*),leftmargin=9mm]
\item Abzas ichida yangi satr
\item Yangi abzas
\item Yangi sahifa
\end{enumerate}
\begin{tcolorbox}[enhanced,breakable,colback=teal!4,colframe=teal!55!black,boxrule=0.45pt,arc=1.2mm,left=2.8mm,right=2.8mm,top=2mm,bottom=2mm]
\textbf{\color{teal}To'g'ri javob:} \textbf{1-B, 2-A, 3-C}
\par\smallskip
\textbf{Izoh:} To'g'ri moslik: Enter - Yangi abzas, Shift+Enter - Abzas ichida yangi satr, Ctrl+Enter - Yangi sahifa.
\end{tcolorbox}
\vspace{2mm}\hrule height .2pt\vspace{2mm}

% DB 0dbedbd5-fa83-e255-a9fe-49f93be626f9 status=published
\needspace{6\baselineskip}
\textbf{\color{ink}94.} {ootnotesizecolor{muted}Y1 | published | 0dbedbd5}
\par
Footerda sahifa raqami bor, lekin birinchi sahifada boshqa ko'rinish kerak. Qaysi imkoniyat?
\begin{enumerate}[label=\Alph*),leftmargin=9mm]
\item Find Next
\item Track Changes faqat
\item Ctrl+C
\item Different First Page
\end{enumerate}
\begin{tcolorbox}[enhanced,breakable,colback=teal!4,colframe=teal!55!black,boxrule=0.45pt,arc=1.2mm,left=2.8mm,right=2.8mm,top=2mm,bottom=2mm]
\textbf{\color{teal}To'g'ri javob:} \textbf{D) Different First Page}
\par\smallskip
\textbf{Izoh:} Birinchi sahifa kolontituli alohida bo'lishi mumkin.
\end{tcolorbox}
\vspace{2mm}\hrule height .2pt\vspace{2mm}

% DB 3a2ceb71-44ea-48e5-85b4-85e4092a5eb9 status=published
\needspace{6\baselineskip}
\textbf{\color{ink}95.} {ootnotesizecolor{muted}Y3 | published | 3a2ceb71}
\par
MS Word hujjatining birinchi sahifasida kolontitul boshqacha bo'lishi, qolgan sahifalarda esa sahifa raqamlari to'g'ri ko'rinishi kerak. Different First Page rejimini sozlash va tekshirish amallarini tartiblang.
\begin{enumerate}[label=\Alph*),leftmargin=9mm]
\item Qolgan sahifalardagi raqamlarni tekshirish
\item Birinchi sahifa kolontitulini moslash
\item Different First Page rejimini yoqish
\item Sahifa raqamlarini Insert orqali qo'shish
\end{enumerate}
\begin{tcolorbox}[enhanced,breakable,colback=teal!4,colframe=teal!55!black,boxrule=0.45pt,arc=1.2mm,left=2.8mm,right=2.8mm,top=2mm,bottom=2mm]
\textbf{\color{teal}To'g'ri javob:} \textbf{D - C - B - A}
\par\smallskip
\textbf{Izoh:} Bosqichlar: Sahifa raqamlarini Insert orqali qo'shish \ensuremath{\rightarrow} Different First Page rejimini yoqish \ensuremath{\rightarrow} Birinchi sahifa kolontitulini moslash \ensuremath{\rightarrow} Qolgan sahifalardagi raqamlarni tekshirish.
\end{tcolorbox}
\vspace{2mm}\hrule height .2pt\vspace{2mm}

% DB 3e4a062b-262e-6001-4d0c-93d2c2e81e12 status=published
\needspace{6\baselineskip}
\textbf{\color{ink}96.} {ootnotesizecolor{muted}Y1 | published | 3e4a062b}
\par
Hujjatning bir bo'limi Landscape, qolganlari Portrait bo'lsin. Nima kerak?
\begin{enumerate}[label=\Alph*),leftmargin=9mm]
\item Barcha faylni rasm qilish
\item Alohida Section Break va o'sha bo'lim orientatsiyasi
\item Faqat Ctrl+B
\item Faqat Save buyrug'i
\end{enumerate}
\begin{tcolorbox}[enhanced,breakable,colback=teal!4,colframe=teal!55!black,boxrule=0.45pt,arc=1.2mm,left=2.8mm,right=2.8mm,top=2mm,bottom=2mm]
\textbf{\color{teal}To'g'ri javob:} \textbf{B) Alohida Section Break va o'sha bo'lim orientatsiyasi}
\par\smallskip
\textbf{Izoh:} Section Break mustaqil sahifa parametrlari beradi.
\end{tcolorbox}
\vspace{2mm}\hrule height .2pt\vspace{2mm}

% DB 68bec829-9e73-eb04-94c4-14ccb0fcc24f status=published
\needspace{6\baselineskip}
\textbf{\color{ink}97.} {ootnotesizecolor{muted}Y1 | published | 68bec829}
\par
Bo'sh satrlar sabab hujjatda ortiqcha betlar chiqdi. Qaysi ko'rish yordam beradi?
\begin{enumerate}[label=\Alph*),leftmargin=9mm]
\item Insert Picture
\item Animation Pane
\item Show/Hide formatlash belgilarini ko'rsatish
\item Equation Editor
\end{enumerate}
\begin{tcolorbox}[enhanced,breakable,colback=teal!4,colframe=teal!55!black,boxrule=0.45pt,arc=1.2mm,left=2.8mm,right=2.8mm,top=2mm,bottom=2mm]
\textbf{\color{teal}To'g'ri javob:} \textbf{C) Show/Hide formatlash belgilarini ko'rsatish}
\par\smallskip
\textbf{Izoh:} Yashirin Enter va bo'shliqlar ko'rinadi.
\end{tcolorbox}
\vspace{2mm}\hrule height .2pt\vspace{2mm}

% DB 85a8b880-3dd4-35f6-cde5-212473799975 status=published
\needspace{6\baselineskip}
\textbf{\color{ink}98.} {ootnotesizecolor{muted}Y1 | published | 85a8b880}
\par
Word hujjatidagi 3\ensuremath{\times}3 jadvalda "Ball" ustuni sarlavhasi va pastdagi raqamlar tagiga chizilishi kerak. Keyin jadvalni ikkinchi ma'lumot satridan oldin ikkiga bo'lish talab etiladi. (1) va (2) qanday buyruqlar?
\begin{enumerate}[label=\Alph*),leftmargin=9mm]
\item (1) Ctrl+U; (2) Table Layout \ensuremath{\rightarrow} Split Table
\item (1) Ctrl+B; (2) Insert Below
\item (1) Ctrl+I; (2) Merge Cells
\item (1) Ctrl+U; (2) Split Cells
\end{enumerate}
\begin{tcolorbox}[enhanced,breakable,colback=teal!4,colframe=teal!55!black,boxrule=0.45pt,arc=1.2mm,left=2.8mm,right=2.8mm,top=2mm,bottom=2mm]
\textbf{\color{teal}To'g'ri javob:} \textbf{A) (1) Ctrl+U; (2) Table Layout \ensuremath{\rightarrow} Split Table}
\par\smallskip
\textbf{Izoh:} Ctrl+U belgilangan matnni tagiga chizadi. Split Table esa kursor joylashgan satr oldidan yangi jadval ajratadi; Split Cells faqat katakni bo'ladi.
\end{tcolorbox}
\vspace{2mm}\hrule height .2pt\vspace{2mm}

% DB af588c73-10b1-f351-b90e-487b06f7993d status=published
\needspace{6\baselineskip}
\textbf{\color{ink}99.} {ootnotesizecolor{muted}Y2 | published | af588c73}
\par
Ushbu vaziyatdagi vositalar va vazifalarni moslashtiring.  1) Heading 1;  2) Normal;  3) Table of Contents;  Izohlar: A) Birinchi darajali sarlavha uslubi;  B) Asosiy matn uslubi;  C) Sarlavhalardan mundarija tuzish;
\textbf{Chap ustun:}\begin{enumerate}[label=\arabic*),leftmargin=9mm]
\item Heading 1
\item Normal
\item Table of Contents
\end{enumerate}
\textbf{O'ng ustun:}\begin{enumerate}[label=\Alph*),leftmargin=9mm]
\item Birinchi darajali sarlavha uslubi
\item Asosiy matn uslubi
\item Sarlavhalardan mundarija tuzish
\end{enumerate}
\begin{tcolorbox}[enhanced,breakable,colback=teal!4,colframe=teal!55!black,boxrule=0.45pt,arc=1.2mm,left=2.8mm,right=2.8mm,top=2mm,bottom=2mm]
\textbf{\color{teal}To'g'ri javob:} \textbf{1-A, 2-B, 3-C}
\par\smallskip
\textbf{Izoh:} To'g'ri moslik: Heading 1 - Birinchi darajali sarlavha uslubi, Normal - Asosiy matn uslubi, Table of Contents - Sarlavhalardan mundarija tuzish.
\end{tcolorbox}
\vspace{2mm}\hrule height .2pt\vspace{2mm}

% DB cb1015fd-1b6b-ad7a-0f84-de96a8d458b0 status=published
\needspace{6\baselineskip}
\textbf{\color{ink}100.} {ootnotesizecolor{muted}Y1 | published | cb1015fd}
\par
Sarlavhalar avtomatik mundarijaga kirmayapti. Sabab qaysi?
\begin{enumerate}[label=\Alph*),leftmargin=9mm]
\item Abzaslar qisqa
\item Hujjatda matn bor
\item Sarlavhalar mos Heading uslubida belgilanmagan
\item Shrift qora
\end{enumerate}
\begin{tcolorbox}[enhanced,breakable,colback=teal!4,colframe=teal!55!black,boxrule=0.45pt,arc=1.2mm,left=2.8mm,right=2.8mm,top=2mm,bottom=2mm]
\textbf{\color{teal}To'g'ri javob:} \textbf{C) Sarlavhalar mos Heading uslubida belgilanmagan}
\par\smallskip
\textbf{Izoh:} TOC Heading uslublari orqali shakllanadi.
\end{tcolorbox}
\vspace{2mm}\hrule height .2pt\vspace{2mm}

% DB 1f97edd0-3442-5114-8b63-4210412dec92 status=published
\needspace{6\baselineskip}
\textbf{\color{ink}101.} {ootnotesizecolor{muted}Y1 | published | 1f97edd0}
\par
Hujjatda Heading 1 bilan 6 bob, Heading 2 bilan 18 kichik bo'lim belgilandi. 3-bob nomi o'zgargach avtomatik mundarija eski nomni ko'rsatmoqda. Nima qilish kerak?
\begin{enumerate}[label=\Alph*),leftmargin=9mm]
\item Table of Contentsni Update Table bilan yangilash
\item Sarlavhalarni faqat Center qilish
\item Barcha sarlavhani oddiy Boldga almashtirish
\item Barcha Page Breaklarni o'chirish
\end{enumerate}
\begin{tcolorbox}[enhanced,breakable,colback=teal!4,colframe=teal!55!black,boxrule=0.45pt,arc=1.2mm,left=2.8mm,right=2.8mm,top=2mm,bottom=2mm]
\textbf{\color{teal}To'g'ri javob:} \textbf{A) Table of Contentsni Update Table bilan yangilash}
\par\smallskip
\textbf{Izoh:} Mundarija sarlavha uslublaridan hosil qilinadi va o'zgarishdan keyin qayta yangilanishi kerak.
\end{tcolorbox}
\vspace{2mm}\hrule height .2pt\vspace{2mm}

% DB 2d9b5e08-7b6d-e80f-82cb-29f9a6b2601e status=published
\needspace{6\baselineskip}
\textbf{\color{ink}102.} {ootnotesizecolor{muted}Y1 | published | 2d9b5e08}
\par
20 betli hisobotning faqat 8-betidagi keng jadval Landscape, boshqa betlar Portrait bo'lsin; keyingi betlarning sahifa raqami saqlansin. Eng to'g'ri usul?
\begin{enumerate}[label=\Alph*),leftmargin=9mm]
\item Jadvalni rasmga aylantirib Headerni o'chirish
\item Jadvaldan oldin 15 marta Enter bosish
\item Hujjatning butun Orientationini Landscape qilish
\item Jadval atrofida Section Break yaratib o'rta bo'limni Landscape qilish, kolontitulni tekshirish
\end{enumerate}
\begin{tcolorbox}[enhanced,breakable,colback=teal!4,colframe=teal!55!black,boxrule=0.45pt,arc=1.2mm,left=2.8mm,right=2.8mm,top=2mm,bottom=2mm]
\textbf{\color{teal}To'g'ri javob:} \textbf{D) Jadval atrofida Section Break yaratib o'rta bo'limni Landscape qilish, kolontitulni tekshirish}
\par\smallskip
\textbf{Izoh:} Alohida bo'limlar turli sahifa parametrlari va kolontitul sozlamalarini qo'llashga imkon beradi.
\end{tcolorbox}
\vspace{2mm}\hrule height .2pt\vspace{2mm}

% DB 4de1de93-de3f-4a41-a630-b73906fb29d2 status=published
\needspace{6\baselineskip}
\textbf{\color{ink}103.} {ootnotesizecolor{muted}Y3 | published | 4de1de93}
\par
MS Word hujjatida avtomatik mundarija mavjud. Bo'lim sarlavhasi o'zgarganidan keyin mundarijada yangilangan nom ko'rinishi uchun amallarni tartiblang.
\begin{enumerate}[label=\Alph*),leftmargin=9mm]
\item Mundarijani yangilash
\item Sarlavhalarga uslublarni belgilash
\item Avtomatik mundarija qo'shish
\item Bo'lim sarlavhasini o'zgartirish
\end{enumerate}
\begin{tcolorbox}[enhanced,breakable,colback=teal!4,colframe=teal!55!black,boxrule=0.45pt,arc=1.2mm,left=2.8mm,right=2.8mm,top=2mm,bottom=2mm]
\textbf{\color{teal}To'g'ri javob:} \textbf{B - C - D - A}
\par\smallskip
\textbf{Izoh:} Bosqichlar: Sarlavhalarga uslublarni belgilash \ensuremath{\rightarrow} Avtomatik mundarija qo'shish \ensuremath{\rightarrow} Bo'lim sarlavhasini o'zgartirish \ensuremath{\rightarrow} Mundarijani yangilash.
\end{tcolorbox}
\vspace{2mm}\hrule height .2pt\vspace{2mm}

% DB 99bbf68e-5be7-48a6-9026-9b948044de24 status=published
\needspace{6\baselineskip}
\textbf{\color{ink}104.} {ootnotesizecolor{muted}Y3 | published | 99bbf68e}
\par
MS Word hujjatidagi sarlavhalardan avtomatik mundarija yaratildi. Keyin sarlavhalardan biri tahrirlandi. Hujjatda Heading uslubini qo'llashdan yangilangan mundarijani chiqarishgacha amallarni tartiblang.
\begin{enumerate}[label=\Alph*),leftmargin=9mm]
\item Sarlavhalarga Heading uslubi berish
\item Mundarija qo'shish
\item Sarlavhalarni tahrirlash
\item Mundarijani Update Table qilish
\end{enumerate}
\begin{tcolorbox}[enhanced,breakable,colback=teal!4,colframe=teal!55!black,boxrule=0.45pt,arc=1.2mm,left=2.8mm,right=2.8mm,top=2mm,bottom=2mm]
\textbf{\color{teal}To'g'ri javob:} \textbf{A - B - C - D}
\par\smallskip
\textbf{Izoh:} Sarlavhalarga Heading uslubi beriladi, avtomatik mundarija yaratiladi, sarlavha o'zgartiriladi va mundarija Update Table bilan yangilanadi. Tartib: 1 \ensuremath{\rightarrow} 2 \ensuremath{\rightarrow} 3 \ensuremath{\rightarrow} 4.
\end{tcolorbox}
\vspace{2mm}\hrule height .2pt\vspace{2mm}

% DB 9ec07bbc-9eb6-2bfa-84f1-1ec53a8c0908 status=published
\needspace{6\baselineskip}
\textbf{\color{ink}105.} {ootnotesizecolor{muted}Y2 | published | 9ec07bbc}
\par
Quyidagi tushunchalarni mazmuniga qarab bog'lang.  1) Section Break;  2) Page Break;  3) Different First Page;  Izohlar: A) Birinchi sahifa kolontitulini ajratish;  B) Mustaqil sahifa parametrlari bo'limi;  C) Matnni navbatdagi sahifaga o'tkazish;
\textbf{Chap ustun:}\begin{enumerate}[label=\arabic*),leftmargin=9mm]
\item Section Break
\item Page Break
\item Different First Page
\end{enumerate}
\textbf{O'ng ustun:}\begin{enumerate}[label=\Alph*),leftmargin=9mm]
\item Birinchi sahifa kolontitulini ajratish
\item Mustaqil sahifa parametrlari bo'limi
\item Matnni navbatdagi sahifaga o'tkazish
\end{enumerate}
\begin{tcolorbox}[enhanced,breakable,colback=teal!4,colframe=teal!55!black,boxrule=0.45pt,arc=1.2mm,left=2.8mm,right=2.8mm,top=2mm,bottom=2mm]
\textbf{\color{teal}To'g'ri javob:} \textbf{1-B, 2-C, 3-A}
\par\smallskip
\textbf{Izoh:} To'g'ri moslik: Section Break - Mustaqil sahifa parametrlari bo'limi, Page Break - Matnni navbatdagi sahifaga o'tkazish, Different First Page - Birinchi sahifa kolontitulini ajratish.
\end{tcolorbox}
\vspace{2mm}\hrule height .2pt\vspace{2mm}

% DB c068ddf6-15b7-7022-5ca1-005c4c18e4c4 status=published
\needspace{6\baselineskip}
\textbf{\color{ink}106.} {ootnotesizecolor{muted}Y2 | published | c068ddf6}
\par
Moslashtiring.  1) Margin;  2) Indent;  3) Section Break;  Ta'riflar: A) Alohida bo'lim parametrlari;  B) Abzas chekinishi;  C) Sahifa hoshiyasi;
\textbf{Chap ustun:}\begin{enumerate}[label=\arabic*),leftmargin=9mm]
\item Margin
\item Indent
\item Section Break
\end{enumerate}
\textbf{O'ng ustun:}\begin{enumerate}[label=\Alph*),leftmargin=9mm]
\item Alohida bo'lim parametrlari
\item Abzas chekinishi
\item Sahifa hoshiyasi
\end{enumerate}
\begin{tcolorbox}[enhanced,breakable,colback=teal!4,colframe=teal!55!black,boxrule=0.45pt,arc=1.2mm,left=2.8mm,right=2.8mm,top=2mm,bottom=2mm]
\textbf{\color{teal}To'g'ri javob:} \textbf{1-C, 2-B, 3-A}
\par\smallskip
\textbf{Izoh:} Mos keladigan ta'riflar: 1-C, 2-B, 3-A.
\end{tcolorbox}
\vspace{2mm}\hrule height .2pt\vspace{2mm}

% DB cbce8857-90d5-7905-4ac2-c2f41289e2c4 status=published
\needspace{6\baselineskip}
\textbf{\color{ink}107.} {ootnotesizecolor{muted}Y2 | published | cbce8857}
\par
Wordda formatlash holati va tuzatish vositasini moslang.
\textbf{Chap ustun:}\begin{enumerate}[label=\arabic*),leftmargin=9mm]
\item Mundarija eski sarlavhani ko'rsatmoqda
\item Faqat bitta bo'lim Landscape bo'lsin
\item Ro'yxat qayta boshlanishi kerak
\item Birinchi sahifa kolontituli boshqacha bo'lsin
\end{enumerate}
\textbf{O'ng ustun:}\begin{enumerate}[label=\Alph*),leftmargin=9mm]
\item Different First Page
\item Restart at 1
\item Update entire table
\item Section Break
\end{enumerate}
\begin{tcolorbox}[enhanced,breakable,colback=teal!4,colframe=teal!55!black,boxrule=0.45pt,arc=1.2mm,left=2.8mm,right=2.8mm,top=2mm,bottom=2mm]
\textbf{\color{teal}To'g'ri javob:} \textbf{1-C, 2-D, 3-B, 4-A}
\par\smallskip
\textbf{Izoh:} 1-C, 2-D, 3-B, 4-A; turli real hujjat muammolarini mos Word sozlamalariga bog'lash.
\end{tcolorbox}
\vspace{2mm}\hrule height .2pt\vspace{2mm}

% DB dd8d8ed0-9187-5f44-2c54-a1dbb3b6d803 status=published
\needspace{6\baselineskip}
\textbf{\color{ink}108.} {ootnotesizecolor{muted}Y1 | published | dd8d8ed0}
\par
30 betli hisobotda 12-betdagi yagona jadval ko'ndalang (Landscape), qolganlari bo'yiga (Portrait) bo'lishi kerak. 13-bet ham Portrait qolishi uchun nima kerak?
\begin{enumerate}[label=\Alph*),leftmargin=9mm]
\item Jadvaldan oldin va keyin Section Break (Next Page) kiritib, faqat o'rtadagi sectionni Landscape qilish
\item Faqat 12-betdan avval Page Break qo'yish
\item 12-bet jadvalini Ctrl+C bilan yangi faylga olish
\item Barcha hujjatning Orientation sozlamasini Landscape qilish
\end{enumerate}
\begin{tcolorbox}[enhanced,breakable,colback=teal!4,colframe=teal!55!black,boxrule=0.45pt,arc=1.2mm,left=2.8mm,right=2.8mm,top=2mm,bottom=2mm]
\textbf{\color{teal}To'g'ri javob:} \textbf{A) Jadvaldan oldin va keyin Section Break (Next Page) kiritib, faqat o'rtadagi sectionni Landscape qilish}
\par\smallskip
\textbf{Izoh:} Section Break har bir qismga alohida Page Setup imkonini beradi. Oddiy Page Break orientatsiyani faqat bitta bet uchun ajratmaydi.
\end{tcolorbox}
\vspace{2mm}\hrule height .2pt\vspace{2mm}

% DB a6bd8408-f965-48eb-9a77-30331de88282 status=review
\needspace{6\baselineskip}
\textbf{\color{ink}109.} {ootnotesizecolor{muted}Y1 | review | a6bd8408}
\par
Enter va Shift+Enter farqi qaysi?
\begin{enumerate}[label=\Alph*),leftmargin=9mm]
\item Enter faqat satr uzadi, Shift+Enter esa yangi abzasni boshlaydi
\item Enter va Shift+Enter yangi abzas yaratadi, farqi faqat tekislanishda
\item Enter yangi abzas, Shift+Enter abzas ichida satr uzilishi yaratadi
\item Enter sahifa oxirini belgilaydi, Shift+Enter matnni keyingi sahifaga ko'chiradi
\end{enumerate}
\begin{tcolorbox}[enhanced,breakable,colback=teal!4,colframe=teal!55!black,boxrule=0.45pt,arc=1.2mm,left=2.8mm,right=2.8mm,top=2mm,bottom=2mm]
\textbf{\color{teal}To'g'ri javob:} \textbf{C) Enter yangi abzas, Shift+Enter abzas ichida satr uzilishi yaratadi}
\par\smallskip
\textbf{Izoh:} Wordda Enter yangi abzas (paragraph) boshlaydi va uning abzas formatlash xususiyatlari ajraladi. Shift+Enter esa bir abzas ichida majburiy satr uzilishini kiritadi; u yangi abzas yoki yangi sahifa yaratmaydi. Bu farq tekislash, ro'yxatlar va oraliqlar uchun ahamiyatli.
\end{tcolorbox}
\vspace{2mm}\hrule height .2pt\vspace{2mm}

% DB bb0dda0f-e300-4120-aac9-7343355dd09b status=review
\needspace{6\baselineskip}
\textbf{\color{ink}110.} {ootnotesizecolor{muted}Y1 | review | bb0dda0f}
\par
30 betli hisobotda 12-betdagi yagona jadval ko'ndalang (Landscape), qolganlari bo'yiga (Portrait) bo'lishi kerak. 13-bet ham Portrait qolishi uchun nima kerak?
\begin{enumerate}[label=\Alph*),leftmargin=9mm]
\item Jadvaldan oldin va keyin Section Break (Next Page) kiritib, faqat o'rtadagi sectionni Landscape qilish
\item Jadvaldan oldin va keyin Page Break kiritib, o'rtadagi sahifaga alohida uslub berish
\item Jadvaldan oldin Section Break kiritib, undan keyingi barcha sahifalarni Landscape qilish
\item Jadvalni 90 daraja burilgan rasmga aylantirib, barcha sahifalarni Portrait qoldirish
\end{enumerate}
\begin{tcolorbox}[enhanced,breakable,colback=teal!4,colframe=teal!55!black,boxrule=0.45pt,arc=1.2mm,left=2.8mm,right=2.8mm,top=2mm,bottom=2mm]
\textbf{\color{teal}To'g'ri javob:} \textbf{A) Jadvaldan oldin va keyin Section Break (Next Page) kiritib, faqat o'rtadagi sectionni Landscape qilish}
\par\smallskip
\textbf{Izoh:} Wordda Orientation sozlamasi section darajasida boshqariladi. Kerakli jadval sahifasidan oldin va keyin Section Break (Next Page) qo'yilib, faqat o'rta bo'lim Landscape qilinadi; keyingi bo'lim yana Portrait bo'lib qoladi. Page Break o'zi mustaqil section yaratmaydi, rasmga aylantirish esa tahrirlanadigan jadval talabini bajarmaydi.
\end{tcolorbox}
\vspace{2mm}\hrule height .2pt\vspace{2mm}
\subsection{MS Word: jadval, rasm, formula, havola va hujjat tuzilishi}

% DB 01b6e6df-71a1-5be4-01f3-fbc3fa490548 status=published
\needspace{6\baselineskip}
\textbf{\color{ink}111.} {ootnotesizecolor{muted}Y1 | published | 01b6e6df}
\par
Wordda jadval katagi nima?
\begin{enumerate}[label=\Alph*),leftmargin=9mm]
\item Hujjat sarlavhasi
\item Sahifa cheti
\item Satr bilan ustun kesishmasi
\item Fayl nomi
\end{enumerate}
\begin{tcolorbox}[enhanced,breakable,colback=teal!4,colframe=teal!55!black,boxrule=0.45pt,arc=1.2mm,left=2.8mm,right=2.8mm,top=2mm,bottom=2mm]
\textbf{\color{teal}To'g'ri javob:} \textbf{C) Satr bilan ustun kesishmasi}
\par\smallskip
\textbf{Izoh:} Katak satr va ustun kesishuvidir.
\end{tcolorbox}
\vspace{2mm}\hrule height .2pt\vspace{2mm}

% DB 2f597528-f732-73cf-cdb6-0120d0efd00f status=published
\needspace{6\baselineskip}
\textbf{\color{ink}112.} {ootnotesizecolor{muted}Y1 | published | 2f597528}
\par
Merge Cells nima qiladi?
\begin{enumerate}[label=\Alph*),leftmargin=9mm]
\item Rasmni qirqadi
\item Barcha sahifalarni chop etadi
\item Faqat yangi fayl yaratadi
\item Bir necha belgilangan katakni birlashtiradi
\end{enumerate}
\begin{tcolorbox}[enhanced,breakable,colback=teal!4,colframe=teal!55!black,boxrule=0.45pt,arc=1.2mm,left=2.8mm,right=2.8mm,top=2mm,bottom=2mm]
\textbf{\color{teal}To'g'ri javob:} \textbf{D) Bir necha belgilangan katakni birlashtiradi}
\par\smallskip
\textbf{Izoh:} Merge Cells kataklarni birlashtiradi.
\end{tcolorbox}
\vspace{2mm}\hrule height .2pt\vspace{2mm}

% DB 3b9107b1-9795-da7a-8671-62313d6881f6 status=published
\needspace{6\baselineskip}
\textbf{\color{ink}113.} {ootnotesizecolor{muted}Y1 | published | 3b9107b1}
\par
Wrap Text qanday vazifa bajaradi?
\begin{enumerate}[label=\Alph*),leftmargin=9mm]
\item Matnning rasm atrofida joylashishini belgilaydi
\item Formulani hisoblaydi
\item Hujjatni shifrlaydi
\item Printer turini tanlaydi
\end{enumerate}
\begin{tcolorbox}[enhanced,breakable,colback=teal!4,colframe=teal!55!black,boxrule=0.45pt,arc=1.2mm,left=2.8mm,right=2.8mm,top=2mm,bottom=2mm]
\textbf{\color{teal}To'g'ri javob:} \textbf{A) Matnning rasm atrofida joylashishini belgilaydi}
\par\smallskip
\textbf{Izoh:} Wrap Text rasm va matn munosabatidir.
\end{tcolorbox}
\vspace{2mm}\hrule height .2pt\vspace{2mm}

% DB 4c4bb3f3-e61f-4f2e-92b6-e37e3b223625 status=published
\needspace{6\baselineskip}
\textbf{\color{ink}114.} {ootnotesizecolor{muted}Y1 | published | 4c4bb3f3}
\par
MS Word dasturida kompyuterdagi fayldan hujjatga rasm qo'shish uchun qaysi ketma-ketlik to'g'ri?
\begin{enumerate}[label=\Alph*),leftmargin=9mm]
\item Insert \ensuremath{\rightarrow} Shapes \ensuremath{\rightarrow} rasmni tanlash
\item Insert \ensuremath{\rightarrow} Pictures \ensuremath{\rightarrow} This Device \ensuremath{\rightarrow} rasmni tanlash
\item Design \ensuremath{\rightarrow} Page Color \ensuremath{\rightarrow} rasmni tanlash
\item Layout \ensuremath{\rightarrow} Breaks \ensuremath{\rightarrow} rasmni tanlash
\end{enumerate}
\begin{tcolorbox}[enhanced,breakable,colback=teal!4,colframe=teal!55!black,boxrule=0.45pt,arc=1.2mm,left=2.8mm,right=2.8mm,top=2mm,bottom=2mm]
\textbf{\color{teal}To'g'ri javob:} \textbf{B) Insert \ensuremath{\rightarrow} Pictures \ensuremath{\rightarrow} This Device \ensuremath{\rightarrow} rasmni tanlash}
\par\smallskip
\textbf{Izoh:} Kompyuterdagi rasm fayli Insert bo'limidagi Pictures buyrug'i orqali, so'ng This Device tanlovi bilan qo'shiladi.
\end{tcolorbox}
\vspace{2mm}\hrule height .2pt\vspace{2mm}

% DB 5da5171c-fc42-35ea-edd2-6756733b07cb status=published
\needspace{6\baselineskip}
\textbf{\color{ink}115.} {ootnotesizecolor{muted}Y1 | published | 5da5171c}
\par
Mail Merge nimani anglatadi?
\begin{enumerate}[label=\Alph*),leftmargin=9mm]
\item Rasm qatlamini kesish
\item Jadvalni filtrlash
\item Faqat bir faylni PDFga o'tkazish
\item Shablon va ma'lumotlar ro'yxatini birlashtirib shaxsiy hujjat chiqarish
\end{enumerate}
\begin{tcolorbox}[enhanced,breakable,colback=teal!4,colframe=teal!55!black,boxrule=0.45pt,arc=1.2mm,left=2.8mm,right=2.8mm,top=2mm,bottom=2mm]
\textbf{\color{teal}To'g'ri javob:} \textbf{D) Shablon va ma'lumotlar ro'yxatini birlashtirib shaxsiy hujjat chiqarish}
\par\smallskip
\textbf{Izoh:} Mail Merge shaxsiylashtirilgan hujjatlar tayyorlaydi.
\end{tcolorbox}
\vspace{2mm}\hrule height .2pt\vspace{2mm}

% DB 5e7c0024-af8c-4d8d-aeb4-e6a1c2da8a9c status=published
\needspace{6\baselineskip}
\textbf{\color{ink}116.} {ootnotesizecolor{muted}Y3 | published | 5e7c0024}
\par
MS Wordda jadval hosil qilib, kataklarga matn kiritilganidan keyin sarlavha uchun tegishli kataklarni birlashtirish kerak. Shu ishning amallarini ketma-ket joylashtiring.
\begin{enumerate}[label=\Alph*),leftmargin=9mm]
\item Insert Table orqali jadval hosil qilish
\item Tegishli kataklarni belgilash
\item Merge Cells qo'llash
\item Kataklarga matn kiritish
\end{enumerate}
\begin{tcolorbox}[enhanced,breakable,colback=teal!4,colframe=teal!55!black,boxrule=0.45pt,arc=1.2mm,left=2.8mm,right=2.8mm,top=2mm,bottom=2mm]
\textbf{\color{teal}To'g'ri javob:} \textbf{A - D - B - C}
\par\smallskip
\textbf{Izoh:} Bosqichlar: Insert Table orqali jadval hosil qilish \ensuremath{\rightarrow} Kataklarga matn kiritish \ensuremath{\rightarrow} Tegishli kataklarni belgilash \ensuremath{\rightarrow} Merge Cells qo'llash.
\end{tcolorbox}
\vspace{2mm}\hrule height .2pt\vspace{2mm}

% DB 6ff87bd2-f560-ddd0-1500-64a1a9ed3c06 status=published
\needspace{6\baselineskip}
\textbf{\color{ink}117.} {ootnotesizecolor{muted}Y2 | published | 6ff87bd2}
\par
Quyidagi buyruqlarni vazifasi bilan moslashtiring.  1) Insert Table;  2) Merge Cells;  3) Split Cells;  Izohlar: A) Kataklarni birlashtirish;  B) Katakni bo'lish;  C) Jadval kiritish;
\textbf{Chap ustun:}\begin{enumerate}[label=\arabic*),leftmargin=9mm]
\item Insert Table
\item Merge Cells
\item Split Cells
\end{enumerate}
\textbf{O'ng ustun:}\begin{enumerate}[label=\Alph*),leftmargin=9mm]
\item Kataklarni birlashtirish
\item Katakni bo'lish
\item Jadval kiritish
\end{enumerate}
\begin{tcolorbox}[enhanced,breakable,colback=teal!4,colframe=teal!55!black,boxrule=0.45pt,arc=1.2mm,left=2.8mm,right=2.8mm,top=2mm,bottom=2mm]
\textbf{\color{teal}To'g'ri javob:} \textbf{1-C, 2-A, 3-B}
\par\smallskip
\textbf{Izoh:} To'g'ri moslik: Insert Table - Jadval kiritish, Merge Cells - Kataklarni birlashtirish, Split Cells - Katakni bo'lish.
\end{tcolorbox}
\vspace{2mm}\hrule height .2pt\vspace{2mm}

% DB f2d5dc6b-05da-57cf-9348-635ec8731abe status=published
\needspace{6\baselineskip}
\textbf{\color{ink}118.} {ootnotesizecolor{muted}Y1 | published | f2d5dc6b}
\par
MS Word vositalarini ular yaratadigan obyektlar bilan moslashtiring: 1) WordArt; 2) ClipArt; 3) Shapes; 4) SmartArt. a) tayyor rasmlar to'plamidan olingan rasm; b) geometrik strelka; c) bezakli "Informatika" yozuvi; d) "Maqsad \ensuremath{\rightarrow} Harakat \ensuremath{\rightarrow} Natija" jarayon diagrammasi.
\begin{enumerate}[label=\Alph*),leftmargin=9mm]
\item 1-a; 2-b; 3-c; 4-d
\item 1-b; 2-c; 3-a; 4-d
\item 1-c; 2-a; 3-b; 4-d
\item 1-a; 2-c; 3-d; 4-b
\end{enumerate}
\begin{tcolorbox}[enhanced,breakable,colback=teal!4,colframe=teal!55!black,boxrule=0.45pt,arc=1.2mm,left=2.8mm,right=2.8mm,top=2mm,bottom=2mm]
\textbf{\color{teal}To'g'ri javob:} \textbf{C) 1-c; 2-a; 3-b; 4-d}
\par\smallskip
\textbf{Izoh:} WordArt bezakli matn (c), ClipArt tayyor rasm (a), Shapes strelka kabi shakl (b), SmartArt esa jarayon diagrammasi (d) yaratadi.
\end{tcolorbox}
\vspace{2mm}\hrule height .2pt\vspace{2mm}

% DB 0a0676c9-d30b-f035-f6ca-d4b348bdbdea status=published
\needspace{6\baselineskip}
\textbf{\color{ink}119.} {ootnotesizecolor{muted}Y1 | published | 0a0676c9}
\par
Wordda rasm ustiga sarlavha sifatida qiyshaytirilgan bezakli matn qo'yish kerak. Diagramma yoki jarayon bosqichlari talab qilinmaydi. Qaysi vosita eng bevosita mos?
\begin{enumerate}[label=\Alph*),leftmargin=9mm]
\item SmartArt
\item Mail Merge
\item Table of Contents
\item WordArt
\end{enumerate}
\begin{tcolorbox}[enhanced,breakable,colback=teal!4,colframe=teal!55!black,boxrule=0.45pt,arc=1.2mm,left=2.8mm,right=2.8mm,top=2mm,bottom=2mm]
\textbf{\color{teal}To'g'ri javob:} \textbf{D) WordArt}
\par\smallskip
\textbf{Izoh:} WordArt bezakli matn ko'rinishlari uchun; SmartArt ko'proq tarkibiy sxema/diagrammalar uchun ishlatiladi. 13-aprel manbasidagi obyektlarni farqlashga o'xshash original vaziyat.
\end{tcolorbox}
\vspace{2mm}\hrule height .2pt\vspace{2mm}

% DB 472d3cff-bf91-a8f0-a8e7-458ee07fe054 status=published
\needspace{6\baselineskip}
\textbf{\color{ink}120.} {ootnotesizecolor{muted}Y1 | published | 472d3cff}
\par
Table of Contents va qo'lda ro'yxatning farqi?
\begin{enumerate}[label=\Alph*),leftmargin=9mm]
\item Qo'lda ro'yxat doim avtomatik yangilanadi
\item Avtomatik TOC sarlavha uslublari bilan yangilanishi mumkin
\item TOC sahifa raqamini o'chiradi
\item TOC faqat rasm qo'shadi
\end{enumerate}
\begin{tcolorbox}[enhanced,breakable,colback=teal!4,colframe=teal!55!black,boxrule=0.45pt,arc=1.2mm,left=2.8mm,right=2.8mm,top=2mm,bottom=2mm]
\textbf{\color{teal}To'g'ri javob:} \textbf{B) Avtomatik TOC sarlavha uslublari bilan yangilanishi mumkin}
\par\smallskip
\textbf{Izoh:} TOC tuzilmalardan hosil qilinadi.
\end{tcolorbox}
\vspace{2mm}\hrule height .2pt\vspace{2mm}

% DB 6414edb3-806c-15e3-6ffa-85679a854929 status=published
\needspace{6\baselineskip}
\textbf{\color{ink}121.} {ootnotesizecolor{muted}Y2 | published | 6414edb3}
\par
Atamalar va ularning aniq ma'nolarini moslashtiring.  1) Wrap Text;  2) Crop;  3) Alt Text;  Izohlar: A) Tasvir chetlarini qirqib ko'rsatish;  B) Matnning rasm atrofiga joylashuvi;  C) Rasmga matnli tavsif yozish;
\textbf{Chap ustun:}\begin{enumerate}[label=\arabic*),leftmargin=9mm]
\item Wrap Text
\item Crop
\item Alt Text
\end{enumerate}
\textbf{O'ng ustun:}\begin{enumerate}[label=\Alph*),leftmargin=9mm]
\item Tasvir chetlarini qirqib ko'rsatish
\item Matnning rasm atrofiga joylashuvi
\item Rasmga matnli tavsif yozish
\end{enumerate}
\begin{tcolorbox}[enhanced,breakable,colback=teal!4,colframe=teal!55!black,boxrule=0.45pt,arc=1.2mm,left=2.8mm,right=2.8mm,top=2mm,bottom=2mm]
\textbf{\color{teal}To'g'ri javob:} \textbf{1-B, 2-A, 3-C}
\par\smallskip
\textbf{Izoh:} To'g'ri moslik: Wrap Text - Matnning rasm atrofiga joylashuvi, Crop - Tasvir chetlarini qirqib ko'rsatish, Alt Text - Rasmga matnli tavsif yozish.
\end{tcolorbox}
\vspace{2mm}\hrule height .2pt\vspace{2mm}

% DB 6c74e018-cd2a-4b78-bd5e-c927d8a749a6 status=published
\needspace{6\baselineskip}
\textbf{\color{ink}122.} {ootnotesizecolor{muted}Y3 | published | 6c74e018}
\par
MS Word hujjatiga rasm qo'shib, Wrap Text rejimi yordamida matnga nisbatan joylashuvini sozlash, rasmni ko'chirish va natijani tekshirish amallarini tartiblang.
\begin{enumerate}[label=\Alph*),leftmargin=9mm]
\item Rasmni hujjatga kiritish
\item Rasmni kerakli joyga ko'chirish
\item Matn bilan o'zaro joylashuvni tekshirish
\item Wrap Text rejimini tanlash
\end{enumerate}
\begin{tcolorbox}[enhanced,breakable,colback=teal!4,colframe=teal!55!black,boxrule=0.45pt,arc=1.2mm,left=2.8mm,right=2.8mm,top=2mm,bottom=2mm]
\textbf{\color{teal}To'g'ri javob:} \textbf{A - D - B - C}
\par\smallskip
\textbf{Izoh:} Bosqichlar: Rasmni hujjatga kiritish \ensuremath{\rightarrow} Wrap Text rejimini tanlash \ensuremath{\rightarrow} Rasmni kerakli joyga ko'chirish \ensuremath{\rightarrow} Matn bilan o'zaro joylashuvni tekshirish.
\end{tcolorbox}
\vspace{2mm}\hrule height .2pt\vspace{2mm}

% DB 7a8e3fcb-2fe2-579a-9194-f4948f651584 status=published
\needspace{6\baselineskip}
\textbf{\color{ink}123.} {ootnotesizecolor{muted}Y1 | published | 7a8e3fcb}
\par
MS Word jadvalining sarlavha qatoridagi \texttt{No.}, \texttt{subject}, \texttt{grade} yozuvlari og'ma ko'rinishga o'tkazildi. So'ng \texttt{subject} ustunining chap tomoniga yangi bo'sh ustun qo'shildi. Qaysi buyruqlar ketma-ketligi bajarilgan?
\begin{enumerate}[label=\Alph*),leftmargin=9mm]
\item Ctrl+I; Split Table
\item Ctrl+B; Insert Right
\item Ctrl+B; Split Table
\item Ctrl+I; Insert Left
\end{enumerate}
\begin{tcolorbox}[enhanced,breakable,colback=teal!4,colframe=teal!55!black,boxrule=0.45pt,arc=1.2mm,left=2.8mm,right=2.8mm,top=2mm,bottom=2mm]
\textbf{\color{teal}To'g'ri javob:} \textbf{D) Ctrl+I; Insert Left}
\par\smallskip
\textbf{Izoh:} Ctrl+I matnni italic (og'ma) qiladi. Tanlangan \texttt{subject} ustunining chapiga yangi ustun qo'shish uchun Table Tools \ensuremath{\rightarrow} Insert Left tanlanadi.
\end{tcolorbox}
\vspace{2mm}\hrule height .2pt\vspace{2mm}

% DB d30b0547-f33e-5ad0-076b-a3d173d5666d status=published
\needspace{6\baselineskip}
\textbf{\color{ink}124.} {ootnotesizecolor{muted}Y1 | published | d30b0547}
\par
Symbol va Equation qanday farqlanadi?
\begin{enumerate}[label=\Alph*),leftmargin=9mm]
\item Symbol faqat PDF yaratadi
\item Symbol alohida belgi, Equation tuzilmali matematik ifoda uchun
\item Equation faqat rasmni bo'yaydi
\item Ikkalasi ham faqat matnni almashtiradi
\end{enumerate}
\begin{tcolorbox}[enhanced,breakable,colback=teal!4,colframe=teal!55!black,boxrule=0.45pt,arc=1.2mm,left=2.8mm,right=2.8mm,top=2mm,bottom=2mm]
\textbf{\color{teal}To'g'ri javob:} \textbf{B) Symbol alohida belgi, Equation tuzilmali matematik ifoda uchun}
\par\smallskip
\textbf{Izoh:} Murakkab kasr va darajani Equation qulay yozadi.
\end{tcolorbox}
\vspace{2mm}\hrule height .2pt\vspace{2mm}

% DB eb574b0a-d5ad-60c8-a591-5aa10f7da69e status=published
\needspace{6\baselineskip}
\textbf{\color{ink}125.} {ootnotesizecolor{muted}Y1 | published | eb574b0a}
\par
Footnote va Endnote farqi qaysi?
\begin{enumerate}[label=\Alph*),leftmargin=9mm]
\item Footnote faqat rasm uchun
\item Ikkisi ham slayd animatsiyasi
\item Footnote sahifa pastida, Endnote odatda yakunda joylashadi
\item Endnote faqat chap margin
\end{enumerate}
\begin{tcolorbox}[enhanced,breakable,colback=teal!4,colframe=teal!55!black,boxrule=0.45pt,arc=1.2mm,left=2.8mm,right=2.8mm,top=2mm,bottom=2mm]
\textbf{\color{teal}To'g'ri javob:} \textbf{C) Footnote sahifa pastida, Endnote odatda yakunda joylashadi}
\par\smallskip
\textbf{Izoh:} Izohlar joyi boshqacha.
\end{tcolorbox}
\vspace{2mm}\hrule height .2pt\vspace{2mm}

% DB f3f8d4bb-bb80-3982-927c-9a003fe34bfc status=published
\needspace{6\baselineskip}
\textbf{\color{ink}126.} {ootnotesizecolor{muted}Y1 | published | f3f8d4bb}
\par
Crop bilan Resize farqi qaysi?
\begin{enumerate}[label=\Alph*),leftmargin=9mm]
\item Crop ko'rinadigan qismini kesadi, Resize o'lchamini o'zgartiradi
\item Crop faqat matn nusxalaydi
\item Resize doim sahifani o'chiradi
\item Ikkalasi ham katak qo'shadi
\end{enumerate}
\begin{tcolorbox}[enhanced,breakable,colback=teal!4,colframe=teal!55!black,boxrule=0.45pt,arc=1.2mm,left=2.8mm,right=2.8mm,top=2mm,bottom=2mm]
\textbf{\color{teal}To'g'ri javob:} \textbf{A) Crop ko'rinadigan qismini kesadi, Resize o'lchamini o'zgartiradi}
\par\smallskip
\textbf{Izoh:} Crop va masshtab o'zgartirish boshqa amal.
\end{tcolorbox}
\vspace{2mm}\hrule height .2pt\vspace{2mm}

% DB 2245c687-3ec8-016a-ba43-c13e74591b34 status=published
\needspace{6\baselineskip}
\textbf{\color{ink}127.} {ootnotesizecolor{muted}Y2 | published | 2245c687}
\par
Amallarni izohi bilan moslang.  1) Equation;  2) Symbol;  3) Shapes;  Izohlar: A) Matematik formula kiritish;  B) Geometrik shakl kiritish;  C) Maxsus belgi kiritish;
\textbf{Chap ustun:}\begin{enumerate}[label=\arabic*),leftmargin=9mm]
\item Equation
\item Symbol
\item Shapes
\end{enumerate}
\textbf{O'ng ustun:}\begin{enumerate}[label=\Alph*),leftmargin=9mm]
\item Matematik formula kiritish
\item Geometrik shakl kiritish
\item Maxsus belgi kiritish
\end{enumerate}
\begin{tcolorbox}[enhanced,breakable,colback=teal!4,colframe=teal!55!black,boxrule=0.45pt,arc=1.2mm,left=2.8mm,right=2.8mm,top=2mm,bottom=2mm]
\textbf{\color{teal}To'g'ri javob:} \textbf{1-A, 2-C, 3-B}
\par\smallskip
\textbf{Izoh:} To'g'ri moslik: Equation - Matematik formula kiritish, Symbol - Maxsus belgi kiritish, Shapes - Geometrik shakl kiritish.
\end{tcolorbox}
\vspace{2mm}\hrule height .2pt\vspace{2mm}

% DB 8084a5db-221e-83fa-697f-041467fe7178 status=published
\needspace{6\baselineskip}
\textbf{\color{ink}128.} {ootnotesizecolor{muted}Y1 | published | 8084a5db}
\par
Hujjatdagi rasmdan faqat o'ng tomon ko'rinsin. Nima qilinadi?
\begin{enumerate}[label=\Alph*),leftmargin=9mm]
\item Find
\item Rotate Text
\item Mail Merge
\item Crop
\end{enumerate}
\begin{tcolorbox}[enhanced,breakable,colback=teal!4,colframe=teal!55!black,boxrule=0.45pt,arc=1.2mm,left=2.8mm,right=2.8mm,top=2mm,bottom=2mm]
\textbf{\color{teal}To'g'ri javob:} \textbf{D) Crop}
\par\smallskip
\textbf{Izoh:} Crop ko'rinadigan qismni chegaralaydi.
\end{tcolorbox}
\vspace{2mm}\hrule height .2pt\vspace{2mm}

% DB 970fab08-8636-43cc-9f78-5f978f60e4de status=published
\needspace{6\baselineskip}
\textbf{\color{ink}129.} {ootnotesizecolor{muted}Y1 | published | 970fab08}
\par
MS Word hujjatidagi atirgul rasmi matn satrlarining \textbf{orqasida} joylashgan; matn tasvir ustida o'qiladi (tasvir yozuvni to'sib qo'ymaydi). Wrap Text buyrug'ining qaysi parametri mos?
\begin{enumerate}[label=\Alph*),leftmargin=9mm]
\item Top and Bottom
\item Tight
\item In Front of Text
\item Behind Text
\end{enumerate}
\par\smallskip\textit{Rasmli savol: }\href{https://github.com/sarvar9417/attestatsiya/blob/main/public/kuzgi-demo/demo_q39_word.svg}{Rasmni onlayn ochish}
\begin{tcolorbox}[enhanced,breakable,colback=teal!4,colframe=teal!55!black,boxrule=0.45pt,arc=1.2mm,left=2.8mm,right=2.8mm,top=2mm,bottom=2mm]
\textbf{\color{teal}To'g'ri javob:} \textbf{D) Behind Text}
\par\smallskip
\textbf{Izoh:} Behind Text tanlansa rasm matn qatlamining orqasiga o'tadi, matn ustidan ko'rinadi. Tight matnni rasm chetlaridan aylanib joylashtiradi.
\end{tcolorbox}
\vspace{2mm}\hrule height .2pt\vspace{2mm}

% DB c5a01d2c-4be1-9996-1f4b-5e989047110f status=published
\needspace{6\baselineskip}
\textbf{\color{ink}130.} {ootnotesizecolor{muted}Y1 | published | c5a01d2c}
\par
Jadvalda ikki katakni bitta sarlavhaga aylantirish uchun?
\begin{enumerate}[label=\Alph*),leftmargin=9mm]
\item Delete Table
\item Crop Picture
\item Merge Cells
\item Split Cells
\end{enumerate}
\begin{tcolorbox}[enhanced,breakable,colback=teal!4,colframe=teal!55!black,boxrule=0.45pt,arc=1.2mm,left=2.8mm,right=2.8mm,top=2mm,bottom=2mm]
\textbf{\color{teal}To'g'ri javob:} \textbf{C) Merge Cells}
\par\smallskip
\textbf{Izoh:} Birlashtirish uchun Merge ishlatiladi.
\end{tcolorbox}
\vspace{2mm}\hrule height .2pt\vspace{2mm}

% DB e2e6e167-3e3e-44e7-b2cf-73636c850b57 status=published
\needspace{6\baselineskip}
\textbf{\color{ink}131.} {ootnotesizecolor{muted}Y3 | published | e2e6e167}
\par
MS Word hujjatining kerakli joyiga Insert Equation yordamida matematik ifoda kiritib, natijani tekshirish uchun amallarni tartiblang.
\begin{enumerate}[label=\Alph*),leftmargin=9mm]
\item Kursorni formula kiritiladigan joyga qo'yish
\item Insert Equationni ochish
\item Ifodani hujjat ichida ko'rib chiqish
\item Matematik ifodani tuzish
\end{enumerate}
\begin{tcolorbox}[enhanced,breakable,colback=teal!4,colframe=teal!55!black,boxrule=0.45pt,arc=1.2mm,left=2.8mm,right=2.8mm,top=2mm,bottom=2mm]
\textbf{\color{teal}To'g'ri javob:} \textbf{A - B - D - C}
\par\smallskip
\textbf{Izoh:} Bosqichlar: Kursorni formula kiritiladigan joyga qo'yish \ensuremath{\rightarrow} Insert Equationni ochish \ensuremath{\rightarrow} Matematik ifodani tuzish \ensuremath{\rightarrow} Ifodani hujjat ichida ko'rib chiqish.
\end{tcolorbox}
\vspace{2mm}\hrule height .2pt\vspace{2mm}

% DB f152c511-c32c-c1f3-db8e-85ec3893d404 status=published
\needspace{6\baselineskip}
\textbf{\color{ink}132.} {ootnotesizecolor{muted}Y1 | published | f152c511}
\par
30 o'quvchiga shaxsiy sertifikat tayyorlash kerak. Nima qulay?
\begin{enumerate}[label=\Alph*),leftmargin=9mm]
\item Mail Merge shablon va ism ro'yxati
\item Word Count
\item Har bir slaydga video qo'shish
\item Faqat Bold ishlatish
\end{enumerate}
\begin{tcolorbox}[enhanced,breakable,colback=teal!4,colframe=teal!55!black,boxrule=0.45pt,arc=1.2mm,left=2.8mm,right=2.8mm,top=2mm,bottom=2mm]
\textbf{\color{teal}To'g'ri javob:} \textbf{A) Mail Merge shablon va ism ro'yxati}
\par\smallskip
\textbf{Izoh:} Mail Merge turli yozuvlarni shablonga joylashtiradi.
\end{tcolorbox}
\vspace{2mm}\hrule height .2pt\vspace{2mm}

% DB 4ca7a4a7-28d2-3a68-77f5-2c1784340869 status=published
\needspace{6\baselineskip}
\textbf{\color{ink}133.} {ootnotesizecolor{muted}Y1 | published | 4ca7a4a7}
\par
Mail Merge ro'yxatida 40 satr bor va 1-satr sarlavha, qolganlari shaxslar. To'g'ri yondashuv?
\begin{enumerate}[label=\Alph*),leftmargin=9mm]
\item Har katak alohida fayl
\item 40 ta sarlavha va 0 yozuv
\item 39 yozuv va maydon nomlari sifatida sarlavha
\item Sarlavha doim birinchi qabul qiluvchi
\end{enumerate}
\begin{tcolorbox}[enhanced,breakable,colback=teal!4,colframe=teal!55!black,boxrule=0.45pt,arc=1.2mm,left=2.8mm,right=2.8mm,top=2mm,bottom=2mm]
\textbf{\color{teal}To'g'ri javob:} \textbf{C) 39 yozuv va maydon nomlari sifatida sarlavha}
\par\smallskip
\textbf{Izoh:} Odatda ustun nomlari maydon, qolgan satrlar yozuvdir.
\end{tcolorbox}
\vspace{2mm}\hrule height .2pt\vspace{2mm}

% DB 522bd647-0f34-4965-b46b-456cfdd2d6bb status=published
\needspace{6\baselineskip}
\textbf{\color{ink}134.} {ootnotesizecolor{muted}Y3 | published | 522bd647}
\par
MS Wordda asosiy xat matni oldindan tayyorlangan, lekin qabul qiluvchilar ro'yxati hali yo'q. Mail Merge orqali har bir kishining ma'lumotlari kiritiladigan xatlarni tayyorlash uchun ro'yxatni yaratish, shablonga ulash, maydonlarni joylashtirish va ko'rib chiqish amallarini tartiblang.
\begin{enumerate}[label=\Alph*),leftmargin=9mm]
\item Merge Field maydonlarini kiritish
\item Preview Results yordamida tekshirish
\item Yaratilgan qabul qiluvchilar ro'yxatini tayyor xat shabloniga ulash
\item Qabul qiluvchilar ro'yxatini yaratish
\end{enumerate}
\begin{tcolorbox}[enhanced,breakable,colback=teal!4,colframe=teal!55!black,boxrule=0.45pt,arc=1.2mm,left=2.8mm,right=2.8mm,top=2mm,bottom=2mm]
\textbf{\color{teal}To'g'ri javob:} \textbf{D - C - A - B}
\par\smallskip
\textbf{Izoh:} Oldin qabul qiluvchilar ro'yxati yaratiladi, u tayyor asosiy xat shabloniga ulanadi, Merge Field maydonlari kiritiladi va Preview Results orqali shaxsiylashtirish tekshiriladi. Tartib: 4 \ensuremath{\rightarrow} 3 \ensuremath{\rightarrow} 1 \ensuremath{\rightarrow} 2.
\end{tcolorbox}
\vspace{2mm}\hrule height .2pt\vspace{2mm}

% DB 5ae3da42-d172-6a22-7f5d-02f4053d84b9 status=published
\needspace{6\baselineskip}
\textbf{\color{ink}135.} {ootnotesizecolor{muted}Y1 | published | 5ae3da42}
\par
Ko'p sahifalik jadvalning har betida ustun nomlari qaytarilsin. Nima sozlanadi?
\begin{enumerate}[label=\Alph*),leftmargin=9mm]
\item Repeat Header Rows
\item Insert Animation
\item Table to Picture
\item Sort Descending
\end{enumerate}
\begin{tcolorbox}[enhanced,breakable,colback=teal!4,colframe=teal!55!black,boxrule=0.45pt,arc=1.2mm,left=2.8mm,right=2.8mm,top=2mm,bottom=2mm]
\textbf{\color{teal}To'g'ri javob:} \textbf{A) Repeat Header Rows}
\par\smallskip
\textbf{Izoh:} Takroriy sarlavha qatori o'qishni yengillashtiradi.
\end{tcolorbox}
\vspace{2mm}\hrule height .2pt\vspace{2mm}

% DB 80c49d70-e23c-e5b9-c6b5-cbd9c0595db5 status=published
\needspace{6\baselineskip}
\textbf{\color{ink}136.} {ootnotesizecolor{muted}Y1 | published | 80c49d70}
\par
Hujjatdagi rasm ustiga matn chiqsin, rasm fon sifatida joylashsin. Qaysi tartib qulay?
\begin{enumerate}[label=\Alph*),leftmargin=9mm]
\item Rasmni In Line qilib qoldirish
\item Behind Text uchun matn o'ralishini sozlash
\item Find All
\item Faqat Share bosish
\end{enumerate}
\begin{tcolorbox}[enhanced,breakable,colback=teal!4,colframe=teal!55!black,boxrule=0.45pt,arc=1.2mm,left=2.8mm,right=2.8mm,top=2mm,bottom=2mm]
\textbf{\color{teal}To'g'ri javob:} \textbf{B) Behind Text uchun matn o'ralishini sozlash}
\par\smallskip
\textbf{Izoh:} Behind Text matnni tasvir ustiga chiqarish imkonini beradi.
\end{tcolorbox}
\vspace{2mm}\hrule height .2pt\vspace{2mm}

% DB bcc20e93-4cf0-3b09-2769-b00f7a80cab8 status=published
\needspace{6\baselineskip}
\textbf{\color{ink}137.} {ootnotesizecolor{muted}Y2 | published | bcc20e93}
\par
Ushbu vaziyatdagi vositalar va vazifalarni moslashtiring.  1) Merge Field;  2) Recipient List;  3) Preview Results;  Izohlar: A) Manbadan dinamik maydon;  B) Birlashtirilgan hujjat natijasini tekshirish;  C) Qabul qiluvchilar jadvali;
\textbf{Chap ustun:}\begin{enumerate}[label=\arabic*),leftmargin=9mm]
\item Merge Field
\item Recipient List
\item Preview Results
\end{enumerate}
\textbf{O'ng ustun:}\begin{enumerate}[label=\Alph*),leftmargin=9mm]
\item Manbadan dinamik maydon
\item Birlashtirilgan hujjat natijasini tekshirish
\item Qabul qiluvchilar jadvali
\end{enumerate}
\begin{tcolorbox}[enhanced,breakable,colback=teal!4,colframe=teal!55!black,boxrule=0.45pt,arc=1.2mm,left=2.8mm,right=2.8mm,top=2mm,bottom=2mm]
\textbf{\color{teal}To'g'ri javob:} \textbf{1-A, 2-C, 3-B}
\par\smallskip
\textbf{Izoh:} To'g'ri moslik: Merge Field - Manbadan dinamik maydon, Recipient List - Qabul qiluvchilar jadvali, Preview Results - Birlashtirilgan hujjat natijasini tekshirish.
\end{tcolorbox}
\vspace{2mm}\hrule height .2pt\vspace{2mm}

% DB c1afb60a-982a-b0ad-be78-b62786990051 status=published
\needspace{6\baselineskip}
\textbf{\color{ink}138.} {ootnotesizecolor{muted}Y2 | published | c1afb60a}
\par
Hujjat obyektlari va ulardan foydalanish vazifalarini moslang. Barcha mosliklar bitta ma'noli bo'lsin.
\textbf{Chap ustun:}\begin{enumerate}[label=\arabic*),leftmargin=9mm]
\item WordArt
\item SmartArt
\item Cross-reference
\item Caption
\end{enumerate}
\textbf{O'ng ustun:}\begin{enumerate}[label=\Alph*),leftmargin=9mm]
\item Obyektga raqamli sarlavha kiritish
\item Hujjat ichidagi boshqa nomerlangan obyektga havola
\item Bezakli matn ko'rinishi
\item Grafik jarayon yoki iyerarxiyani ko'rsatish
\end{enumerate}
\begin{tcolorbox}[enhanced,breakable,colback=teal!4,colframe=teal!55!black,boxrule=0.45pt,arc=1.2mm,left=2.8mm,right=2.8mm,top=2mm,bottom=2mm]
\textbf{\color{teal}To'g'ri javob:} \textbf{1-C, 2-D, 3-B, 4-A}
\par\smallskip
\textbf{Izoh:} WordArt-C; SmartArt-D; Cross-reference-B; Caption-A. Rasm va interfeysni ajratish asosidagi 13-aprel uslubiga mos.
\end{tcolorbox}
\vspace{2mm}\hrule height .2pt\vspace{2mm}

% DB c8c9b04e-b1cb-66c8-c065-b3438c9aee97 status=published
\needspace{6\baselineskip}
\textbf{\color{ink}139.} {ootnotesizecolor{muted}Y1 | published | c8c9b04e}
\par
125 ta sertifikat uchun Word maketi va Excel ro'yxati tayyor. Excel ro'yxatidagi "Sinf" ustuni keyin o'zgargan. Qaysi jarayon o'zgargan ma'lumotlar bilan sertifikatlarni qayta yaratishga yordam beradi?
\begin{enumerate}[label=\Alph*),leftmargin=9mm]
\item PowerPointda 125 ta mustaqil slayd yaratish
\item Excelda faqat Sort qilish va avvalgi PDFni o'zgartirmasdan tarqatish
\item Word Mail Merge'da Excel manbasini tekshirish/yangilash, Merge Fields ni bog'lash, Preview Results orqali tekshirish va Finish \& Merge qilish
\item Word Find \& Replace bilan 125 ta ismni bittadan almashtirish
\end{enumerate}
\begin{tcolorbox}[enhanced,breakable,colback=teal!4,colframe=teal!55!black,boxrule=0.45pt,arc=1.2mm,left=2.8mm,right=2.8mm,top=2mm,bottom=2mm]
\textbf{\color{teal}To'g'ri javob:} \textbf{C) Word Mail Merge'da Excel manbasini tekshirish/yangilash, Merge Fields ni bog'lash, Preview Results orqali tekshirish va Finish \& Merge qilish}
\par\smallskip
\textbf{Izoh:} Mail Merge ma'lumot manbasidagi dolzarb yozuvlar asosida alohida sertifikatlar hosil qiladi. Manbaga ulanishni yangilash, ustun-maydon mosligini va oldindan ko'rishni tekshirish zarur.
\end{tcolorbox}
\vspace{2mm}\hrule height .2pt\vspace{2mm}

% DB cf8a08e4-979e-dff6-bde5-0fb8a0e1a04e status=published
\needspace{6\baselineskip}
\textbf{\color{ink}140.} {ootnotesizecolor{muted}Y1 | published | cf8a08e4}
\par
Kimyoviy ifodada indeks va kasrli formulalar kerak. Eng mos vosita?
\begin{enumerate}[label=\Alph*),leftmargin=9mm]
\item WordArt faqat
\item Equation
\item Orientation
\item Footer
\end{enumerate}
\begin{tcolorbox}[enhanced,breakable,colback=teal!4,colframe=teal!55!black,boxrule=0.45pt,arc=1.2mm,left=2.8mm,right=2.8mm,top=2mm,bottom=2mm]
\textbf{\color{teal}To'g'ri javob:} \textbf{B) Equation}
\par\smallskip
\textbf{Izoh:} Equation murakkab matematik tuzilma yaratadi.
\end{tcolorbox}
\vspace{2mm}\hrule height .2pt\vspace{2mm}

% DB 005e0fb4-a459-f25c-befc-e9fe87181e09 status=published
\needspace{6\baselineskip}
\textbf{\color{ink}141.} {ootnotesizecolor{muted}Y2 | published | 005e0fb4}
\par
Moslashtiring.  1) Merge Cells;  2) Crop;  3) Hyperlink;  Ta'riflar: A) Matn yoki obyektni manzilga bog'lash;  B) Rasmning ko'rinadigan qismini chegaralash;  C) Kataklarni birlashtirish;
\textbf{Chap ustun:}\begin{enumerate}[label=\arabic*),leftmargin=9mm]
\item Merge Cells
\item Crop
\item Hyperlink
\end{enumerate}
\textbf{O'ng ustun:}\begin{enumerate}[label=\Alph*),leftmargin=9mm]
\item Matn yoki obyektni manzilga bog'lash
\item Rasmning ko'rinadigan qismini chegaralash
\item Kataklarni birlashtirish
\end{enumerate}
\begin{tcolorbox}[enhanced,breakable,colback=teal!4,colframe=teal!55!black,boxrule=0.45pt,arc=1.2mm,left=2.8mm,right=2.8mm,top=2mm,bottom=2mm]
\textbf{\color{teal}To'g'ri javob:} \textbf{1-C, 2-B, 3-A}
\par\smallskip
\textbf{Izoh:} Mos keladigan ta'riflar: 1-C, 2-B, 3-A.
\end{tcolorbox}
\vspace{2mm}\hrule height .2pt\vspace{2mm}

% DB 35dbc6d7-2989-17be-a8c4-1b5b5c24a6ea status=published
\needspace{6\baselineskip}
\textbf{\color{ink}142.} {ootnotesizecolor{muted}Y1 | published | 35dbc6d7}
\par
Wordda 7-rasm uchun Caption va unga matndan Cross-reference bor. Oldiga yana bir rasm joylashtirilib nomerlash yangilandi. Matndagi havola eski raqamni ko'rsatmoqda. Keyingi amal?
\begin{enumerate}[label=\Alph*),leftmargin=9mm]
\item Word Countni qayta hisoblash
\item Havolani faqat ko'k rang qilish
\item Rasmga Crop qo'llash
\item Cross-reference maydonlarini yangilab mos rasm raqamini tekshirish
\end{enumerate}
\begin{tcolorbox}[enhanced,breakable,colback=teal!4,colframe=teal!55!black,boxrule=0.45pt,arc=1.2mm,left=2.8mm,right=2.8mm,top=2mm,bottom=2mm]
\textbf{\color{teal}To'g'ri javob:} \textbf{D) Cross-reference maydonlarini yangilab mos rasm raqamini tekshirish}
\par\smallskip
\textbf{Izoh:} Raqamlangan obyektlar joyi o'zgargach havola maydonlarini yangilash lozim.
\end{tcolorbox}
\vspace{2mm}\hrule height .2pt\vspace{2mm}

% DB 4b395dd0-aedd-a0ad-b567-e9b4eef824d1 status=published
\needspace{6\baselineskip}
\textbf{\color{ink}143.} {ootnotesizecolor{muted}Y3 | published | 4b395dd0}
\par
Wordda Excel ro'yxatidan foydalanib xatosiz shaxsiylashtirilgan sertifikatlar yaratish uchun quyidagi amallarni mantiqan bog'liq ketma-ketlikda joylang.
\begin{enumerate}[label=\Alph*),leftmargin=9mm]
\item Merge Field maydonlarini qo'yish
\item Excel ma'lumotlar manbasini tanlash
\item Preview Results bilan qiymatlarni nazorat qilish
\item Finish \& Merge orqali yakuniy hujjatlarni yaratish
\end{enumerate}
\begin{tcolorbox}[enhanced,breakable,colback=teal!4,colframe=teal!55!black,boxrule=0.45pt,arc=1.2mm,left=2.8mm,right=2.8mm,top=2mm,bottom=2mm]
\textbf{\color{teal}To'g'ri javob:} \textbf{B - A - C - D}
\par\smallskip
\textbf{Izoh:} Ma'lumot manbasi biriktirilgach ustunlarga mos Merge Field maydonlari ishlaydi; yakuniy yaratishdan avval natija tekshiriladi.
\end{tcolorbox}
\vspace{2mm}\hrule height .2pt\vspace{2mm}

% DB 7bc6e6d6-bf03-4378-a005-54a70b974d83 status=published
\needspace{6\baselineskip}
\textbf{\color{ink}144.} {ootnotesizecolor{muted}Y3 | published | 7bc6e6d6}
\par
MS Wordda ma'lumotlar manbasi hali tayyorlanmagan, ammo qabul qiluvchilar uchun umumiy xat yaratish rejalashtirilgan. Ro'yxatni tayyorlash, unga bog'langan xat shablonini yaratish va Merge Fieldlar natijasini tekshirish tartibini belgilang.
\begin{enumerate}[label=\Alph*),leftmargin=9mm]
\item Preview Results bilan tekshirish
\item Ma'lumotlar manbasini tayyorlash
\item Merge Fieldlarni joylashtirish
\item Tayyor ma'lumotlar manbasiga bog'langan xat shablonini yaratish
\end{enumerate}
\begin{tcolorbox}[enhanced,breakable,colback=teal!4,colframe=teal!55!black,boxrule=0.45pt,arc=1.2mm,left=2.8mm,right=2.8mm,top=2mm,bottom=2mm]
\textbf{\color{teal}To'g'ri javob:} \textbf{B - D - C - A}
\par\smallskip
\textbf{Izoh:} Ma'lumotlar manbasi avval tayyorlanadi, keyin unga bog'langan shablon yaratiladi. So'ng Merge Field maydonlari kiritilib, Preview Results bilan tekshiriladi. Tartib: 2 \ensuremath{\rightarrow} 4 \ensuremath{\rightarrow} 3 \ensuremath{\rightarrow} 1.
\end{tcolorbox}
\vspace{2mm}\hrule height .2pt\vspace{2mm}

% DB 813fb43d-17bf-7f9c-7368-8d1fe26287d9 status=published
\needspace{6\baselineskip}
\textbf{\color{ink}145.} {ootnotesizecolor{muted}Y1 | published | 813fb43d}
\par
Maktab 40 ta sertifikat chiqaradi. Ro'yxatda 40 ta satr, har satrda ism va sinf bor. Word shablonda ism/sinf avtomatik almashsin. Qaysi yechim mos?
\begin{enumerate}[label=\Alph*),leftmargin=9mm]
\item Har bir sertifikatni alohida qo'lda yozish
\item Slide Master bilan Word hujjatini yaratish
\item Recipient List + Merge Field (Ism, Sinf) + Preview Results + Finish Merge
\item Merge Cells orqali 40 katakni birlashtirish
\end{enumerate}
\begin{tcolorbox}[enhanced,breakable,colback=teal!4,colframe=teal!55!black,boxrule=0.45pt,arc=1.2mm,left=2.8mm,right=2.8mm,top=2mm,bottom=2mm]
\textbf{\color{teal}To'g'ri javob:} \textbf{C) Recipient List + Merge Field (Ism, Sinf) + Preview Results + Finish Merge}
\par\smallskip
\textbf{Izoh:} Mail Merge shablon va ma'lumot manbasini bog'laydi; Preview Results bo'sh yoki noto'g'ri maydonlarni tekshirishga yordam beradi.
\end{tcolorbox}
\vspace{2mm}\hrule height .2pt\vspace{2mm}

% DB b3bd02a9-09cd-be84-c862-c24e4d11d9db status=published
\needspace{6\baselineskip}
\textbf{\color{ink}146.} {ootnotesizecolor{muted}Y2 | published | b3bd02a9}
\par
Quyidagi tushunchalarni mazmuniga qarab bog'lang.  1) Caption;  2) Cross-reference;  3) Bookmark;  Izohlar: A) Rasm yoki jadvalga raqamli izoh;  B) Hujjatning nomlangan manzili;  C) Boshqa elementga ichki havola;
\textbf{Chap ustun:}\begin{enumerate}[label=\arabic*),leftmargin=9mm]
\item Caption
\item Cross-reference
\item Bookmark
\end{enumerate}
\textbf{O'ng ustun:}\begin{enumerate}[label=\Alph*),leftmargin=9mm]
\item Rasm yoki jadvalga raqamli izoh
\item Hujjatning nomlangan manzili
\item Boshqa elementga ichki havola
\end{enumerate}
\begin{tcolorbox}[enhanced,breakable,colback=teal!4,colframe=teal!55!black,boxrule=0.45pt,arc=1.2mm,left=2.8mm,right=2.8mm,top=2mm,bottom=2mm]
\textbf{\color{teal}To'g'ri javob:} \textbf{1-A, 2-C, 3-B}
\par\smallskip
\textbf{Izoh:} To'g'ri moslik: Caption - Rasm yoki jadvalga raqamli izoh, Cross-reference - Boshqa elementga ichki havola, Bookmark - Hujjatning nomlangan manzili.
\end{tcolorbox}
\vspace{2mm}\hrule height .2pt\vspace{2mm}

% DB e2e2c9e6-56a5-429d-b6a2-c1a2efd7697a status=published
\needspace{6\baselineskip}
\textbf{\color{ink}147.} {ootnotesizecolor{muted}Y3 | published | e2e2c9e6}
\par
MS Word hujjatida rasmga avtomatik raqamlangan Caption qo'yish, matnda shu rasmga Cross-reference havolasi yaratish, so'ng rasm raqamlanishi o'zgarganidan keyin havoladagi raqamni yangilash kerak. Quyidagi amallarni shu natijaga erishish uchun mantiqiy ketma-ketlikda joylashtiring.
\begin{enumerate}[label=\Alph*),leftmargin=9mm]
\item Cross-reference maydonini yangilash
\item Jadval yoki rasmga Caption qo'shish
\item Matndan Cross-reference qo'shish
\item Rasm raqamlanishini o'zgartirish
\end{enumerate}
\begin{tcolorbox}[enhanced,breakable,colback=teal!4,colframe=teal!55!black,boxrule=0.45pt,arc=1.2mm,left=2.8mm,right=2.8mm,top=2mm,bottom=2mm]
\textbf{\color{teal}To'g'ri javob:} \textbf{B - C - D - A}
\par\smallskip
\textbf{Izoh:} Avval rasmga Caption qo'yiladi, keyin shu rasmga Cross-reference yaratiladi. Rasm raqamlanishi o'zgargach, havola yangi raqamni ko'rsatishi uchun Cross-reference maydoni yangilanadi. To'g'ri tartib: 2 \ensuremath{\rightarrow} 3 \ensuremath{\rightarrow} 4 \ensuremath{\rightarrow} 1.
\end{tcolorbox}
\vspace{2mm}\hrule height .2pt\vspace{2mm}

% DB abed67b4-f4d5-44f4-a313-2f82909fb405 status=review
\needspace{6\baselineskip}
\textbf{\color{ink}148.} {ootnotesizecolor{muted}Y1 | review | abed67b4}
\par
Maktab 40 ta sertifikat chiqaradi. Ro'yxatda 40 ta satr, har satrda ism va sinf bor. Word shablonda ism/sinf avtomatik almashsin. Qaysi yechim mos?
\begin{enumerate}[label=\Alph*),leftmargin=9mm]
\item Insert \ensuremath{\rightarrow} Table \ensuremath{\rightarrow} Merge Cells \ensuremath{\rightarrow} Copy \ensuremath{\rightarrow} Save As PDF
\item References \ensuremath{\rightarrow} Insert Citation \ensuremath{\rightarrow} Update Field \ensuremath{\rightarrow} Print Current Page
\item Recipient List + Merge Field (Ism, Sinf) + Preview Results + Finish Merge
\item Review \ensuremath{\rightarrow} Track Changes \ensuremath{\rightarrow} Accept All \ensuremath{\rightarrow} Print 40 Copies
\end{enumerate}
\begin{tcolorbox}[enhanced,breakable,colback=teal!4,colframe=teal!55!black,boxrule=0.45pt,arc=1.2mm,left=2.8mm,right=2.8mm,top=2mm,bottom=2mm]
\textbf{\color{teal}To'g'ri javob:} \textbf{C) Recipient List + Merge Field (Ism, Sinf) + Preview Results + Finish Merge}
\par\smallskip
\textbf{Izoh:} Word Mail Merge ma'lumot manbasidagi 40 ta yozuvni shablon bilan bog'laydi; Recipient List va Merge Field yordamida ism hamda sinf maydonlari qo'shiladi, Preview Results orqali tekshiriladi va Finish Merge bilan alohida sertifikatlar tayyorlanadi. Jadval kataklarini birlashtirish yoki oddiy nusxalash satrlar bo'yicha avtomatik almashtirish emas.
\end{tcolorbox}
\vspace{2mm}\hrule height .2pt\vspace{2mm}
\section{Microsoft Excel}
\subsection{MS Excel: elektron jadval tuzilishi, kiritish va formatlash}

% DB 136362ca-2f19-23bb-7f96-90e183a59823 status=published
\needspace{6\baselineskip}
\textbf{\color{ink}149.} {ootnotesizecolor{muted}Y2 | published | 136362ca}
\par
Quyidagi buyruqlarni vazifasi bilan moslashtiring.  1) A1;  2) C5;  3) B2:D4;  Izohlar: A) Birinchi ustunning birinchi satri;  B) Ko'p katakli to'rtburchak diapazon;  C) Uchinchi ustunning beshinchi satri;
\textbf{Chap ustun:}\begin{enumerate}[label=\arabic*),leftmargin=9mm]
\item A1
\item C5
\item B2:D4
\end{enumerate}
\textbf{O'ng ustun:}\begin{enumerate}[label=\Alph*),leftmargin=9mm]
\item Birinchi ustunning birinchi satri
\item Ko'p katakli to'rtburchak diapazon
\item Uchinchi ustunning beshinchi satri
\end{enumerate}
\begin{tcolorbox}[enhanced,breakable,colback=teal!4,colframe=teal!55!black,boxrule=0.45pt,arc=1.2mm,left=2.8mm,right=2.8mm,top=2mm,bottom=2mm]
\textbf{\color{teal}To'g'ri javob:} \textbf{1-A, 2-C, 3-B}
\par\smallskip
\textbf{Izoh:} To'g'ri moslik: A1 - Birinchi ustunning birinchi satri, C5 - Uchinchi ustunning beshinchi satri, B2:D4 - Ko'p katakli to'rtburchak diapazon.
\end{tcolorbox}
\vspace{2mm}\hrule height .2pt\vspace{2mm}

% DB 8c8f05d3-508c-ae8b-5436-3731d3622280 status=published
\needspace{6\baselineskip}
\textbf{\color{ink}150.} {ootnotesizecolor{muted}Y1 | published | 8c8f05d3}
\par
Excelda B4 manzili nimani anglatadi?
\begin{enumerate}[label=\Alph*),leftmargin=9mm]
\item 4 ta ustunli diapazon
\item B varaqning 4-fayli
\item B ustuni va 4-satr kesishmasi
\item 4-ustun B-satr
\end{enumerate}
\begin{tcolorbox}[enhanced,breakable,colback=teal!4,colframe=teal!55!black,boxrule=0.45pt,arc=1.2mm,left=2.8mm,right=2.8mm,top=2mm,bottom=2mm]
\textbf{\color{teal}To'g'ri javob:} \textbf{C) B ustuni va 4-satr kesishmasi}
\par\smallskip
\textbf{Izoh:} Katak manzilida avval ustun harfi, so'ng satr raqami.
\end{tcolorbox}
\vspace{2mm}\hrule height .2pt\vspace{2mm}

% DB 8e26379d-053d-6deb-32c9-8b7a3b257d6b status=published
\needspace{6\baselineskip}
\textbf{\color{ink}151.} {ootnotesizecolor{muted}Y1 | published | 8e26379d}
\par
A1:A5 diapazonida nechta katak bor?
\begin{enumerate}[label=\Alph*),leftmargin=9mm]
\item 10 ta
\item 4 ta
\item 6 ta
\item 5 ta
\end{enumerate}
\begin{tcolorbox}[enhanced,breakable,colback=teal!4,colframe=teal!55!black,boxrule=0.45pt,arc=1.2mm,left=2.8mm,right=2.8mm,top=2mm,bottom=2mm]
\textbf{\color{teal}To'g'ri javob:} \textbf{D) 5 ta}
\par\smallskip
\textbf{Izoh:} 1-dan 5-gacha beshta satr.
\end{tcolorbox}
\vspace{2mm}\hrule height .2pt\vspace{2mm}

% DB c61d5170-ef74-5e3a-898f-a4a4b57be955 status=published
\needspace{6\baselineskip}
\textbf{\color{ink}152.} {ootnotesizecolor{muted}Y1 | published | c61d5170}
\par
Microsoft Excelda ustun va satr kesishmasida joylashgan jadvalning asosiy elementi nima?
\begin{enumerate}[label=\Alph*),leftmargin=9mm]
\item Ustun
\item Diagramma
\item Formula
\item Katak
\end{enumerate}
\begin{tcolorbox}[enhanced,breakable,colback=teal!4,colframe=teal!55!black,boxrule=0.45pt,arc=1.2mm,left=2.8mm,right=2.8mm,top=2mm,bottom=2mm]
\textbf{\color{teal}To'g'ri javob:} \textbf{D) Katak}
\par\smallskip
\textbf{Izoh:} Elektron jadvalning satr va ustun kesishmasida joylashgan asosiy elementi katak (cell) deyiladi.
\end{tcolorbox}
\vspace{2mm}\hrule height .2pt\vspace{2mm}

% DB ca132419-4ebf-4fb2-b790-7d3b24df8048 status=published
\needspace{6\baselineskip}
\textbf{\color{ink}153.} {ootnotesizecolor{muted}Y3 | published | ca132419}
\par
MS Excelda yangi ish kitobini yaratib, ustun sarlavhalari va qiymatlarni kiritish hamda faylni XLSX formatida saqlash uchun amallarni tartiblang.
\begin{enumerate}[label=\Alph*),leftmargin=9mm]
\item Faylni XLSX ko'rinishida saqlash
\item Ish kitobini yaratish
\item Qiymatlarni kataklarga kiritish
\item Ish varag'iga sarlavhalarni yozish
\end{enumerate}
\begin{tcolorbox}[enhanced,breakable,colback=teal!4,colframe=teal!55!black,boxrule=0.45pt,arc=1.2mm,left=2.8mm,right=2.8mm,top=2mm,bottom=2mm]
\textbf{\color{teal}To'g'ri javob:} \textbf{B - D - C - A}
\par\smallskip
\textbf{Izoh:} Bosqichlar: Ish kitobini yaratish \ensuremath{\rightarrow} Ish varag'iga sarlavhalarni yozish \ensuremath{\rightarrow} Qiymatlarni kataklarga kiritish \ensuremath{\rightarrow} Faylni XLSX ko'rinishida saqlash.
\end{tcolorbox}
\vspace{2mm}\hrule height .2pt\vspace{2mm}

% DB cc3e84f0-0e5d-27cb-5457-6e45973c7641 status=published
\needspace{6\baselineskip}
\textbf{\color{ink}154.} {ootnotesizecolor{muted}Y1 | published | cc3e84f0}
\par
Name Box odatda nimani ko'rsatadi?
\begin{enumerate}[label=\Alph*),leftmargin=9mm]
\item Printer nomini
\item Diagramma rangini
\item Faol katak manzilini
\item Shrift o'lchamini
\end{enumerate}
\begin{tcolorbox}[enhanced,breakable,colback=teal!4,colframe=teal!55!black,boxrule=0.45pt,arc=1.2mm,left=2.8mm,right=2.8mm,top=2mm,bottom=2mm]
\textbf{\color{teal}To'g'ri javob:} \textbf{C) Faol katak manzilini}
\par\smallskip
\textbf{Izoh:} Name Box katak yoki tanlangan diapazon nomi/manzilini ko'rsatishi mumkin.
\end{tcolorbox}
\vspace{2mm}\hrule height .2pt\vspace{2mm}

% DB e02df9aa-0ecd-537d-9c54-662be8304b55 status=published
\needspace{6\baselineskip}
\textbf{\color{ink}155.} {ootnotesizecolor{muted}Y1 | published | e02df9aa}
\par
Elektron jadval katagiga kiritilgan qiymat son ham, formula ham bo'lmasa, u qanday turdagi ma'lumot sifatida qabul qilinadi?
\begin{enumerate}[label=\Alph*),leftmargin=9mm]
\item Funksiya
\item Formula
\item Matn
\item Vaqt
\end{enumerate}
\begin{tcolorbox}[enhanced,breakable,colback=teal!4,colframe=teal!55!black,boxrule=0.45pt,arc=1.2mm,left=2.8mm,right=2.8mm,top=2mm,bottom=2mm]
\textbf{\color{teal}To'g'ri javob:} \textbf{C) Matn}
\par\smallskip
\textbf{Izoh:} Son yoki \texttt{=} bilan boshlanuvchi formula sifatida tanilmagan ketma-ketlik matn sifatida qabul qilinadi.
\end{tcolorbox}
\vspace{2mm}\hrule height .2pt\vspace{2mm}

% DB eddb337b-0211-48cd-af75-e5566fa341d3 status=published
\needspace{6\baselineskip}
\textbf{\color{ink}156.} {ootnotesizecolor{muted}Y1 | published | eddb337b}
\par
Excelda workbook nima?
\begin{enumerate}[label=\Alph*),leftmargin=9mm]
\item Jadval formulasini yoki qiymatini saqlovchi alohida elektron katak
\item Bir yoki bir necha ish varag'idan iborat ish kitobi
\item Ustun va satrlardan tashkil topgan bitta ish varag'i
\item Jadval ma'lumotlari asosida yaratiladigan diagramma va grafik obyekti
\end{enumerate}
\begin{tcolorbox}[enhanced,breakable,colback=teal!4,colframe=teal!55!black,boxrule=0.45pt,arc=1.2mm,left=2.8mm,right=2.8mm,top=2mm,bottom=2mm]
\textbf{\color{teal}To'g'ri javob:} \textbf{B) Bir yoki bir necha ish varag'idan iborat ish kitobi}
\par\smallskip
\textbf{Izoh:} Excel workbook - tarkibida bir yoki bir nechta worksheet saqlaydigan ish kitobidir. Katak bitta qiymat yoki formulani saqlaydi, worksheet ish kitobining bitta varag'i, chart esa jadval ma'lumotidan hosil bo'ladigan diagrammadir.
\end{tcolorbox}
\vspace{2mm}\hrule height .2pt\vspace{2mm}

% DB 36335e9f-c600-4a3b-b7b1-b05516cd2a76 status=published
\needspace{6\baselineskip}
\textbf{\color{ink}157.} {ootnotesizecolor{muted}Y1 | published | 36335e9f}
\par
Excelda katakka \texttt{=A1+B1} formulasi kiritildi. Odatdagi ko'rinishda Formula Bar bilan natija ko'rsatilgan katak orasidagi farq qaysi javobda to'g'ri?
\begin{enumerate}[label=\Alph*),leftmargin=9mm]
\item Formula Bar katakning rangini, katak esa shrift turini ko'rsatadi
\item Formula Bar ifodani, katak hisoblangan natijani ko'rsatishi mumkin
\item Formula Bar katak manzilini, katak esa faqat funksiyalar ro'yxatini ko'rsatadi
\item Formula Bar natijani yashiradi, katak esa formulalarni faqat matn ko'rinishida chiqaradi
\end{enumerate}
\begin{tcolorbox}[enhanced,breakable,colback=teal!4,colframe=teal!55!black,boxrule=0.45pt,arc=1.2mm,left=2.8mm,right=2.8mm,top=2mm,bottom=2mm]
\textbf{\color{teal}To'g'ri javob:} \textbf{B) Formula Bar ifodani, katak hisoblangan natijani ko'rsatishi mumkin}
\par\smallskip
\textbf{Izoh:} Excelda Formula Bar tanlangan katakdagi formula ifodasini ko'rsatishi, jadval katagi esa shu formula hisoblangan natijasini aks ettirishi mumkin. Formula ko'rinishini maxsus sozlama bilan ham o'zgartirish mumkin.
\end{tcolorbox}
\vspace{2mm}\hrule height .2pt\vspace{2mm}

% DB 3e69b502-bd7d-4892-b855-0861e93b93f1 status=published
\needspace{6\baselineskip}
\textbf{\color{ink}158.} {ootnotesizecolor{muted}Y3 | published | 3e69b502}
\par
MS Excelda jadvalning birinchi katagidan oxirgi katagigacha bo'lgan diapazonga yagona chegara formati qo'llash kerak. Diapazonni aniqlash, belgilash va formatlash amallarini tartiblang.
\begin{enumerate}[label=\Alph*),leftmargin=9mm]
\item Birinchi katakni tanlash
\item Diapazonning oxirgi katagini aniqlash
\item Formatni tanlangan diapazonga qo'llash
\item Kerakli diapazonni belgilash
\end{enumerate}
\begin{tcolorbox}[enhanced,breakable,colback=teal!4,colframe=teal!55!black,boxrule=0.45pt,arc=1.2mm,left=2.8mm,right=2.8mm,top=2mm,bottom=2mm]
\textbf{\color{teal}To'g'ri javob:} \textbf{A - B - D - C}
\par\smallskip
\textbf{Izoh:} Bosqichlar: Birinchi katakni tanlash \ensuremath{\rightarrow} Diapazonning oxirgi katagini aniqlash \ensuremath{\rightarrow} Kerakli diapazonni belgilash \ensuremath{\rightarrow} Formatni tanlangan diapazonga qo'llash.
\end{tcolorbox}
\vspace{2mm}\hrule height .2pt\vspace{2mm}

% DB 451af33b-3b10-d53a-3dd5-a7d7b49b1edf status=published
\needspace{6\baselineskip}
\textbf{\color{ink}159.} {ootnotesizecolor{muted}Y1 | published | 451af33b}
\par
Currency va Percentage formatlari qanday farqlanadi?
\begin{enumerate}[label=\Alph*),leftmargin=9mm]
\item Currency pul ko'rinishi, Percentage foizli ko'rinish
\item Ikkisi faqat sonni o'chiradi
\item Percentage sana formatidir
\item Currency matnni avtomatik tarjima qiladi
\end{enumerate}
\begin{tcolorbox}[enhanced,breakable,colback=teal!4,colframe=teal!55!black,boxrule=0.45pt,arc=1.2mm,left=2.8mm,right=2.8mm,top=2mm,bottom=2mm]
\textbf{\color{teal}To'g'ri javob:} \textbf{A) Currency pul ko'rinishi, Percentage foizli ko'rinish}
\par\smallskip
\textbf{Izoh:} Son formati qiymat ko'rinishiga ta'sir qiladi.
\end{tcolorbox}
\vspace{2mm}\hrule height .2pt\vspace{2mm}

% DB 715e89a4-f00c-4d4c-ba38-5149897c2d95 status=published
\needspace{6\baselineskip}
\textbf{\color{ink}160.} {ootnotesizecolor{muted}Y1 | published | 715e89a4}
\par
Excel katagida sonli 0,25 qiymati bor. Percentage formati berilgach u 25\% bo'lib ko'rindi. Odatda katakning ichki qiymatiga nima bo'ladi?
\begin{enumerate}[label=\Alph*),leftmargin=9mm]
\item Odatda saqlangan qiymatni emas, ko'rsatilishini o'zgartiradi
\item Format qiymatni har doim 100 baravar oshirib qayta saqlaydi
\item Format qiymatni matnga aylantirib, hisoblashlarni avtomatik to'xtatadi
\item Format foiz yozuvini faqat qo'lda kiritilgan matnga o'tkazadi
\end{enumerate}
\begin{tcolorbox}[enhanced,breakable,colback=teal!4,colframe=teal!55!black,boxrule=0.45pt,arc=1.2mm,left=2.8mm,right=2.8mm,top=2mm,bottom=2mm]
\textbf{\color{teal}To'g'ri javob:} \textbf{A) Odatda saqlangan qiymatni emas, ko'rsatilishini o'zgartiradi}
\par\smallskip
\textbf{Izoh:} Percentage formati sonning ko'rinishini 100 ga ko'paytirilgan foiz shaklida ko'rsatadi; 0,25 ning saqlangan sonli qiymati odatda o'zgarmaydi. Faqat ko'rinish va formatlash boshqacha bo'ladi.
\end{tcolorbox}
\vspace{2mm}\hrule height .2pt\vspace{2mm}

% DB e06e75b4-f921-5473-973d-52538ebf7c14 status=published
\needspace{6\baselineskip}
\textbf{\color{ink}161.} {ootnotesizecolor{muted}Y2 | published | e06e75b4}
\par
Atamalar va ularning aniq ma'nolarini moslashtiring.  1) Workbook;  2) Worksheet;  3) Cell;  Izohlar: A) Satr va ustun kesishgan katak;  B) Satr-ustunlar to'ri;  C) Bir yoki bir nechta ish varag'i bor fayl;
\textbf{Chap ustun:}\begin{enumerate}[label=\arabic*),leftmargin=9mm]
\item Workbook
\item Worksheet
\item Cell
\end{enumerate}
\textbf{O'ng ustun:}\begin{enumerate}[label=\Alph*),leftmargin=9mm]
\item Satr va ustun kesishgan katak
\item Satr-ustunlar to'ri
\item Bir yoki bir nechta ish varag'i bor fayl
\end{enumerate}
\begin{tcolorbox}[enhanced,breakable,colback=teal!4,colframe=teal!55!black,boxrule=0.45pt,arc=1.2mm,left=2.8mm,right=2.8mm,top=2mm,bottom=2mm]
\textbf{\color{teal}To'g'ri javob:} \textbf{1-C, 2-B, 3-A}
\par\smallskip
\textbf{Izoh:} To'g'ri moslik: Workbook - Bir yoki bir nechta ish varag'i bor fayl, Worksheet - Satr-ustunlar to'ri, Cell - Satr va ustun kesishgan katak.
\end{tcolorbox}
\vspace{2mm}\hrule height .2pt\vspace{2mm}

% DB ea1b04c4-14fa-7286-c675-d6c9a007354f status=published
\needspace{6\baselineskip}
\textbf{\color{ink}162.} {ootnotesizecolor{muted}Y1 | published | ea1b04c4}
\par
A1:B3 va A1:A3 diapazonlari qanday farq qiladi?
\begin{enumerate}[label=\Alph*),leftmargin=9mm]
\item Ikkalasida 3 katak
\item Birinchida 5, ikkinchida 4 katak
\item Birinchida 2, ikkinchida 6 katak
\item Birinchida 6, ikkinchida 3 katak
\end{enumerate}
\begin{tcolorbox}[enhanced,breakable,colback=teal!4,colframe=teal!55!black,boxrule=0.45pt,arc=1.2mm,left=2.8mm,right=2.8mm,top=2mm,bottom=2mm]
\textbf{\color{teal}To'g'ri javob:} \textbf{D) Birinchida 6, ikkinchida 3 katak}
\par\smallskip
\textbf{Izoh:} Kataklar soni ustunlar \ensuremath{\times} satrlar.
\end{tcolorbox}
\vspace{2mm}\hrule height .2pt\vspace{2mm}

% DB 043cec72-42ff-4df3-b893-eb7d562b7f55 status=published
\needspace{6\baselineskip}
\textbf{\color{ink}163.} {ootnotesizecolor{muted}Y3 | published | 043cec72}
\par
MS Excelda B2 katagiga 0.25 qiymatini kiritib, uni foiz formatida 25\% ko'rinishida chiqarish kerak. Katakni tanlashdan yakuniy natijani tekshirishgacha amallarni tartiblang.
\begin{enumerate}[label=\Alph*),leftmargin=9mm]
\item 25\% ko'rinishini tekshirish
\item Percentage formatini berish
\item Katakni tanlash
\item Katakka 0.25 sonini kiritish
\end{enumerate}
\begin{tcolorbox}[enhanced,breakable,colback=teal!4,colframe=teal!55!black,boxrule=0.45pt,arc=1.2mm,left=2.8mm,right=2.8mm,top=2mm,bottom=2mm]
\textbf{\color{teal}To'g'ri javob:} \textbf{C - D - B - A}
\par\smallskip
\textbf{Izoh:} Birinchi navbatda B2 katagi tanlanadi, unga 0.25 qiymati kiritiladi, Percentage formatiga o'tkaziladi va natija 25\% ekanligi tekshiriladi. To'g'ri tartib: 3 \ensuremath{\rightarrow} 4 \ensuremath{\rightarrow} 2 \ensuremath{\rightarrow} 1. Avvalgi 4 \ensuremath{\rightarrow} 3 ketma-ketligi xato edi.
\end{tcolorbox}
\vspace{2mm}\hrule height .2pt\vspace{2mm}

% DB 0bd2d525-d932-3ebc-d2ba-352026388a02 status=published
\needspace{6\baselineskip}
\textbf{\color{ink}164.} {ootnotesizecolor{muted}Y1 | published | 0bd2d525}
\par
Jadvalda nomlar va ballar bor. Nomlar matn, ballar son bo'lishi uchun?
\begin{enumerate}[label=\Alph*),leftmargin=9mm]
\item Faqat diagramma qo'shish
\item Barcha qiymatga \% format berish
\item Har ustunda mos turdagi qiymat kiritish
\item Har qatorni birlashtirish
\end{enumerate}
\begin{tcolorbox}[enhanced,breakable,colback=teal!4,colframe=teal!55!black,boxrule=0.45pt,arc=1.2mm,left=2.8mm,right=2.8mm,top=2mm,bottom=2mm]
\textbf{\color{teal}To'g'ri javob:} \textbf{C) Har ustunda mos turdagi qiymat kiritish}
\par\smallskip
\textbf{Izoh:} Tartibli jadval tahlilni osonlashtiradi.
\end{tcolorbox}
\vspace{2mm}\hrule height .2pt\vspace{2mm}

% DB 34584eab-2166-0e40-22aa-30d1a990cf59 status=published
\needspace{6\baselineskip}
\textbf{\color{ink}165.} {ootnotesizecolor{muted}Y1 | published | 34584eab}
\par
C7 katakni tanlash uchun nima qilinadi?
\begin{enumerate}[label=\Alph*),leftmargin=9mm]
\item C varaqni o'chirish
\item C ustuni, 7-satrdagi katakka bosiladi
\item 7-ustun, C-satrga
\item 7 ta worksheet yaratish
\end{enumerate}
\begin{tcolorbox}[enhanced,breakable,colback=teal!4,colframe=teal!55!black,boxrule=0.45pt,arc=1.2mm,left=2.8mm,right=2.8mm,top=2mm,bottom=2mm]
\textbf{\color{teal}To'g'ri javob:} \textbf{B) C ustuni, 7-satrdagi katakka bosiladi}
\par\smallskip
\textbf{Izoh:} Harf ustun, raqam satr.
\end{tcolorbox}
\vspace{2mm}\hrule height .2pt\vspace{2mm}

% DB 3b449782-8524-3dd6-ad20-f47234d727ce status=published
\needspace{6\baselineskip}
\textbf{\color{ink}166.} {ootnotesizecolor{muted}Y2 | published | 3b449782}
\par
Amallarni izohi bilan moslang.  1) Text;  2) Number;  3) Date;  Izohlar: A) Son qiymati;  B) Sana qiymati;  C) Belgili yozuv;
\textbf{Chap ustun:}\begin{enumerate}[label=\arabic*),leftmargin=9mm]
\item Text
\item Number
\item Date
\end{enumerate}
\textbf{O'ng ustun:}\begin{enumerate}[label=\Alph*),leftmargin=9mm]
\item Son qiymati
\item Sana qiymati
\item Belgili yozuv
\end{enumerate}
\begin{tcolorbox}[enhanced,breakable,colback=teal!4,colframe=teal!55!black,boxrule=0.45pt,arc=1.2mm,left=2.8mm,right=2.8mm,top=2mm,bottom=2mm]
\textbf{\color{teal}To'g'ri javob:} \textbf{1-C, 2-A, 3-B}
\par\smallskip
\textbf{Izoh:} To'g'ri moslik: Text - Belgili yozuv, Number - Son qiymati, Date - Sana qiymati.
\end{tcolorbox}
\vspace{2mm}\hrule height .2pt\vspace{2mm}

% DB a1397866-c1ee-9937-c0c4-5a4fc6ac8232 status=published
\needspace{6\baselineskip}
\textbf{\color{ink}167.} {ootnotesizecolor{muted}Y1 | published | a1397866}
\par
A1 da 0.4 bor. Percentage formati qo'llansa qanday ko'rinadi?
\begin{enumerate}[label=\Alph*),leftmargin=9mm]
\item 0.4\%
\item 4000\%
\item 4\%
\item 40\%
\end{enumerate}
\begin{tcolorbox}[enhanced,breakable,colback=teal!4,colframe=teal!55!black,boxrule=0.45pt,arc=1.2mm,left=2.8mm,right=2.8mm,top=2mm,bottom=2mm]
\textbf{\color{teal}To'g'ri javob:} \textbf{D) 40\%}
\par\smallskip
\textbf{Izoh:} 0.4\ensuremath{\times}100\%=40\%.
\end{tcolorbox}
\vspace{2mm}\hrule height .2pt\vspace{2mm}

% DB c38ac91d-273d-d9c1-d4d7-22f4ede0446f status=published
\needspace{6\baselineskip}
\textbf{\color{ink}168.} {ootnotesizecolor{muted}Y1 | published | c38ac91d}
\par
A2:D5 oralig'ida nechta katak mavjud?
\begin{enumerate}[label=\Alph*),leftmargin=9mm]
\item 12 ta
\item 8 ta
\item 16 ta
\item 20 ta
\end{enumerate}
\begin{tcolorbox}[enhanced,breakable,colback=teal!4,colframe=teal!55!black,boxrule=0.45pt,arc=1.2mm,left=2.8mm,right=2.8mm,top=2mm,bottom=2mm]
\textbf{\color{teal}To'g'ri javob:} \textbf{C) 16 ta}
\par\smallskip
\textbf{Izoh:} 4 ustun \ensuremath{\times} 4 satr = 16.
\end{tcolorbox}
\vspace{2mm}\hrule height .2pt\vspace{2mm}

% DB 2b844ab9-d45c-b6b3-5ef0-9e157af988bb status=published
\needspace{6\baselineskip}
\textbf{\color{ink}169.} {ootnotesizecolor{muted}Y1 | published | 2b844ab9}
\par
A1:A10 diapazonini belgilash nechta ustunni qamrab oladi?
\begin{enumerate}[label=\Alph*),leftmargin=9mm]
\item 1 ta ustun
\item 2 ta ustun
\item 10 ta ustun
\item 0 ta ustun
\end{enumerate}
\begin{tcolorbox}[enhanced,breakable,colback=teal!4,colframe=teal!55!black,boxrule=0.45pt,arc=1.2mm,left=2.8mm,right=2.8mm,top=2mm,bottom=2mm]
\textbf{\color{teal}To'g'ri javob:} \textbf{A) 1 ta ustun}
\par\smallskip
\textbf{Izoh:} Faqat A ustuni tanlangan.
\end{tcolorbox}
\vspace{2mm}\hrule height .2pt\vspace{2mm}

% DB 3331e87c-5494-34c0-fa87-50cc19cc9ed9 status=published
\needspace{6\baselineskip}
\textbf{\color{ink}170.} {ootnotesizecolor{muted}Y1 | published | 3331e87c}
\par
Excelda B2:B6 sonlar 4,6,8,10,12. Nechta katak va jami yig'indi?
\begin{enumerate}[label=\Alph*),leftmargin=9mm]
\item 5 katak, yig'indi 40
\item 6 katak, yig'indi 42
\item 5 katak, yig'indi 30
\item 4 katak, yig'indi 40
\end{enumerate}
\begin{tcolorbox}[enhanced,breakable,colback=teal!4,colframe=teal!55!black,boxrule=0.45pt,arc=1.2mm,left=2.8mm,right=2.8mm,top=2mm,bottom=2mm]
\textbf{\color{teal}To'g'ri javob:} \textbf{A) 5 katak, yig'indi 40}
\par\smallskip
\textbf{Izoh:} 4+6+8+10+12=40.
\end{tcolorbox}
\vspace{2mm}\hrule height .2pt\vspace{2mm}

% DB 3ef1fd82-8450-23ff-4963-6683ec841eb5 status=published
\needspace{6\baselineskip}
\textbf{\color{ink}171.} {ootnotesizecolor{muted}Y1 | published | 3ef1fd82}
\par
A1:C3 oralig'idagi 9 katakdan 2 tasi bo'sh. Band kataklar soni?
\begin{enumerate}[label=\Alph*),leftmargin=9mm]
\item 11 ta
\item 7 ta
\item 6 ta
\item 9 ta
\end{enumerate}
\begin{tcolorbox}[enhanced,breakable,colback=teal!4,colframe=teal!55!black,boxrule=0.45pt,arc=1.2mm,left=2.8mm,right=2.8mm,top=2mm,bottom=2mm]
\textbf{\color{teal}To'g'ri javob:} \textbf{B) 7 ta}
\par\smallskip
\textbf{Izoh:} 9-2=7.
\end{tcolorbox}
\vspace{2mm}\hrule height .2pt\vspace{2mm}

% DB 6013c973-1a6a-41aa-ac4b-bb89d3c48119 status=published
\needspace{6\baselineskip}
\textbf{\color{ink}172.} {ootnotesizecolor{muted}Y3 | published | 6013c973}
\par
MS Excel jadvalidagi sarlavha satri pastga aylantirganda ham ko'rinib turishi kerak. Freeze Panes yordamida sozlash va natijani tekshirish amallarini tartiblang.
\begin{enumerate}[label=\Alph*),leftmargin=9mm]
\item View bo'limida Freeze Panes tanlash
\item Sarlavha joyida qolganini tekshirish
\item Aylantirib ko'rish
\item Jadvalda sarlavha satrini aniqlash
\end{enumerate}
\begin{tcolorbox}[enhanced,breakable,colback=teal!4,colframe=teal!55!black,boxrule=0.45pt,arc=1.2mm,left=2.8mm,right=2.8mm,top=2mm,bottom=2mm]
\textbf{\color{teal}To'g'ri javob:} \textbf{D - A - C - B}
\par\smallskip
\textbf{Izoh:} Bosqichlar: Jadvalda sarlavha satrini aniqlash \ensuremath{\rightarrow} View bo'limida Freeze Panes tanlash \ensuremath{\rightarrow} Aylantirib ko'rish \ensuremath{\rightarrow} Sarlavha joyida qolganini tekshirish.
\end{tcolorbox}
\vspace{2mm}\hrule height .2pt\vspace{2mm}

% DB 73f33406-e43b-6830-7c9d-ecccf43aa6d7 status=published
\needspace{6\baselineskip}
\textbf{\color{ink}173.} {ootnotesizecolor{muted}Y2 | published | 73f33406}
\par
Ushbu vaziyatdagi vositalar va vazifalarni moslashtiring.  1) Percentage;  2) Currency;  3) General;  Izohlar: A) Ulushning foizli ko'rinishi;  B) Umumiy sukut format;  C) Pul birligidagi ko'rinish;
\textbf{Chap ustun:}\begin{enumerate}[label=\arabic*),leftmargin=9mm]
\item Percentage
\item Currency
\item General
\end{enumerate}
\textbf{O'ng ustun:}\begin{enumerate}[label=\Alph*),leftmargin=9mm]
\item Ulushning foizli ko'rinishi
\item Umumiy sukut format
\item Pul birligidagi ko'rinish
\end{enumerate}
\begin{tcolorbox}[enhanced,breakable,colback=teal!4,colframe=teal!55!black,boxrule=0.45pt,arc=1.2mm,left=2.8mm,right=2.8mm,top=2mm,bottom=2mm]
\textbf{\color{teal}To'g'ri javob:} \textbf{1-A, 2-C, 3-B}
\par\smallskip
\textbf{Izoh:} To'g'ri moslik: Percentage - Ulushning foizli ko'rinishi, Currency - Pul birligidagi ko'rinish, General - Umumiy sukut format.
\end{tcolorbox}
\vspace{2mm}\hrule height .2pt\vspace{2mm}

% DB d53d3cfc-f971-a72c-8c98-90cd2084dbdb status=published
\needspace{6\baselineskip}
\textbf{\color{ink}174.} {ootnotesizecolor{muted}Y1 | published | d53d3cfc}
\par
Katakda \#\#\#\#\# ko'rindi. Ehtimoliy oddiy sabab?
\begin{enumerate}[label=\Alph*),leftmargin=9mm]
\item Formulada kalit so'z yo'q
\item Bu albatta sheet himoyasi
\item Har doim \#DIV/0! xatosi
\item Ustun eni son yoki sanani ko'rsatishga tor
\end{enumerate}
\begin{tcolorbox}[enhanced,breakable,colback=teal!4,colframe=teal!55!black,boxrule=0.45pt,arc=1.2mm,left=2.8mm,right=2.8mm,top=2mm,bottom=2mm]
\textbf{\color{teal}To'g'ri javob:} \textbf{D) Ustun eni son yoki sanani ko'rsatishga tor}
\par\smallskip
\textbf{Izoh:} \#\#\#\#\# ko'pincha ko'rsatish kengligi bilan bog'liq.
\end{tcolorbox}
\vspace{2mm}\hrule height .2pt\vspace{2mm}

% DB 598bd40a-0bc3-4920-84cf-790574121068 status=published
\needspace{6\baselineskip}
\textbf{\color{ink}175.} {ootnotesizecolor{muted}Y3 | published | 598bd40a}
\par
Excel ish kitobidagi Sheet1 varag'i nomini "Natijalar" deb o'zgartirib, muvaffaqiyatli bajarilganini tasdiqlash amallarini tartiblang.
\begin{enumerate}[label=\Alph*),leftmargin=9mm]
\item Ma'noli yangi nomni yozish
\item Nom o'zgarganini tasdiqlash
\item Rename amalini ishga tushirish
\item Ish varag'i tabini tanlash
\end{enumerate}
\begin{tcolorbox}[enhanced,breakable,colback=teal!4,colframe=teal!55!black,boxrule=0.45pt,arc=1.2mm,left=2.8mm,right=2.8mm,top=2mm,bottom=2mm]
\textbf{\color{teal}To'g'ri javob:} \textbf{D - C - A - B}
\par\smallskip
\textbf{Izoh:} Bosqichlar: Ish varag'i tabini tanlash \ensuremath{\rightarrow} Rename amalini ishga tushirish \ensuremath{\rightarrow} Ma'noli yangi nomni yozish \ensuremath{\rightarrow} Nom o'zgarganini tasdiqlash.
\end{tcolorbox}
\vspace{2mm}\hrule height .2pt\vspace{2mm}

% DB 62da1340-b118-c7ce-6226-6916952049b5 status=published
\needspace{6\baselineskip}
\textbf{\color{ink}176.} {ootnotesizecolor{muted}Y2 | published | 62da1340}
\par
Moslashtiring.  1) Workbook;  2) Worksheet;  3) Range;  Ta'riflar: A) Kataklar oralig'i;  B) Ish kitobi;  C) Ish varag'i;
\textbf{Chap ustun:}\begin{enumerate}[label=\arabic*),leftmargin=9mm]
\item Workbook
\item Worksheet
\item Range
\end{enumerate}
\textbf{O'ng ustun:}\begin{enumerate}[label=\Alph*),leftmargin=9mm]
\item Kataklar oralig'i
\item Ish kitobi
\item Ish varag'i
\end{enumerate}
\begin{tcolorbox}[enhanced,breakable,colback=teal!4,colframe=teal!55!black,boxrule=0.45pt,arc=1.2mm,left=2.8mm,right=2.8mm,top=2mm,bottom=2mm]
\textbf{\color{teal}To'g'ri javob:} \textbf{1-B, 2-C, 3-A}
\par\smallskip
\textbf{Izoh:} Mos keladigan ta'riflar: 1-B, 2-C, 3-A.
\end{tcolorbox}
\vspace{2mm}\hrule height .2pt\vspace{2mm}

% DB 9937250d-36e5-401a-5907-614d16e5e226 status=published
\needspace{6\baselineskip}
\textbf{\color{ink}177.} {ootnotesizecolor{muted}Y1 | published | 9937250d}
\par
Excelda B3:E7 oralig'i tanlandi. Undan birinchi va oxirgi satrlar olib tashlansa qolgan maydonda nechta katak bo'ladi?
\begin{enumerate}[label=\Alph*),leftmargin=9mm]
\item 20 ta
\item 12 ta
\item 8 ta
\item 15 ta
\end{enumerate}
\begin{tcolorbox}[enhanced,breakable,colback=teal!4,colframe=teal!55!black,boxrule=0.45pt,arc=1.2mm,left=2.8mm,right=2.8mm,top=2mm,bottom=2mm]
\textbf{\color{teal}To'g'ri javob:} \textbf{B) 12 ta}
\par\smallskip
\textbf{Izoh:} B:E to'rtta ustun, 4-6-satrlar uchta satr: 4\ensuremath{\times}3=12.
\end{tcolorbox}
\vspace{2mm}\hrule height .2pt\vspace{2mm}

% DB a0858d98-adeb-44db-a7f4-a935782c6e7d status=published
\needspace{6\baselineskip}
\textbf{\color{ink}178.} {ootnotesizecolor{muted}Y3 | published | a0858d98}
\par
MS Excelda yangi ma'lumotlar jadvalini tayyorlashda ustunlarga mos sarlavha berish, ma'lumotlarni kiritish, formatni tekshirish va ish kitobini saqlash bosqichlarini tartiblang.
\begin{enumerate}[label=\Alph*),leftmargin=9mm]
\item Qiymat turi va formatini tekshirish
\item Jadvalni saqlash
\item Ustunlarga mos sarlavha berish
\item Kataklarga ma'lumotlarni kiritish
\end{enumerate}
\begin{tcolorbox}[enhanced,breakable,colback=teal!4,colframe=teal!55!black,boxrule=0.45pt,arc=1.2mm,left=2.8mm,right=2.8mm,top=2mm,bottom=2mm]
\textbf{\color{teal}To'g'ri javob:} \textbf{C - D - A - B}
\par\smallskip
\textbf{Izoh:} Jadval tuzilishi uchun avval sarlavhalar yoziladi, so'ng qiymatlar kiritiladi, turi va formati tekshiriladi, oxirida ish kitobi saqlanadi. Tartib: 3 \ensuremath{\rightarrow} 4 \ensuremath{\rightarrow} 1 \ensuremath{\rightarrow} 2.
\end{tcolorbox}
\vspace{2mm}\hrule height .2pt\vspace{2mm}

% DB e696099f-116f-e2fd-3db9-f8361ff642ea status=published
\needspace{6\baselineskip}
\textbf{\color{ink}179.} {ootnotesizecolor{muted}Y1 | published | e696099f}
\par
Katakda 0.045 son saqlangan. U Percentage formatida bitta kasr xonasi bilan ko'rsatildi. Ekrandagi ko'rinish va asl qiymat qaysi?
\begin{enumerate}[label=\Alph*),leftmargin=9mm]
\item 4.5\%, saqlangan qiymat 4.5
\item 45.0\%, saqlangan qiymat 45
\item 4.5\%, saqlangan qiymat 0.045
\item 0.045\%, saqlangan qiymat 0.045
\end{enumerate}
\begin{tcolorbox}[enhanced,breakable,colback=teal!4,colframe=teal!55!black,boxrule=0.45pt,arc=1.2mm,left=2.8mm,right=2.8mm,top=2mm,bottom=2mm]
\textbf{\color{teal}To'g'ri javob:} \textbf{C) 4.5\%, saqlangan qiymat 0.045}
\par\smallskip
\textbf{Izoh:} 0.045\ensuremath{\times}100=4.5\%, format saqlangan sonni o'zgartirmaydi.
\end{tcolorbox}
\vspace{2mm}\hrule height .2pt\vspace{2mm}

% DB e69fcca4-c609-6853-a852-4ebac0043a32 status=published
\needspace{6\baselineskip}
\textbf{\color{ink}180.} {ootnotesizecolor{muted}Y2 | published | e69fcca4}
\par
Quyidagi tushunchalarni mazmuniga qarab bog'lang.  1) Freeze Panes;  2) Wrap Text;  3) Merge \& Center;  Izohlar: A) Sarlavhani aylantirganda ko'rinishda ushlab turish;  B) Tanlangan kataklarni birlashtirib markazlash;  C) Katakdagi uzun matnni qatorlarga ajratish;
\textbf{Chap ustun:}\begin{enumerate}[label=\arabic*),leftmargin=9mm]
\item Freeze Panes
\item Wrap Text
\item Merge \& Center
\end{enumerate}
\textbf{O'ng ustun:}\begin{enumerate}[label=\Alph*),leftmargin=9mm]
\item Sarlavhani aylantirganda ko'rinishda ushlab turish
\item Tanlangan kataklarni birlashtirib markazlash
\item Katakdagi uzun matnni qatorlarga ajratish
\end{enumerate}
\begin{tcolorbox}[enhanced,breakable,colback=teal!4,colframe=teal!55!black,boxrule=0.45pt,arc=1.2mm,left=2.8mm,right=2.8mm,top=2mm,bottom=2mm]
\textbf{\color{teal}To'g'ri javob:} \textbf{1-A, 2-C, 3-B}
\par\smallskip
\textbf{Izoh:} To'g'ri moslik: Freeze Panes - Sarlavhani aylantirganda ko'rinishda ushlab turish, Wrap Text - Katakdagi uzun matnni qatorlarga ajratish, Merge \& Center - Tanlangan kataklarni birlashtirib markazlash.
\end{tcolorbox}
\vspace{2mm}\hrule height .2pt\vspace{2mm}

% DB 79e953a9-1ae4-437c-b275-4715358b3408 status=review
\needspace{6\baselineskip}
\textbf{\color{ink}181.} {ootnotesizecolor{muted}Y1 | review | 79e953a9}
\par
Katakda \#\#\#\#\# ko'rindi. Ehtimoliy oddiy sabab?
\begin{enumerate}[label=\Alph*),leftmargin=9mm]
\item Katak formulasi nolga bo'lish xatosini qaytarayotgan bo'lishi mumkin
\item Katakdagi matn kodlash turi bilan mos kelmayotgan bo'lishi mumkin
\item Katakda manzil topilmaganligi sababli \#REF! xatosi yashirilgan bo'lishi mumkin
\item Ustun eni son yoki sanani ko'rsatishga tor
\end{enumerate}
\begin{tcolorbox}[enhanced,breakable,colback=teal!4,colframe=teal!55!black,boxrule=0.45pt,arc=1.2mm,left=2.8mm,right=2.8mm,top=2mm,bottom=2mm]
\textbf{\color{teal}To'g'ri javob:} \textbf{D) Ustun eni son yoki sanani ko'rsatishga tor}
\par\smallskip
\textbf{Izoh:} Excelda katakda \#\#\#\#\# ko'rinishi ko'pincha ustun eni son yoki sana/vaqt qiymatini ko'rsatishga yetmasligini bildiradi. Ustunni kengaytirish odatiy tekshiruvdir. Nolga bo'lish va noto'g'ri manzil kabi formula xatolari esa odatda \#DIV/0! yoki \#REF! ko'rinishida chiqadi; ayrim sana/vaqt qiymatlari ham maxsus holatlar keltirishi mumkin.
\end{tcolorbox}
\vspace{2mm}\hrule height .2pt\vspace{2mm}
\subsection{MS Excel: formulalar, ustuvorlik va katak havolalari}

% DB 74297d88-47f2-48c7-a252-72a85e33192c status=published
\needspace{6\baselineskip}
\textbf{\color{ink}182.} {ootnotesizecolor{muted}Y2 | published | 74297d88}
\par
Excel katak manzillarini nisbiy yoki mutlaq murojaat turiga ko'ra moslashtiring.
\textbf{Chap ustun:}\begin{enumerate}[label=\arabic*),leftmargin=9mm]
\item A1
\item \$A\$1
\item A\$1
\end{enumerate}
\textbf{O'ng ustun:}\begin{enumerate}[label=\Alph*),leftmargin=9mm]
\item Nisbiy manzil
\item To'liq mutlaq manzil
\item Satr mutlaq manzil
\end{enumerate}
\begin{tcolorbox}[enhanced,breakable,colback=teal!4,colframe=teal!55!black,boxrule=0.45pt,arc=1.2mm,left=2.8mm,right=2.8mm,top=2mm,bottom=2mm]
\textbf{\color{teal}To'g'ri javob:} \textbf{1-A, 2-B, 3-C}
\par\smallskip
\textbf{Izoh:} A1 - ustun va satr nisbiy; \$A\$1 - ustun va satr mutlaq; A\$1 - faqat satr mutlaq.
\end{tcolorbox}
\vspace{2mm}\hrule height .2pt\vspace{2mm}

% DB 88e7f4d7-88af-5684-bebe-b760cc145d10 status=published
\needspace{6\baselineskip}
\textbf{\color{ink}183.} {ootnotesizecolor{muted}Y1 | published | 88e7f4d7}
\par
=(7+2)*3 natija?
\begin{enumerate}[label=\Alph*),leftmargin=9mm]
\item 13
\item 27
\item 21
\item 9
\end{enumerate}
\begin{tcolorbox}[enhanced,breakable,colback=teal!4,colframe=teal!55!black,boxrule=0.45pt,arc=1.2mm,left=2.8mm,right=2.8mm,top=2mm,bottom=2mm]
\textbf{\color{teal}To'g'ri javob:} \textbf{B) 27}
\par\smallskip
\textbf{Izoh:} Qavsdagi 9 avval hisoblanadi.
\end{tcolorbox}
\vspace{2mm}\hrule height .2pt\vspace{2mm}

% DB 8be5a32a-440b-642a-5a4f-19b91b88a4d0 status=published
\needspace{6\baselineskip}
\textbf{\color{ink}184.} {ootnotesizecolor{muted}Y1 | published | 8be5a32a}
\par
=A\$2 havolasida nima mutlaq?
\begin{enumerate}[label=\Alph*),leftmargin=9mm]
\item Ikkalasi
\item Hech biri
\item 2-satr
\item A ustun
\end{enumerate}
\begin{tcolorbox}[enhanced,breakable,colback=teal!4,colframe=teal!55!black,boxrule=0.45pt,arc=1.2mm,left=2.8mm,right=2.8mm,top=2mm,bottom=2mm]
\textbf{\color{teal}To'g'ri javob:} \textbf{C) 2-satr}
\par\smallskip
\textbf{Izoh:} \$ belgi satr oldida.
\end{tcolorbox}
\vspace{2mm}\hrule height .2pt\vspace{2mm}

% DB 8fa79405-fc7d-231a-7cf9-396da2cac92a status=published
\needspace{6\baselineskip}
\textbf{\color{ink}185.} {ootnotesizecolor{muted}Y1 | published | 8fa79405}
\par
=A1 havolasi qanday tur?
\begin{enumerate}[label=\Alph*),leftmargin=9mm]
\item To'liq mutlaq
\item Nisbiy
\item Faqat satri mutlaq
\item Faqat ustuni mutlaq
\end{enumerate}
\begin{tcolorbox}[enhanced,breakable,colback=teal!4,colframe=teal!55!black,boxrule=0.45pt,arc=1.2mm,left=2.8mm,right=2.8mm,top=2mm,bottom=2mm]
\textbf{\color{teal}To'g'ri javob:} \textbf{B) Nisbiy}
\par\smallskip
\textbf{Izoh:} \$ belgisi yo'q.
\end{tcolorbox}
\vspace{2mm}\hrule height .2pt\vspace{2mm}

% DB 9406b688-aa19-4909-879d-55421859aadd status=published
\needspace{6\baselineskip}
\textbf{\color{ink}186.} {ootnotesizecolor{muted}Y3 | published | 9406b688}
\par
MS Excelda C3 katakka =A1+\$B\$2 formula yozilib, u D5 katakka ko'chiriladi. Nisbiy A1 havolasi B3 ga o'tishi, mutlaq \$B\$2 havolasi esa o'zgarmasligi kerak. Ko'chirilgan formulani =B3+\$B\$2 tarzida olish uchun amallarni tartiblang.
\begin{enumerate}[label=\Alph*),leftmargin=9mm]
\item Boshlang'ich formulani yozish
\item \$B\$2 mutlaq havolasi o'zgarmaganini tekshirib, yakuniy formulani yozish
\item Nisbiy A1 havolasini B3 ga o'zgartirish
\item C3 dan D5 gacha siljish (+1 ustun, +2 satr) ekanini aniqlash
\end{enumerate}
\begin{tcolorbox}[enhanced,breakable,colback=teal!4,colframe=teal!55!black,boxrule=0.45pt,arc=1.2mm,left=2.8mm,right=2.8mm,top=2mm,bottom=2mm]
\textbf{\color{teal}To'g'ri javob:} \textbf{A - D - C - B}
\par\smallskip
\textbf{Izoh:} Avval C3 dagi boshlang'ich formula yoziladi. C3 dan D5 ga ko'chish +1 ustun va +2 satr siljishdir. Shu sabab A1 nisbiy havolasi B3 ga o'tadi, \$B\$2 esa mutlaq bo'lgani uchun saqlanadi. Natija: =B3+\$B\$2. Tartib: 1 \ensuremath{\rightarrow} 4 \ensuremath{\rightarrow} 3 \ensuremath{\rightarrow} 2.
\end{tcolorbox}
\vspace{2mm}\hrule height .2pt\vspace{2mm}

% DB b535a0f3-58cb-2563-c8ad-55c4555919e7 status=published
\needspace{6\baselineskip}
\textbf{\color{ink}187.} {ootnotesizecolor{muted}Y1 | published | b535a0f3}
\par
Excelda =\$A\$1 havolasi qanday tur?
\begin{enumerate}[label=\Alph*),leftmargin=9mm]
\item Mutlaq ustun va mutlaq satr
\item Nisbiy ustun va nisbiy satr
\item Nisbiy ustun, mutlaq satr
\item Mutlaq ustun, nisbiy satr
\end{enumerate}
\begin{tcolorbox}[enhanced,breakable,colback=teal!4,colframe=teal!55!black,boxrule=0.45pt,arc=1.2mm,left=2.8mm,right=2.8mm,top=2mm,bottom=2mm]
\textbf{\color{teal}To'g'ri javob:} \textbf{A) Mutlaq ustun va mutlaq satr}
\par\smallskip
\textbf{Izoh:} \$A\$1 ikkala koordinatani mahkamlaydi.
\end{tcolorbox}
\vspace{2mm}\hrule height .2pt\vspace{2mm}

% DB 14f4a3eb-3552-4d69-aecf-259720c6e0af status=published
\needspace{6\baselineskip}
\textbf{\color{ink}188.} {ootnotesizecolor{muted}Y1 | published | 14f4a3eb}
\par
Excel jadvalida A1=2, B1=4 va C1=x. Formulalar: A2=\texttt{(B1-A1)/2}, B2=\texttt{2-A1/2}, C2=\texttt{(C1-A1)*2-4}. A2:C2 qiymatlaridan tuzilgan doiraviy diagrammada A2 va B2 sektorlarining har biri doiraning 1/4 qismini, C2 sektori esa 1/2 qismini egallaydi. C1 katagiga qaysi butun son yozilgan?
\begin{enumerate}[label=\Alph*),leftmargin=9mm]
\item 7
\item 5
\item 6
\item 9
\end{enumerate}
\begin{tcolorbox}[enhanced,breakable,colback=teal!4,colframe=teal!55!black,boxrule=0.45pt,arc=1.2mm,left=2.8mm,right=2.8mm,top=2mm,bottom=2mm]
\textbf{\color{teal}To'g'ri javob:} \textbf{B) 5}
\par\smallskip
\textbf{Izoh:} A2=(4-2)/2=1 va B2=2-2/2=1. Diagrammada C2 ulardan ikki baravar, ya'ni 2. 2(x-2)-4=2 tenglamadan x=5.
\end{tcolorbox}
\vspace{2mm}\hrule height .2pt\vspace{2mm}

% DB 155d0a3c-4430-3901-baf5-6e17de805e63 status=published
\needspace{6\baselineskip}
\textbf{\color{ink}189.} {ootnotesizecolor{muted}Y1 | published | 155d0a3c}
\par
Excel formulasida \texttt{=SUM(A1:A5} yozilgan va yakuniy qavs qo'yilmagan. Quyidagi fikrlarning qaysi biri eng aniq?
\begin{enumerate}[label=\Alph*),leftmargin=9mm]
\item Bu xatolik qavs jufti yetishmasligi bilan bog'liq; Excel formulani to'g'rilashni taklif qilishi mumkin
\item Bu har doim \#REF! xatosini keltiradi
\item Bu formula baribir barcha versiyalarda majburiy to'g'ri
\item Bu avtomatik 5 ga ko'paytirish deb talqin qilinadi
\end{enumerate}
\begin{tcolorbox}[enhanced,breakable,colback=teal!4,colframe=teal!55!black,boxrule=0.45pt,arc=1.2mm,left=2.8mm,right=2.8mm,top=2mm,bottom=2mm]
\textbf{\color{teal}To'g'ri javob:} \textbf{A) Bu xatolik qavs jufti yetishmasligi bilan bog'liq; Excel formulani to'g'rilashni taklif qilishi mumkin}
\par\smallskip
\textbf{Izoh:} Sintaktik qavs yetishmaydi; \#REF! esa odatda yaroqsiz katak havolasi bilan bog'liq.
\end{tcolorbox}
\vspace{2mm}\hrule height .2pt\vspace{2mm}

% DB 1a0ebc5f-c211-5da0-b3f3-c6acca2e250c status=published
\needspace{6\baselineskip}
\textbf{\color{ink}190.} {ootnotesizecolor{muted}Y1 | published | 1a0ebc5f}
\par
Excelda A1=5, B1=0, C1=10, A2=-3, B2=6, C2=5. E5 dagi \texttt{=SUM(\$A\$1:\$B\$1)+\$A2*C2} formula D4 ga ko'chirilsa, natija qancha bo'ladi?
\begin{enumerate}[label=\Alph*),leftmargin=9mm]
\item 0
\item 5
\item 10
\item 15
\end{enumerate}
\begin{tcolorbox}[enhanced,breakable,colback=teal!4,colframe=teal!55!black,boxrule=0.45pt,arc=1.2mm,left=2.8mm,right=2.8mm,top=2mm,bottom=2mm]
\textbf{\color{teal}To'g'ri javob:} \textbf{B) 5}
\par\smallskip
\textbf{Izoh:} Bir katak chapga va yuqoriga ko'chirganda formula \texttt{=SUM(\$A\$1:\$B\$1)+\$A1*B1} bo'ladi: 5+5\ensuremath{\times}0=5.
\end{tcolorbox}
\vspace{2mm}\hrule height .2pt\vspace{2mm}

% DB 2303afb4-c7cc-bf5c-1d52-d36d0121ddda status=published
\needspace{6\baselineskip}
\textbf{\color{ink}191.} {ootnotesizecolor{muted}Y1 | published | 2303afb4}
\par
=B2 va =\$B\$2 nusxalanganda qanday farq?
\begin{enumerate}[label=\Alph*),leftmargin=9mm]
\item \$B\$2 foiz deb qabul qilinadi
\item Ikkalasi doim siljiydi
\item B2 siljiydi, \$B\$2 siljimaydi
\item Ikkalasi doim o'chadi
\end{enumerate}
\begin{tcolorbox}[enhanced,breakable,colback=teal!4,colframe=teal!55!black,boxrule=0.45pt,arc=1.2mm,left=2.8mm,right=2.8mm,top=2mm,bottom=2mm]
\textbf{\color{teal}To'g'ri javob:} \textbf{C) B2 siljiydi, \$B\$2 siljimaydi}
\par\smallskip
\textbf{Izoh:} Mutlaq havola ko'chirishda barqaror.
\end{tcolorbox}
\vspace{2mm}\hrule height .2pt\vspace{2mm}

% DB 303ef09b-1030-52f6-97ae-08731ae831f8 status=published
\needspace{6\baselineskip}
\textbf{\color{ink}192.} {ootnotesizecolor{muted}Y1 | published | 303ef09b}
\par
MS Excelda A1:A5 kataklariga mos ravishda 1, 2, 3, 4, 5 sonlari kiritilgan. B1 katakdagi \texttt{=\$A\$1+A2} formula B5 gacha avtomatik to'ldirilsa, B5 dagi natijani toping.
\begin{enumerate}[label=\Alph*),leftmargin=9mm]
\item 7
\item 1
\item 5
\item 6
\end{enumerate}
\begin{tcolorbox}[enhanced,breakable,colback=teal!4,colframe=teal!55!black,boxrule=0.45pt,arc=1.2mm,left=2.8mm,right=2.8mm,top=2mm,bottom=2mm]
\textbf{\color{teal}To'g'ri javob:} \textbf{B) 1}
\par\smallskip
\textbf{Izoh:} B5 dagi formula \texttt{=\$A\$1+A6} bo'ladi. \texttt{\$A\$1} mutlaq havola bo'lib 1 ni saqlaydi; bo'sh A6 arifmetik amalda 0 sifatida olinadi. Natija 1.
\end{tcolorbox}
\vspace{2mm}\hrule height .2pt\vspace{2mm}

% DB 3b6afbf3-c67b-02da-ece3-6c53dc01f6f1 status=published
\needspace{6\baselineskip}
\textbf{\color{ink}193.} {ootnotesizecolor{muted}Y1 | published | 3b6afbf3}
\par
\#DIV/0! qachon chiqishi mumkin?
\begin{enumerate}[label=\Alph*),leftmargin=9mm]
\item Nolga bo'lganda
\item Slayd o'chganda
\item Shrift kattalashganda
\item Hujjat sahifasi tor bo'lganda
\end{enumerate}
\begin{tcolorbox}[enhanced,breakable,colback=teal!4,colframe=teal!55!black,boxrule=0.45pt,arc=1.2mm,left=2.8mm,right=2.8mm,top=2mm,bottom=2mm]
\textbf{\color{teal}To'g'ri javob:} \textbf{A) Nolga bo'lganda}
\par\smallskip
\textbf{Izoh:} Nolga bo'lish noaniq.
\end{tcolorbox}
\vspace{2mm}\hrule height .2pt\vspace{2mm}

% DB 40c9648b-8341-4de2-bc5c-f56c9a9cd25f status=published
\needspace{6\baselineskip}
\textbf{\color{ink}194.} {ootnotesizecolor{muted}Y2 | published | 40c9648b}
\par
Excel katak murojaatlarini ustun va satrning nisbiy yoki mutlaq holatiga ko'ra moslashtiring.
\textbf{Chap ustun:}\begin{enumerate}[label=\arabic*),leftmargin=9mm]
\item \$A\$1
\item A\$1
\item \$A1
\end{enumerate}
\textbf{O'ng ustun:}\begin{enumerate}[label=\Alph*),leftmargin=9mm]
\item Ustuni mutlaq, satri nisbiy
\item Ikkala koordinatasi mutlaq
\item Satr mutlaq, ustun nisbiy
\end{enumerate}
\begin{tcolorbox}[enhanced,breakable,colback=teal!4,colframe=teal!55!black,boxrule=0.45pt,arc=1.2mm,left=2.8mm,right=2.8mm,top=2mm,bottom=2mm]
\textbf{\color{teal}To'g'ri javob:} \textbf{1-B, 2-C, 3-A}
\par\smallskip
\textbf{Izoh:} \$A\$1 - ustun va satr mutlaq; A\$1 - satr mutlaq va ustun nisbiy; \$A1 - ustun mutlaq va satr nisbiy.
\end{tcolorbox}
\vspace{2mm}\hrule height .2pt\vspace{2mm}

% DB 4b52cefe-1b0a-416c-a60b-edabe92fe6f2 status=published
\needspace{6\baselineskip}
\textbf{\color{ink}195.} {ootnotesizecolor{muted}Y2 | published | 4b52cefe}
\par
Excel formulalaridagi katak murojaatlarini qaysi koordinata mahkamlanganiga ko'ra moslashtiring.
\textbf{Chap ustun:}\begin{enumerate}[label=\arabic*),leftmargin=9mm]
\item =\$A1
\item =A\$1
\item =\$A\$1
\end{enumerate}
\textbf{O'ng ustun:}\begin{enumerate}[label=\Alph*),leftmargin=9mm]
\item Ustun mahkamlangan
\item Satr mahkamlangan
\item Ikkalasi mahkamlangan
\end{enumerate}
\begin{tcolorbox}[enhanced,breakable,colback=teal!4,colframe=teal!55!black,boxrule=0.45pt,arc=1.2mm,left=2.8mm,right=2.8mm,top=2mm,bottom=2mm]
\textbf{\color{teal}To'g'ri javob:} \textbf{1-A, 2-B, 3-C}
\par\smallskip
\textbf{Izoh:} =\$A1 havolasida A ustuni mahkam, satri nisbiy; =A\$1 havolasida 1-satr mahkam, ustuni nisbiy; =\$A\$1 da ikkalasi mahkam.
\end{tcolorbox}
\vspace{2mm}\hrule height .2pt\vspace{2mm}

% DB 4eb2106b-77a3-5018-84c6-7b2116845bc2 status=published
\needspace{6\baselineskip}
\textbf{\color{ink}196.} {ootnotesizecolor{muted}Y1 | published | 4eb2106b}
\par
Excelda B1 katakda umumiy soliq stavkasi 0,13, B3:B6 da ishchilar maoshi bor. C3 ga \texttt{=B3*B1} yozilib C6 gacha to'ldirilganda faqat birinchi natija to'g'ri chiqdi. To'g'ri formulani toping.
\begin{enumerate}[label=\Alph*),leftmargin=9mm]
\item \texttt{=B3*B1*100\%}
\item \texttt{=B3*\$B1}
\item B1 ni 13\% deb qayta yozish
\item \texttt{=B3*\$B\$1}
\end{enumerate}
\begin{tcolorbox}[enhanced,breakable,colback=teal!4,colframe=teal!55!black,boxrule=0.45pt,arc=1.2mm,left=2.8mm,right=2.8mm,top=2mm,bottom=2mm]
\textbf{\color{teal}To'g'ri javob:} \textbf{D) \texttt{=B3*\$B\$1}}
\par\smallskip
\textbf{Izoh:} Maosh havolasi qator bo'yicha o'zgarishi, stavka esa doim B1 da qolishi kerak. Shu sabab B1 mutlaq havola: \texttt{=B3*\$B\$1}.
\end{tcolorbox}
\vspace{2mm}\hrule height .2pt\vspace{2mm}

% DB 7df2815b-2e0d-7fd8-0fd1-f76b84ca92f1 status=published
\needspace{6\baselineskip}
\textbf{\color{ink}197.} {ootnotesizecolor{muted}Y1 | published | 7df2815b}
\par
=B\$3 va =\$B3 farqi?
\begin{enumerate}[label=\Alph*),leftmargin=9mm]
\item Ikkalasida ham satr mahkam
\item Ikkalasida ustun mahkam
\item Hech qaysisi havola emas
\item Birinchida satr, ikkinchida ustun mahkam
\end{enumerate}
\begin{tcolorbox}[enhanced,breakable,colback=teal!4,colframe=teal!55!black,boxrule=0.45pt,arc=1.2mm,left=2.8mm,right=2.8mm,top=2mm,bottom=2mm]
\textbf{\color{teal}To'g'ri javob:} \textbf{D) Birinchida satr, ikkinchida ustun mahkam}
\par\smallskip
\textbf{Izoh:} \$ dan keyingi komponent mahkam.
\end{tcolorbox}
\vspace{2mm}\hrule height .2pt\vspace{2mm}

% DB 7dfd9e40-a6b6-4202-a01c-c24e52d1b1c7 status=published
\needspace{6\baselineskip}
\textbf{\color{ink}198.} {ootnotesizecolor{muted}Y1 | published | 7dfd9e40}
\par
Excelda \texttt{A1}, \texttt{A\$1} va \texttt{\$A1} havolalari formula nusxalanganda turlicha o'zgaradi. Nisbiy va aralash havolani farqlashning asosiy mezoni qaysi?
\begin{enumerate}[label=\Alph*),leftmargin=9mm]
\item Ustun harflari alifbo bo'yicha qanday tartibda yozilgani
\item Katakning matn yoki raqam formatida saqlangani
\item Formuladagi havolaning bir nechta varaqda takrorlanishi
\item Qaysi koordinata oldida \$ borligi
\end{enumerate}
\begin{tcolorbox}[enhanced,breakable,colback=teal!4,colframe=teal!55!black,boxrule=0.45pt,arc=1.2mm,left=2.8mm,right=2.8mm,top=2mm,bottom=2mm]
\textbf{\color{teal}To'g'ri javob:} \textbf{D) Qaysi koordinata oldida \$ borligi}
\par\smallskip
\textbf{Izoh:} Dollar belgisi satr yoki ustun koordinatasini nusxalashda mahkamlaydi. A1 nisbiy, A\$1 satri mahkamlangan, \$A1 esa ustuni mahkamlangan aralash havoladir.
\end{tcolorbox}
\vspace{2mm}\hrule height .2pt\vspace{2mm}

% DB 9200cd04-23cc-4fb8-bdc3-3a80bcfde578 status=published
\needspace{6\baselineskip}
\textbf{\color{ink}199.} {ootnotesizecolor{muted}Y2 | published | 9200cd04}
\par
Quyidagi formulalarni katak havolasi turiga ko'ra moslashtiring. 1) =B2+3; 2) =B\$2+3; 3) =\$B2+3; Izohlar: A) Faqat ustun qat'iy; B) Ustun ham, satr ham nisbiy; C) Faqat satr qat'iy;
\textbf{Chap ustun:}\begin{enumerate}[label=\arabic*),leftmargin=9mm]
\item =B2+3
\item =B\$2+3
\item =\$B2+3
\end{enumerate}
\textbf{O'ng ustun:}\begin{enumerate}[label=\Alph*),leftmargin=9mm]
\item Faqat ustun qat'iy
\item Ustun ham, satr ham nisbiy
\item Faqat satr qat'iy
\end{enumerate}
\begin{tcolorbox}[enhanced,breakable,colback=teal!4,colframe=teal!55!black,boxrule=0.45pt,arc=1.2mm,left=2.8mm,right=2.8mm,top=2mm,bottom=2mm]
\textbf{\color{teal}To'g'ri javob:} \textbf{1-B, 2-C, 3-A}
\par\smallskip
\textbf{Izoh:} B2 havolasida ustun ham, satr ham nisbiy. B\$2 da faqat 2-satr mahkamlangan; \$B2 da faqat B ustuni mahkamlangan. Qisqa izohlarning har biri aynan bitta formulaga mos.
\end{tcolorbox}
\vspace{2mm}\hrule height .2pt\vspace{2mm}

% DB 96b6a635-cb5f-d07f-e68a-7af7843920de status=published
\needspace{6\baselineskip}
\textbf{\color{ink}200.} {ootnotesizecolor{muted}Y2 | published | 96b6a635}
\par
Atamalar va ularning aniq ma'nolarini moslashtiring.  1) \#DIV/0!;  2) \#REF!;  3) \#NAME?;  Izohlar: A) Tanimagan funksiya yoki nom;  B) Nolga bo'lish;  C) Yaroqsiz katak manzili;
\textbf{Chap ustun:}\begin{enumerate}[label=\arabic*),leftmargin=9mm]
\item \#DIV/0!
\item \#REF!
\item \#NAME?
\end{enumerate}
\textbf{O'ng ustun:}\begin{enumerate}[label=\Alph*),leftmargin=9mm]
\item Tanimagan funksiya yoki nom
\item Nolga bo'lish
\item Yaroqsiz katak manzili
\end{enumerate}
\begin{tcolorbox}[enhanced,breakable,colback=teal!4,colframe=teal!55!black,boxrule=0.45pt,arc=1.2mm,left=2.8mm,right=2.8mm,top=2mm,bottom=2mm]
\textbf{\color{teal}To'g'ri javob:} \textbf{1-B, 2-C, 3-A}
\par\smallskip
\textbf{Izoh:} To'g'ri moslik: \#DIV/0! - Nolga bo'lish, \#REF! - Yaroqsiz katak manzili, \#NAME? - Tanimagan funksiya yoki nom.
\end{tcolorbox}
\vspace{2mm}\hrule height .2pt\vspace{2mm}

% DB c25dd49c-ae18-55b1-be39-39fba31f3607 status=published
\needspace{6\baselineskip}
\textbf{\color{ink}201.} {ootnotesizecolor{muted}Y1 | published | c25dd49c}
\par
Excelda soliq stavkasi 12\%. Jadval natijalari: Muhammad - 945 va 113,4; Dudiyor - 687 va 82,44; Orif - 1024 va 122,88; Shohjahon - 1246 va 159,62; Mehriddin - 964 va 115,68. Qaysi xodimning solig'i xato hisoblangan?
\begin{enumerate}[label=\Alph*),leftmargin=9mm]
\item Orif
\item Mehriddin
\item Shohjahon
\item Dudiyor
\end{enumerate}
\begin{tcolorbox}[enhanced,breakable,colback=teal!4,colframe=teal!55!black,boxrule=0.45pt,arc=1.2mm,left=2.8mm,right=2.8mm,top=2mm,bottom=2mm]
\textbf{\color{teal}To'g'ri javob:} \textbf{C) Shohjahon}
\par\smallskip
\textbf{Izoh:} 1246\ensuremath{\times}12\%=149,52, jadvalda esa 159,62 yozilgan. Qolgan ko'rsatilgan natijalar 12\% ga teng.
\end{tcolorbox}
\vspace{2mm}\hrule height .2pt\vspace{2mm}

% DB c2a72d1e-96c4-f979-9681-7cbe29d5a5b4 status=published
\needspace{6\baselineskip}
\textbf{\color{ink}202.} {ootnotesizecolor{muted}Y1 | published | c2a72d1e}
\par
A1=12, B1=0. C1 katakdagi \texttt{=IFERROR(A1/B1,"XATO")} nimani qaytaradi?
\begin{enumerate}[label=\Alph*),leftmargin=9mm]
\item 0 sonini
\item \texttt{\#DIV/0!}ni bevosita
\item 12 sonini
\item \texttt{XATO} matnini
\end{enumerate}
\begin{tcolorbox}[enhanced,breakable,colback=teal!4,colframe=teal!55!black,boxrule=0.45pt,arc=1.2mm,left=2.8mm,right=2.8mm,top=2mm,bottom=2mm]
\textbf{\color{teal}To'g'ri javob:} \textbf{D) \texttt{XATO} matnini}
\par\smallskip
\textbf{Izoh:} A1/B1 nolga bo'linadi, lekin IFERROR ikkinchi argumentdagi \texttt{XATO} matnini qaytaradi.
\end{tcolorbox}
\vspace{2mm}\hrule height .2pt\vspace{2mm}

% DB f670dec6-7281-5b2b-bd69-5aa3ba279615 status=published
\needspace{6\baselineskip}
\textbf{\color{ink}203.} {ootnotesizecolor{muted}Y1 | published | f670dec6}
\par
Excel jadvalida \texttt{A2=13}, \texttt{C1=9}, \texttt{A3=2}, \texttt{C2=25}. B2 katakchadagi formula \texttt{=\$A2+\$C1}. Formula B2 dan C3 ga ko'chirilsa, C3 ning qiymati nechaga teng bo'ladi?
\begin{enumerate}[label=\Alph*),leftmargin=9mm]
\item 20
\item 25
\item 22
\item 27
\end{enumerate}
\begin{tcolorbox}[enhanced,breakable,colback=teal!4,colframe=teal!55!black,boxrule=0.45pt,arc=1.2mm,left=2.8mm,right=2.8mm,top=2mm,bottom=2mm]
\textbf{\color{teal}To'g'ri javob:} \textbf{D) 27}
\par\smallskip
\textbf{Izoh:} Bir ustun o'ngga va bir qator pastga ko'chirilganda \texttt{\$A2} \ensuremath{\rightarrow} \texttt{\$A3}, \texttt{\$C1} \ensuremath{\rightarrow} \texttt{\$C2}. Demak, \texttt{C3=\$A3+\$C2=2+25=27}.
\end{tcolorbox}
\vspace{2mm}\hrule height .2pt\vspace{2mm}

% DB f73203bd-b562-55ad-bea9-8422923e4fc4 status=published
\needspace{6\baselineskip}
\textbf{\color{ink}204.} {ootnotesizecolor{muted}Y1 | published | f73203bd}
\par
Elektron jadvalning qiymatlari quyidagicha: \newline  \newline |   | A | B | C | D | E | \newline |---|---:|---:|---:|---:|---:| \newline | 1 | 1 | 10 | 100 | 1000 | 10000 | \newline | 2 | 2 | 20 | - | 2000 | 20000 | \newline | 3 | 3 | \texttt{=A\$2+D\$3} | 300 | 3000 | 30000 | \newline | 4 | 4 | 40 | 400 | 4000 | 40000 | \newline  \newline B3 katakchadagi formula C2 katakchaga ko'chirildi. C2 dagi formulaning qiymatini toping.
\begin{enumerate}[label=\Alph*),leftmargin=9mm]
\item 3020
\item 3002
\item 30020
\item 3200
\end{enumerate}
\begin{tcolorbox}[enhanced,breakable,colback=teal!4,colframe=teal!55!black,boxrule=0.45pt,arc=1.2mm,left=2.8mm,right=2.8mm,top=2mm,bottom=2mm]
\textbf{\color{teal}To'g'ri javob:} \textbf{C) 30020}
\par\smallskip
\textbf{Izoh:} B3 dan C2 ga ko'chirish - bir ustun o'ngga va bir qator yuqoriga. Aralash havolalarda qatorlar \texttt{\$} bilan mahkamlangan, ustunlar esa siljiydi: \texttt{=A\$2+D\$3} \ensuremath{\rightarrow} \texttt{=B\$2+E\$3}. Natija 20+30000=30020.
\end{tcolorbox}
\vspace{2mm}\hrule height .2pt\vspace{2mm}

% DB fbd8ec23-bc93-4074-a053-7e933f71e958 status=published
\needspace{6\baselineskip}
\textbf{\color{ink}205.} {ootnotesizecolor{muted}Y3 | published | fbd8ec23}
\par
MS Excel katagiga hisoblash formulasi kiritildi. Qavslar va arifmetik amallar ustuvorligiga rioya qilib hisob natijasini topish va tekshirish jarayonini tartiblang.
\begin{enumerate}[label=\Alph*),leftmargin=9mm]
\item Arifmetik ustuvorlikni aniqlash
\item Qolgan amallarni bajarib natijani tekshirish
\item Qavs ichini hisoblash
\item Katakda formulani kiritish
\end{enumerate}
\begin{tcolorbox}[enhanced,breakable,colback=teal!4,colframe=teal!55!black,boxrule=0.45pt,arc=1.2mm,left=2.8mm,right=2.8mm,top=2mm,bottom=2mm]
\textbf{\color{teal}To'g'ri javob:} \textbf{D - A - C - B}
\par\smallskip
\textbf{Izoh:} Bosqichlar: Katakda formulani kiritish \ensuremath{\rightarrow} Arifmetik ustuvorlikni aniqlash \ensuremath{\rightarrow} Qavs ichini hisoblash \ensuremath{\rightarrow} Qolgan amallarni bajarib natijani tekshirish.
\end{tcolorbox}
\vspace{2mm}\hrule height .2pt\vspace{2mm}

% DB fe7da103-6d94-5a8a-b130-e20981925182 status=published
\needspace{6\baselineskip}
\textbf{\color{ink}206.} {ootnotesizecolor{muted}Y1 | published | fe7da103}
\par
Excelda B1:K1 kataklariga 0 dan 9 gacha, A2:A5 kataklariga 6 dan 9 gacha sonlar yozilgan. B2 dagi formula B2:K5 diapazoniga ko'chirilganda 60 dan 99 gacha sonlar jadvali hosil bo'lishi uchun B2 ga qaysi formula yozilishi kerak?
\begin{enumerate}[label=\Alph*),leftmargin=9mm]
\item =\$A2*10+B\$1
\item =A\$2*10+\$B1
\item =\$A2*10+\$B1
\item =A2*10+B1
\end{enumerate}
\begin{tcolorbox}[enhanced,breakable,colback=teal!4,colframe=teal!55!black,boxrule=0.45pt,arc=1.2mm,left=2.8mm,right=2.8mm,top=2mm,bottom=2mm]
\textbf{\color{teal}To'g'ri javob:} \textbf{A) =\$A2*10+B\$1}
\par\smallskip
\textbf{Izoh:} O'nlar raqami doim A ustunidan olinadi, shuning uchun ustun \texttt{\$A} bilan qotiriladi. Birlar raqami doim 1-qatoridan olinadi, shuning uchun \texttt{B\$1}. Formula: \texttt{=\$A2*10+B\$1}.
\end{tcolorbox}
\vspace{2mm}\hrule height .2pt\vspace{2mm}

% DB 119a9688-9111-4e6d-c59d-203698eabd7d status=published
\needspace{6\baselineskip}
\textbf{\color{ink}207.} {ootnotesizecolor{muted}Y1 | published | 119a9688}
\par
B2 formulasi =A2*2. Uni B4 ga nusxalasak?
\begin{enumerate}[label=\Alph*),leftmargin=9mm]
\item =A4*2
\item =C4*2
\item =A2*2
\item =A3*4
\end{enumerate}
\begin{tcolorbox}[enhanced,breakable,colback=teal!4,colframe=teal!55!black,boxrule=0.45pt,arc=1.2mm,left=2.8mm,right=2.8mm,top=2mm,bottom=2mm]
\textbf{\color{teal}To'g'ri javob:} \textbf{A) =A4*2}
\par\smallskip
\textbf{Izoh:} Ikki satr pastga siljiydi.
\end{tcolorbox}
\vspace{2mm}\hrule height .2pt\vspace{2mm}

% DB 2f8363ca-08f3-f955-e6f5-58782af8878d status=published
\needspace{6\baselineskip}
\textbf{\color{ink}208.} {ootnotesizecolor{muted}Y1 | published | 2f8363ca}
\par
C3 formulasi =B3+\$F\$1. Uni C5 ga nusxalasak?
\begin{enumerate}[label=\Alph*),leftmargin=9mm]
\item =D5+\$F\$1
\item =B5+\$F\$1
\item =B3+\$F\$3
\item =B5+\$F\$3
\end{enumerate}
\begin{tcolorbox}[enhanced,breakable,colback=teal!4,colframe=teal!55!black,boxrule=0.45pt,arc=1.2mm,left=2.8mm,right=2.8mm,top=2mm,bottom=2mm]
\textbf{\color{teal}To'g'ri javob:} \textbf{B) =B5+\$F\$1}
\par\smallskip
\textbf{Izoh:} Nisbiy satr siljiydi, mutlaq qoladi.
\end{tcolorbox}
\vspace{2mm}\hrule height .2pt\vspace{2mm}

% DB 345e7d43-de0e-5158-86d3-684db7b2ff7d status=published
\needspace{6\baselineskip}
\textbf{\color{ink}209.} {ootnotesizecolor{muted}Y1 | published | 345e7d43}
\par
MS Excelda A1 katakdagi \texttt{=C\$2*\$D\$7-\$B5} formula D3 katakka nusxalandi. D3 dagi formula qaysi?
\begin{enumerate}[label=\Alph*),leftmargin=9mm]
\item \texttt{=F\$2*\$D\$7-\$B7}
\item \texttt{=F\$4*\$D\$7-\$B7}
\item \texttt{=C\$2*\$G\$7-\$E7}
\item \texttt{=F\$2*\$G\$9-\$B5}
\end{enumerate}
\begin{tcolorbox}[enhanced,breakable,colback=teal!4,colframe=teal!55!black,boxrule=0.45pt,arc=1.2mm,left=2.8mm,right=2.8mm,top=2mm,bottom=2mm]
\textbf{\color{teal}To'g'ri javob:} \textbf{A) \texttt{=F\$2*\$D\$7-\$B7}}
\par\smallskip
\textbf{Izoh:} A1 dan D3 ga siljish +3 ustun va +2 satr. C\$2 dagi ustun nisbiy bo'lib F ga siljiydi, \$D\$7 o'zgarmaydi, \$B5 dagi satr 7 ga siljiydi. Natija \texttt{=F\$2*\$D\$7-\$B7}.
\end{tcolorbox}
\vspace{2mm}\hrule height .2pt\vspace{2mm}

% DB 393f6959-2ac6-57e8-2a6b-479ab0b9c863 status=published
\needspace{6\baselineskip}
\textbf{\color{ink}210.} {ootnotesizecolor{muted}Y1 | published | 393f6959}
\par
=5*(2+4)-7 natija?
\begin{enumerate}[label=\Alph*),leftmargin=9mm]
\item 53
\item 23
\item 13
\item 17
\end{enumerate}
\begin{tcolorbox}[enhanced,breakable,colback=teal!4,colframe=teal!55!black,boxrule=0.45pt,arc=1.2mm,left=2.8mm,right=2.8mm,top=2mm,bottom=2mm]
\textbf{\color{teal}To'g'ri javob:} \textbf{B) 23}
\par\smallskip
\textbf{Izoh:} Qavsdan so'ng ko'paytirish: 30-7=23.
\end{tcolorbox}
\vspace{2mm}\hrule height .2pt\vspace{2mm}

% DB 70619879-6b49-69ba-f97c-3ed2f3867bf0 status=published
\needspace{6\baselineskip}
\textbf{\color{ink}211.} {ootnotesizecolor{muted}Y2 | published | 70619879}
\par
Amallarni izohi bilan moslang.  1) SUM(A1:A3);  2) A1\textasciicircum{}2;  3) A1\&B1;  Izohlar: A) Diapazon yig'indisi;  B) Kvadrat daraja;  C) Matnlarni birlashtirish;
\textbf{Chap ustun:}\begin{enumerate}[label=\arabic*),leftmargin=9mm]
\item SUM(A1:A3)
\item A1\textasciicircum{}2
\item A1\&B1
\end{enumerate}
\textbf{O'ng ustun:}\begin{enumerate}[label=\Alph*),leftmargin=9mm]
\item Diapazon yig'indisi
\item Kvadrat daraja
\item Matnlarni birlashtirish
\end{enumerate}
\begin{tcolorbox}[enhanced,breakable,colback=teal!4,colframe=teal!55!black,boxrule=0.45pt,arc=1.2mm,left=2.8mm,right=2.8mm,top=2mm,bottom=2mm]
\textbf{\color{teal}To'g'ri javob:} \textbf{1-A, 2-B, 3-C}
\par\smallskip
\textbf{Izoh:} To'g'ri moslik: SUM(A1:A3) - Diapazon yig'indisi, A1\textasciicircum{}2 - Kvadrat daraja, A1\&B1 - Matnlarni birlashtirish.
\end{tcolorbox}
\vspace{2mm}\hrule height .2pt\vspace{2mm}

% DB 76b7b44c-efb6-8ba2-76c7-afa61805b81e status=published
\needspace{6\baselineskip}
\textbf{\color{ink}212.} {ootnotesizecolor{muted}Y1 | published | 76b7b44c}
\par
B2 formulasi =A\$1+\$C2. Uni D5 ga ko'chirsak?
\begin{enumerate}[label=\Alph*),leftmargin=9mm]
\item =A\$1+\$C2
\item =C\$4+\$E5
\item =C\$1+\$C5
\item =A\$4+\$C5
\end{enumerate}
\begin{tcolorbox}[enhanced,breakable,colback=teal!4,colframe=teal!55!black,boxrule=0.45pt,arc=1.2mm,left=2.8mm,right=2.8mm,top=2mm,bottom=2mm]
\textbf{\color{teal}To'g'ri javob:} \textbf{C) =C\$1+\$C5}
\par\smallskip
\textbf{Izoh:} Ikki ustun va uch satr siljishi qo'llanadi.
\end{tcolorbox}
\vspace{2mm}\hrule height .2pt\vspace{2mm}

% DB fb804c62-d436-46f8-be52-e397e2a69e1e status=published
\needspace{6\baselineskip}
\textbf{\color{ink}213.} {ootnotesizecolor{muted}Y3 | published | fb804c62}
\par
MS Excel formulasi xatolik kodini qaytardi. Xatoning sababini havola yoki bo'luvchidan qidirib, formulani tuzatish va natijani qayta tekshirish bosqichlarini tartiblang.
\begin{enumerate}[label=\Alph*),leftmargin=9mm]
\item Tuzatilgan natijani qayta hisoblash
\item Havola yoki bo'luvchini tekshirish
\item Formula bilan xato chiqishini kuzatish
\item Xato kodini o'qish
\end{enumerate}
\begin{tcolorbox}[enhanced,breakable,colback=teal!4,colframe=teal!55!black,boxrule=0.45pt,arc=1.2mm,left=2.8mm,right=2.8mm,top=2mm,bottom=2mm]
\textbf{\color{teal}To'g'ri javob:} \textbf{C - D - B - A}
\par\smallskip
\textbf{Izoh:} Bosqichlar: Formula bilan xato chiqishini kuzatish \ensuremath{\rightarrow} Xato kodini o'qish \ensuremath{\rightarrow} Havola yoki bo'luvchini tekshirish \ensuremath{\rightarrow} Tuzatilgan natijani qayta hisoblash.
\end{tcolorbox}
\vspace{2mm}\hrule height .2pt\vspace{2mm}

% DB 1d0b3fcc-7ee6-fac0-425d-765f43641cb0 status=published
\needspace{6\baselineskip}
\textbf{\color{ink}214.} {ootnotesizecolor{muted}Y1 | published | 1d0b3fcc}
\par
A1=4, A2=2. =A1/A2*3 natija?
\begin{enumerate}[label=\Alph*),leftmargin=9mm]
\item 2/3
\item 12
\item 1
\item 6
\end{enumerate}
\begin{tcolorbox}[enhanced,breakable,colback=teal!4,colframe=teal!55!black,boxrule=0.45pt,arc=1.2mm,left=2.8mm,right=2.8mm,top=2mm,bottom=2mm]
\textbf{\color{teal}To'g'ri javob:} \textbf{D) 6}
\par\smallskip
\textbf{Izoh:} Bo'lish va ko'paytirish chapdan o'ngga: 4/2*3=6.
\end{tcolorbox}
\vspace{2mm}\hrule height .2pt\vspace{2mm}

% DB 24deb8ad-4d70-eacb-3b72-dabe277617b9 status=published
\needspace{6\baselineskip}
\textbf{\color{ink}215.} {ootnotesizecolor{muted}Y1 | published | 24deb8ad}
\par
D4 katakda \texttt{=\$B\$2+A3*C\$1} bor. Uni F7 katakka nusxalash kerak. F7 formulasini toping.
\begin{enumerate}[label=\Alph*),leftmargin=9mm]
\item \texttt{=\$B\$2+A3*C\$1}
\item \texttt{=\$B\$2+C3*E\$1}
\item \texttt{=\$B\$2+C6*E\$1}
\item \texttt{=\$D\$5+C6*E\$4}
\end{enumerate}
\begin{tcolorbox}[enhanced,breakable,colback=teal!4,colframe=teal!55!black,boxrule=0.45pt,arc=1.2mm,left=2.8mm,right=2.8mm,top=2mm,bottom=2mm]
\textbf{\color{teal}To'g'ri javob:} \textbf{C) \texttt{=\$B\$2+C6*E\$1}}
\par\smallskip
\textbf{Izoh:} D4\ensuremath{\rightarrow}F7: ustun +2, satr +3; A3\ensuremath{\rightarrow}C6; C\$1\ensuremath{\rightarrow}E\$1; \$B\$2 o'zgarmaydi.
\end{tcolorbox}
\vspace{2mm}\hrule height .2pt\vspace{2mm}

% DB 3545a66c-92be-44c4-82d5-9f865b46ec81 status=published
\needspace{6\baselineskip}
\textbf{\color{ink}216.} {ootnotesizecolor{muted}Y3 | published | 3545a66c}
\par
Excel jadvalidagi narxlarga bitta katakda saqlangan soliq stavkasi qo'llanadi. Stavkani mutlaq havola orqali barcha satrlarga qo'llash ishlarini tartiblang.
\begin{enumerate}[label=\Alph*),leftmargin=9mm]
\item Stavka katagini \$ bilan mutlaqlash
\item Formulani boshqa satrlarga nusxalash
\item Soliq stavkasini alohida katakka yozish
\item Narxlar diapazonini kiritish
\end{enumerate}
\begin{tcolorbox}[enhanced,breakable,colback=teal!4,colframe=teal!55!black,boxrule=0.45pt,arc=1.2mm,left=2.8mm,right=2.8mm,top=2mm,bottom=2mm]
\textbf{\color{teal}To'g'ri javob:} \textbf{D - C - A - B}
\par\smallskip
\textbf{Izoh:} Bosqichlar: Narxlar diapazonini kiritish \ensuremath{\rightarrow} Soliq stavkasini alohida katakka yozish \ensuremath{\rightarrow} Stavka katagini \$ bilan mutlaqlash \ensuremath{\rightarrow} Formulani boshqa satrlarga nusxalash.
\end{tcolorbox}
\vspace{2mm}\hrule height .2pt\vspace{2mm}

% DB 57a05572-d9b0-5451-ae1b-6c13fd475721 status=published
\needspace{6\baselineskip}
\textbf{\color{ink}217.} {ootnotesizecolor{muted}Y2 | published | 57a05572}
\par
Excel xatolik kodlarini ularning aniq sabab yoki misoli bilan moslang.
\textbf{Chap ustun:}\begin{enumerate}[label=\arabic*),leftmargin=9mm]
\item \#VALUE!
\item \#DIV/0!
\item \#NAME?
\item \#NUM!
\item \#REF!
\item \#NULL!
\end{enumerate}
\textbf{O'ng ustun:}\begin{enumerate}[label=\Alph*),leftmargin=9mm]
\item Hisoblashda mos kelmaydigan ma'lumot turi ishlatilgan, masalan matnni songa qo'shish
\item Formula nolga bo'lishga uringan
\item Formula noma'lum yoki xato yozilgan nomni ishlatgan, masalan \texttt{SU()}
\item Sonli argument yaroqsiz, masalan haqiqiy sonlar sohasida \texttt{SQRT(-9)}
\item Formula o'chirilgan katak yoki ustunga yaroqsiz havola qilgan
\item Kesishmaydigan ikki oraliq orasida kesishish operatori - bo'sh joy - qo'llangan
\end{enumerate}
\begin{tcolorbox}[enhanced,breakable,colback=teal!4,colframe=teal!55!black,boxrule=0.45pt,arc=1.2mm,left=2.8mm,right=2.8mm,top=2mm,bottom=2mm]
\textbf{\color{teal}To'g'ri javob:} \textbf{1-A, 2-B, 3-C, 4-D, 5-E, 6-F}
\par\smallskip
\textbf{Izoh:} \#VALUE! - tur mos emas; \#DIV/0! - nolga bo'lish; \#NAME? - noma'lum nom; \#NUM! - yaroqsiz sonli argument; \#REF! - yaroqsiz havola; \#NULL! - mavjud bo'lmagan oraliqlar kesishmasi.
\end{tcolorbox}
\vspace{2mm}\hrule height .2pt\vspace{2mm}

% DB 669fc812-2d94-68f7-a302-0982c10a4c9f status=published
\needspace{6\baselineskip}
\textbf{\color{ink}218.} {ootnotesizecolor{muted}Y1 | published | 669fc812}
\par
A1=18, B1=0. =A1/B1 nima beradi?
\begin{enumerate}[label=\Alph*),leftmargin=9mm]
\item 0
\item 18
\item \#REF!
\item \#DIV/0!
\end{enumerate}
\begin{tcolorbox}[enhanced,breakable,colback=teal!4,colframe=teal!55!black,boxrule=0.45pt,arc=1.2mm,left=2.8mm,right=2.8mm,top=2mm,bottom=2mm]
\textbf{\color{teal}To'g'ri javob:} \textbf{D) \#DIV/0!}
\par\smallskip
\textbf{Izoh:} Nolga bo'lib bo'lmaydi.
\end{tcolorbox}
\vspace{2mm}\hrule height .2pt\vspace{2mm}

% DB ab4743c6-233e-369e-5db7-2fde96117caf status=published
\needspace{6\baselineskip}
\textbf{\color{ink}219.} {ootnotesizecolor{muted}Y1 | published | ab4743c6}
\par
Excelda E4 katakka \texttt{=C\$2*\$D3+\$B\$1} yozilgan. F6 ga nusxalansa, yangi formula qaysi?
\begin{enumerate}[label=\Alph*),leftmargin=9mm]
\item \texttt{=D\$2*\$D5+\$C\$3}
\item \texttt{=D\$2*\$D5+\$B\$1}
\item \texttt{=D\$4*\$E5+\$B\$1}
\item \texttt{=C\$2*\$D3+\$B\$1}
\end{enumerate}
\begin{tcolorbox}[enhanced,breakable,colback=teal!4,colframe=teal!55!black,boxrule=0.45pt,arc=1.2mm,left=2.8mm,right=2.8mm,top=2mm,bottom=2mm]
\textbf{\color{teal}To'g'ri javob:} \textbf{B) \texttt{=D\$2*\$D5+\$B\$1}}
\par\smallskip
\textbf{Izoh:} E4\ensuremath{\rightarrow}F6: ustun +1, satr +2; C\$2\ensuremath{\rightarrow}D\$2; \$D3\ensuremath{\rightarrow}\$D5; \$B\$1 mutlaq qoladi.
\end{tcolorbox}
\vspace{2mm}\hrule height .2pt\vspace{2mm}

% DB acd9b6ac-556f-3f93-320a-dde3d09d0d58 status=published
\needspace{6\baselineskip}
\textbf{\color{ink}220.} {ootnotesizecolor{muted}Y1 | published | acd9b6ac}
\par
A1=6 va A2=3. =A1/A2+A2\textasciicircum{}2 natija?
\begin{enumerate}[label=\Alph*),leftmargin=9mm]
\item 9
\item 5
\item 11
\item 15
\end{enumerate}
\begin{tcolorbox}[enhanced,breakable,colback=teal!4,colframe=teal!55!black,boxrule=0.45pt,arc=1.2mm,left=2.8mm,right=2.8mm,top=2mm,bottom=2mm]
\textbf{\color{teal}To'g'ri javob:} \textbf{C) 11}
\par\smallskip
\textbf{Izoh:} 6/3+9=11.
\end{tcolorbox}
\vspace{2mm}\hrule height .2pt\vspace{2mm}

% DB ad1d4868-2293-17a3-12e8-5787094006ca status=published
\needspace{6\baselineskip}
\textbf{\color{ink}221.} {ootnotesizecolor{muted}Y1 | published | ad1d4868}
\par
E3 formulasi =C\$2+\$D3. Uni H5 ga ko'chirish natijasi?
\begin{enumerate}[label=\Alph*),leftmargin=9mm]
\item =F\$2+\$D5
\item =H\$2+\$D8
\item =C\$2+\$D3
\item =F\$4+\$G5
\end{enumerate}
\begin{tcolorbox}[enhanced,breakable,colback=teal!4,colframe=teal!55!black,boxrule=0.45pt,arc=1.2mm,left=2.8mm,right=2.8mm,top=2mm,bottom=2mm]
\textbf{\color{teal}To'g'ri javob:} \textbf{A) =F\$2+\$D5}
\par\smallskip
\textbf{Izoh:} Uch ustun va ikki satr siljishi; satr 2, ustun D mahkam.
\end{tcolorbox}
\vspace{2mm}\hrule height .2pt\vspace{2mm}

% DB b952d7e5-350c-8b9f-e2ae-065048c0eec0 status=published
\needspace{6\baselineskip}
\textbf{\color{ink}222.} {ootnotesizecolor{muted}Y2 | published | b952d7e5}
\par
Excel xatolik kodlarini sabab bilan moslang. Bu savolda xatoliklar alohida, o'zaro farqli holatlar sifatida berilgan.
\textbf{Chap ustun:}\begin{enumerate}[label=\arabic*),leftmargin=9mm]
\item \#DIV/0!
\item \#REF!
\item \#NAME?
\item \#VALUE!
\end{enumerate}
\textbf{O'ng ustun:}\begin{enumerate}[label=\Alph*),leftmargin=9mm]
\item Noto'g'ri turdagi matnni hisoblashga qo'shish
\item Formula nomi noto'g'ri yozilgan
\item Nolga bo'lish
\item O'chirilgan katakka yaroqsiz havola
\end{enumerate}
\begin{tcolorbox}[enhanced,breakable,colback=teal!4,colframe=teal!55!black,boxrule=0.45pt,arc=1.2mm,left=2.8mm,right=2.8mm,top=2mm,bottom=2mm]
\textbf{\color{teal}To'g'ri javob:} \textbf{1-C, 2-D, 3-B, 4-A}
\par\smallskip
\textbf{Izoh:} 1-C, 2-D, 3-B, 4-A. 13-apreldagi xatoliklar moslashtirishiga o'xshash; manbadagi izohlar tekshirilgan va aniq vaziyatlarga almashtirilgan.
\end{tcolorbox}
\vspace{2mm}\hrule height .2pt\vspace{2mm}

% DB 424438e9-5b9e-4326-707a-25ecdc904e12 status=published
\needspace{6\baselineskip}
\textbf{\color{ink}223.} {ootnotesizecolor{muted}Y1 | published | 424438e9}
\par
D4 katakdagi =B\$2*\$C3+\$D\$1 formulasi F6 ga ko'chirildi. Yangi formula qaysi?
\begin{enumerate}[label=\Alph*),leftmargin=9mm]
\item =D\$2*\$C5+\$D\$1
\item =B\$2*\$C3+\$D\$1
\item =D4*\$E5+\$D1
\item =D\$4*\$E5+\$F\$3
\end{enumerate}
\begin{tcolorbox}[enhanced,breakable,colback=teal!4,colframe=teal!55!black,boxrule=0.45pt,arc=1.2mm,left=2.8mm,right=2.8mm,top=2mm,bottom=2mm]
\textbf{\color{teal}To'g'ri javob:} \textbf{A) =D\$2*\$C5+\$D\$1}
\par\smallskip
\textbf{Izoh:} D4\ensuremath{\rightarrow}F6 siljish +2 ustun, +2 satr: B\$2\ensuremath{\rightarrow}D\$2, \$C3\ensuremath{\rightarrow}\$C5, \$D\$1 o'zgarmaydi.
\end{tcolorbox}
\vspace{2mm}\hrule height .2pt\vspace{2mm}

% DB 47875ffc-5890-44c4-aa10-5446a15deba1 status=published
\needspace{6\baselineskip}
\textbf{\color{ink}224.} {ootnotesizecolor{muted}Y3 | published | 47875ffc}
\par
E4 katakka hali formula yozilmagan, G6 katak ham bo'sh. Quyidagi topshiriqda formulani ko'chirib, natijasini tekshirish uchun bajariladigan bog'liq amallarni tartiblang.
\begin{enumerate}[label=\Alph*),leftmargin=9mm]
\item G6 katakdagi hosil bo'lgan havolalarni ko'chirish qoidalariga ko'ra tekshirish
\item E4 katakka \texttt{=\$B\$2+A3*C\$1} formulasini kiritish
\item E4 dagi formulani G6 katakka nusxalash
\item Tekshirilgan G6 formulasining qiymatini manba kataklar bilan solishtirish
\end{enumerate}
\begin{tcolorbox}[enhanced,breakable,colback=teal!4,colframe=teal!55!black,boxrule=0.45pt,arc=1.2mm,left=2.8mm,right=2.8mm,top=2mm,bottom=2mm]
\textbf{\color{teal}To'g'ri javob:} \textbf{B - C - A - D}
\par\smallskip
\textbf{Izoh:} Avval E4 da formula bo'lishi kerak; keyin u G6 ga nusxalanadi. E4\ensuremath{\rightarrow}G6 siljish (+2 ustun, +2 satr): \texttt{=\$B\$2+C5*E\$1}. Shundan keyin qiymat manba kataklardagi sonlar asosida tekshiriladi. Havolalarni tekshirish hisoblangan natijaga ishonishdan oldingi nazorat bosqichidir.
\end{tcolorbox}
\vspace{2mm}\hrule height .2pt\vspace{2mm}

% DB 512def30-e649-e58a-78d3-2cd370338a5f status=published
\needspace{6\baselineskip}
\textbf{\color{ink}225.} {ootnotesizecolor{muted}Y1 | published | 512def30}
\par
B2=80, C2="Ha". Formula =IF(AND(B2\textgreater{}=60,C2="Ha"),B2*0.1,0) hisoblanganda qanday natija chiqadi?
\begin{enumerate}[label=\Alph*),leftmargin=9mm]
\item 80
\item 8
\item 6
\item 0
\end{enumerate}
\begin{tcolorbox}[enhanced,breakable,colback=teal!4,colframe=teal!55!black,boxrule=0.45pt,arc=1.2mm,left=2.8mm,right=2.8mm,top=2mm,bottom=2mm]
\textbf{\color{teal}To'g'ri javob:} \textbf{B) 8}
\par\smallskip
\textbf{Izoh:} AND sharti rost (80\textgreater{}=60 va Ha), demak 80\ensuremath{\times}0.1=8.
\end{tcolorbox}
\vspace{2mm}\hrule height .2pt\vspace{2mm}

% DB 999205dc-7c28-4a09-af38-e1451d4c07bf status=published
\needspace{6\baselineskip}
\textbf{\color{ink}226.} {ootnotesizecolor{muted}Y3 | published | 999205dc}
\par
MS Excelda C3 katakdagi =A1+\$B\$2 formulani D5 katakka nusxalash kerak. Nisbiy havola A1 dan B3 ga siljishi, mutlaq \$B\$2 esa saqlanishi lozim. Ko'chirish va yakuniy formulani tekshirish amallarini tartiblang.
\begin{enumerate}[label=\Alph*),leftmargin=9mm]
\item Yakuniy nusxa formulasidagi mutlaq koordinatalar o'zgarmaganini tekshirish
\item Formula joylashgan boshlang'ich katakni topish
\item Ko'chirishning satr/ustun siljishini hisoblash
\item Nisbiy koordinatalarni siljitib yangi formulani yozish
\end{enumerate}
\begin{tcolorbox}[enhanced,breakable,colback=teal!4,colframe=teal!55!black,boxrule=0.45pt,arc=1.2mm,left=2.8mm,right=2.8mm,top=2mm,bottom=2mm]
\textbf{\color{teal}To'g'ri javob:} \textbf{B - C - D - A}
\par\smallskip
\textbf{Izoh:} Avval asl formula katagi C3 aniqlanadi, D5 ga siljish (+1 ustun, +2 satr) hisoblanadi. Nisbiy havola B3 ga siljitilib =B3+\$B\$2 yoziladi, so'ng mutlaq \$B\$2 o'zgarmagani tekshiriladi. Tartib: 2 \ensuremath{\rightarrow} 3 \ensuremath{\rightarrow} 4 \ensuremath{\rightarrow} 1.
\end{tcolorbox}
\vspace{2mm}\hrule height .2pt\vspace{2mm}

% DB a5108c98-663e-4489-a351-dafb2ba8ea8f status=published
\needspace{6\baselineskip}
\textbf{\color{ink}227.} {ootnotesizecolor{muted}Y3 | published | a5108c98}
\par
MS Excel jadvalidagi birinchi satr formulasini Fill Handle bilan pastga ko'chirib, barcha natijalarni tekshirish uchun amallarni tartiblang.
\begin{enumerate}[label=\Alph*),leftmargin=9mm]
\item Bir necha natijani qo'lda nazorat qilish
\item Hisoblash uchun manzillarni belgilash
\item Birinchi satrda formulani sinash
\item Fill Handle bilan pastga ko'chirish
\end{enumerate}
\begin{tcolorbox}[enhanced,breakable,colback=teal!4,colframe=teal!55!black,boxrule=0.45pt,arc=1.2mm,left=2.8mm,right=2.8mm,top=2mm,bottom=2mm]
\textbf{\color{teal}To'g'ri javob:} \textbf{B - C - D - A}
\par\smallskip
\textbf{Izoh:} Bosqichlar: Hisoblash uchun manzillarni belgilash \ensuremath{\rightarrow} Birinchi satrda formulani sinash \ensuremath{\rightarrow} Fill Handle bilan pastga ko'chirish \ensuremath{\rightarrow} Bir necha natijani qo'lda nazorat qilish.
\end{tcolorbox}
\vspace{2mm}\hrule height .2pt\vspace{2mm}

% DB b9e6fee1-9ca3-fbc0-c604-e12ccd134ed8 status=published
\needspace{6\baselineskip}
\textbf{\color{ink}228.} {ootnotesizecolor{muted}Y1 | published | b9e6fee1}
\par
Excelda D5 katakka \texttt{=\$B\$1+\$A2*B2} yozilgan. Uni C6 ga ko'chirganda formuladagi havolalar qanday bo'ladi va qiymat nechaga teng? A1:C3 da A3=4, B1=0, B2=7, A2=-2, C3=8.
\begin{enumerate}[label=\Alph*),leftmargin=9mm]
\item \texttt{=\$B\$1+\$A3*A3}, natija 16
\item \texttt{=\$B\$1+\$A2*B2}, natija -14
\item \texttt{=\$A\$1+\$A3*A3}, natija 16
\item \texttt{=\$B\$2+\$A3*A3}, natija 23
\end{enumerate}
\begin{tcolorbox}[enhanced,breakable,colback=teal!4,colframe=teal!55!black,boxrule=0.45pt,arc=1.2mm,left=2.8mm,right=2.8mm,top=2mm,bottom=2mm]
\textbf{\color{teal}To'g'ri javob:} \textbf{A) \texttt{=\$B\$1+\$A3*A3}, natija 16}
\par\smallskip
\textbf{Izoh:} D5\ensuremath{\rightarrow}C6 da ustun siljishi -1, satr +1. \$B\$1 o'zgarmaydi, \$A2\ensuremath{\rightarrow}\$A3, B2\ensuremath{\rightarrow}A3. 0+4\ensuremath{\times}4=16.
\end{tcolorbox}
\vspace{2mm}\hrule height .2pt\vspace{2mm}
\subsection{MS Excel: matematik, statistik, mantiqiy va matn funksiyalari}

% DB 0040e844-a908-4a03-be13-0fe06fa55522 status=published
\needspace{6\baselineskip}
\textbf{\color{ink}229.} {ootnotesizecolor{muted}Y3 | published | 0040e844}
\par
O'quvchilarning ballari yozilgan Excel diapazoni uchun AVERAGE funksiyasida o'rtacha qiymatni hisoblash va tekshirish amallarini tartiblang.
\begin{enumerate}[label=\Alph*),leftmargin=9mm]
\item Ballar diapazonini aniqlash
\item Diapazon argumentini kiritish
\item AVERAGE funksiyasini tanlash
\item Hosil bo'lgan o'rtacha qiymatni tekshirish
\end{enumerate}
\begin{tcolorbox}[enhanced,breakable,colback=teal!4,colframe=teal!55!black,boxrule=0.45pt,arc=1.2mm,left=2.8mm,right=2.8mm,top=2mm,bottom=2mm]
\textbf{\color{teal}To'g'ri javob:} \textbf{A - C - B - D}
\par\smallskip
\textbf{Izoh:} Bosqichlar: Ballar diapazonini aniqlash \ensuremath{\rightarrow} AVERAGE funksiyasini tanlash \ensuremath{\rightarrow} Diapazon argumentini kiritish \ensuremath{\rightarrow} Hosil bo'lgan o'rtacha qiymatni tekshirish.
\end{tcolorbox}
\vspace{2mm}\hrule height .2pt\vspace{2mm}

% DB 68549b73-71cf-3f2c-3070-5c02b0f6caa4 status=published
\needspace{6\baselineskip}
\textbf{\color{ink}230.} {ootnotesizecolor{muted}Y2 | published | 68549b73}
\par
Quyidagi buyruqlarni vazifasi bilan moslashtiring.  1) SUM;  2) AVERAGE;  3) MAX;  Izohlar: A) Yig'indi;  B) Eng katta qiymat;  C) O'rtacha arifmetik;
\textbf{Chap ustun:}\begin{enumerate}[label=\arabic*),leftmargin=9mm]
\item SUM
\item AVERAGE
\item MAX
\end{enumerate}
\textbf{O'ng ustun:}\begin{enumerate}[label=\Alph*),leftmargin=9mm]
\item Yig'indi
\item Eng katta qiymat
\item O'rtacha arifmetik
\end{enumerate}
\begin{tcolorbox}[enhanced,breakable,colback=teal!4,colframe=teal!55!black,boxrule=0.45pt,arc=1.2mm,left=2.8mm,right=2.8mm,top=2mm,bottom=2mm]
\textbf{\color{teal}To'g'ri javob:} \textbf{1-A, 2-C, 3-B}
\par\smallskip
\textbf{Izoh:} To'g'ri moslik: SUM - Yig'indi, AVERAGE - O'rtacha arifmetik, MAX - Eng katta qiymat.
\end{tcolorbox}
\vspace{2mm}\hrule height .2pt\vspace{2mm}

% DB 6beb2430-8ce7-f6c1-b076-b210394e08c6 status=published
\needspace{6\baselineskip}
\textbf{\color{ink}231.} {ootnotesizecolor{muted}Y1 | published | 6beb2430}
\par
MAX qanday qiymat qaytaradi?
\begin{enumerate}[label=\Alph*),leftmargin=9mm]
\item Eng kichik sonni
\item Eng katta sonli qiymatni
\item Barcha matnni
\item Kataklar sonini
\end{enumerate}
\begin{tcolorbox}[enhanced,breakable,colback=teal!4,colframe=teal!55!black,boxrule=0.45pt,arc=1.2mm,left=2.8mm,right=2.8mm,top=2mm,bottom=2mm]
\textbf{\color{teal}To'g'ri javob:} \textbf{B) Eng katta sonli qiymatni}
\par\smallskip
\textbf{Izoh:} MAX maksimal sonni topadi.
\end{tcolorbox}
\vspace{2mm}\hrule height .2pt\vspace{2mm}

% DB 7fba5aa6-734e-46c1-b70a-95e462d60b1a status=published
\needspace{6\baselineskip}
\textbf{\color{ink}232.} {ootnotesizecolor{muted}Y1 | published | 7fba5aa6}
\par
Excelda \texttt{=IF(A1\textgreater{}=60,"O'tdi","O'tmadi")} formulasi bahoga qarab ikkita matndan birini qaytaradi. IF funksiyasining asosiy vazifasi nima?
\begin{enumerate}[label=\Alph*),leftmargin=9mm]
\item Shart rost/yolg'onligiga qarab natija tanlash
\item Shartga mos keladigan kataklar sonini hisoblash
\item Shartga mos kelmaydigan qatorlarni vaqtincha yashirish
\item Shart qanday bo'lishidan qat'i nazar sonlarni qo'shish
\end{enumerate}
\begin{tcolorbox}[enhanced,breakable,colback=teal!4,colframe=teal!55!black,boxrule=0.45pt,arc=1.2mm,left=2.8mm,right=2.8mm,top=2mm,bottom=2mm]
\textbf{\color{teal}To'g'ri javob:} \textbf{A) Shart rost/yolg'onligiga qarab natija tanlash}
\par\smallskip
\textbf{Izoh:} IF mantiqiy shartni baholab, u rost yoki yolg'on bo'lishiga qarab berilgan ikki natijadan birini qaytaradi. COUNTIF sanash, Filter esa ko'rsatishni cheklash vazifasini bajaradi.
\end{tcolorbox}
\vspace{2mm}\hrule height .2pt\vspace{2mm}

% DB 90f497c4-d3c0-f863-4975-6556747d7fbd status=published
\needspace{6\baselineskip}
\textbf{\color{ink}233.} {ootnotesizecolor{muted}Y1 | published | 90f497c4}
\par
SUM funksiyasi nima qiladi?
\begin{enumerate}[label=\Alph*),leftmargin=9mm]
\item Yangi varaq yaratadi
\item Diagrammani o'chiradi
\item Matnni kesadi
\item Sonli qiymatlarni jamlaydi
\end{enumerate}
\begin{tcolorbox}[enhanced,breakable,colback=teal!4,colframe=teal!55!black,boxrule=0.45pt,arc=1.2mm,left=2.8mm,right=2.8mm,top=2mm,bottom=2mm]
\textbf{\color{teal}To'g'ri javob:} \textbf{D) Sonli qiymatlarni jamlaydi}
\par\smallskip
\textbf{Izoh:} SUM yig'indi funksiyasidir.
\end{tcolorbox}
\vspace{2mm}\hrule height .2pt\vspace{2mm}

% DB e8fb0fbf-c9ca-5a06-8e25-b7f4251ddc1f status=published
\needspace{6\baselineskip}
\textbf{\color{ink}234.} {ootnotesizecolor{muted}Y1 | published | e8fb0fbf}
\par
MS Excelda B3=10, C3=20 va D3=30. Qaysi formula natijasi 1 ga teng?
\begin{enumerate}[label=\Alph*),leftmargin=9mm]
\item \texttt{=SIGN(B3+C3+D3)}
\item \texttt{=GCD(B3,ABS(C3),D3)}
\item \texttt{=ROUND(B3,C3,D3)}
\item \texttt{=LCM(B3,C3,D3)}
\end{enumerate}
\begin{tcolorbox}[enhanced,breakable,colback=teal!4,colframe=teal!55!black,boxrule=0.45pt,arc=1.2mm,left=2.8mm,right=2.8mm,top=2mm,bottom=2mm]
\textbf{\color{teal}To'g'ri javob:} \textbf{A) \texttt{=SIGN(B3+C3+D3)}}
\par\smallskip
\textbf{Izoh:} B3+C3+D3=60 musbat son, shuning uchun SIGN(60)=1. GCD=10, LCM=60; ROUND esa bu ko'rinishda uchta argument bilan noto'g'ri yozilgan.
\end{tcolorbox}
\vspace{2mm}\hrule height .2pt\vspace{2mm}

% DB e92fc773-2c1a-4c0c-a194-19eac3dde5f4 status=published
\needspace{6\baselineskip}
\textbf{\color{ink}235.} {ootnotesizecolor{muted}Y1 | published | e92fc773}
\par
Excelda B2:B6 kataklariga 2, 4, 6, 8 va 10 yozildi. \texttt{=AVERAGE(B2:B6)} natijasi qanday statistik ko'rsatkichni ifodalaydi?
\begin{enumerate}[label=\Alph*),leftmargin=9mm]
\item Sonli qiymatlarning arifmetik o'rtachasi
\item Sonli qiymatlarning eng kattasi
\item Sonli qiymatlarning umumiy yig'indisi
\item Sonli qiymatlarning umumiy soni
\end{enumerate}
\begin{tcolorbox}[enhanced,breakable,colback=teal!4,colframe=teal!55!black,boxrule=0.45pt,arc=1.2mm,left=2.8mm,right=2.8mm,top=2mm,bottom=2mm]
\textbf{\color{teal}To'g'ri javob:} \textbf{A) Sonli qiymatlarning arifmetik o'rtachasi}
\par\smallskip
\textbf{Izoh:} AVERAGE sonli qiymatlarning arifmetik o'rtachasini hisoblaydi: (2+4+6+8+10)/5=6. MAX, SUM va COUNT boshqa natijalarni beradi.
\end{tcolorbox}
\vspace{2mm}\hrule height .2pt\vspace{2mm}

% DB 58f8bbbf-ed70-41d1-b4a7-77db5007a54a status=published
\needspace{6\baselineskip}
\textbf{\color{ink}236.} {ootnotesizecolor{muted}Y3 | published | 58f8bbbf}
\par
Excelda ballga ko'ra shart rost bo'lganda bitta, yolg'on bo'lganda boshqa natija chiqaruvchi IF formulasini tuzib, ikki xil qiymatda sinash amallarini tartiblang.
\begin{enumerate}[label=\Alph*),leftmargin=9mm]
\item IF funksiyasi shartini yozish
\item Ikki xil kirish qiymatida sinash
\item TRUE va FALSE natijalarini berish
\item Shart bajarilishi kerak bo'lgan katakni aniqlash
\end{enumerate}
\begin{tcolorbox}[enhanced,breakable,colback=teal!4,colframe=teal!55!black,boxrule=0.45pt,arc=1.2mm,left=2.8mm,right=2.8mm,top=2mm,bottom=2mm]
\textbf{\color{teal}To'g'ri javob:} \textbf{D - A - C - B}
\par\smallskip
\textbf{Izoh:} Bosqichlar: Shart bajarilishi kerak bo'lgan katakni aniqlash \ensuremath{\rightarrow} IF funksiyasi shartini yozish \ensuremath{\rightarrow} TRUE va FALSE natijalarini berish \ensuremath{\rightarrow} Ikki xil kirish qiymatida sinash.
\end{tcolorbox}
\vspace{2mm}\hrule height .2pt\vspace{2mm}

% DB 718928fc-a98c-75e3-4f8c-4029291b9b17 status=published
\needspace{6\baselineskip}
\textbf{\color{ink}237.} {ootnotesizecolor{muted}Y2 | published | 718928fc}
\par
Atamalar va ularning aniq ma'nolarini moslashtiring.  1) COUNT;  2) COUNTA;  3) COUNTBLANK;  Izohlar: A) Bo'sh kataklar soni;  B) Bo'sh bo'lmaganlar soni;  C) Sonli kataklar soni;
\textbf{Chap ustun:}\begin{enumerate}[label=\arabic*),leftmargin=9mm]
\item COUNT
\item COUNTA
\item COUNTBLANK
\end{enumerate}
\textbf{O'ng ustun:}\begin{enumerate}[label=\Alph*),leftmargin=9mm]
\item Bo'sh kataklar soni
\item Bo'sh bo'lmaganlar soni
\item Sonli kataklar soni
\end{enumerate}
\begin{tcolorbox}[enhanced,breakable,colback=teal!4,colframe=teal!55!black,boxrule=0.45pt,arc=1.2mm,left=2.8mm,right=2.8mm,top=2mm,bottom=2mm]
\textbf{\color{teal}To'g'ri javob:} \textbf{1-C, 2-B, 3-A}
\par\smallskip
\textbf{Izoh:} To'g'ri moslik: COUNT - Sonli kataklar soni, COUNTA - Bo'sh bo'lmaganlar soni, COUNTBLANK - Bo'sh kataklar soni.
\end{tcolorbox}
\vspace{2mm}\hrule height .2pt\vspace{2mm}

% DB 89ec6c7c-f222-5f60-baa9-1bbbdfacff4b status=published
\needspace{6\baselineskip}
\textbf{\color{ink}238.} {ootnotesizecolor{muted}Y1 | published | 89ec6c7c}
\par
MS Excel 2021 da \texttt{=ROUNDDOWN(POWER(343;1/3)/ROUND(2,459;1);1)} formulasi qanday natija qaytaradi?
\begin{enumerate}[label=\Alph*),leftmargin=9mm]
\item 2,8
\item 3
\item 2
\item 1
\end{enumerate}
\begin{tcolorbox}[enhanced,breakable,colback=teal!4,colframe=teal!55!black,boxrule=0.45pt,arc=1.2mm,left=2.8mm,right=2.8mm,top=2mm,bottom=2mm]
\textbf{\color{teal}To'g'ri javob:} \textbf{A) 2,8}
\par\smallskip
\textbf{Izoh:} POWER(343;1/3)=7; ROUND(2,459;1)=2,5; 7/2,5=2,8. ROUNDDOWN(...;1) natijani bir kasr xonasigacha pastga yaxlitlab, 2,8 ni beradi.
\end{tcolorbox}
\vspace{2mm}\hrule height .2pt\vspace{2mm}

% DB 939fc993-7f18-1dd8-17ca-35882f0f5692 status=published
\needspace{6\baselineskip}
\textbf{\color{ink}239.} {ootnotesizecolor{muted}Y1 | published | 939fc993}
\par
SUM va COUNTIF farqi?
\begin{enumerate}[label=\Alph*),leftmargin=9mm]
\item SUM bo'sh katakni topadi
\item COUNTIF hamma sonni qo'shadi
\item Ikkalasi faqat matn formatlaydi
\item SUM qiymatlarni qo'shadi, COUNTIF shartga moslarni sanaydi
\end{enumerate}
\begin{tcolorbox}[enhanced,breakable,colback=teal!4,colframe=teal!55!black,boxrule=0.45pt,arc=1.2mm,left=2.8mm,right=2.8mm,top=2mm,bottom=2mm]
\textbf{\color{teal}To'g'ri javob:} \textbf{D) SUM qiymatlarni qo'shadi, COUNTIF shartga moslarni sanaydi}
\par\smallskip
\textbf{Izoh:} COUNTIF natija sanog'idir.
\end{tcolorbox}
\vspace{2mm}\hrule height .2pt\vspace{2mm}

% DB 9c69b073-3459-2c5b-5b8e-1f94d2a426f0 status=published
\needspace{6\baselineskip}
\textbf{\color{ink}240.} {ootnotesizecolor{muted}Y1 | published | 9c69b073}
\par
AND va OR farqi?
\begin{enumerate}[label=\Alph*),leftmargin=9mm]
\item Ikkalasi faqat satrni ko'chiradi
\item AND kamida bitta, OR barcha shart rost bo'lsa TRUE
\item AND barcha, OR kamida bitta shart rost bo'lsa TRUE
\item Ikkalasi faqat qo'shadi
\end{enumerate}
\begin{tcolorbox}[enhanced,breakable,colback=teal!4,colframe=teal!55!black,boxrule=0.45pt,arc=1.2mm,left=2.8mm,right=2.8mm,top=2mm,bottom=2mm]
\textbf{\color{teal}To'g'ri javob:} \textbf{C) AND barcha, OR kamida bitta shart rost bo'lsa TRUE}
\par\smallskip
\textbf{Izoh:} Mantiqiy birlashtirish mezonlari boshqacha.
\end{tcolorbox}
\vspace{2mm}\hrule height .2pt\vspace{2mm}

% DB c0ae14d1-07fc-445f-b884-b723f8abbcc1 status=published
\needspace{6\baselineskip}
\textbf{\color{ink}241.} {ootnotesizecolor{muted}Y1 | published | c0ae14d1}
\par
Excelda bitta oraliqda sonli kataklar, matnli kataklar va bo'sh kataklar bor. COUNT va COUNTA funksiyalarining asosiy farqi qaysi javobda to'g'ri ko'rsatilgan?
\begin{enumerate}[label=\Alph*),leftmargin=9mm]
\item COUNT matn kiritilgan kataklarni, COUNTA esa faqat sonlarni sanaydi
\item COUNT sonlarni, COUNTA bo'sh bo'lmagan kataklarni sanaydi
\item COUNT barcha bo'sh kataklarni, COUNTA esa formulalarni sanaydi
\item COUNT barcha qiymatlarni qo'shadi, COUNTA sonlarning o'rtachasini hisoblaydi
\end{enumerate}
\begin{tcolorbox}[enhanced,breakable,colback=teal!4,colframe=teal!55!black,boxrule=0.45pt,arc=1.2mm,left=2.8mm,right=2.8mm,top=2mm,bottom=2mm]
\textbf{\color{teal}To'g'ri javob:} \textbf{B) COUNT sonlarni, COUNTA bo'sh bo'lmagan kataklarni sanaydi}
\par\smallskip
\textbf{Izoh:} COUNT oraliqdagi son sifatida saqlangan qiymatlarni hisoblaydi; COUNTA esa bo'sh bo'lmagan kataklarni hisoblaydi. Bu funksiyalar yig'indi yoki o'rtachani topmaydi.
\end{tcolorbox}
\vspace{2mm}\hrule height .2pt\vspace{2mm}

% DB dd332a53-f833-5df0-953d-dc8e9511c92f status=published
\needspace{6\baselineskip}
\textbf{\color{ink}242.} {ootnotesizecolor{muted}Y1 | published | dd332a53}
\par
Excelda \texttt{SUM(A1:B1)=12}, \texttt{SUM(D1:E1)=7} va \texttt{AVERAGE(A1:E1)=5}. C1 katak qiymatini toping.
\begin{enumerate}[label=\Alph*),leftmargin=9mm]
\item 0
\item 4
\item 6
\item 5
\end{enumerate}
\begin{tcolorbox}[enhanced,breakable,colback=teal!4,colframe=teal!55!black,boxrule=0.45pt,arc=1.2mm,left=2.8mm,right=2.8mm,top=2mm,bottom=2mm]
\textbf{\color{teal}To'g'ri javob:} \textbf{C) 6}
\par\smallskip
\textbf{Izoh:} Besh katakning jami \texttt{5\ensuremath{\cdot}5=25}. Ma'lum kataklar yig'indisi \texttt{12+7=19}; demak, \texttt{C1=25-19=6}.
\end{tcolorbox}
\vspace{2mm}\hrule height .2pt\vspace{2mm}

% DB e1832efc-c109-4033-bffd-260bad97cfb7 status=published
\needspace{6\baselineskip}
\textbf{\color{ink}243.} {ootnotesizecolor{muted}Y1 | published | e1832efc}
\par
Excel jadvalida qiymatlar berilgan mezonga mos kelganda nechta yozuv borligini va shu yozuvlardagi sonlarning yig'indisini alohida hisoblash talab qilinadi. COUNTIF va SUMIF farqi qaysi?
\begin{enumerate}[label=\Alph*),leftmargin=9mm]
\item COUNTIF shartga mos sonlarni qo'shadi, SUMIF esa yozuvlarni sanaydi
\item COUNTIF qiymatlarni saralaydi, SUMIF esa mos qatorlarni vaqtincha yashiradi
\item COUNTIF sanaydi, SUMIF shartga mos qiymatlarni yig'adi
\item COUNTIF mos qiymatlarning o'rtachasini, SUMIF esa eng kattasini topadi
\end{enumerate}
\begin{tcolorbox}[enhanced,breakable,colback=teal!4,colframe=teal!55!black,boxrule=0.45pt,arc=1.2mm,left=2.8mm,right=2.8mm,top=2mm,bottom=2mm]
\textbf{\color{teal}To'g'ri javob:} \textbf{C) COUNTIF sanaydi, SUMIF shartga mos qiymatlarni yig'adi}
\par\smallskip
\textbf{Izoh:} COUNTIF mezonga mos yozuvlar sonini hisoblaydi, SUMIF esa mezonga mos qiymatlarni yig'adi. Boshqa variantlar ularning vazifalarini almashtirgan yoki AVERAGEIF/MAXIFS kabi hisoblarni chalkashtirgan.
\end{tcolorbox}
\vspace{2mm}\hrule height .2pt\vspace{2mm}

% DB 0c759b3e-36cd-5647-5294-0ab0c8140d9e status=published
\needspace{6\baselineskip}
\textbf{\color{ink}244.} {ootnotesizecolor{muted}Y2 | published | 0c759b3e}
\par
Amallarni izohi bilan moslang.  1) IF;  2) AND;  3) OR;  Izohlar: A) Kamida bitta shart bajarilishi;  B) Shartga qarab natija tanlash;  C) Barcha shartlar bajarilishi;
\textbf{Chap ustun:}\begin{enumerate}[label=\arabic*),leftmargin=9mm]
\item IF
\item AND
\item OR
\end{enumerate}
\textbf{O'ng ustun:}\begin{enumerate}[label=\Alph*),leftmargin=9mm]
\item Kamida bitta shart bajarilishi
\item Shartga qarab natija tanlash
\item Barcha shartlar bajarilishi
\end{enumerate}
\begin{tcolorbox}[enhanced,breakable,colback=teal!4,colframe=teal!55!black,boxrule=0.45pt,arc=1.2mm,left=2.8mm,right=2.8mm,top=2mm,bottom=2mm]
\textbf{\color{teal}To'g'ri javob:} \textbf{1-B, 2-C, 3-A}
\par\smallskip
\textbf{Izoh:} To'g'ri moslik: IF - Shartga qarab natija tanlash, AND - Barcha shartlar bajarilishi, OR - Kamida bitta shart bajarilishi.
\end{tcolorbox}
\vspace{2mm}\hrule height .2pt\vspace{2mm}

% DB 1e2bce66-8bcc-456a-82c6-0d98bbd65a78 status=published
\needspace{6\baselineskip}
\textbf{\color{ink}245.} {ootnotesizecolor{muted}Y1 | published | 1e2bce66}
\par
MS Excelda \texttt{=SUMM(A1:C3)} natijasi 15, \texttt{=MAKS(B1:C3)} natijasi 8. Jadval: \newline  \newline |   | A | B | C | \newline |---|---:|---:|---:| \newline |1|?|-1|0| \newline |2|-1|?|2| \newline |3|6|-2|1| \newline  \newline \texttt{=PROIZVED(A1;B2)} natijasini toping.
\begin{enumerate}[label=\Alph*),leftmargin=9mm]
\item 16
\item -16
\item 12
\item 14
\end{enumerate}
\begin{tcolorbox}[enhanced,breakable,colback=teal!4,colframe=teal!55!black,boxrule=0.45pt,arc=1.2mm,left=2.8mm,right=2.8mm,top=2mm,bottom=2mm]
\textbf{\color{teal}To'g'ri javob:} \textbf{A) 16}
\par\smallskip
\textbf{Izoh:} Ma'lum kataklar yig'indisi 5, demak A1+B2=10. B1:C3 diapazonidagi maksimum 8 bo'lishi uchun B2=8; shunda A1=2 va A1\ensuremath{\cdot}B2=16. [Yutuqli\_test, test kodi 12; Telegram xabarlar 6446 va 6556.]
\end{tcolorbox}
\vspace{2mm}\hrule height .2pt\vspace{2mm}

% DB 33906778-191b-87d3-a582-a50611d068a5 status=published
\needspace{6\baselineskip}
\textbf{\color{ink}246.} {ootnotesizecolor{muted}Y1 | published | 33906778}
\par
A1:A4=2,4,6,8. MAX-MIN natija?
\begin{enumerate}[label=\Alph*),leftmargin=9mm]
\item 2
\item 6
\item 8
\item 4
\end{enumerate}
\begin{tcolorbox}[enhanced,breakable,colback=teal!4,colframe=teal!55!black,boxrule=0.45pt,arc=1.2mm,left=2.8mm,right=2.8mm,top=2mm,bottom=2mm]
\textbf{\color{teal}To'g'ri javob:} \textbf{B) 6}
\par\smallskip
\textbf{Izoh:} 8-2=6.
\end{tcolorbox}
\vspace{2mm}\hrule height .2pt\vspace{2mm}

% DB 3c0fe183-8128-5c97-9264-6e529ea0a410 status=published
\needspace{6\baselineskip}
\textbf{\color{ink}247.} {ootnotesizecolor{muted}Y1 | published | 3c0fe183}
\par
A1:A4=2,4,6,8. AVERAGE natija?
\begin{enumerate}[label=\Alph*),leftmargin=9mm]
\item 5
\item 20
\item 2
\item 8
\end{enumerate}
\begin{tcolorbox}[enhanced,breakable,colback=teal!4,colframe=teal!55!black,boxrule=0.45pt,arc=1.2mm,left=2.8mm,right=2.8mm,top=2mm,bottom=2mm]
\textbf{\color{teal}To'g'ri javob:} \textbf{A) 5}
\par\smallskip
\textbf{Izoh:} 20/4=5.
\end{tcolorbox}
\vspace{2mm}\hrule height .2pt\vspace{2mm}

% DB 45b5b2af-fbbf-5eb5-9b52-39d6a179710a status=published
\needspace{6\baselineskip}
\textbf{\color{ink}248.} {ootnotesizecolor{muted}Y1 | published | 45b5b2af}
\par
Excelda avval \texttt{X1=10, Y1=20, Z1=30}, keyin esa \texttt{X1=50, Y1=100, Z1=200} qiymatlari berildi. Qaysi formula ikkala holatda ham bir xil natija qaytaradi?
\begin{enumerate}[label=\Alph*),leftmargin=9mm]
\item =IF(X1+Y1\textgreater{}=Z1; "Yes"; "No")
\item =IF(Z1\textgreater{}Y1+X1; "Yes"; "No")
\item =IF(Y1-Z1\textless{}X1; "Yes"; "No")
\item =IF(Y1+X1=Z1; "Yes"; "No")
\end{enumerate}
\begin{tcolorbox}[enhanced,breakable,colback=teal!4,colframe=teal!55!black,boxrule=0.45pt,arc=1.2mm,left=2.8mm,right=2.8mm,top=2mm,bottom=2mm]
\textbf{\color{teal}To'g'ri javob:} \textbf{C) =IF(Y1-Z1\textless{}X1; "Yes"; "No")}
\par\smallskip
\textbf{Izoh:} C variantida dastlab \texttt{20-30\textless{}10}, keyin \texttt{100-200\textless{}50}; ikkala tengsizlik ham rost va formula ikki safar ham \texttt{Yes} qaytaradi. Qolgan formulalar natijasi o'zgaradi.
\end{tcolorbox}
\vspace{2mm}\hrule height .2pt\vspace{2mm}

% DB 60b773ad-13a8-4daf-05f1-608f9f4e8cf4 status=published
\needspace{6\baselineskip}
\textbf{\color{ink}249.} {ootnotesizecolor{muted}Y1 | published | 60b773ad}
\par
Excelda quyidagi baholar jadvali bor: A2:B4 = (40,2), (60,3), (80,5). D2=60. \texttt{=VLOOKUP(D2,A2:B4,2,FALSE)} qanday qiymat qaytaradi?
\begin{enumerate}[label=\Alph*),leftmargin=9mm]
\item 2
\item 5
\item \texttt{\#N/A}
\item 3
\end{enumerate}
\begin{tcolorbox}[enhanced,breakable,colback=teal!4,colframe=teal!55!black,boxrule=0.45pt,arc=1.2mm,left=2.8mm,right=2.8mm,top=2mm,bottom=2mm]
\textbf{\color{teal}To'g'ri javob:} \textbf{D) 3}
\par\smallskip
\textbf{Izoh:} FALSE aynan 60 ni qidiradi va shu satrning 2-ustunidagi 3 qiymatini oladi.
\end{tcolorbox}
\vspace{2mm}\hrule height .2pt\vspace{2mm}

% DB 6771fe57-b559-9b5f-4bb3-96bd1a972400 status=published
\needspace{6\baselineskip}
\textbf{\color{ink}250.} {ootnotesizecolor{muted}Y1 | published | 6771fe57}
\par
A1=7, B1=-9. Excelda \texttt{=IF(AND(A1\textgreater{}0,B1\textless{}0),ABS(B1)+A1,0)} qiymat nechaga teng?
\begin{enumerate}[label=\Alph*),leftmargin=9mm]
\item 0
\item 16
\item 2
\item 9
\end{enumerate}
\begin{tcolorbox}[enhanced,breakable,colback=teal!4,colframe=teal!55!black,boxrule=0.45pt,arc=1.2mm,left=2.8mm,right=2.8mm,top=2mm,bottom=2mm]
\textbf{\color{teal}To'g'ri javob:} \textbf{B) 16}
\par\smallskip
\textbf{Izoh:} Ikkala shart rost, ABS(-9)+7=16.
\end{tcolorbox}
\vspace{2mm}\hrule height .2pt\vspace{2mm}

% DB 7b00bc02-40b2-442b-9cdd-4dee53e6c660 status=published
\needspace{6\baselineskip}
\textbf{\color{ink}251.} {ootnotesizecolor{muted}Y3 | published | 7b00bc02}
\par
Excel jadvalidagi belgilangan ustunda berilgan mezonga mos kataklar sonini COUNTIF bilan hisoblab, natijani tekshirish bosqichlarini tartiblang.
\begin{enumerate}[label=\Alph*),leftmargin=9mm]
\item Hisoblanadigan ustunni belgilash
\item Natijani qo'lda sanab tekshirish
\item Mezon argumentini berish
\item COUNTIF funksiyasi uchun diapazonni yozish
\end{enumerate}
\begin{tcolorbox}[enhanced,breakable,colback=teal!4,colframe=teal!55!black,boxrule=0.45pt,arc=1.2mm,left=2.8mm,right=2.8mm,top=2mm,bottom=2mm]
\textbf{\color{teal}To'g'ri javob:} \textbf{A - D - C - B}
\par\smallskip
\textbf{Izoh:} Bosqichlar: Hisoblanadigan ustunni belgilash \ensuremath{\rightarrow} COUNTIF funksiyasi uchun diapazonni yozish \ensuremath{\rightarrow} Mezon argumentini berish \ensuremath{\rightarrow} Natijani qo'lda sanab tekshirish.
\end{tcolorbox}
\vspace{2mm}\hrule height .2pt\vspace{2mm}

% DB 7ecae3e5-5e18-7689-8665-1e0d63a4f984 status=published
\needspace{6\baselineskip}
\textbf{\color{ink}252.} {ootnotesizecolor{muted}Y1 | published | 7ecae3e5}
\par
A1=40, B1=75. =OR(A1\textgreater{}=60,B1\textgreater{}=60) natija?
\begin{enumerate}[label=\Alph*),leftmargin=9mm]
\item TRUE
\item 60
\item 115
\item FALSE
\end{enumerate}
\begin{tcolorbox}[enhanced,breakable,colback=teal!4,colframe=teal!55!black,boxrule=0.45pt,arc=1.2mm,left=2.8mm,right=2.8mm,top=2mm,bottom=2mm]
\textbf{\color{teal}To'g'ri javob:} \textbf{A) TRUE}
\par\smallskip
\textbf{Izoh:} Kamida bitta shart rost.
\end{tcolorbox}
\vspace{2mm}\hrule height .2pt\vspace{2mm}

% DB 93ca9f5d-2d85-4181-8a53-76a2a5ca4635 status=published
\needspace{6\baselineskip}
\textbf{\color{ink}253.} {ootnotesizecolor{muted}Y1 | published | 93ca9f5d}
\par
MS Excelda A1=5 va B1=3. Quyidagi formula natijasini toping: \newline  \newline \texttt{=PSTR(2024;1;A1-B1)*LEVSIMV(2023;ZNAK(B1+A1))}
\begin{enumerate}[label=\Alph*),leftmargin=9mm]
\item 24
\item 20
\item 40
\item 2
\end{enumerate}
\begin{tcolorbox}[enhanced,breakable,colback=teal!4,colframe=teal!55!black,boxrule=0.45pt,arc=1.2mm,left=2.8mm,right=2.8mm,top=2mm,bottom=2mm]
\textbf{\color{teal}To'g'ri javob:} \textbf{C) 40}
\par\smallskip
\textbf{Izoh:} PSTR(2024;1;2)="20". ZNAK(8)=1 va LEVSIMV(2023;1)="2". Sonlarga aylantirilganda 20\ensuremath{\cdot}2=40. [Yutuqli\_test, test kodi 12; Telegram xabarlar 6446 va 6556.]
\end{tcolorbox}
\vspace{2mm}\hrule height .2pt\vspace{2mm}

% DB b14fcc1b-36f6-4f40-9dac-c94a7e0b68a6 status=published
\needspace{6\baselineskip}
\textbf{\color{ink}254.} {ootnotesizecolor{muted}Y1 | published | b14fcc1b}
\par
MS Excel. Agar A1 = -4, B1 = 5, B2 = 6 bo'lsa, quyidagi formula natijasi -30 ga teng bo'lishi uchun A2 katakka kiritilishi kerak bo'lgan qiymatni aniqlang. \newline  \newline \texttt{=ESLI(I(A1+B2\textless{}A2*B1;A1*A2\textless{}0);A1*B2-15+A2;A1*B1+5-A2)}
\begin{enumerate}[label=\Alph*),leftmargin=9mm]
\item 7
\item 5
\item 9
\item 8
\end{enumerate}
\begin{tcolorbox}[enhanced,breakable,colback=teal!4,colframe=teal!55!black,boxrule=0.45pt,arc=1.2mm,left=2.8mm,right=2.8mm,top=2mm,bottom=2mm]
\textbf{\color{teal}To'g'ri javob:} \textbf{C) 9}
\par\smallskip
\textbf{Izoh:} A2 musbat bo'lganda ikkala shart rost, natija -4\ensuremath{\times}6-15+A2=-39+A2. -30 uchun A2=9.
\end{tcolorbox}
\vspace{2mm}\hrule height .2pt\vspace{2mm}

% DB e725988a-b33f-4742-4bcd-d1ead089efc2 status=published
\needspace{6\baselineskip}
\textbf{\color{ink}255.} {ootnotesizecolor{muted}Y1 | published | e725988a}
\par
A1:A4=2,4,6,8. SUM natija?
\begin{enumerate}[label=\Alph*),leftmargin=9mm]
\item 4
\item 8
\item 5
\item 20
\end{enumerate}
\begin{tcolorbox}[enhanced,breakable,colback=teal!4,colframe=teal!55!black,boxrule=0.45pt,arc=1.2mm,left=2.8mm,right=2.8mm,top=2mm,bottom=2mm]
\textbf{\color{teal}To'g'ri javob:} \textbf{D) 20}
\par\smallskip
\textbf{Izoh:} Yig'indi 20.
\end{tcolorbox}
\vspace{2mm}\hrule height .2pt\vspace{2mm}

% DB 0ff26eab-cc2c-dc33-b73c-96243f19112b status=published
\needspace{6\baselineskip}
\textbf{\color{ink}256.} {ootnotesizecolor{muted}Y1 | published | 0ff26eab}
\par
A1:A5 kataklarda 12, bo'sh, "Baho", 0 va 8 yozilgan. \texttt{=COUNTA(A1:A5)-COUNT(A1:A5)} natija?
\begin{enumerate}[label=\Alph*),leftmargin=9mm]
\item 0
\item 1
\item 2
\item 3
\end{enumerate}
\begin{tcolorbox}[enhanced,breakable,colback=teal!4,colframe=teal!55!black,boxrule=0.45pt,arc=1.2mm,left=2.8mm,right=2.8mm,top=2mm,bottom=2mm]
\textbf{\color{teal}To'g'ri javob:} \textbf{B) 1}
\par\smallskip
\textbf{Izoh:} COUNTA=4 (12, matn, 0, 8), COUNT=3 (12, 0, 8); farq 1.
\end{tcolorbox}
\vspace{2mm}\hrule height .2pt\vspace{2mm}

% DB 1d8d7408-d9b0-c361-d6c6-25e245e1e4ea status=published
\needspace{6\baselineskip}
\textbf{\color{ink}257.} {ootnotesizecolor{muted}Y2 | published | 1d8d7408}
\par
Ushbu vaziyatdagi vositalar va vazifalarni moslashtiring.  1) LEFT;  2) MID;  3) RIGHT;  Izohlar: A) Matn oxiridan belgilar;  B) Matn o'rtasidan belgilar;  C) Matn boshidan belgilar;
\textbf{Chap ustun:}\begin{enumerate}[label=\arabic*),leftmargin=9mm]
\item LEFT
\item MID
\item RIGHT
\end{enumerate}
\textbf{O'ng ustun:}\begin{enumerate}[label=\Alph*),leftmargin=9mm]
\item Matn oxiridan belgilar
\item Matn o'rtasidan belgilar
\item Matn boshidan belgilar
\end{enumerate}
\begin{tcolorbox}[enhanced,breakable,colback=teal!4,colframe=teal!55!black,boxrule=0.45pt,arc=1.2mm,left=2.8mm,right=2.8mm,top=2mm,bottom=2mm]
\textbf{\color{teal}To'g'ri javob:} \textbf{1-C, 2-B, 3-A}
\par\smallskip
\textbf{Izoh:} To'g'ri moslik: LEFT - Matn boshidan belgilar, MID - Matn o'rtasidan belgilar, RIGHT - Matn oxiridan belgilar.
\end{tcolorbox}
\vspace{2mm}\hrule height .2pt\vspace{2mm}

% DB 40d21a98-90f1-0148-62d6-1cb58800f2f3 status=published
\needspace{6\baselineskip}
\textbf{\color{ink}258.} {ootnotesizecolor{muted}Y1 | published | 40d21a98}
\par
=RIGHT("Excel",3) natija?
\begin{enumerate}[label=\Alph*),leftmargin=9mm]
\item xel
\item el
\item cel
\item Exc
\end{enumerate}
\begin{tcolorbox}[enhanced,breakable,colback=teal!4,colframe=teal!55!black,boxrule=0.45pt,arc=1.2mm,left=2.8mm,right=2.8mm,top=2mm,bottom=2mm]
\textbf{\color{teal}To'g'ri javob:} \textbf{C) cel}
\par\smallskip
\textbf{Izoh:} Excelning oxirgi uch belgisi cel.
\end{tcolorbox}
\vspace{2mm}\hrule height .2pt\vspace{2mm}

% DB 5941f1b7-1718-1924-f23f-386fe94139ef status=published
\needspace{6\baselineskip}
\textbf{\color{ink}259.} {ootnotesizecolor{muted}Y1 | published | 5941f1b7}
\par
=LEFT("PowerPoint",5) natija?
\begin{enumerate}[label=\Alph*),leftmargin=9mm]
\item owe
\item Power
\item PowerP
\item Point
\end{enumerate}
\begin{tcolorbox}[enhanced,breakable,colback=teal!4,colframe=teal!55!black,boxrule=0.45pt,arc=1.2mm,left=2.8mm,right=2.8mm,top=2mm,bottom=2mm]
\textbf{\color{teal}To'g'ri javob:} \textbf{B) Power}
\par\smallskip
\textbf{Izoh:} Birinchi besh belgi olinadi.
\end{tcolorbox}
\vspace{2mm}\hrule height .2pt\vspace{2mm}

% DB 6a7894a3-9797-c184-eb21-9d5e83dead0b status=published
\needspace{6\baselineskip}
\textbf{\color{ink}260.} {ootnotesizecolor{muted}Y1 | published | 6a7894a3}
\par
=INT(-2.3) natija?
\begin{enumerate}[label=\Alph*),leftmargin=9mm]
\item -2
\item 2
\item -3
\item 3
\end{enumerate}
\begin{tcolorbox}[enhanced,breakable,colback=teal!4,colframe=teal!55!black,boxrule=0.45pt,arc=1.2mm,left=2.8mm,right=2.8mm,top=2mm,bottom=2mm]
\textbf{\color{teal}To'g'ri javob:} \textbf{C) -3}
\par\smallskip
\textbf{Izoh:} INT pastki butun songa yo'naltiradi.
\end{tcolorbox}
\vspace{2mm}\hrule height .2pt\vspace{2mm}

% DB 8c8cefa9-5e54-4def-877a-aa877a7b91fc status=published
\needspace{6\baselineskip}
\textbf{\color{ink}261.} {ootnotesizecolor{muted}Y3 | published | 8c8cefa9}
\par
MS Excelda matn yozilgan katakdagi belgilarning boshidan kerakli miqdorini LEFT funksiyasi bilan ajratib olish va tekshirish amallarini tartiblang.
\begin{enumerate}[label=\Alph*),leftmargin=9mm]
\item LEFT funksiyasiga manzilni berish
\item Hosil bo'lgan satrni tekshirish
\item Olinadigan belgilar sonini belgilash
\item Matn katagini aniqlash
\end{enumerate}
\begin{tcolorbox}[enhanced,breakable,colback=teal!4,colframe=teal!55!black,boxrule=0.45pt,arc=1.2mm,left=2.8mm,right=2.8mm,top=2mm,bottom=2mm]
\textbf{\color{teal}To'g'ri javob:} \textbf{D - A - C - B}
\par\smallskip
\textbf{Izoh:} Bosqichlar: Matn katagini aniqlash \ensuremath{\rightarrow} LEFT funksiyasiga manzilni berish \ensuremath{\rightarrow} Olinadigan belgilar sonini belgilash \ensuremath{\rightarrow} Hosil bo'lgan satrni tekshirish.
\end{tcolorbox}
\vspace{2mm}\hrule height .2pt\vspace{2mm}

% DB 910dd755-9523-3b20-954a-54a9d038898e status=published
\needspace{6\baselineskip}
\textbf{\color{ink}262.} {ootnotesizecolor{muted}Y3 | published | 910dd755}
\par
A2:B4 jadvalining birinchi ustunida 60, 70, 80 sonlari bor. D2 katakda esa \texttt{"60a"} matni yozilgan. E2 katak bo'sh. Exact-match VLOOKUP bilan qidirishdagi xatoni aniqlash va tuzatish ketma-ketligini joylang.
\begin{enumerate}[label=\Alph*),leftmargin=9mm]
\item E2 katakda \texttt{\#N/A} xatosi chiqqanini kuzatish
\item D2 dagi \texttt{"60a"} matni jadvaldagi \texttt{60} soniga teng emasligini aniqlash
\item D2 ni \texttt{60} soniga tuzatib, E2 formulasini qayta hisoblatish
\item E2 ga \texttt{=VLOOKUP(D2,A2:B4,2,FALSE)} formulasini kiritish
\end{enumerate}
\begin{tcolorbox}[enhanced,breakable,colback=teal!4,colframe=teal!55!black,boxrule=0.45pt,arc=1.2mm,left=2.8mm,right=2.8mm,top=2mm,bottom=2mm]
\textbf{\color{teal}To'g'ri javob:} \textbf{D - A - B - C}
\par\smallskip
\textbf{Izoh:} E2 ga aniq qidiruv formulasi kiritiladi; \texttt{"60a"} A ustunidagi \texttt{60} ga mos kelmagani uchun \texttt{\#N/A} kuzatiladi. Xatoga sabab bo'lgan qidiruv kaliti tekshiriladi va D2 raqamli 60 ga tuzatiladi. \texttt{FALSE} exact match talab qiladi.
\end{tcolorbox}
\vspace{2mm}\hrule height .2pt\vspace{2mm}

% DB e0b3efef-71f6-e9d5-484f-0d8910064d8e status=published
\needspace{6\baselineskip}
\textbf{\color{ink}263.} {ootnotesizecolor{muted}Y1 | published | e0b3efef}
\par
=MID("Formula",2,3) natija?
\begin{enumerate}[label=\Alph*),leftmargin=9mm]
\item For
\item mula
\item rmu
\item orm
\end{enumerate}
\begin{tcolorbox}[enhanced,breakable,colback=teal!4,colframe=teal!55!black,boxrule=0.45pt,arc=1.2mm,left=2.8mm,right=2.8mm,top=2mm,bottom=2mm]
\textbf{\color{teal}To'g'ri javob:} \textbf{D) orm}
\par\smallskip
\textbf{Izoh:} 2-belgidan uchta belgi: orm.
\end{tcolorbox}
\vspace{2mm}\hrule height .2pt\vspace{2mm}

% DB 2ed04742-6420-c508-acdc-e12d3fbc24cf status=published
\needspace{6\baselineskip}
\textbf{\color{ink}264.} {ootnotesizecolor{muted}Y1 | published | 2ed04742}
\par
A1:A4 qiymatlar: 5, 0, "Test", 10. \texttt{=IF(COUNT(A1:A4)=3,AVERAGE(A1:A4),-1)} natija?
\begin{enumerate}[label=\Alph*),leftmargin=9mm]
\item 5
\item 3.75
\item -1
\item 15
\end{enumerate}
\begin{tcolorbox}[enhanced,breakable,colback=teal!4,colframe=teal!55!black,boxrule=0.45pt,arc=1.2mm,left=2.8mm,right=2.8mm,top=2mm,bottom=2mm]
\textbf{\color{teal}To'g'ri javob:} \textbf{A) 5}
\par\smallskip
\textbf{Izoh:} COUNT=3 (matn e'tiborga olinmaydi), AVERAGE=(5+0+10)/3=5.
\end{tcolorbox}
\vspace{2mm}\hrule height .2pt\vspace{2mm}

% DB 3894bbbd-7b1a-847c-a512-f1aa36cae48f status=published
\needspace{6\baselineskip}
\textbf{\color{ink}265.} {ootnotesizecolor{muted}Y1 | published | 3894bbbd}
\par
A2:A4 = Daftar, Kitob, Daftar; B2:B4 = 3, 7, 5. =SUMIF(A2:A4,"Daftar",B2:B4)-COUNTIF(A2:A4,"Daftar") natija?
\begin{enumerate}[label=\Alph*),leftmargin=9mm]
\item 6
\item 2
\item 13
\item 8
\end{enumerate}
\begin{tcolorbox}[enhanced,breakable,colback=teal!4,colframe=teal!55!black,boxrule=0.45pt,arc=1.2mm,left=2.8mm,right=2.8mm,top=2mm,bottom=2mm]
\textbf{\color{teal}To'g'ri javob:} \textbf{A) 6}
\par\smallskip
\textbf{Izoh:} SUMIF 3+5=8; COUNTIF 2; 8-2=6.
\end{tcolorbox}
\vspace{2mm}\hrule height .2pt\vspace{2mm}

% DB 5f57dd30-12a0-f49a-5fda-783dc4814d34 status=published
\needspace{6\baselineskip}
\textbf{\color{ink}266.} {ootnotesizecolor{muted}Y1 | published | 5f57dd30}
\par
A2:A6 da sinf nomlari: A, B, A, A, B. B2:B6 da ballar: 4, 6, 9, 2, 8. \texttt{=SUMIF(A2:A6,"A",B2:B6)-COUNTIF(A2:A6,"A")} natijasi?
\begin{enumerate}[label=\Alph*),leftmargin=9mm]
\item 9
\item 11
\item 12
\item 15
\end{enumerate}
\begin{tcolorbox}[enhanced,breakable,colback=teal!4,colframe=teal!55!black,boxrule=0.45pt,arc=1.2mm,left=2.8mm,right=2.8mm,top=2mm,bottom=2mm]
\textbf{\color{teal}To'g'ri javob:} \textbf{C) 12}
\par\smallskip
\textbf{Izoh:} A guruhi bo'yicha balllar 4+9+2=15, COUNTIF=3; 15-3=12.
\end{tcolorbox}
\vspace{2mm}\hrule height .2pt\vspace{2mm}

% DB 64c58a81-703e-48e0-a2b1-2b64d1b53a59 status=published
\needspace{6\baselineskip}
\textbf{\color{ink}267.} {ootnotesizecolor{muted}Y3 | published | 64c58a81}
\par
MS Excelda B2:B18 kataklardagi o'quvchilar ballarining o'rtacha qiymatini hisoblab, natijani qo'lda tekshirish kerak. AVERAGE formulasini tuzish amallarini tartiblang.
\begin{enumerate}[label=\Alph*),leftmargin=9mm]
\item Funksiya va argumentlarni yozish
\item Ballar joylashgan B2:B18 diapazonini belgilash
\item Natijani qo'lda tekshirish
\item Belgilangan ballar diapazoni uchun AVERAGE funksiyasini tanlash
\end{enumerate}
\begin{tcolorbox}[enhanced,breakable,colback=teal!4,colframe=teal!55!black,boxrule=0.45pt,arc=1.2mm,left=2.8mm,right=2.8mm,top=2mm,bottom=2mm]
\textbf{\color{teal}To'g'ri javob:} \textbf{B - D - A - C}
\par\smallskip
\textbf{Izoh:} Avval B2:B18 diapazoni belgilanadi, unga tegishli AVERAGE funksiyasi tanlanadi, argumentlari kiritiladi va natija tekshiriladi. Tartib: 2 \ensuremath{\rightarrow} 4 \ensuremath{\rightarrow} 1 \ensuremath{\rightarrow} 3.
\end{tcolorbox}
\vspace{2mm}\hrule height .2pt\vspace{2mm}

% DB c7faec3f-da80-2e2a-cc8c-d28f3ca5766c status=published
\needspace{6\baselineskip}
\textbf{\color{ink}268.} {ootnotesizecolor{muted}Y2 | published | c7faec3f}
\par
Quyidagi tushunchalarni mazmuniga qarab bog'lang.  1) VLOOKUP;  2) HLOOKUP;  3) COUNTIF;  Izohlar: A) Shartga mos kataklarni sanash;  B) Birinchi ustunda vertikal qidiruv;  C) Birinchi satrda gorizontal qidiruv;
\textbf{Chap ustun:}\begin{enumerate}[label=\arabic*),leftmargin=9mm]
\item VLOOKUP
\item HLOOKUP
\item COUNTIF
\end{enumerate}
\textbf{O'ng ustun:}\begin{enumerate}[label=\Alph*),leftmargin=9mm]
\item Shartga mos kataklarni sanash
\item Birinchi ustunda vertikal qidiruv
\item Birinchi satrda gorizontal qidiruv
\end{enumerate}
\begin{tcolorbox}[enhanced,breakable,colback=teal!4,colframe=teal!55!black,boxrule=0.45pt,arc=1.2mm,left=2.8mm,right=2.8mm,top=2mm,bottom=2mm]
\textbf{\color{teal}To'g'ri javob:} \textbf{1-B, 2-C, 3-A}
\par\smallskip
\textbf{Izoh:} To'g'ri moslik: VLOOKUP - Birinchi ustunda vertikal qidiruv, HLOOKUP - Birinchi satrda gorizontal qidiruv, COUNTIF - Shartga mos kataklarni sanash.
\end{tcolorbox}
\vspace{2mm}\hrule height .2pt\vspace{2mm}

% DB d1454ff2-635b-cd85-9f51-7b1bdc898600 status=published
\needspace{6\baselineskip}
\textbf{\color{ink}269.} {ootnotesizecolor{muted}Y1 | published | d1454ff2}
\par
A1:A5 kataklarda 5, 0, bo'sh, "Ha", 7 bor. =IF(COUNT(A1:A5)=3,AVERAGE(A1:A5),0) natijasi qancha?
\begin{enumerate}[label=\Alph*),leftmargin=9mm]
\item 0
\item 3
\item 6
\item 4
\end{enumerate}
\begin{tcolorbox}[enhanced,breakable,colback=teal!4,colframe=teal!55!black,boxrule=0.45pt,arc=1.2mm,left=2.8mm,right=2.8mm,top=2mm,bottom=2mm]
\textbf{\color{teal}To'g'ri javob:} \textbf{D) 4}
\par\smallskip
\textbf{Izoh:} COUNT=3 (5,0,7), AVERAGE=(5+0+7)/3=4; shart rost.
\end{tcolorbox}
\vspace{2mm}\hrule height .2pt\vspace{2mm}

% DB df1f37b4-065d-34d3-103a-76173d30c13c status=published
\needspace{6\baselineskip}
\textbf{\color{ink}270.} {ootnotesizecolor{muted}Y2 | published | df1f37b4}
\par
Moslashtiring.  1) COUNT;  2) COUNTA;  3) COUNTIF;  Ta'riflar: A) Shartga mos kataklar soni;  B) Sonli kataklar soni;  C) Bo'sh bo'lmagan kataklar soni;
\textbf{Chap ustun:}\begin{enumerate}[label=\arabic*),leftmargin=9mm]
\item COUNT
\item COUNTA
\item COUNTIF
\end{enumerate}
\textbf{O'ng ustun:}\begin{enumerate}[label=\Alph*),leftmargin=9mm]
\item Shartga mos kataklar soni
\item Sonli kataklar soni
\item Bo'sh bo'lmagan kataklar soni
\end{enumerate}
\begin{tcolorbox}[enhanced,breakable,colback=teal!4,colframe=teal!55!black,boxrule=0.45pt,arc=1.2mm,left=2.8mm,right=2.8mm,top=2mm,bottom=2mm]
\textbf{\color{teal}To'g'ri javob:} \textbf{1-B, 2-C, 3-A}
\par\smallskip
\textbf{Izoh:} Mos keladigan ta'riflar: 1-B, 2-C, 3-A.
\end{tcolorbox}
\vspace{2mm}\hrule height .2pt\vspace{2mm}

% DB ed517b78-58e6-9aea-c1c8-fb20959dc042 status=published
\needspace{6\baselineskip}
\textbf{\color{ink}271.} {ootnotesizecolor{muted}Y2 | published | ed517b78}
\par
Funksiya va natijani moslang. A1:A4 = \{5; 0; "T"; 15\}; B1=6.
\textbf{Chap ustun:}\begin{enumerate}[label=\arabic*),leftmargin=9mm]
\item COUNT(A1:A4)
\item COUNTA(A1:A4)
\item MAX(A1:A4)
\item SUM(A1:A4)+B1
\end{enumerate}
\textbf{O'ng ustun:}\begin{enumerate}[label=\Alph*),leftmargin=9mm]
\item 26
\item 4
\item 15
\item 3
\end{enumerate}
\begin{tcolorbox}[enhanced,breakable,colback=teal!4,colframe=teal!55!black,boxrule=0.45pt,arc=1.2mm,left=2.8mm,right=2.8mm,top=2mm,bottom=2mm]
\textbf{\color{teal}To'g'ri javob:} \textbf{1-D, 2-B, 3-C, 4-A}
\par\smallskip
\textbf{Izoh:} COUNT=3; COUNTA=4; MAX=15; SUM=20+6=26.
\end{tcolorbox}
\vspace{2mm}\hrule height .2pt\vspace{2mm}

% DB f719f505-4778-4f00-a97c-1090784492b3 status=published
\needspace{6\baselineskip}
\textbf{\color{ink}272.} {ootnotesizecolor{muted}Y3 | published | f719f505}
\par
MS Excelda mahsulot kodi bo'yicha jadvaldan narxni faqat aniq moslik asosida qaytaruvchi VLOOKUP formulasining argumentlarini mantiqiy ketma-ketlikda tuzing.
\begin{enumerate}[label=\Alph*),leftmargin=9mm]
\item Qaytariladigan ustun raqamini ko'rsatish
\item Izlanadigan kalitni aniqlash
\item VLOOKUP uchun vertikal jadval diapazonini berish
\item Aniq moslik uchun FALSE parametrini tanlash
\end{enumerate}
\begin{tcolorbox}[enhanced,breakable,colback=teal!4,colframe=teal!55!black,boxrule=0.45pt,arc=1.2mm,left=2.8mm,right=2.8mm,top=2mm,bottom=2mm]
\textbf{\color{teal}To'g'ri javob:} \textbf{B - C - A - D}
\par\smallskip
\textbf{Izoh:} Bosqichlar: Izlanadigan kalitni aniqlash \ensuremath{\rightarrow} VLOOKUP uchun vertikal jadval diapazonini berish \ensuremath{\rightarrow} Qaytariladigan ustun raqamini ko'rsatish \ensuremath{\rightarrow} Aniq moslik uchun FALSE parametrini tanlash.
\end{tcolorbox}
\vspace{2mm}\hrule height .2pt\vspace{2mm}

% DB 224820ab-02c6-4791-8121-d5ebc971c598 status=review
\needspace{6\baselineskip}
\textbf{\color{ink}273.} {ootnotesizecolor{muted}Y1 | review | 224820ab}
\par
SUM va COUNTIF farqi?
\begin{enumerate}[label=\Alph*),leftmargin=9mm]
\item SUM qiymatlarni sanaydi, COUNTIF shartga moslarni qo'shadi
\item SUM shartni tekshiradi, COUNTIF jami qiymatni hisoblaydi
\item SUM faqat raqam formatlaydi, COUNTIF kataklarni saralaydi
\item SUM qiymatlarni qo'shadi, COUNTIF shartga moslarni sanaydi
\end{enumerate}
\begin{tcolorbox}[enhanced,breakable,colback=teal!4,colframe=teal!55!black,boxrule=0.45pt,arc=1.2mm,left=2.8mm,right=2.8mm,top=2mm,bottom=2mm]
\textbf{\color{teal}To'g'ri javob:} \textbf{D) SUM qiymatlarni qo'shadi, COUNTIF shartga moslarni sanaydi}
\par\smallskip
\textbf{Izoh:} SUM berilgan sonli qiymatlarning yig'indisini hisoblaydi. COUNTIF esa ko'rsatilgan diapazonda belgilangan shartga mos keladigan kataklar sonini qaytaradi. Shart asosida qiymatlarni yig'ish uchun alohida SUMIF ishlatiladi; SUM va COUNTIF vazifalarini almashtirmaslik kerak.
\end{tcolorbox}
\vspace{2mm}\hrule height .2pt\vspace{2mm}
\subsection{MS Excel: saralash, filtr, diagramma va ma'lumot tahlili}

% DB 5201a7da-08d8-20bc-e534-c95d58ac68d0 status=published
\needspace{6\baselineskip}
\textbf{\color{ink}274.} {ootnotesizecolor{muted}Y1 | published | 5201a7da}
\par
Excelda Sort nimani bajaradi?
\begin{enumerate}[label=\Alph*),leftmargin=9mm]
\item Barcha qatorlarni o'chiradi
\item Matnni shifrlaydi
\item Yozuvlarni mezon bo'yicha tartiblaydi
\item Slaydga animatsiya qo'shadi
\end{enumerate}
\begin{tcolorbox}[enhanced,breakable,colback=teal!4,colframe=teal!55!black,boxrule=0.45pt,arc=1.2mm,left=2.8mm,right=2.8mm,top=2mm,bottom=2mm]
\textbf{\color{teal}To'g'ri javob:} \textbf{C) Yozuvlarni mezon bo'yicha tartiblaydi}
\par\smallskip
\textbf{Izoh:} Sort tartibni o'zgartiradi.
\end{tcolorbox}
\vspace{2mm}\hrule height .2pt\vspace{2mm}

% DB 5fecdc5c-5075-8687-ff61-37b9e4dc7ca5 status=published
\needspace{6\baselineskip}
\textbf{\color{ink}275.} {ootnotesizecolor{muted}Y2 | published | 5fecdc5c}
\par
Quyidagi buyruqlarni vazifasi bilan moslashtiring.  1) Sort;  2) Filter;  3) Conditional Formatting;  Izohlar: A) Katakni shartga ko'ra bezash;  B) Faqat mos qatorlarni ko'rsatish;  C) Qatorlarni tartiblash;
\textbf{Chap ustun:}\begin{enumerate}[label=\arabic*),leftmargin=9mm]
\item Sort
\item Filter
\item Conditional Formatting
\end{enumerate}
\textbf{O'ng ustun:}\begin{enumerate}[label=\Alph*),leftmargin=9mm]
\item Katakni shartga ko'ra bezash
\item Faqat mos qatorlarni ko'rsatish
\item Qatorlarni tartiblash
\end{enumerate}
\begin{tcolorbox}[enhanced,breakable,colback=teal!4,colframe=teal!55!black,boxrule=0.45pt,arc=1.2mm,left=2.8mm,right=2.8mm,top=2mm,bottom=2mm]
\textbf{\color{teal}To'g'ri javob:} \textbf{1-C, 2-B, 3-A}
\par\smallskip
\textbf{Izoh:} To'g'ri moslik: Sort - Qatorlarni tartiblash, Filter - Faqat mos qatorlarni ko'rsatish, Conditional Formatting - Katakni shartga ko'ra bezash.
\end{tcolorbox}
\vspace{2mm}\hrule height .2pt\vspace{2mm}

% DB 62189a80-aeab-404f-958e-1e00734aea05 status=published
\needspace{6\baselineskip}
\textbf{\color{ink}276.} {ootnotesizecolor{muted}Y1 | published | 62189a80}
\par
Excelda 200 nafar o'quvchi ro'yxatidan faqat 9-sinfdagilar vaqtincha ko'rinishi kerak. Qolgan yozuvlar o'chirilmasin. Filter nima qiladi?
\begin{enumerate}[label=\Alph*),leftmargin=9mm]
\item Mos kelgan yozuvlarni o'chirib, qolganlarini jadvalda saqlaydi
\item Barcha yozuvlarni alifbo bo'yicha qayta joylashtirib, tartiblaydi
\item Har bir yozuvni boshqa varaqqa ko'chirib, aslini yangilaydi
\item Mezonga mos yozuvlarni ko'rsatib, qolganlarini vaqtincha yashiradi
\end{enumerate}
\begin{tcolorbox}[enhanced,breakable,colback=teal!4,colframe=teal!55!black,boxrule=0.45pt,arc=1.2mm,left=2.8mm,right=2.8mm,top=2mm,bottom=2mm]
\textbf{\color{teal}To'g'ri javob:} \textbf{D) Mezonga mos yozuvlarni ko'rsatib, qolganlarini vaqtincha yashiradi}
\par\smallskip
\textbf{Izoh:} Filter shartga mos satrlarni ko'rsatib, boshqa satrlarni vaqtincha yashiradi; ma'lumotlarni o'chirmaydi. Sort esa tartibni o'zgartiradi.
\end{tcolorbox}
\vspace{2mm}\hrule height .2pt\vspace{2mm}

% DB 7788d7cd-8558-cc9b-fb11-bf0d5562c47b status=published
\needspace{6\baselineskip}
\textbf{\color{ink}277.} {ootnotesizecolor{muted}Y1 | published | 7788d7cd}
\par
Line diagrammasi odatda nimaga qulay?
\begin{enumerate}[label=\Alph*),leftmargin=9mm]
\item Vaqt bo'yicha tendensiyani ko'rsatishga
\item Faqat matn tarjima qilishga
\item Wordda jadvalni bo'lishga
\item Slayd raqamini qo'shishga
\end{enumerate}
\begin{tcolorbox}[enhanced,breakable,colback=teal!4,colframe=teal!55!black,boxrule=0.45pt,arc=1.2mm,left=2.8mm,right=2.8mm,top=2mm,bottom=2mm]
\textbf{\color{teal}To'g'ri javob:} \textbf{A) Vaqt bo'yicha tendensiyani ko'rsatishga}
\par\smallskip
\textbf{Izoh:} Line trend tasvirlaydi.
\end{tcolorbox}
\vspace{2mm}\hrule height .2pt\vspace{2mm}

% DB f4e3b11b-5ed7-4ae8-adc5-0e0ce870ecb8 status=published
\needspace{6\baselineskip}
\textbf{\color{ink}278.} {ootnotesizecolor{muted}Y3 | published | f4e3b11b}
\par
MS Excel jadvalidagi faqat belgilangan mezonga mos satrlarni ko'rsatish uchun Data Filter qo'llash va ko'rinib qolgan satrlarni tekshirish amallarini tartiblang.
\begin{enumerate}[label=\Alph*),leftmargin=9mm]
\item Data Filter buyrug'ini yoqish
\item Ko'ringan satrlarni tahlil qilish
\item Mezonni tanlash
\item Sarlavhali jadvalni tanlash
\end{enumerate}
\begin{tcolorbox}[enhanced,breakable,colback=teal!4,colframe=teal!55!black,boxrule=0.45pt,arc=1.2mm,left=2.8mm,right=2.8mm,top=2mm,bottom=2mm]
\textbf{\color{teal}To'g'ri javob:} \textbf{D - A - C - B}
\par\smallskip
\textbf{Izoh:} Bosqichlar: Sarlavhali jadvalni tanlash \ensuremath{\rightarrow} Data Filter buyrug'ini yoqish \ensuremath{\rightarrow} Mezonni tanlash \ensuremath{\rightarrow} Ko'ringan satrlarni tahlil qilish.
\end{tcolorbox}
\vspace{2mm}\hrule height .2pt\vspace{2mm}

% DB 146f0067-e35a-41fc-bd9a-ae9073b1b7d4 status=published
\needspace{6\baselineskip}
\textbf{\color{ink}279.} {ootnotesizecolor{muted}Y1 | published | 146f0067}
\par
Jadvalda baholarni kattadan kichikka joylash va faqat muayyan sinfning qatorlarini ko'rsatish kerak. Sort va Filter orasidagi farqni toping.
\begin{enumerate}[label=\Alph*),leftmargin=9mm]
\item Sort tartiblaydi, Filter qatorlarni vaqtincha yashiradi
\item Sort qatorlarni o'chiradi, Filter esa alohida nusxa ko'chiradi
\item Sort qatorlarni yashiradi, Filter esa barcha qatorni saralaydi
\item Sort ustunlarni birlashtiradi, Filter esa qiymatlarni avtomatik hisoblaydi
\end{enumerate}
\begin{tcolorbox}[enhanced,breakable,colback=teal!4,colframe=teal!55!black,boxrule=0.45pt,arc=1.2mm,left=2.8mm,right=2.8mm,top=2mm,bottom=2mm]
\textbf{\color{teal}To'g'ri javob:} \textbf{A) Sort tartiblaydi, Filter qatorlarni vaqtincha yashiradi}
\par\smallskip
\textbf{Izoh:} Sort yozuvlarning tartibini o'zgartiradi, Filter mezonga mos bo'lmagan qatorlarni ko'rinishdan vaqtincha yashiradi. Hech biri o'z vazifasi bo'yicha yozuvlarni majburan o'chirmaydi.
\end{tcolorbox}
\vspace{2mm}\hrule height .2pt\vspace{2mm}

% DB 1dc6fe77-ca41-5c2e-ba27-8ae37439db6f status=published
\needspace{6\baselineskip}
\textbf{\color{ink}280.} {ootnotesizecolor{muted}Y1 | published | 1dc6fe77}
\par
Diagrammada Damas avtomobili ishlab chiqarilishi I chorakda 700 ta, III chorakda 650 ta va IV chorakda 650 ta. Shu uch chorak bo'yicha Damas ishlab chiqarish ulushlarini toping.
\begin{enumerate}[label=\Alph*),leftmargin=9mm]
\item 35\%; 32,5\%; 32,5\%
\item 30\%; 35\%; 35\%
\item 33,3\%; 33,3\%; 33,3\%
\item 32,5\%; 35\%; 32,5\%
\end{enumerate}
\begin{tcolorbox}[enhanced,breakable,colback=teal!4,colframe=teal!55!black,boxrule=0.45pt,arc=1.2mm,left=2.8mm,right=2.8mm,top=2mm,bottom=2mm]
\textbf{\color{teal}To'g'ri javob:} \textbf{A) 35\%; 32,5\%; 32,5\%}
\par\smallskip
\textbf{Izoh:} Jami 700+650+650=2000. Ulushlar: 700/2000=35\%, 650/2000=32,5\% va 650/2000=32,5\%. Manbadagi doiraviy diagramma variantlaridagi yaxlitlash nomuvofiqligi aniq sonli variantlar bilan tuzatildi.
\end{tcolorbox}
\vspace{2mm}\hrule height .2pt\vspace{2mm}

% DB 29ec2b83-f0ab-4c2d-ab50-c1cf661b81a4 status=published
\needspace{6\baselineskip}
\textbf{\color{ink}281.} {ootnotesizecolor{muted}Y1 | published | 29ec2b83}
\par
Hududlar bo'yicha bir yillik savdo hajmi va yil davomida oyma-oy savdo o'zgarishi ko'rsatiladi. Column va Line diagrammalari qachon maqsadga muvofiq?
\begin{enumerate}[label=\Alph*),leftmargin=9mm]
\item Column ikki sonli o'zgaruvchi aloqasini, Line esa ulushlarni ko'rsatadi
\item Column toifalar taqqoslovi, Line dinamik trend uchun
\item Column umumiy ulushni foizlarda, Line esa matn belgilarini ko'rsatadi
\item Column faqat alohida yozuvlarni, Line esa fayl nomlarini ko'rsatadi
\end{enumerate}
\begin{tcolorbox}[enhanced,breakable,colback=teal!4,colframe=teal!55!black,boxrule=0.45pt,arc=1.2mm,left=2.8mm,right=2.8mm,top=2mm,bottom=2mm]
\textbf{\color{teal}To'g'ri javob:} \textbf{B) Column toifalar taqqoslovi, Line dinamik trend uchun}
\par\smallskip
\textbf{Izoh:} Column kategoriyalar orasida qiymatlarni taqqoslashga, Line esa ketma-ket davrlardagi o'zgarish tendensiyasini kuzatishga mos. Tanlov ma'lumot va tahlil maqsadiga bog'liq.
\end{tcolorbox}
\vspace{2mm}\hrule height .2pt\vspace{2mm}

% DB 80699016-79e0-5a5f-8476-716bfd724830 status=published
\needspace{6\baselineskip}
\textbf{\color{ink}282.} {ootnotesizecolor{muted}Y1 | published | 80699016}
\par
Diagrammadagi Damas ishlab chiqarish sonlari: I chorak - 700 ta, III chorak - 650 ta, IV chorak - 650 ta. Shu uch chorak bo'yicha Damas avtomobillarining ulushlarini I, III va IV chorak tartibida toping.
\begin{enumerate}[label=\Alph*),leftmargin=9mm]
\item 35\%; 32,5\%; 32,5\%
\item 32,5\%; 35\%; 32,5\%
\item 33\%; 35\%; 32\%
\item 25\%; 30\%; 45\%
\end{enumerate}
\begin{tcolorbox}[enhanced,breakable,colback=teal!4,colframe=teal!55!black,boxrule=0.45pt,arc=1.2mm,left=2.8mm,right=2.8mm,top=2mm,bottom=2mm]
\textbf{\color{teal}To'g'ri javob:} \textbf{A) 35\%; 32,5\%; 32,5\%}
\par\smallskip
\textbf{Izoh:} Jami 700+650+650=2000 ta. I chorak: 700/2000=35\%; III va IV choraklarning har biri 650/2000=32,5\%. Asl PDFdagi doiraviy diagramma variantlarida aniq 35\%; 32,5\%; 32,5\% javobi yo'q edi; savol aniq sonli variantlar bilan tiklandi.
\end{tcolorbox}
\vspace{2mm}\hrule height .2pt\vspace{2mm}

% DB a7ec6dcc-1524-493b-94da-dbafad5ed8ba status=published
\needspace{6\baselineskip}
\textbf{\color{ink}283.} {ootnotesizecolor{muted}Y3 | published | a7ec6dcc}
\par
MS Excelda bir necha davrdagi ko'rsatkichlarning o'zgarishini chiziqli diagrammada aks ettirib, o'qlar va sarlavhani nomlash amallarini tartiblang.
\begin{enumerate}[label=\Alph*),leftmargin=9mm]
\item Line grafik turini tanlash
\item Manba ma'lumotlarini tekshirish
\item O'q nomlari va sarlavhani yozish
\item Diagramma uchun mos diapazonni belgilash
\end{enumerate}
\begin{tcolorbox}[enhanced,breakable,colback=teal!4,colframe=teal!55!black,boxrule=0.45pt,arc=1.2mm,left=2.8mm,right=2.8mm,top=2mm,bottom=2mm]
\textbf{\color{teal}To'g'ri javob:} \textbf{B - D - A - C}
\par\smallskip
\textbf{Izoh:} Bosqichlar: Manba ma'lumotlarini tekshirish \ensuremath{\rightarrow} Diagramma uchun mos diapazonni belgilash \ensuremath{\rightarrow} Line grafik turini tanlash \ensuremath{\rightarrow} O'q nomlari va sarlavhani yozish.
\end{tcolorbox}
\vspace{2mm}\hrule height .2pt\vspace{2mm}

% DB b73ff9f8-e26d-3878-c259-6d2d97ceaa79 status=published
\needspace{6\baselineskip}
\textbf{\color{ink}284.} {ootnotesizecolor{muted}Y2 | published | b73ff9f8}
\par
Atamalar va ularning aniq ma'nolarini moslashtiring.  1) Column chart;  2) Line chart;  3) Pie chart;  Izohlar: A) Toifalar qiymatini ustunlarda taqqoslash;  B) Bitta butun ulushlari;  C) Ketma-ket qiymatlar dinamikasi;
\textbf{Chap ustun:}\begin{enumerate}[label=\arabic*),leftmargin=9mm]
\item Column chart
\item Line chart
\item Pie chart
\end{enumerate}
\textbf{O'ng ustun:}\begin{enumerate}[label=\Alph*),leftmargin=9mm]
\item Toifalar qiymatini ustunlarda taqqoslash
\item Bitta butun ulushlari
\item Ketma-ket qiymatlar dinamikasi
\end{enumerate}
\begin{tcolorbox}[enhanced,breakable,colback=teal!4,colframe=teal!55!black,boxrule=0.45pt,arc=1.2mm,left=2.8mm,right=2.8mm,top=2mm,bottom=2mm]
\textbf{\color{teal}To'g'ri javob:} \textbf{1-A, 2-C, 3-B}
\par\smallskip
\textbf{Izoh:} To'g'ri moslik: Column chart - Toifalar qiymatini ustunlarda taqqoslash, Line chart - Ketma-ket qiymatlar dinamikasi, Pie chart - Bitta butun ulushlari.
\end{tcolorbox}
\vspace{2mm}\hrule height .2pt\vspace{2mm}

% DB bea73c95-742d-382b-bb00-560a2b44d883 status=published
\needspace{6\baselineskip}
\textbf{\color{ink}285.} {ootnotesizecolor{muted}Y1 | published | bea73c95}
\par
Filtrdan so'ng jadvalda faqat balli kamida 70 va tumani "Navoiy" bo'lganlar ko'rinsin. Qaysi shart to'g'ri?
\begin{enumerate}[label=\Alph*),leftmargin=9mm]
\item Ball \textless{} 70 AND Tuman = Navoiy
\item Faqat ballni kamayish tartibida Sort qilish
\item Ball \ensuremath{\geq} 70 AND Tuman = Navoiy
\item Ball \ensuremath{\geq} 70 OR Tuman = Navoiy
\end{enumerate}
\begin{tcolorbox}[enhanced,breakable,colback=teal!4,colframe=teal!55!black,boxrule=0.45pt,arc=1.2mm,left=2.8mm,right=2.8mm,top=2mm,bottom=2mm]
\textbf{\color{teal}To'g'ri javob:} \textbf{C) Ball \ensuremath{\geq} 70 AND Tuman = Navoiy}
\par\smallskip
\textbf{Izoh:} Ikkala talab bir vaqtda bajarilishi shart; Sort filtr o'rnini bosmaydi.
\end{tcolorbox}
\vspace{2mm}\hrule height .2pt\vspace{2mm}

% DB ea21ba4a-b578-4b35-a0e2-1b6f9a324f90 status=published
\needspace{6\baselineskip}
\textbf{\color{ink}286.} {ootnotesizecolor{muted}Y1 | published | ea21ba4a}
\par
Sinflarning umumiy o'quvchi sonidagi ulushi va o'quvchilar bo'yi bilan vazni o'rtasidagi bog'liqlik turli grafiklarda tasvirlanadi. Pie va Scatter farqini toping.
\begin{enumerate}[label=\Alph*),leftmargin=9mm]
\item Pie vaqt davridagi trendni, Scatter esa tartiblangan reytingni ko'rsatadi
\item Pie sonli juftliklarni, Scatter esa yakuniy tarkib ulushini ko'rsatadi
\item Pie ulush, Scatter ikki sonli o'zgaruvchi aloqasini ko'rsatadi
\item Pie filtrlash natijasini, Scatter esa barcha satr tartibini ko'rsatadi
\end{enumerate}
\begin{tcolorbox}[enhanced,breakable,colback=teal!4,colframe=teal!55!black,boxrule=0.45pt,arc=1.2mm,left=2.8mm,right=2.8mm,top=2mm,bottom=2mm]
\textbf{\color{teal}To'g'ri javob:} \textbf{C) Pie ulush, Scatter ikki sonli o'zgaruvchi aloqasini ko'rsatadi}
\par\smallskip
\textbf{Izoh:} Pie umumiy qiymat tarkibidagi ulushlarni tasvirlaydi. Scatter esa ikkita sonli o'zgaruvchi orasidagi bog'lanishni nuqtalar orqali ko'rsatadi.
\end{tcolorbox}
\vspace{2mm}\hrule height .2pt\vspace{2mm}

% DB 21384f62-a4d5-4d72-d08c-6244a78b9253 status=published
\needspace{6\baselineskip}
\textbf{\color{ink}287.} {ootnotesizecolor{muted}Y1 | published | 21384f62}
\par
5 o'quvchi ballarini eng kattadan kichikka tartiblash uchun?
\begin{enumerate}[label=\Alph*),leftmargin=9mm]
\item Merge Cells
\item Data Validation faqat
\item Descending Sort
\item Ascending Sort
\end{enumerate}
\begin{tcolorbox}[enhanced,breakable,colback=teal!4,colframe=teal!55!black,boxrule=0.45pt,arc=1.2mm,left=2.8mm,right=2.8mm,top=2mm,bottom=2mm]
\textbf{\color{teal}To'g'ri javob:} \textbf{C) Descending Sort}
\par\smallskip
\textbf{Izoh:} Kamayish tartibi descending.
\end{tcolorbox}
\vspace{2mm}\hrule height .2pt\vspace{2mm}

% DB 5ee4bbd8-8c9e-4b8b-864b-90ad70db8c2b status=published
\needspace{6\baselineskip}
\textbf{\color{ink}288.} {ootnotesizecolor{muted}Y3 | published | 5ee4bbd8}
\par
MS Excelda ball kiritiladigan kataklarga 0 dan 100 gacha qiymatlarni qabul qiluvchi Data Validation cheklovi qo'yish va uni sinash amallarini tartiblang.
\begin{enumerate}[label=\Alph*),leftmargin=9mm]
\item Chegaradan tashqari qiymatni sinash
\item 0 dan 100 gacha cheklov qo'yish
\item Data Validation oynasini ochish
\item Ball kataklarini tanlash
\end{enumerate}
\begin{tcolorbox}[enhanced,breakable,colback=teal!4,colframe=teal!55!black,boxrule=0.45pt,arc=1.2mm,left=2.8mm,right=2.8mm,top=2mm,bottom=2mm]
\textbf{\color{teal}To'g'ri javob:} \textbf{D - C - B - A}
\par\smallskip
\textbf{Izoh:} Bosqichlar: Ball kataklarini tanlash \ensuremath{\rightarrow} Data Validation oynasini ochish \ensuremath{\rightarrow} 0 dan 100 gacha cheklov qo'yish \ensuremath{\rightarrow} Chegaradan tashqari qiymatni sinash.
\end{tcolorbox}
\vspace{2mm}\hrule height .2pt\vspace{2mm}

% DB 7b83e332-6344-22e8-1f90-482eb9931c5e status=published
\needspace{6\baselineskip}
\textbf{\color{ink}289.} {ootnotesizecolor{muted}Y1 | published | 7b83e332}
\par
Jadvaldan faqat Toshkent yozuvlari ko'rinsin. Nima qilinadi?
\begin{enumerate}[label=\Alph*),leftmargin=9mm]
\item Barcha boshqa qatorni o'chirish
\item Slayd o'tishini yoqish
\item Wordda Replace All
\item Matn filtri: Toshkent
\end{enumerate}
\begin{tcolorbox}[enhanced,breakable,colback=teal!4,colframe=teal!55!black,boxrule=0.45pt,arc=1.2mm,left=2.8mm,right=2.8mm,top=2mm,bottom=2mm]
\textbf{\color{teal}To'g'ri javob:} \textbf{D) Matn filtri: Toshkent}
\par\smallskip
\textbf{Izoh:} Filtr yozuvlarni yo'qotmaydi.
\end{tcolorbox}
\vspace{2mm}\hrule height .2pt\vspace{2mm}

% DB 8032a35f-8aa0-afe1-be31-af5a98c6d9b5 status=published
\needspace{6\baselineskip}
\textbf{\color{ink}290.} {ootnotesizecolor{muted}Y1 | published | 8032a35f}
\par
12 oy savdosi o'zgarishini ko'rsatish uchun qaysi diagramma mos?
\begin{enumerate}[label=\Alph*),leftmargin=9mm]
\item Line
\item WordArt
\item Pie faqat
\item Doughnut faqat
\end{enumerate}
\begin{tcolorbox}[enhanced,breakable,colback=teal!4,colframe=teal!55!black,boxrule=0.45pt,arc=1.2mm,left=2.8mm,right=2.8mm,top=2mm,bottom=2mm]
\textbf{\color{teal}To'g'ri javob:} \textbf{A) Line}
\par\smallskip
\textbf{Izoh:} Vaqt ketma-ketligi trendni ifodalaydi.
\end{tcolorbox}
\vspace{2mm}\hrule height .2pt\vspace{2mm}

% DB a6ada9f7-cdee-0645-b889-50a98598615e status=published
\needspace{6\baselineskip}
\textbf{\color{ink}291.} {ootnotesizecolor{muted}Y1 | published | a6ada9f7}
\par
Excelda A sinf 24, B sinf 36, C sinf 60 ta kitob to'pladi. Jami bo'yicha doiraviy diagrammada C sinf ulushi qanchaga teng?
\begin{enumerate}[label=\Alph*),leftmargin=9mm]
\item 60\%
\item 33\%
\item 25\%
\item 50\%
\end{enumerate}
\begin{tcolorbox}[enhanced,breakable,colback=teal!4,colframe=teal!55!black,boxrule=0.45pt,arc=1.2mm,left=2.8mm,right=2.8mm,top=2mm,bottom=2mm]
\textbf{\color{teal}To'g'ri javob:} \textbf{D) 50\%}
\par\smallskip
\textbf{Izoh:} Jami 24+36+60=120; C sinf ulushi 60/120=50\%. Diagrammadagi nisbiy ulush masalasi.
\end{tcolorbox}
\vspace{2mm}\hrule height .2pt\vspace{2mm}

% DB c841756a-9651-bcc3-919e-17e9406abb9b status=published
\needspace{6\baselineskip}
\textbf{\color{ink}292.} {ootnotesizecolor{muted}Y2 | published | c841756a}
\par
Amallarni izohi bilan moslang.  1) Chart title;  2) Legend;  3) Axis title;  Izohlar: A) O'lchov o'qi yorlig'i;  B) Grafik mazmuni sarlavhasi;  C) Ma'lumot seriyalari kaliti;
\textbf{Chap ustun:}\begin{enumerate}[label=\arabic*),leftmargin=9mm]
\item Chart title
\item Legend
\item Axis title
\end{enumerate}
\textbf{O'ng ustun:}\begin{enumerate}[label=\Alph*),leftmargin=9mm]
\item O'lchov o'qi yorlig'i
\item Grafik mazmuni sarlavhasi
\item Ma'lumot seriyalari kaliti
\end{enumerate}
\begin{tcolorbox}[enhanced,breakable,colback=teal!4,colframe=teal!55!black,boxrule=0.45pt,arc=1.2mm,left=2.8mm,right=2.8mm,top=2mm,bottom=2mm]
\textbf{\color{teal}To'g'ri javob:} \textbf{1-B, 2-C, 3-A}
\par\smallskip
\textbf{Izoh:} To'g'ri moslik: Chart title - Grafik mazmuni sarlavhasi, Legend - Ma'lumot seriyalari kaliti, Axis title - O'lchov o'qi yorlig'i.
\end{tcolorbox}
\vspace{2mm}\hrule height .2pt\vspace{2mm}

% DB fbd8fa68-e9f8-5855-b465-52fca2c1d34d status=published
\needspace{6\baselineskip}
\textbf{\color{ink}293.} {ootnotesizecolor{muted}Y1 | published | fbd8fa68}
\par
Sotuv diagrammasida sentabrda Rice=40 va Barley=60, oktabrda Rice=100 va Barley=60, noyabrda Rice=100 va Barley=60. Uch davr bo'yicha Rice+Barley jami sotuvining sentabr, oktabr va noyabr ulushlarini toping.
\begin{enumerate}[label=\Alph*),leftmargin=9mm]
\item 50\%; 26\%; 24\%
\item 33,3\%; 33,3\%; 33,3\%
\item 50\%; 25\%; 25\%
\item 23,8\%; 38,1\%; 38,1\%
\end{enumerate}
\begin{tcolorbox}[enhanced,breakable,colback=teal!4,colframe=teal!55!black,boxrule=0.45pt,arc=1.2mm,left=2.8mm,right=2.8mm,top=2mm,bottom=2mm]
\textbf{\color{teal}To'g'ri javob:} \textbf{D) 23,8\%; 38,1\%; 38,1\%}
\par\smallskip
\textbf{Izoh:} Davrlar bo'yicha Rice+Barley: 100, 160, 160; jami 420. Ulushlar 100/420\ensuremath{\approx}23,8\%, 160/420\ensuremath{\approx}38,1\% va 160/420\ensuremath{\approx}38,1\%. Asl rasmli variantlar tekshiriladigan aniq foizlarga aylantirildi.
\end{tcolorbox}
\vspace{2mm}\hrule height .2pt\vspace{2mm}

% DB 364aa4d2-5b87-eb1e-09fc-c8b202771846 status=published
\needspace{6\baselineskip}
\textbf{\color{ink}294.} {ootnotesizecolor{muted}Y2 | published | 364aa4d2}
\par
Excel diagrammalarining vazifalari bilan moslang.
\textbf{Chap ustun:}\begin{enumerate}[label=\arabic*),leftmargin=9mm]
\item Column
\item Line
\item Pie
\item Scatter (XY)
\end{enumerate}
\textbf{O'ng ustun:}\begin{enumerate}[label=\Alph*),leftmargin=9mm]
\item Ikki sonli o'zgaruvchi orasidagi tarqoqlikni ko'rsatish
\item Toifalar qiymatlarini yonma-yon qiyoslash
\item Vaqt bo'yicha yo'nalishni ko'rsatish
\item Bitta butunning qismlari ulushini ifodalash
\end{enumerate}
\begin{tcolorbox}[enhanced,breakable,colback=teal!4,colframe=teal!55!black,boxrule=0.45pt,arc=1.2mm,left=2.8mm,right=2.8mm,top=2mm,bottom=2mm]
\textbf{\color{teal}To'g'ri javob:} \textbf{1-B, 2-C, 3-D, 4-A}
\par\smallskip
\textbf{Izoh:} Column-B, Line-C, Pie-D, Scatter-A. Demo 2026 diagramma talqini shakliga asoslangan yangi moslashtirish.
\end{tcolorbox}
\vspace{2mm}\hrule height .2pt\vspace{2mm}

% DB 39fe09d6-888f-58ad-47dc-f4ed502dc62e status=published
\needspace{6\baselineskip}
\textbf{\color{ink}295.} {ootnotesizecolor{muted}Y1 | published | 39fe09d6}
\par
Ro'yxatda Filter yoqilib 100 yozuvdan 23 tasi ko'rindi. Qolgan 77 tasi?
\begin{enumerate}[label=\Alph*),leftmargin=9mm]
\item Yashirilgan, ma'lumotlar saqlangan
\item Rasmga aylangan
\item Yangi workbookga ko'chgan
\item Albatta o'chirilgan
\end{enumerate}
\begin{tcolorbox}[enhanced,breakable,colback=teal!4,colframe=teal!55!black,boxrule=0.45pt,arc=1.2mm,left=2.8mm,right=2.8mm,top=2mm,bottom=2mm]
\textbf{\color{teal}To'g'ri javob:} \textbf{A) Yashirilgan, ma'lumotlar saqlangan}
\par\smallskip
\textbf{Izoh:} Filter qatorlarni vaqtincha yashiradi.
\end{tcolorbox}
\vspace{2mm}\hrule height .2pt\vspace{2mm}

% DB 8d935e13-833e-54c6-1108-16324def8da1 status=published
\needspace{6\baselineskip}
\textbf{\color{ink}296.} {ootnotesizecolor{muted}Y2 | published | 8d935e13}
\par
Ushbu vaziyatdagi vositalar va vazifalarni moslashtiring.  1) Data Validation;  2) Verification;  3) What-if analysis;  Izohlar: A) Modelning taxminiy shartini o'zgartirib ko'rish;  B) Kiritilgan qiymatni manba bilan taqqoslash;  C) Kiritiladigan qiymatni qoida bilan cheklash;
\textbf{Chap ustun:}\begin{enumerate}[label=\arabic*),leftmargin=9mm]
\item Data Validation
\item Verification
\item What-if analysis
\end{enumerate}
\textbf{O'ng ustun:}\begin{enumerate}[label=\Alph*),leftmargin=9mm]
\item Modelning taxminiy shartini o'zgartirib ko'rish
\item Kiritilgan qiymatni manba bilan taqqoslash
\item Kiritiladigan qiymatni qoida bilan cheklash
\end{enumerate}
\begin{tcolorbox}[enhanced,breakable,colback=teal!4,colframe=teal!55!black,boxrule=0.45pt,arc=1.2mm,left=2.8mm,right=2.8mm,top=2mm,bottom=2mm]
\textbf{\color{teal}To'g'ri javob:} \textbf{1-C, 2-B, 3-A}
\par\smallskip
\textbf{Izoh:} To'g'ri moslik: Data Validation - Kiritiladigan qiymatni qoida bilan cheklash, Verification - Kiritilgan qiymatni manba bilan taqqoslash, What-if analysis - Modelning taxminiy shartini o'zgartirib ko'rish.
\end{tcolorbox}
\vspace{2mm}\hrule height .2pt\vspace{2mm}

% DB a61f322b-110a-626f-443d-14d6b248de7a status=published
\needspace{6\baselineskip}
\textbf{\color{ink}297.} {ootnotesizecolor{muted}Y1 | published | a61f322b}
\par
Excel jadvalidagi (Fan, Ball) qatorlar: (Informatika, 87), (Matematika, 79), (Informatika, 65), (Informatika, 92), (Matematika, 84), (Informatika, 80). Fan="Informatika" VA Ball\ensuremath{\geq}80 filtri yoqilganda nechta satr ko'rinadi va ularning ballari yig'indisi nechaga teng?
\begin{enumerate}[label=\Alph*),leftmargin=9mm]
\item 3 ta satr; yig'indi 259
\item 4 ta satr; yig'indi 324
\item 2 ta satr; yig'indi 179
\item 3 ta satr; yig'indi 244
\end{enumerate}
\begin{tcolorbox}[enhanced,breakable,colback=teal!4,colframe=teal!55!black,boxrule=0.45pt,arc=1.2mm,left=2.8mm,right=2.8mm,top=2mm,bottom=2mm]
\textbf{\color{teal}To'g'ri javob:} \textbf{A) 3 ta satr; yig'indi 259}
\par\smallskip
\textbf{Izoh:} Informatika satrlari 87,65,92,80. \ensuremath{\geq}80 mezoniga 87,92,80 mos: 3 ta satr, jami 87+92+80=259. Matematika qatorlari filtrdan o'tmaydi. Yig'indini barcha yozuvlarga emas, filtrga mos ko'rinadigan yozuvlarga hisoblash kerak.
\end{tcolorbox}
\vspace{2mm}\hrule height .2pt\vspace{2mm}

% DB c0a176f0-5d93-4401-b6bc-1b0136db3d0f status=published
\needspace{6\baselineskip}
\textbf{\color{ink}298.} {ootnotesizecolor{muted}Y3 | published | c0a176f0}
\par
MS Excelda ma'lumotlar jadvalini tanlangan ustun bo'yicha o'sish tartibida saralab, barcha ustunlar bir qatorda qolganini tekshirish amallarini tartiblang.
\begin{enumerate}[label=\Alph*),leftmargin=9mm]
\item Ascending tartibini qo'llash
\item Sort mezoni ustunini tanlash
\item Manba jadvalni tayyorlash
\item Butun qatorlar birga siljiganini tekshirish
\end{enumerate}
\begin{tcolorbox}[enhanced,breakable,colback=teal!4,colframe=teal!55!black,boxrule=0.45pt,arc=1.2mm,left=2.8mm,right=2.8mm,top=2mm,bottom=2mm]
\textbf{\color{teal}To'g'ri javob:} \textbf{C - B - A - D}
\par\smallskip
\textbf{Izoh:} Bosqichlar: Manba jadvalni tayyorlash \ensuremath{\rightarrow} Sort mezoni ustunini tanlash \ensuremath{\rightarrow} Ascending tartibini qo'llash \ensuremath{\rightarrow} Butun qatorlar birga siljiganini tekshirish.
\end{tcolorbox}
\vspace{2mm}\hrule height .2pt\vspace{2mm}

% DB cd772e39-d561-8aab-7986-a6f200445980 status=published
\needspace{6\baselineskip}
\textbf{\color{ink}299.} {ootnotesizecolor{muted}Y1 | published | cd772e39}
\par
50 balldan kichik kataklar qizil bo'lsin. Nima kerak?
\begin{enumerate}[label=\Alph*),leftmargin=9mm]
\item Insert Pie
\item Conditional Formatting, less than 50
\item Sort descending
\item AutoSum
\end{enumerate}
\begin{tcolorbox}[enhanced,breakable,colback=teal!4,colframe=teal!55!black,boxrule=0.45pt,arc=1.2mm,left=2.8mm,right=2.8mm,top=2mm,bottom=2mm]
\textbf{\color{teal}To'g'ri javob:} \textbf{B) Conditional Formatting, less than 50}
\par\smallskip
\textbf{Izoh:} Shartli formatlash rang qo'llaydi.
\end{tcolorbox}
\vspace{2mm}\hrule height .2pt\vspace{2mm}

% DB ed0d594c-a445-23cf-343b-ad9b86b99017 status=published
\needspace{6\baselineskip}
\textbf{\color{ink}300.} {ootnotesizecolor{muted}Y1 | published | ed0d594c}
\par
Diagrammada ikkita seriya bor: 2025 va 2026. Qaysi element farqlashga yordam beradi?
\begin{enumerate}[label=\Alph*),leftmargin=9mm]
\item Page Number
\item Headers
\item Legend
\item Footnote
\end{enumerate}
\begin{tcolorbox}[enhanced,breakable,colback=teal!4,colframe=teal!55!black,boxrule=0.45pt,arc=1.2mm,left=2.8mm,right=2.8mm,top=2mm,bottom=2mm]
\textbf{\color{teal}To'g'ri javob:} \textbf{C) Legend}
\par\smallskip
\textbf{Izoh:} Legenda seriya belgilari va ranglarini sharhlaydi.
\end{tcolorbox}
\vspace{2mm}\hrule height .2pt\vspace{2mm}

% DB 0bdd9d71-0a7f-0d56-75ac-12dd4eb83d96 status=published
\needspace{6\baselineskip}
\textbf{\color{ink}301.} {ootnotesizecolor{muted}Y1 | published | 0bdd9d71}
\par
Uch sinf natijasi 18, 24 va 12. Excel filtri qiymati 15 yoki undan yuqori sinflarni ko'rsatadi. Ko'rinayotgan qiymatlar yig'indisi?
\begin{enumerate}[label=\Alph*),leftmargin=9mm]
\item 30
\item 54
\item 36
\item 42
\end{enumerate}
\begin{tcolorbox}[enhanced,breakable,colback=teal!4,colframe=teal!55!black,boxrule=0.45pt,arc=1.2mm,left=2.8mm,right=2.8mm,top=2mm,bottom=2mm]
\textbf{\color{teal}To'g'ri javob:} \textbf{D) 42}
\par\smallskip
\textbf{Izoh:} Filtrdan so'ng 18 va 24 ko'rinadi, 12 yashiriladi; 18+24=42.
\end{tcolorbox}
\vspace{2mm}\hrule height .2pt\vspace{2mm}

% DB 1e17b68d-85c1-83b6-fa98-7083e9bfdf5c status=published
\needspace{6\baselineskip}
\textbf{\color{ink}302.} {ootnotesizecolor{muted}Y2 | published | 1e17b68d}
\par
Quyidagi tushunchalarni mazmuniga qarab bog'lang.  1) Ascending;  2) Descending;  3) Custom Filter;  Izohlar: A) Kamayish tartibi;  B) Bir necha mezon bilan cheklash;  C) O'sish tartibi;
\textbf{Chap ustun:}\begin{enumerate}[label=\arabic*),leftmargin=9mm]
\item Ascending
\item Descending
\item Custom Filter
\end{enumerate}
\textbf{O'ng ustun:}\begin{enumerate}[label=\Alph*),leftmargin=9mm]
\item Kamayish tartibi
\item Bir necha mezon bilan cheklash
\item O'sish tartibi
\end{enumerate}
\begin{tcolorbox}[enhanced,breakable,colback=teal!4,colframe=teal!55!black,boxrule=0.45pt,arc=1.2mm,left=2.8mm,right=2.8mm,top=2mm,bottom=2mm]
\textbf{\color{teal}To'g'ri javob:} \textbf{1-C, 2-A, 3-B}
\par\smallskip
\textbf{Izoh:} To'g'ri moslik: Ascending - O'sish tartibi, Descending - Kamayish tartibi, Custom Filter - Bir necha mezon bilan cheklash.
\end{tcolorbox}
\vspace{2mm}\hrule height .2pt\vspace{2mm}

% DB 70d4d41b-c8cf-b919-3fb8-dbcd796b0d90 status=published
\needspace{6\baselineskip}
\textbf{\color{ink}303.} {ootnotesizecolor{muted}Y1 | published | 70d4d41b}
\par
Ustunli diagrammada uch filial savdosi A=80, B=120, C=100. Agar har birining soni 20 tadan oshsa, qaysi fikr o'zgarmaydi?
\begin{enumerate}[label=\Alph*),leftmargin=9mm]
\item B filialning mutlaq savdo hajmi
\item Filiallar orasidagi absolyut farq va kamayish tartibi
\item Ularning umumiy yig'indisi
\item A filialning jami savdodagi foiz ulushi
\end{enumerate}
\begin{tcolorbox}[enhanced,breakable,colback=teal!4,colframe=teal!55!black,boxrule=0.45pt,arc=1.2mm,left=2.8mm,right=2.8mm,top=2mm,bottom=2mm]
\textbf{\color{teal}To'g'ri javob:} \textbf{B) Filiallar orasidagi absolyut farq va kamayish tartibi}
\par\smallskip
\textbf{Izoh:} Hamma ustunga bir xil 20 qo'shilganda ayirmalar saqlanadi va tartib ham o'zgarmaydi; ulush, yig'indi va qiymatlar o'zgaradi.
\end{tcolorbox}
\vspace{2mm}\hrule height .2pt\vspace{2mm}

% DB 973913f2-699c-5f8e-ee87-32414b63153e status=published
\needspace{6\baselineskip}
\textbf{\color{ink}304.} {ootnotesizecolor{muted}Y1 | published | 973913f2}
\par
2024-yilda 80 ta, 2025-yilda 100 ta, 2026-yilda 90 ta buyurtma bo'lgan. 2024-yilga nisbatan 2026-yil foiz o'zgarishi nechaga teng?
\begin{enumerate}[label=\Alph*),leftmargin=9mm]
\item 10\% kamayish
\item 25\% o'sish
\item 12.5\% o'sish
\item 10\% o'sish
\end{enumerate}
\begin{tcolorbox}[enhanced,breakable,colback=teal!4,colframe=teal!55!black,boxrule=0.45pt,arc=1.2mm,left=2.8mm,right=2.8mm,top=2mm,bottom=2mm]
\textbf{\color{teal}To'g'ri javob:} \textbf{C) 12.5\% o'sish}
\par\smallskip
\textbf{Izoh:} (90-80)/80\ensuremath{\times}100\%=12.5\%; o'rtadagi 100 birlamchi asos emas.
\end{tcolorbox}
\vspace{2mm}\hrule height .2pt\vspace{2mm}

% DB 9f6ab428-81ee-497d-3d14-49a8470eb06e status=published
\needspace{6\baselineskip}
\textbf{\color{ink}305.} {ootnotesizecolor{muted}Y3 | published | 9f6ab428}
\par
Faqat "Informatika" fanidan 70 balldan yuqori natijalarni diagrammada ko'rsatishning quyidagi amallarini tartiblang. Amal o'rni yakuniy, tekshiriladigan natijaga tayanadi.
\begin{enumerate}[label=\Alph*),leftmargin=9mm]
\item Saralangan manbadan diagramma qurish
\item Ma'lumotlarning Fan va Ball ustunlarini tekshirish
\item Fan=Informatika hamda Ball\textgreater{}70 filtr shartini qo'yish
\item Diagramma ustunlarini ko'rinadigan natijalar bilan solishtirish
\end{enumerate}
\begin{tcolorbox}[enhanced,breakable,colback=teal!4,colframe=teal!55!black,boxrule=0.45pt,arc=1.2mm,left=2.8mm,right=2.8mm,top=2mm,bottom=2mm]
\textbf{\color{teal}To'g'ri javob:} \textbf{B - C - A - D}
\par\smallskip
\textbf{Izoh:} Avval ma'lumot va mezon tekshiriladi, keyin filtrlanadi; diagramma shu ko'rinadigan yozuvlarga tayangan holda quriladi.
\end{tcolorbox}
\vspace{2mm}\hrule height .2pt\vspace{2mm}

% DB db8817c6-dfe0-49a2-ba4f-396a365e6a55 status=published
\needspace{6\baselineskip}
\textbf{\color{ink}306.} {ootnotesizecolor{muted}Y3 | published | db8817c6}
\par
Excel formulasining to'g'riligini tekshirish uchun modelga sinov qiymatlarini kiritib, qo'lda hisoblangan natija bilan solishtirish va zarur bo'lsa tuzatish amallarini tartiblang.
\begin{enumerate}[label=\Alph*),leftmargin=9mm]
\item Formula natijasi bilan solishtirish
\item Xato chiqsa formulani tuzatib qayta sinash
\item Modelga sinov qiymatlarini kiritish
\item Kutilgan natijani qo'lda aniqlash
\end{enumerate}
\begin{tcolorbox}[enhanced,breakable,colback=teal!4,colframe=teal!55!black,boxrule=0.45pt,arc=1.2mm,left=2.8mm,right=2.8mm,top=2mm,bottom=2mm]
\textbf{\color{teal}To'g'ri javob:} \textbf{C - D - A - B}
\par\smallskip
\textbf{Izoh:} Bosqichlar: Modelga sinov qiymatlarini kiritish \ensuremath{\rightarrow} Kutilgan natijani qo'lda aniqlash \ensuremath{\rightarrow} Formula natijasi bilan solishtirish \ensuremath{\rightarrow} Xato chiqsa formulani tuzatib qayta sinash.
\end{tcolorbox}
\vspace{2mm}\hrule height .2pt\vspace{2mm}

% DB e6110eed-1623-45a0-b542-b63d60c8f4c9 status=published
\needspace{6\baselineskip}
\textbf{\color{ink}307.} {ootnotesizecolor{muted}Y3 | published | e6110eed}
\par
MS Excelda sarlavhali jadvaldan faqat muayyan mezonga mos qatorlarni ko'rsatish uchun Filter buyrug'ini ishlatish kerak. Jadvalni tayyorlash, filtrlash va natijani tahlil qilish amallarini tartiblang.
\begin{enumerate}[label=\Alph*),leftmargin=9mm]
\item Butun jadvalda Filter qo'llash
\item Jadvaldagi sarlavha va ma'lumotlarni tekshirish
\item Kerakli mezonni berish
\item Mos qatorlar natijasini tahlil qilish
\end{enumerate}
\begin{tcolorbox}[enhanced,breakable,colback=teal!4,colframe=teal!55!black,boxrule=0.45pt,arc=1.2mm,left=2.8mm,right=2.8mm,top=2mm,bottom=2mm]
\textbf{\color{teal}To'g'ri javob:} \textbf{B - A - C - D}
\par\smallskip
\textbf{Izoh:} Avval sarlavha va ma'lumotlar tekshiriladi, keyin jadvalda Filter yoqiladi, mezon tanlanadi va qolgan qatorlar tahlil qilinadi. Tartib: 2 \ensuremath{\rightarrow} 1 \ensuremath{\rightarrow} 3 \ensuremath{\rightarrow} 4.
\end{tcolorbox}
\vspace{2mm}\hrule height .2pt\vspace{2mm}

% DB eb9a9e17-0d98-247c-d651-90730a06c28f status=published
\needspace{6\baselineskip}
\textbf{\color{ink}308.} {ootnotesizecolor{muted}Y2 | published | eb9a9e17}
\par
Moslashtiring.  1) Sort;  2) Filter;  3) Conditional Formatting;  Ta'riflar: A) Ko'rinishni shartga bog'lash;  B) Qatorlarni tartiblash;  C) Shartga mos kelmagan qatorlarni yashirish;
\textbf{Chap ustun:}\begin{enumerate}[label=\arabic*),leftmargin=9mm]
\item Sort
\item Filter
\item Conditional Formatting
\end{enumerate}
\textbf{O'ng ustun:}\begin{enumerate}[label=\Alph*),leftmargin=9mm]
\item Ko'rinishni shartga bog'lash
\item Qatorlarni tartiblash
\item Shartga mos kelmagan qatorlarni yashirish
\end{enumerate}
\begin{tcolorbox}[enhanced,breakable,colback=teal!4,colframe=teal!55!black,boxrule=0.45pt,arc=1.2mm,left=2.8mm,right=2.8mm,top=2mm,bottom=2mm]
\textbf{\color{teal}To'g'ri javob:} \textbf{1-B, 2-C, 3-A}
\par\smallskip
\textbf{Izoh:} Mos keladigan ta'riflar: 1-B, 2-C, 3-A.
\end{tcolorbox}
\vspace{2mm}\hrule height .2pt\vspace{2mm}

% DB 49b90af6-241b-4dba-be55-18b9337cefc3 status=review
\needspace{6\baselineskip}
\textbf{\color{ink}309.} {ootnotesizecolor{muted}Y1 | review | 49b90af6}
\par
50 balldan kichik kataklar qizil bo'lsin. Nima kerak?
\begin{enumerate}[label=\Alph*),leftmargin=9mm]
\item Data Validation \ensuremath{\rightarrow} Whole number \ensuremath{\rightarrow} Less than 50
\item Conditional Formatting, less than 50
\item Sort \& Filter \ensuremath{\rightarrow} Sort Smallest to Largest
\item Home \ensuremath{\rightarrow} Fill Color \ensuremath{\rightarrow} Red (all cells)
\end{enumerate}
\begin{tcolorbox}[enhanced,breakable,colback=teal!4,colframe=teal!55!black,boxrule=0.45pt,arc=1.2mm,left=2.8mm,right=2.8mm,top=2mm,bottom=2mm]
\textbf{\color{teal}To'g'ri javob:} \textbf{B) Conditional Formatting, less than 50}
\par\smallskip
\textbf{Izoh:} Excelda 50 dan kichik qiymatlarga avtomatik qizil rang berish uchun Conditional Formatting qoidasida Less Than 50 sharti belgilanadi. Data Validation kiritishni cheklaydi, Sort \& Filter tartibni o'zgartiradi, oddiy Fill Color esa shartga qarab avtomatik yangilanmaydi.
\end{tcolorbox}
\vspace{2mm}\hrule height .2pt\vspace{2mm}
\section{Microsoft PowerPoint}
\subsection{MS PowerPoint: taqdimot yaratish, slayd dizayni va obyektlar}

% DB 07210290-6049-99d7-5761-a6af60180b86 status=published
\needspace{6\baselineskip}
\textbf{\color{ink}310.} {ootnotesizecolor{muted}Y1 | published | 07210290}
\par
Theme nima?
\begin{enumerate}[label=\Alph*),leftmargin=9mm]
\item Printer ruxsati
\item Bitta harf uchun izoh
\item Fayl hajmi birligi
\item Rang, shrift va dizaynning umumiy uyg'un to'plami
\end{enumerate}
\begin{tcolorbox}[enhanced,breakable,colback=teal!4,colframe=teal!55!black,boxrule=0.45pt,arc=1.2mm,left=2.8mm,right=2.8mm,top=2mm,bottom=2mm]
\textbf{\color{teal}To'g'ri javob:} \textbf{D) Rang, shrift va dizaynning umumiy uyg'un to'plami}
\par\smallskip
\textbf{Izoh:} Theme umumiy taqdimot ko'rinishini birlashtiradi.
\end{tcolorbox}
\vspace{2mm}\hrule height .2pt\vspace{2mm}

% DB 14ccf817-3ac6-4bd6-ad86-d94891838e23 status=published
\needspace{6\baselineskip}
\textbf{\color{ink}311.} {ootnotesizecolor{muted}Y1 | published | 14ccf817}
\par
PowerPointda ko'plab slaydlardagi sarlavha shriftini va umumiy joylashuvini markazdan boshqarish kerak. Slide Master uchun eng to'g'ri tavsif qaysi?
\begin{enumerate}[label=\Alph*),leftmargin=9mm]
\item Taqdimotdagi bitta obyektga animatsiya berishni boshqaruvchi maxsus panel
\item Slaydlar orasidagi o'tish effektlari vaqtini belgilovchi maxsus menyu
\item Takroriy slayd dizayni va joylashuvlarini boshqaruvchi asos
\item Taqdimotning har bir slaydi uchun alohida ovoz yozuvchi vosita
\end{enumerate}
\begin{tcolorbox}[enhanced,breakable,colback=teal!4,colframe=teal!55!black,boxrule=0.45pt,arc=1.2mm,left=2.8mm,right=2.8mm,top=2mm,bottom=2mm]
\textbf{\color{teal}To'g'ri javob:} \textbf{C) Takroriy slayd dizayni va joylashuvlarini boshqaruvchi asos}
\par\smallskip
\textbf{Izoh:} Slide Master maket va umumiy dizayn uslublarini markazdan boshqaradi; slaydga qo'lda berilgan formatlar ayrim joylarda ustuvor bo'lishi mumkin. Animatsiya va Transition boshqa boshqaruv vositalari.
\end{tcolorbox}
\vspace{2mm}\hrule height .2pt\vspace{2mm}

% DB 20f0b927-0fe3-4088-a6da-c7df3895a2f2 status=published
\needspace{6\baselineskip}
\textbf{\color{ink}312.} {ootnotesizecolor{muted}Y3 | published | 20f0b927}
\par
Muayyan auditoriyaga mo'ljallangan PowerPoint taqdimotini ishlab chiqishda auditoriya, mazmun rejasi, slayd maketlari va obyektlarni joylashtirish ishlarini tartiblang.
\begin{enumerate}[label=\Alph*),leftmargin=9mm]
\item Taqdimot mazmunini rejalash
\item Auditoriyani aniqlash
\item Mos slayd Layoutlarini tanlash
\item Slaydlarga obyektlarni joylashtirish
\end{enumerate}
\begin{tcolorbox}[enhanced,breakable,colback=teal!4,colframe=teal!55!black,boxrule=0.45pt,arc=1.2mm,left=2.8mm,right=2.8mm,top=2mm,bottom=2mm]
\textbf{\color{teal}To'g'ri javob:} \textbf{B - A - C - D}
\par\smallskip
\textbf{Izoh:} Bosqichlar: Auditoriyani aniqlash \ensuremath{\rightarrow} Taqdimot mazmunini rejalash \ensuremath{\rightarrow} Mos slayd Layoutlarini tanlash \ensuremath{\rightarrow} Slaydlarga obyektlarni joylashtirish.
\end{tcolorbox}
\vspace{2mm}\hrule height .2pt\vspace{2mm}

% DB 668d0b30-cab3-586f-a999-92487c9ba448 status=published
\needspace{6\baselineskip}
\textbf{\color{ink}313.} {ootnotesizecolor{muted}Y1 | published | 668d0b30}
\par
MS PowerPointda slayd kolontituliga joriy sanani joylashtirish uchun qaysi yo'l to'g'ri?
\begin{enumerate}[label=\Alph*),leftmargin=9mm]
\item Insert \ensuremath{\rightarrow} Text \ensuremath{\rightarrow} Date \& Time; Slide bo'limida Date and time ni yoqish
\item Animations \ensuremath{\rightarrow} Animation Pane \ensuremath{\rightarrow} Date \& Time
\item Insert \ensuremath{\rightarrow} Text \ensuremath{\rightarrow} Text Box orqali sanani qo'lda yozish
\item Animations \ensuremath{\rightarrow} Object \ensuremath{\rightarrow} Slide \ensuremath{\rightarrow} Date and Time
\end{enumerate}
\begin{tcolorbox}[enhanced,breakable,colback=teal!4,colframe=teal!55!black,boxrule=0.45pt,arc=1.2mm,left=2.8mm,right=2.8mm,top=2mm,bottom=2mm]
\textbf{\color{teal}To'g'ri javob:} \textbf{A) Insert \ensuremath{\rightarrow} Text \ensuremath{\rightarrow} Date \& Time; Slide bo'limida Date and time ni yoqish}
\par\smallskip
\textbf{Izoh:} Insert lentasidagi Text guruhidan Date \& Time (Header and Footer oynasi) ochilib, Slide bo'limida Date and time belgilanadi.
\end{tcolorbox}
\vspace{2mm}\hrule height .2pt\vspace{2mm}

% DB a48bfa62-b7ff-f382-0f49-08711ef826c0 status=published
\needspace{6\baselineskip}
\textbf{\color{ink}314.} {ootnotesizecolor{muted}Y2 | published | a48bfa62}
\par
Quyidagi buyruqlarni vazifasi bilan moslashtiring.  1) New Slide;  2) Duplicate Slide;  3) Slide Sorter;  Izohlar: A) Slaydlar tartibini ko'rib almashtirish;  B) Slayd nusxasi;  C) Yangi slayd qo'shish;
\textbf{Chap ustun:}\begin{enumerate}[label=\arabic*),leftmargin=9mm]
\item New Slide
\item Duplicate Slide
\item Slide Sorter
\end{enumerate}
\textbf{O'ng ustun:}\begin{enumerate}[label=\Alph*),leftmargin=9mm]
\item Slaydlar tartibini ko'rib almashtirish
\item Slayd nusxasi
\item Yangi slayd qo'shish
\end{enumerate}
\begin{tcolorbox}[enhanced,breakable,colback=teal!4,colframe=teal!55!black,boxrule=0.45pt,arc=1.2mm,left=2.8mm,right=2.8mm,top=2mm,bottom=2mm]
\textbf{\color{teal}To'g'ri javob:} \textbf{1-C, 2-B, 3-A}
\par\smallskip
\textbf{Izoh:} To'g'ri moslik: New Slide - Yangi slayd qo'shish, Duplicate Slide - Slayd nusxasi, Slide Sorter - Slaydlar tartibini ko'rib almashtirish.
\end{tcolorbox}
\vspace{2mm}\hrule height .2pt\vspace{2mm}

% DB cb5bd226-45f3-b84b-c5f1-3f4d764fefef status=published
\needspace{6\baselineskip}
\textbf{\color{ink}315.} {ootnotesizecolor{muted}Y1 | published | cb5bd226}
\par
Slide Layout nima?
\begin{enumerate}[label=\Alph*),leftmargin=9mm]
\item Slayd faylining kengaytmasi
\item Faqat shrift turi
\item Slaydda matn va obyekt joylari maketi
\item Slayd namoyishining tezligi
\end{enumerate}
\begin{tcolorbox}[enhanced,breakable,colback=teal!4,colframe=teal!55!black,boxrule=0.45pt,arc=1.2mm,left=2.8mm,right=2.8mm,top=2mm,bottom=2mm]
\textbf{\color{teal}To'g'ri javob:} \textbf{C) Slaydda matn va obyekt joylari maketi}
\par\smallskip
\textbf{Izoh:} Layout slayd ichki tuzilmasini belgilaydi.
\end{tcolorbox}
\vspace{2mm}\hrule height .2pt\vspace{2mm}

% DB e2aa2b76-e216-a0dc-ac66-b0dd23e3501c status=published
\needspace{6\baselineskip}
\textbf{\color{ink}316.} {ootnotesizecolor{muted}Y1 | published | e2aa2b76}
\par
PowerPointning asosiy vazifasi?
\begin{enumerate}[label=\Alph*),leftmargin=9mm]
\item Elektron jadval formulasini hisoblash
\item Slaydli taqdimot yaratish va namoyish etish
\item Hujjatni antivirusdan o'tkazish
\item Operatsion tizimni o'rnatish
\end{enumerate}
\begin{tcolorbox}[enhanced,breakable,colback=teal!4,colframe=teal!55!black,boxrule=0.45pt,arc=1.2mm,left=2.8mm,right=2.8mm,top=2mm,bottom=2mm]
\textbf{\color{teal}To'g'ri javob:} \textbf{B) Slaydli taqdimot yaratish va namoyish etish}
\par\smallskip
\textbf{Izoh:} PowerPoint taqdimot dasturidir.
\end{tcolorbox}
\vspace{2mm}\hrule height .2pt\vspace{2mm}

% DB 030ab3fc-72af-cb11-74dd-d0c4a63c6d9d status=published
\needspace{6\baselineskip}
\textbf{\color{ink}317.} {ootnotesizecolor{muted}Y1 | published | 030ab3fc}
\par
Theme va Layout qanday farqlanadi?
\begin{enumerate}[label=\Alph*),leftmargin=9mm]
\item Theme faqat tezlikni, Layout shifrlashni bildiradi
\item Layout faqat animatsiya uchun
\item Ikkalasi ham faqat PDF yaratadi
\item Theme dizayn uslubi, Layout obyekt joylari tuzilishi
\end{enumerate}
\begin{tcolorbox}[enhanced,breakable,colback=teal!4,colframe=teal!55!black,boxrule=0.45pt,arc=1.2mm,left=2.8mm,right=2.8mm,top=2mm,bottom=2mm]
\textbf{\color{teal}To'g'ri javob:} \textbf{D) Theme dizayn uslubi, Layout obyekt joylari tuzilishi}
\par\smallskip
\textbf{Izoh:} Dizayn va joylashuv tushunchalari boshqa.
\end{tcolorbox}
\vspace{2mm}\hrule height .2pt\vspace{2mm}

% DB 0f81b5ab-dbe1-9207-580d-2b1e919e6f11 status=published
\needspace{6\baselineskip}
\textbf{\color{ink}318.} {ootnotesizecolor{muted}Y1 | published | 0f81b5ab}
\par
Group va Align farqi?
\begin{enumerate}[label=\Alph*),leftmargin=9mm]
\item Group birlashtirib boshqaradi, Align tekislaydi
\item Ikkalasi PDFga eksport qiladi
\item Group faqat font hajmini belgilaydi
\item Align hamma obyektni o'chiradi
\end{enumerate}
\begin{tcolorbox}[enhanced,breakable,colback=teal!4,colframe=teal!55!black,boxrule=0.45pt,arc=1.2mm,left=2.8mm,right=2.8mm,top=2mm,bottom=2mm]
\textbf{\color{teal}To'g'ri javob:} \textbf{A) Group birlashtirib boshqaradi, Align tekislaydi}
\par\smallskip
\textbf{Izoh:} Guruhlash va tekislash turli funksiyalar.
\end{tcolorbox}
\vspace{2mm}\hrule height .2pt\vspace{2mm}

% DB 67563f81-96f2-b441-b07a-b3757b4ad5c1 status=published
\needspace{6\baselineskip}
\textbf{\color{ink}319.} {ootnotesizecolor{muted}Y1 | published | 67563f81}
\par
Bring Forward va Send Backward farqi?
\begin{enumerate}[label=\Alph*),leftmargin=9mm]
\item Slaydni saqlash formatini tanlaydi
\item Obyektning ustma-ustlik tartibini oldinga yoki orqaga o'tkazadi
\item Printerga ko'rsatma yuboradi
\item Taqdimot ovozini oshiradi
\end{enumerate}
\begin{tcolorbox}[enhanced,breakable,colback=teal!4,colframe=teal!55!black,boxrule=0.45pt,arc=1.2mm,left=2.8mm,right=2.8mm,top=2mm,bottom=2mm]
\textbf{\color{teal}To'g'ri javob:} \textbf{B) Obyektning ustma-ustlik tartibini oldinga yoki orqaga o'tkazadi}
\par\smallskip
\textbf{Izoh:} Layer tartibi obyektlarning qaysi biri ustida ko'rinishini belgilaydi.
\end{tcolorbox}
\vspace{2mm}\hrule height .2pt\vspace{2mm}

% DB 98b3e6af-a785-f0b1-00de-5b2de7eb5423 status=published
\needspace{6\baselineskip}
\textbf{\color{ink}320.} {ootnotesizecolor{muted}Y1 | published | 98b3e6af}
\par
Presenter View va Slide Show ekranining farqi?
\begin{enumerate}[label=\Alph*),leftmargin=9mm]
\item Ma'ruzachida qaydlar ko'rinishi, tinglovchida joriy slayd bo'lishi mumkin
\item Tinglovchi har doim xom formulani ko'radi
\item Presenter View faqat Wordda
\item Ikkalasi ham printer sozlamasi
\end{enumerate}
\begin{tcolorbox}[enhanced,breakable,colback=teal!4,colframe=teal!55!black,boxrule=0.45pt,arc=1.2mm,left=2.8mm,right=2.8mm,top=2mm,bottom=2mm]
\textbf{\color{teal}To'g'ri javob:} \textbf{A) Ma'ruzachida qaydlar ko'rinishi, tinglovchida joriy slayd bo'lishi mumkin}
\par\smallskip
\textbf{Izoh:} Presenter View taqdimotchiga yordam beradi.
\end{tcolorbox}
\vspace{2mm}\hrule height .2pt\vspace{2mm}

% DB 9985bb62-397b-8ad5-4e48-1fcf73966231 status=published
\needspace{6\baselineskip}
\textbf{\color{ink}321.} {ootnotesizecolor{muted}Y2 | published | 9985bb62}
\par
Atamalar va ularning aniq ma'nolarini moslashtiring.  1) Theme;  2) Layout;  3) Slide Master;  Izohlar: A) Rang-shrift umumiy uslubi;  B) Slayd joy ajratgichlari;  C) Ko'p slaydga umumiy andoza;
\textbf{Chap ustun:}\begin{enumerate}[label=\arabic*),leftmargin=9mm]
\item Theme
\item Layout
\item Slide Master
\end{enumerate}
\textbf{O'ng ustun:}\begin{enumerate}[label=\Alph*),leftmargin=9mm]
\item Rang-shrift umumiy uslubi
\item Slayd joy ajratgichlari
\item Ko'p slaydga umumiy andoza
\end{enumerate}
\begin{tcolorbox}[enhanced,breakable,colback=teal!4,colframe=teal!55!black,boxrule=0.45pt,arc=1.2mm,left=2.8mm,right=2.8mm,top=2mm,bottom=2mm]
\textbf{\color{teal}To'g'ri javob:} \textbf{1-A, 2-B, 3-C}
\par\smallskip
\textbf{Izoh:} To'g'ri moslik: Theme - Rang-shrift umumiy uslubi, Layout - Slayd joy ajratgichlari, Slide Master - Ko'p slaydga umumiy andoza.
\end{tcolorbox}
\vspace{2mm}\hrule height .2pt\vspace{2mm}

% DB a26d7b61-2213-4a13-98c5-8e0503032027 status=published
\needspace{6\baselineskip}
\textbf{\color{ink}322.} {ootnotesizecolor{muted}Y3 | published | a26d7b61}
\par
PowerPoint taqdimotidagi barcha slaydlarda bir xil logotip ko'rinishi uchun uni Slide Master orqali kiritish va natijani tekshirish amallarini tartiblang.
\begin{enumerate}[label=\Alph*),leftmargin=9mm]
\item Slide Master rejimini ochish
\item Umumiy logotipni kiritish
\item Master rejimidan chiqish
\item Slaydlarda logotipni tekshirish
\end{enumerate}
\begin{tcolorbox}[enhanced,breakable,colback=teal!4,colframe=teal!55!black,boxrule=0.45pt,arc=1.2mm,left=2.8mm,right=2.8mm,top=2mm,bottom=2mm]
\textbf{\color{teal}To'g'ri javob:} \textbf{A - B - C - D}
\par\smallskip
\textbf{Izoh:} Bosqichlar: Slide Master rejimini ochish \ensuremath{\rightarrow} Umumiy logotipni kiritish \ensuremath{\rightarrow} Master rejimidan chiqish \ensuremath{\rightarrow} Slaydlarda logotipni tekshirish.
\end{tcolorbox}
\vspace{2mm}\hrule height .2pt\vspace{2mm}

% DB 050e0063-1663-4b46-adfc-f2e6cfcd9b41 status=published
\needspace{6\baselineskip}
\textbf{\color{ink}323.} {ootnotesizecolor{muted}Y2 | published | 050e0063}
\par
Amallarni izohi bilan moslang.  1) Align;  2) Group;  3) Bring to Front;  Izohlar: A) Obyektlarni birgalikda boshqarish;  B) Obyektni oldingi qatlamga chiqarish;  C) Obyektlarni tekislash;
\textbf{Chap ustun:}\begin{enumerate}[label=\arabic*),leftmargin=9mm]
\item Align
\item Group
\item Bring to Front
\end{enumerate}
\textbf{O'ng ustun:}\begin{enumerate}[label=\Alph*),leftmargin=9mm]
\item Obyektlarni birgalikda boshqarish
\item Obyektni oldingi qatlamga chiqarish
\item Obyektlarni tekislash
\end{enumerate}
\begin{tcolorbox}[enhanced,breakable,colback=teal!4,colframe=teal!55!black,boxrule=0.45pt,arc=1.2mm,left=2.8mm,right=2.8mm,top=2mm,bottom=2mm]
\textbf{\color{teal}To'g'ri javob:} \textbf{1-C, 2-A, 3-B}
\par\smallskip
\textbf{Izoh:} To'g'ri moslik: Align - Obyektlarni tekislash, Group - Obyektlarni birgalikda boshqarish, Bring to Front - Obyektni oldingi qatlamga chiqarish.
\end{tcolorbox}
\vspace{2mm}\hrule height .2pt\vspace{2mm}

% DB 2dce883a-bacf-4688-8f28-89578b02eae9 status=published
\needspace{6\baselineskip}
\textbf{\color{ink}324.} {ootnotesizecolor{muted}Y3 | published | 2dce883a}
\par
PowerPoint slaydiga rasm qo'shib, uning o'lchamini moslash, Align orqali tekislash va namoyishda tekshirish jarayonini tartiblang.
\begin{enumerate}[label=\Alph*),leftmargin=9mm]
\item Rasm hajmini moslash
\item Align bilan tekislash
\item Slide Show rejimida tekshirish
\item Slaydga rasm joylashtirish
\end{enumerate}
\begin{tcolorbox}[enhanced,breakable,colback=teal!4,colframe=teal!55!black,boxrule=0.45pt,arc=1.2mm,left=2.8mm,right=2.8mm,top=2mm,bottom=2mm]
\textbf{\color{teal}To'g'ri javob:} \textbf{D - A - B - C}
\par\smallskip
\textbf{Izoh:} Bosqichlar: Slaydga rasm joylashtirish \ensuremath{\rightarrow} Rasm hajmini moslash \ensuremath{\rightarrow} Align bilan tekislash \ensuremath{\rightarrow} Slide Show rejimida tekshirish.
\end{tcolorbox}
\vspace{2mm}\hrule height .2pt\vspace{2mm}

% DB 48709bed-2d6c-1be0-acb2-efaaccba9314 status=published
\needspace{6\baselineskip}
\textbf{\color{ink}325.} {ootnotesizecolor{muted}Y1 | published | 48709bed}
\par
Slayddagi matn va rasm uchun tayyor joy ajratilishi kerak. Qaysi Layout?
\begin{enumerate}[label=\Alph*),leftmargin=9mm]
\item Blank faqat
\item Section Header faqat
\item Title and Content
\item Title Slide faqat
\end{enumerate}
\begin{tcolorbox}[enhanced,breakable,colback=teal!4,colframe=teal!55!black,boxrule=0.45pt,arc=1.2mm,left=2.8mm,right=2.8mm,top=2mm,bottom=2mm]
\textbf{\color{teal}To'g'ri javob:} \textbf{C) Title and Content}
\par\smallskip
\textbf{Izoh:} Bu maket sarlavha va kontent joyini beradi.
\end{tcolorbox}
\vspace{2mm}\hrule height .2pt\vspace{2mm}

% DB 9b14244e-d898-f1b3-de6b-017fddc4f475 status=published
\needspace{6\baselineskip}
\textbf{\color{ink}326.} {ootnotesizecolor{muted}Y1 | published | 9b14244e}
\par
Taqdimotga video qo'shilgandan keyin uni avtomatik boshlash va faqat slayd ichida ko'rsatish kerak. Qaysi juft sozlama muhim?
\begin{enumerate}[label=\Alph*),leftmargin=9mm]
\item Transitions \ensuremath{\rightarrow} Duration va Page Margins
\item Review \ensuremath{\rightarrow} Track Changes va Compare
\item Insert \ensuremath{\rightarrow} Table va Word Count
\item Playback \ensuremath{\rightarrow} Start Automatically va video o'lchami/joylashuvini tekshirish
\end{enumerate}
\begin{tcolorbox}[enhanced,breakable,colback=teal!4,colframe=teal!55!black,boxrule=0.45pt,arc=1.2mm,left=2.8mm,right=2.8mm,top=2mm,bottom=2mm]
\textbf{\color{teal}To'g'ri javob:} \textbf{D) Playback \ensuremath{\rightarrow} Start Automatically va video o'lchami/joylashuvini tekshirish}
\par\smallskip
\textbf{Izoh:} Video ishga tushirish xulqi Playbackda belgilanadi; ko'rinish slayd ichidagi obyekt xossasi bilan boshqariladi.
\end{tcolorbox}
\vspace{2mm}\hrule height .2pt\vspace{2mm}

% DB b5e454f3-0640-f09f-4730-8224c271ed5b status=published
\needspace{6\baselineskip}
\textbf{\color{ink}327.} {ootnotesizecolor{muted}Y1 | published | b5e454f3}
\par
Yangi titul slayd uchun mos Layout qaysi?
\begin{enumerate}[label=\Alph*),leftmargin=9mm]
\item Two Content faqat
\item Title Slide
\item Blank faqat
\item Comparison faqat
\end{enumerate}
\begin{tcolorbox}[enhanced,breakable,colback=teal!4,colframe=teal!55!black,boxrule=0.45pt,arc=1.2mm,left=2.8mm,right=2.8mm,top=2mm,bottom=2mm]
\textbf{\color{teal}To'g'ri javob:} \textbf{B) Title Slide}
\par\smallskip
\textbf{Izoh:} Title Slide bosh sahifa uchun mo'ljallangan.
\end{tcolorbox}
\vspace{2mm}\hrule height .2pt\vspace{2mm}

% DB e2de47f9-ef4c-c89f-09bc-1b0cae065885 status=published
\needspace{6\baselineskip}
\textbf{\color{ink}328.} {ootnotesizecolor{muted}Y1 | published | e2de47f9}
\par
Taqdimotning har bir slaydida logotip bir xil ko'rinsin. Nima qulay?
\begin{enumerate}[label=\Alph*),leftmargin=9mm]
\item Faqat F5
\item Faqat Spell Check
\item Har safar qo'lda Crop
\item Slide Master
\end{enumerate}
\begin{tcolorbox}[enhanced,breakable,colback=teal!4,colframe=teal!55!black,boxrule=0.45pt,arc=1.2mm,left=2.8mm,right=2.8mm,top=2mm,bottom=2mm]
\textbf{\color{teal}To'g'ri javob:} \textbf{D) Slide Master}
\par\smallskip
\textbf{Izoh:} Master takroriy dizaynni nazorat qiladi.
\end{tcolorbox}
\vspace{2mm}\hrule height .2pt\vspace{2mm}

% DB facbb0c0-bc15-483f-bc02-823c64e16617 status=published
\needspace{6\baselineskip}
\textbf{\color{ink}329.} {ootnotesizecolor{muted}Y1 | published | facbb0c0}
\par
PowerPoint dasturida \textbf{"Zapis' ekrana" (Screen Recording)} funksiyasi bilan ishlashning to'g'ri ketma-ketligini tanlang. (Manba: Word hujjatidagi 1-rasm.)
\begin{enumerate}[label=\Alph*),leftmargin=9mm]
\item Vstavka \ensuremath{\rightarrow} Zvuk \ensuremath{\rightarrow} Zapis' ekrana \ensuremath{\rightarrow} Zapisat'
\item Glavnaya \ensuremath{\rightarrow} Video \ensuremath{\rightarrow} Zapis' ekrana \ensuremath{\rightarrow} Vybrat' oblast' \ensuremath{\rightarrow} Zapisat'
\item Glavnaya \ensuremath{\rightarrow} Ob'ekt \ensuremath{\rightarrow} Zapis' ekrana \ensuremath{\rightarrow} Vybrat' oblast' \ensuremath{\rightarrow} Zapisat'
\item Vstavka \ensuremath{\rightarrow} Zapis' ekrana \ensuremath{\rightarrow} Vybrat' oblast' \ensuremath{\rightarrow} Zvuk \ensuremath{\rightarrow} Zapisat'
\end{enumerate}
\par\smallskip\textit{Rasmli savol: }\href{https://github.com/sarvar9417/attestatsiya/blob/main/public/kuzgi-demo/demo_q01_powerpoint.svg}{Rasmni onlayn ochish}
\begin{tcolorbox}[enhanced,breakable,colback=teal!4,colframe=teal!55!black,boxrule=0.45pt,arc=1.2mm,left=2.8mm,right=2.8mm,top=2mm,bottom=2mm]
\textbf{\color{teal}To'g'ri javob:} \textbf{D) Vstavka \ensuremath{\rightarrow} Zapis' ekrana \ensuremath{\rightarrow} Vybrat' oblast' \ensuremath{\rightarrow} Zvuk \ensuremath{\rightarrow} Zapisat'}
\par\smallskip
\textbf{Izoh:} PowerPoint oynasida Insert (Vstavka) lentasidagi Screen Recording buyrug'i bosiladi, yoziladigan soha tanlanadi, kerak bo'lsa ovoz sozlanadi va Record orqali yozuv boshlanadi.
\end{tcolorbox}
\vspace{2mm}\hrule height .2pt\vspace{2mm}

% DB 070a2f11-f88d-495e-b109-8d6188cf4ee3 status=published
\needspace{6\baselineskip}
\textbf{\color{ink}330.} {ootnotesizecolor{muted}Y3 | published | 070a2f11}
\par
PowerPoint taqdimotidagi menyu obyektidan ma'lum slaydga o'tuvchi ichki gipermurojaat yaratish va uni namoyishda tekshirish amallarini tartiblang.
\begin{enumerate}[label=\Alph*),leftmargin=9mm]
\item Nishon slaydni belgilash
\item Insert Hyperlink oynasini ochish
\item Taqdimotda yo'naltiruvchi obyektni tanlash
\item Namoyishda havolani tekshirish
\end{enumerate}
\begin{tcolorbox}[enhanced,breakable,colback=teal!4,colframe=teal!55!black,boxrule=0.45pt,arc=1.2mm,left=2.8mm,right=2.8mm,top=2mm,bottom=2mm]
\textbf{\color{teal}To'g'ri javob:} \textbf{C - B - A - D}
\par\smallskip
\textbf{Izoh:} Bosqichlar: Taqdimotda yo'naltiruvchi obyektni tanlash \ensuremath{\rightarrow} Insert Hyperlink oynasini ochish \ensuremath{\rightarrow} Nishon slaydni belgilash \ensuremath{\rightarrow} Namoyishda havolani tekshirish.
\end{tcolorbox}
\vspace{2mm}\hrule height .2pt\vspace{2mm}

% DB 45232c4c-22bc-0c54-487b-9e44e981bc8c status=published
\needspace{6\baselineskip}
\textbf{\color{ink}331.} {ootnotesizecolor{muted}Y1 | published | 45232c4c}
\par
Diagramma slaydda nimani aniq ko'rsatishi kerak?
\begin{enumerate}[label=\Alph*),leftmargin=9mm]
\item Sarlavha, o'q va birliklari bilan ma'no
\item Faqat slayd raqami
\item Faqat audio boshi
\item Faqat rangli bezak
\end{enumerate}
\begin{tcolorbox}[enhanced,breakable,colback=teal!4,colframe=teal!55!black,boxrule=0.45pt,arc=1.2mm,left=2.8mm,right=2.8mm,top=2mm,bottom=2mm]
\textbf{\color{teal}To'g'ri javob:} \textbf{A) Sarlavha, o'q va birliklari bilan ma'no}
\par\smallskip
\textbf{Izoh:} Vizual ma'lumot kontekst bilan o'qiladi.
\end{tcolorbox}
\vspace{2mm}\hrule height .2pt\vspace{2mm}

% DB 591c1897-1d08-613d-270c-71ae5260c8fb status=published
\needspace{6\baselineskip}
\textbf{\color{ink}332.} {ootnotesizecolor{muted}Y3 | published | 591c1897}
\par
Yangi PowerPoint taqdimoti yaratilyapti. Talab: umumiy sarlavha formati \textbf{slaydlar yaratilishidan oldin} Slide Master'da sozlansin, so'ng 12 slayd qo'shilib yakuniy namoyish tekshirilsin. Amallarni tartiblang.
\begin{enumerate}[label=\Alph*),leftmargin=9mm]
\item Kerakli maketlar asosida 12 ta slayd yaratib mazmunni joylash
\item Yangi bo'sh taqdimot ochish
\item Hali slaydlar ko'paytirilmasidan Slide Master'da sarlavha shrifti va maketni sozlash
\item Slide Show rejimida sarlavha ko'rinishi bir xilligini tekshirish
\end{enumerate}
\begin{tcolorbox}[enhanced,breakable,colback=teal!4,colframe=teal!55!black,boxrule=0.45pt,arc=1.2mm,left=2.8mm,right=2.8mm,top=2mm,bottom=2mm]
\textbf{\color{teal}To'g'ri javob:} \textbf{B - C - A - D}
\par\smallskip
\textbf{Izoh:} Bo'sh taqdimot ochiladi \ensuremath{\rightarrow} barcha yangi slaydlar uchun Slide Master maketi avval sozlanadi \ensuremath{\rightarrow} 12 slayd shu maketlardan yaratiladi \ensuremath{\rightarrow} yakunda Slide Show'da ko'rinish tekshiriladi. Savolning maxsus sharti maketni slaydlar yaratilishidan oldin sozlashni talab qiladi.
\end{tcolorbox}
\vspace{2mm}\hrule height .2pt\vspace{2mm}

% DB 5f451a7a-fb86-dd90-6998-69b19eb2a189 status=published
\needspace{6\baselineskip}
\textbf{\color{ink}333.} {ootnotesizecolor{muted}Y1 | published | 5f451a7a}
\par
12 slaydda bir xil rang va shrift uslubini saqlash kerak. Nima tanlanadi?
\begin{enumerate}[label=\Alph*),leftmargin=9mm]
\item Theme yoki Slide Master sozlamalari
\item Har slaydga tasodifiy bezak
\item Faqat Crop
\item Faqat CTRL+P
\end{enumerate}
\begin{tcolorbox}[enhanced,breakable,colback=teal!4,colframe=teal!55!black,boxrule=0.45pt,arc=1.2mm,left=2.8mm,right=2.8mm,top=2mm,bottom=2mm]
\textbf{\color{teal}To'g'ri javob:} \textbf{A) Theme yoki Slide Master sozlamalari}
\par\smallskip
\textbf{Izoh:} Umumiy dizayn muvofiqligi bir markazda boshqariladi.
\end{tcolorbox}
\vspace{2mm}\hrule height .2pt\vspace{2mm}

% DB d2e2cdd9-ab62-7604-5367-651cc945264b status=published
\needspace{6\baselineskip}
\textbf{\color{ink}334.} {ootnotesizecolor{muted}Y1 | published | d2e2cdd9}
\par
Matn katta auditoriya uchun mayda bo'lib qoldi. Nima muhim?
\begin{enumerate}[label=\Alph*),leftmargin=9mm]
\item Shaklni maydaroq qilish
\item Fon va matnni bir xil rang qilish
\item Barcha matnni bitta slaydga joylash
\item Kattaroq o'qiladigan shrift va qisqaroq mazmun
\end{enumerate}
\begin{tcolorbox}[enhanced,breakable,colback=teal!4,colframe=teal!55!black,boxrule=0.45pt,arc=1.2mm,left=2.8mm,right=2.8mm,top=2mm,bottom=2mm]
\textbf{\color{teal}To'g'ri javob:} \textbf{D) Kattaroq o'qiladigan shrift va qisqaroq mazmun}
\par\smallskip
\textbf{Izoh:} Ko'rinarlilik taqdimot sifatining birinchi mezoni.
\end{tcolorbox}
\vspace{2mm}\hrule height .2pt\vspace{2mm}

% DB e0a046ff-5371-a01b-f9a3-090f2732dfb9 status=published
\needspace{6\baselineskip}
\textbf{\color{ink}335.} {ootnotesizecolor{muted}Y2 | published | e0a046ff}
\par
Ushbu vaziyatdagi vositalar va vazifalarni moslashtiring.  1) Hyperlink;  2) Speaker Notes;  3) Presenter View;  Izohlar: A) Slayd yoki veb-manzilga o'tish;  B) Ma'ruzachi uchun izohlar;  C) Ma'ruzachi uchun boshqaruv ko'rinishi;
\textbf{Chap ustun:}\begin{enumerate}[label=\arabic*),leftmargin=9mm]
\item Hyperlink
\item Speaker Notes
\item Presenter View
\end{enumerate}
\textbf{O'ng ustun:}\begin{enumerate}[label=\Alph*),leftmargin=9mm]
\item Slayd yoki veb-manzilga o'tish
\item Ma'ruzachi uchun izohlar
\item Ma'ruzachi uchun boshqaruv ko'rinishi
\end{enumerate}
\begin{tcolorbox}[enhanced,breakable,colback=teal!4,colframe=teal!55!black,boxrule=0.45pt,arc=1.2mm,left=2.8mm,right=2.8mm,top=2mm,bottom=2mm]
\textbf{\color{teal}To'g'ri javob:} \textbf{1-A, 2-B, 3-C}
\par\smallskip
\textbf{Izoh:} To'g'ri moslik: Hyperlink - Slayd yoki veb-manzilga o'tish, Speaker Notes - Ma'ruzachi uchun izohlar, Presenter View - Ma'ruzachi uchun boshqaruv ko'rinishi.
\end{tcolorbox}
\vspace{2mm}\hrule height .2pt\vspace{2mm}

% DB fec8ad1a-1290-810e-df83-60363afb32e7 status=published
\needspace{6\baselineskip}
\textbf{\color{ink}336.} {ootnotesizecolor{muted}Y1 | published | fec8ad1a}
\par
Taqdimot 10 slayddan iborat; uchinchi slaydni yettinchi o'ringa o'tkazish uchun?
\begin{enumerate}[label=\Alph*),leftmargin=9mm]
\item Excel Filter qo'llash
\item Slide Sorterda tortib joylash
\item Footer qo'yish
\item Faylni DOCX qilish
\end{enumerate}
\begin{tcolorbox}[enhanced,breakable,colback=teal!4,colframe=teal!55!black,boxrule=0.45pt,arc=1.2mm,left=2.8mm,right=2.8mm,top=2mm,bottom=2mm]
\textbf{\color{teal}To'g'ri javob:} \textbf{B) Slide Sorterda tortib joylash}
\par\smallskip
\textbf{Izoh:} Slaydlar tartibi Sorterda qulay o'zgartiriladi.
\end{tcolorbox}
\vspace{2mm}\hrule height .2pt\vspace{2mm}

% DB 07836d91-188a-6fbb-1ffb-37cfe5d06fb0 status=published
\needspace{6\baselineskip}
\textbf{\color{ink}337.} {ootnotesizecolor{muted}Y1 | published | 07836d91}
\par
Video joylashtirilgan interaktiv PPTXni PDF sifatida yuborishdi. Qaysi ikkita muhim cheklovni tekshirish lozim?
\begin{enumerate}[label=\Alph*),leftmargin=9mm]
\item PDF videoni barcha qurilmada avtomatik bog'lab beradi
\item PDF barcha PPTX animatsiyalarini doim aynan ko'rsatadi
\item Animatsiya/navigatsiyaning statiklashishi va videoning PDFda PPTXdagi kabi ishlamasligi
\item Slide Master PDFni har doim tahrirlanuvchi PPTXga qaytaradi
\end{enumerate}
\begin{tcolorbox}[enhanced,breakable,colback=teal!4,colframe=teal!55!black,boxrule=0.45pt,arc=1.2mm,left=2.8mm,right=2.8mm,top=2mm,bottom=2mm]
\textbf{\color{teal}To'g'ri javob:} \textbf{C) Animatsiya/navigatsiyaning statiklashishi va videoning PDFda PPTXdagi kabi ishlamasligi}
\par\smallskip
\textbf{Izoh:} PDF statik hujjat; interaktiv elementlar va multimedia mosligini alohida nazorat qilish kerak.
\end{tcolorbox}
\vspace{2mm}\hrule height .2pt\vspace{2mm}

% DB 645da3d7-29f2-d785-b402-884ea6ea7277 status=published
\needspace{6\baselineskip}
\textbf{\color{ink}338.} {ootnotesizecolor{muted}Y2 | published | 645da3d7}
\par
Quyidagi tushunchalarni mazmuniga qarab bog'lang.  1) Insert Audio;  2) Insert Video;  3) Screen Recording;  Izohlar: A) Video faylini kiritish;  B) Ekrandagi harakatni video sifatida yozish;  C) Ovoz faylini kiritish;
\textbf{Chap ustun:}\begin{enumerate}[label=\arabic*),leftmargin=9mm]
\item Insert Audio
\item Insert Video
\item Screen Recording
\end{enumerate}
\textbf{O'ng ustun:}\begin{enumerate}[label=\Alph*),leftmargin=9mm]
\item Video faylini kiritish
\item Ekrandagi harakatni video sifatida yozish
\item Ovoz faylini kiritish
\end{enumerate}
\begin{tcolorbox}[enhanced,breakable,colback=teal!4,colframe=teal!55!black,boxrule=0.45pt,arc=1.2mm,left=2.8mm,right=2.8mm,top=2mm,bottom=2mm]
\textbf{\color{teal}To'g'ri javob:} \textbf{1-C, 2-A, 3-B}
\par\smallskip
\textbf{Izoh:} To'g'ri moslik: Insert Audio - Ovoz faylini kiritish, Insert Video - Video faylini kiritish, Screen Recording - Ekrandagi harakatni video sifatida yozish.
\end{tcolorbox}
\vspace{2mm}\hrule height .2pt\vspace{2mm}

% DB 8e48d2b9-e8d8-1ec9-ad80-70a76cb5785a status=published
\needspace{6\baselineskip}
\textbf{\color{ink}339.} {ootnotesizecolor{muted}Y2 | published | 8e48d2b9}
\par
Moslashtiring.  1) Theme;  2) Layout;  3) Slide Master;  Ta'riflar: A) Joy ajratgichlar tuzilishi;  B) Umumiy taqdimot dizayn asosi;  C) Umumiy dizayn uslubi;
\textbf{Chap ustun:}\begin{enumerate}[label=\arabic*),leftmargin=9mm]
\item Theme
\item Layout
\item Slide Master
\end{enumerate}
\textbf{O'ng ustun:}\begin{enumerate}[label=\Alph*),leftmargin=9mm]
\item Joy ajratgichlar tuzilishi
\item Umumiy taqdimot dizayn asosi
\item Umumiy dizayn uslubi
\end{enumerate}
\begin{tcolorbox}[enhanced,breakable,colback=teal!4,colframe=teal!55!black,boxrule=0.45pt,arc=1.2mm,left=2.8mm,right=2.8mm,top=2mm,bottom=2mm]
\textbf{\color{teal}To'g'ri javob:} \textbf{1-C, 2-A, 3-B}
\par\smallskip
\textbf{Izoh:} Mos keladigan ta'riflar: 1-C, 2-A, 3-B.
\end{tcolorbox}
\vspace{2mm}\hrule height .2pt\vspace{2mm}

% DB 97ef9947-94db-4598-a36e-602a2a7d15af status=published
\needspace{6\baselineskip}
\textbf{\color{ink}340.} {ootnotesizecolor{muted}Y3 | published | 97ef9947}
\par
PowerPoint taqdimotiga videofayl joylashtirib, Playback boshlanish sozlamasini tanlash va boshqa kompyuterda ishlashini tekshirish amallarini tartiblang.
\begin{enumerate}[label=\Alph*),leftmargin=9mm]
\item Boshqa kompyuterda ishlashini tekshirish
\item Playback boshlash usulini belgilash
\item Insert Video orqali faylni tanlash
\item Slaydga videoni joylashtirish
\end{enumerate}
\begin{tcolorbox}[enhanced,breakable,colback=teal!4,colframe=teal!55!black,boxrule=0.45pt,arc=1.2mm,left=2.8mm,right=2.8mm,top=2mm,bottom=2mm]
\textbf{\color{teal}To'g'ri javob:} \textbf{C - D - B - A}
\par\smallskip
\textbf{Izoh:} Bosqichlar: Insert Video orqali faylni tanlash \ensuremath{\rightarrow} Slaydga videoni joylashtirish \ensuremath{\rightarrow} Playback boshlash usulini belgilash \ensuremath{\rightarrow} Boshqa kompyuterda ishlashini tekshirish.
\end{tcolorbox}
\vspace{2mm}\hrule height .2pt\vspace{2mm}

% DB a21affa6-665a-3259-ddb3-14f930c551d7 status=published
\needspace{6\baselineskip}
\textbf{\color{ink}341.} {ootnotesizecolor{muted}Y1 | published | a21affa6}
\par
50 slaydli taqdimotning har birida logotip bir xil ko'rinsin, ammo 3 ta slaydning Layouti boshqacha. Qaysi usul yaxshiroq?
\begin{enumerate}[label=\Alph*),leftmargin=9mm]
\item Har slaydga logotipni qo'lda 50 marta joylashtirish
\item Slide Masterdagi umumiy dizaynni sozlash, maxsus Layoutlarni alohida boshqarish
\item Transitions Apply To All bilan logotip qo'shish
\item Barcha slaydni JPGga aylantirish
\end{enumerate}
\begin{tcolorbox}[enhanced,breakable,colback=teal!4,colframe=teal!55!black,boxrule=0.45pt,arc=1.2mm,left=2.8mm,right=2.8mm,top=2mm,bottom=2mm]
\textbf{\color{teal}To'g'ri javob:} \textbf{B) Slide Masterdagi umumiy dizaynni sozlash, maxsus Layoutlarni alohida boshqarish}
\par\smallskip
\textbf{Izoh:} Master umumiy obyektni, Layout esa slayd ichidagi joylashuvni markazlashtiradi.
\end{tcolorbox}
\vspace{2mm}\hrule height .2pt\vspace{2mm}

% DB b7879ba6-da24-2642-6db7-2e20c166e880 status=published
\needspace{6\baselineskip}
\textbf{\color{ink}342.} {ootnotesizecolor{muted}Y1 | published | b7879ba6}
\par
PowerPointda 50 ta slayd Slide Master asosida yaratilgan. Master sarlavha shrifti o'zgartirildi, ammo 8 slayddagi sarlavhalar avvalgidek qoldi - ular qo'lda formatlangan. Qaysi yechim umumiy dizaynni qayta tiklaydi?
\begin{enumerate}[label=\Alph*),leftmargin=9mm]
\item Faqat Rehearse Timings ni qayta ishga tushirish
\item Barcha slaydlarga Transition \ensuremath{\rightarrow} Fade berish
\item Master'dagi umumiy shriftni sozlab, lokal formatlashni bekor qilish yoki kerakli slaydlarda Reset/Layout bilan maketni qayta qo'llash
\item Animation Pane'da har bir effektni o'chirish
\end{enumerate}
\begin{tcolorbox}[enhanced,breakable,colback=teal!4,colframe=teal!55!black,boxrule=0.45pt,arc=1.2mm,left=2.8mm,right=2.8mm,top=2mm,bottom=2mm]
\textbf{\color{teal}To'g'ri javob:} \textbf{C) Master'dagi umumiy shriftni sozlab, lokal formatlashni bekor qilish yoki kerakli slaydlarda Reset/Layout bilan maketni qayta qo'llash}
\par\smallskip
\textbf{Izoh:} Slide Master umumiy formatni belgilaydi, lekin ayrim slaydlardagi qo'lda qo'yilgan lokal formatlar ustuvor bo'lishi mumkin. Muammo qayta Master sozlash bilan birga lokal formatlashni tozalash, Reset yoki layoutni qayta qo'llash orqali hal etiladi.
\end{tcolorbox}
\vspace{2mm}\hrule height .2pt\vspace{2mm}

% DB eecd741a-1d91-4370-95c5-e43b8f0e6f58 status=published
\needspace{6\baselineskip}
\textbf{\color{ink}343.} {ootnotesizecolor{muted}Y3 | published | eecd741a}
\par
PowerPointda muayyan auditoriya uchun mavzuli taqdimot ishlab chiqishda maqsadni aniqlashdan slaydlar rejasini tuzish, Theme/Layout tanlash va namoyishni tekshirishgacha amallarni tartiblang.
\begin{enumerate}[label=\Alph*),leftmargin=9mm]
\item Auditoriya va maqsadni aniqlash
\item Mos Layout/Theme tanlash
\item Obyektlarni joylashtirib namoyishni tekshirish
\item Slaydlar rejasini tuzish
\end{enumerate}
\begin{tcolorbox}[enhanced,breakable,colback=teal!4,colframe=teal!55!black,boxrule=0.45pt,arc=1.2mm,left=2.8mm,right=2.8mm,top=2mm,bottom=2mm]
\textbf{\color{teal}To'g'ri javob:} \textbf{A - D - B - C}
\par\smallskip
\textbf{Izoh:} Auditoriya va taqdimot maqsadi aniqlanadi, slaydlar rejasi tuziladi, mos dizayn tanlanadi, obyektlar joylashtirilib natija namoyishda tekshiriladi. Tartib: 1 \ensuremath{\rightarrow} 4 \ensuremath{\rightarrow} 2 \ensuremath{\rightarrow} 3.
\end{tcolorbox}
\vspace{2mm}\hrule height .2pt\vspace{2mm}
\subsection{MS PowerPoint: animatsiya, transition va namoyish boshqaruvi}

% DB 04d94c39-f014-501b-ae7e-b2188f4f0dc5 status=published
\needspace{6\baselineskip}
\textbf{\color{ink}344.} {ootnotesizecolor{muted}Y1 | published | 04d94c39}
\par
PowerPoint animatsiya galereyasidagi guruhlarni moslang: 1) yashil yulduzli effektlar; 2) sariq yulduzli effektlar; 3) qizil yulduzli effektlar; 4) chiziq, yoy, burilish va shakl trayektoriyalari. a) Exit; b) Motion Paths; c) Entrance; d) Emphasis.
\begin{enumerate}[label=\Alph*),leftmargin=9mm]
\item 1-a; 2-b; 3-c; 4-d
\item 1-a; 2-d; 3-b; 4-c
\item 1-c; 2-d; 3-a; 4-b
\item 1-c; 2-b; 3-a; 4-d
\end{enumerate}
\begin{tcolorbox}[enhanced,breakable,colback=teal!4,colframe=teal!55!black,boxrule=0.45pt,arc=1.2mm,left=2.8mm,right=2.8mm,top=2mm,bottom=2mm]
\textbf{\color{teal}To'g'ri javob:} \textbf{C) 1-c; 2-d; 3-a; 4-b}
\par\smallskip
\textbf{Izoh:} PowerPointda yashil belgilar Entrance, sariq belgilar Emphasis, qizil belgilar Exit, chiziq va shaklli trayektoriyalar esa Motion Paths guruhidir.
\end{tcolorbox}
\vspace{2mm}\hrule height .2pt\vspace{2mm}

% DB 09c50277-b998-1a85-8ffc-7d21e6b9b084 status=published
\needspace{6\baselineskip}
\textbf{\color{ink}345.} {ootnotesizecolor{muted}Y2 | published | 09c50277}
\par
Quyidagi buyruqlarni vazifasi bilan moslashtiring.  1) Entrance;  2) Emphasis;  3) Exit;  Izohlar: A) Obyektni paydo qilish;  B) Ko'rinib turgan obyektga urg'u berish;  C) Obyektni chiqib ketgandek ko'rsatish;
\textbf{Chap ustun:}\begin{enumerate}[label=\arabic*),leftmargin=9mm]
\item Entrance
\item Emphasis
\item Exit
\end{enumerate}
\textbf{O'ng ustun:}\begin{enumerate}[label=\Alph*),leftmargin=9mm]
\item Obyektni paydo qilish
\item Ko'rinib turgan obyektga urg'u berish
\item Obyektni chiqib ketgandek ko'rsatish
\end{enumerate}
\begin{tcolorbox}[enhanced,breakable,colback=teal!4,colframe=teal!55!black,boxrule=0.45pt,arc=1.2mm,left=2.8mm,right=2.8mm,top=2mm,bottom=2mm]
\textbf{\color{teal}To'g'ri javob:} \textbf{1-A, 2-B, 3-C}
\par\smallskip
\textbf{Izoh:} To'g'ri moslik: Entrance - Obyektni paydo qilish, Emphasis - Ko'rinib turgan obyektga urg'u berish, Exit - Obyektni chiqib ketgandek ko'rsatish.
\end{tcolorbox}
\vspace{2mm}\hrule height .2pt\vspace{2mm}

% DB 564ed242-b9e2-0be5-e448-b90723773448 status=published
\needspace{6\baselineskip}
\textbf{\color{ink}346.} {ootnotesizecolor{muted}Y1 | published | 564ed242}
\par
PowerPointda Transition nima?
\begin{enumerate}[label=\Alph*),leftmargin=9mm]
\item Bir slayddan keyingisiga o'tish effekti
\item Fayl kengaytmasi
\item Word jadvaldagi satr
\item Bitta katakning qiymati
\end{enumerate}
\begin{tcolorbox}[enhanced,breakable,colback=teal!4,colframe=teal!55!black,boxrule=0.45pt,arc=1.2mm,left=2.8mm,right=2.8mm,top=2mm,bottom=2mm]
\textbf{\color{teal}To'g'ri javob:} \textbf{A) Bir slayddan keyingisiga o'tish effekti}
\par\smallskip
\textbf{Izoh:} Transition slaydlar almashinuvidir.
\end{tcolorbox}
\vspace{2mm}\hrule height .2pt\vspace{2mm}

% DB 7a1c4db9-9a05-5a89-8f16-62563175d09c status=published
\needspace{6\baselineskip}
\textbf{\color{ink}347.} {ootnotesizecolor{muted}Y1 | published | 7a1c4db9}
\par
PowerPoint Transition sozlamasida \texttt{Duration=02,00}, \texttt{On Mouse Click} belgilangan va \texttt{After=00:10,00} ham belgilangan. Slayd qachon keyingi slaydga o'tadi?
\begin{enumerate}[label=\Alph*),leftmargin=9mm]
\item Sichqoncha bosilganda yoki 2 sekunddan so'ng
\item Faqat sichqoncha bosilganda
\item Faqat klaviatura tugmasi bosilganda
\item Sichqoncha bosilganda yoki 10 sekunddan so'ng
\end{enumerate}
\begin{tcolorbox}[enhanced,breakable,colback=teal!4,colframe=teal!55!black,boxrule=0.45pt,arc=1.2mm,left=2.8mm,right=2.8mm,top=2mm,bottom=2mm]
\textbf{\color{teal}To'g'ri javob:} \textbf{D) Sichqoncha bosilganda yoki 10 sekunddan so'ng}
\par\smallskip
\textbf{Izoh:} \texttt{After=10 s} avtomatik o'tish vaqtini, \texttt{On Mouse Click} qo'lda o'tishni bildiradi. \texttt{Duration=2 s} esa o'tish effektining davomiyligi, kutish vaqti emas.
\end{tcolorbox}
\vspace{2mm}\hrule height .2pt\vspace{2mm}

% DB 7e39c6f1-4a1f-5cc6-65c0-025cf132300b status=published
\needspace{6\baselineskip}
\textbf{\color{ink}348.} {ootnotesizecolor{muted}Y1 | published | 7e39c6f1}
\par
Animation nima?
\begin{enumerate}[label=\Alph*),leftmargin=9mm]
\item Excel funksiyasi
\item Slayd ichidagi obyektga berilgan effekt
\item Printer turini tanlash
\item Faqat slaydlar orasidagi o'tish
\end{enumerate}
\begin{tcolorbox}[enhanced,breakable,colback=teal!4,colframe=teal!55!black,boxrule=0.45pt,arc=1.2mm,left=2.8mm,right=2.8mm,top=2mm,bottom=2mm]
\textbf{\color{teal}To'g'ri javob:} \textbf{B) Slayd ichidagi obyektga berilgan effekt}
\par\smallskip
\textbf{Izoh:} Animation obyekt harakatidir.
\end{tcolorbox}
\vspace{2mm}\hrule height .2pt\vspace{2mm}

% DB 9b9b29b6-c766-b750-fdbc-8b2a5daede9e status=published
\needspace{6\baselineskip}
\textbf{\color{ink}349.} {ootnotesizecolor{muted}Y1 | published | 9b9b29b6}
\par
Entrance toifasi nimani bildiradi?
\begin{enumerate}[label=\Alph*),leftmargin=9mm]
\item Hujjatni shifrlashni
\item Slayd almashtirishni
\item Obyektning paydo bo'lish effekti
\item Obyektni doimiy o'chirish
\end{enumerate}
\begin{tcolorbox}[enhanced,breakable,colback=teal!4,colframe=teal!55!black,boxrule=0.45pt,arc=1.2mm,left=2.8mm,right=2.8mm,top=2mm,bottom=2mm]
\textbf{\color{teal}To'g'ri javob:} \textbf{C) Obyektning paydo bo'lish effekti}
\par\smallskip
\textbf{Izoh:} Entrance obyektni kirita ko'rsatadi.
\end{tcolorbox}
\vspace{2mm}\hrule height .2pt\vspace{2mm}

% DB a759d3d7-bca3-88ea-7180-b4dd269b2390 status=published
\needspace{6\baselineskip}
\textbf{\color{ink}350.} {ootnotesizecolor{muted}Y1 | published | a759d3d7}
\par
Apply To All nima qiladi?
\begin{enumerate}[label=\Alph*),leftmargin=9mm]
\item Barcha slaydni o'chiradi
\item Tanlangan slayd o'tishini barcha slaydlarga qo'llaydi
\item Faqat Office litsenziyasini o'zgartiradi
\item Faqat bitta obyektni kesadi
\end{enumerate}
\begin{tcolorbox}[enhanced,breakable,colback=teal!4,colframe=teal!55!black,boxrule=0.45pt,arc=1.2mm,left=2.8mm,right=2.8mm,top=2mm,bottom=2mm]
\textbf{\color{teal}To'g'ri javob:} \textbf{B) Tanlangan slayd o'tishini barcha slaydlarga qo'llaydi}
\par\smallskip
\textbf{Izoh:} Transitionsdagi Apply To All sozlamani umumlashtiradi.
\end{tcolorbox}
\vspace{2mm}\hrule height .2pt\vspace{2mm}

% DB f19d91ca-7b7b-4a37-bfe4-99ec15bc9421 status=published
\needspace{6\baselineskip}
\textbf{\color{ink}351.} {ootnotesizecolor{muted}Y3 | published | f19d91ca}
\par
PowerPoint slaydidagi obyektga Entrance animatsiyasi qo'shib, Preview orqali natijani ko'rish uchun bajariladigan amallarni tartiblang.
\begin{enumerate}[label=\Alph*),leftmargin=9mm]
\item Entrance effektini tanlash
\item Animatsiya uchun obyektni tanlash
\item Preview orqali ko'rish
\item Animations menyusini ochish
\end{enumerate}
\begin{tcolorbox}[enhanced,breakable,colback=teal!4,colframe=teal!55!black,boxrule=0.45pt,arc=1.2mm,left=2.8mm,right=2.8mm,top=2mm,bottom=2mm]
\textbf{\color{teal}To'g'ri javob:} \textbf{B - D - A - C}
\par\smallskip
\textbf{Izoh:} Bosqichlar: Animatsiya uchun obyektni tanlash \ensuremath{\rightarrow} Animations menyusini ochish \ensuremath{\rightarrow} Entrance effektini tanlash \ensuremath{\rightarrow} Preview orqali ko'rish.
\end{tcolorbox}
\vspace{2mm}\hrule height .2pt\vspace{2mm}

% DB 055dedf2-964d-61fc-71ea-5f8e1e9afcb9 status=published
\needspace{6\baselineskip}
\textbf{\color{ink}352.} {ootnotesizecolor{muted}Y1 | published | 055dedf2}
\par
F5 va Shift+F5 farqi?
\begin{enumerate}[label=\Alph*),leftmargin=9mm]
\item Ikkisi faqat Wordda ishlaydi
\item F5 faqat slayd o'chiradi
\item Shift+F5 faylni PDF qiladi
\item F5 boshidan, Shift+F5 joriy slayddan namoyishni boshlaydi
\end{enumerate}
\begin{tcolorbox}[enhanced,breakable,colback=teal!4,colframe=teal!55!black,boxrule=0.45pt,arc=1.2mm,left=2.8mm,right=2.8mm,top=2mm,bottom=2mm]
\textbf{\color{teal}To'g'ri javob:} \textbf{D) F5 boshidan, Shift+F5 joriy slayddan namoyishni boshlaydi}
\par\smallskip
\textbf{Izoh:} Standart PowerPoint tezkor tugmalari.
\end{tcolorbox}
\vspace{2mm}\hrule height .2pt\vspace{2mm}

% DB 0652538f-aa89-dfc6-5aaf-468ab17408e5 status=published
\needspace{6\baselineskip}
\textbf{\color{ink}353.} {ootnotesizecolor{muted}Y1 | published | 0652538f}
\par
Animation Pane'da 1-effekt "On Click", 2-effekt "With Previous", 3-effekt "After Previous" qilib berilgan. Ikkinchi effekt qachon ishga tushadi?
\begin{enumerate}[label=\Alph*),leftmargin=9mm]
\item Faqat uchinchi marta bosilganda
\item Birinchi effekt bilan bir vaqtda
\item Uchinchi effektdan keyin
\item Slayd ochilganidan 10 daqiqa keyin
\end{enumerate}
\begin{tcolorbox}[enhanced,breakable,colback=teal!4,colframe=teal!55!black,boxrule=0.45pt,arc=1.2mm,left=2.8mm,right=2.8mm,top=2mm,bottom=2mm]
\textbf{\color{teal}To'g'ri javob:} \textbf{B) Birinchi effekt bilan bir vaqtda}
\par\smallskip
\textbf{Izoh:} With Previous - oldingi effekt bilan birga boshlanadi; Delay alohida sozlama.
\end{tcolorbox}
\vspace{2mm}\hrule height .2pt\vspace{2mm}

% DB 2955446d-764a-891b-6aa7-93bc51bf1400 status=published
\needspace{6\baselineskip}
\textbf{\color{ink}354.} {ootnotesizecolor{muted}Y1 | published | 2955446d}
\par
On Click va After (timing) farqi?
\begin{enumerate}[label=\Alph*),leftmargin=9mm]
\item After faqat jadvalga
\item Ikkisi ham faqat Excel
\item On Click har doim avtomatik
\item Birinchisi foydalanuvchi bosishiga, ikkinchisi belgilangan vaqtga tayanadi
\end{enumerate}
\begin{tcolorbox}[enhanced,breakable,colback=teal!4,colframe=teal!55!black,boxrule=0.45pt,arc=1.2mm,left=2.8mm,right=2.8mm,top=2mm,bottom=2mm]
\textbf{\color{teal}To'g'ri javob:} \textbf{D) Birinchisi foydalanuvchi bosishiga, ikkinchisi belgilangan vaqtga tayanadi}
\par\smallskip
\textbf{Izoh:} Boshqaruv shartlari har xil.
\end{tcolorbox}
\vspace{2mm}\hrule height .2pt\vspace{2mm}

% DB 5373afe9-b537-5a13-9ded-d6205c51a089 status=published
\needspace{6\baselineskip}
\textbf{\color{ink}355.} {ootnotesizecolor{muted}Y1 | published | 5373afe9}
\par
PowerPoint Animation Pane chizig'ida Picture 2 ning kirish animatsiyasi 3-9 s, chiqishi 6-10 s; Picture 4 ning kirishi 0-5 s, chiqishi 11-14 s; Picture 6 ning kirishi 9-12 s, chiqishi 12-17 s. Kirish va chiqish davomiyliklari yig'indisi eng katta bo'lgan shaklni toping.
\begin{enumerate}[label=\Alph*),leftmargin=9mm]
\item Picture 6 - 8 s
\item Picture 4 - 6,5 s
\item Picture 2 - 10 s
\item Barchasi teng
\end{enumerate}
\begin{tcolorbox}[enhanced,breakable,colback=teal!4,colframe=teal!55!black,boxrule=0.45pt,arc=1.2mm,left=2.8mm,right=2.8mm,top=2mm,bottom=2mm]
\textbf{\color{teal}To'g'ri javob:} \textbf{C) Picture 2 - 10 s}
\par\smallskip
\textbf{Izoh:} Picture 2: (9-3)+(10-6)=6+4=10 s. Picture 4 va Picture 6 ning yig'indisi 8 s dan.
\end{tcolorbox}
\vspace{2mm}\hrule height .2pt\vspace{2mm}

% DB 5672ac57-35ea-f504-94cd-e5d4f321c22b status=published
\needspace{6\baselineskip}
\textbf{\color{ink}356.} {ootnotesizecolor{muted}Y1 | published | 5672ac57}
\par
O'qituvchi 12 slaydli taqdimotni har safar 15 soniyadan keyin avtomatik almashtirmoqchi. Qaysi sozlama mos?
\begin{enumerate}[label=\Alph*),leftmargin=9mm]
\item Transitions \ensuremath{\rightarrow} Advance Slide \ensuremath{\rightarrow} After: 00:15 (zarur slaydlarda)
\item Animations \ensuremath{\rightarrow} Emphasis \ensuremath{\rightarrow} Spin
\item Insert \ensuremath{\rightarrow} Symbols \ensuremath{\rightarrow} Equation
\item Home \ensuremath{\rightarrow} Replace \ensuremath{\rightarrow} All
\end{enumerate}
\begin{tcolorbox}[enhanced,breakable,colback=teal!4,colframe=teal!55!black,boxrule=0.45pt,arc=1.2mm,left=2.8mm,right=2.8mm,top=2mm,bottom=2mm]
\textbf{\color{teal}To'g'ri javob:} \textbf{A) Transitions \ensuremath{\rightarrow} Advance Slide \ensuremath{\rightarrow} After: 00:15 (zarur slaydlarda)}
\par\smallskip
\textbf{Izoh:} Automatic slayd o'tishi Advance Slide / After maydoni orqali sozlanadi; bosish bilan o'tish alohida yoqilishi mumkin.
\end{tcolorbox}
\vspace{2mm}\hrule height .2pt\vspace{2mm}

% DB 621309bf-18ba-48d0-adcd-8342fb8c634a status=published
\needspace{6\baselineskip}
\textbf{\color{ink}357.} {ootnotesizecolor{muted}Y3 | published | 621309bf}
\par
PowerPoint slaydiga almashinish (Transition) effektini berib, Duration vaqtini sozlash va almashinishni tekshirish amallarini tartiblang.
\begin{enumerate}[label=\Alph*),leftmargin=9mm]
\item Slayd almashinuvini tekshirish
\item Durationni sozlash
\item Transitions bo'limidan effekt berish
\item Slaydni tanlash
\end{enumerate}
\begin{tcolorbox}[enhanced,breakable,colback=teal!4,colframe=teal!55!black,boxrule=0.45pt,arc=1.2mm,left=2.8mm,right=2.8mm,top=2mm,bottom=2mm]
\textbf{\color{teal}To'g'ri javob:} \textbf{D - C - B - A}
\par\smallskip
\textbf{Izoh:} Bosqichlar: Slaydni tanlash \ensuremath{\rightarrow} Transitions bo'limidan effekt berish \ensuremath{\rightarrow} Durationni sozlash \ensuremath{\rightarrow} Slayd almashinuvini tekshirish.
\end{tcolorbox}
\vspace{2mm}\hrule height .2pt\vspace{2mm}

% DB 70fe38b7-cd71-596a-ba76-56aba3accdb8 status=published
\needspace{6\baselineskip}
\textbf{\color{ink}358.} {ootnotesizecolor{muted}Y1 | published | 70fe38b7}
\par
PowerPoint Animation Pane vaqt chizig'ida Star kirish animatsiyasi 0-3 s, chiqish animatsiyasi 4-6,5 s davom etadi. Circle va Square uchun kirish-chiqish davomiyliklari bundan qisqa. Eng uzun jami vaqtni toping.
\begin{enumerate}[label=\Alph*),leftmargin=9mm]
\item Circle - 6 s
\item Star - 5,5 s
\item Circle - 4,5 s
\item Square - 4,5 s
\end{enumerate}
\begin{tcolorbox}[enhanced,breakable,colback=teal!4,colframe=teal!55!black,boxrule=0.45pt,arc=1.2mm,left=2.8mm,right=2.8mm,top=2mm,bottom=2mm]
\textbf{\color{teal}To'g'ri javob:} \textbf{B) Star - 5,5 s}
\par\smallskip
\textbf{Izoh:} Star kirishi 3 s, chiqishi 2,5 s; jami 5,5 s. Bu Circle va Square animatsiyalari jamidan uzun.
\end{tcolorbox}
\vspace{2mm}\hrule height .2pt\vspace{2mm}

% DB 7cd37be9-aadc-e493-794f-92e386a8b14c status=published
\needspace{6\baselineskip}
\textbf{\color{ink}359.} {ootnotesizecolor{muted}Y2 | published | 7cd37be9}
\par
Atamalar va ularning aniq ma'nolarini moslashtiring.  1) On Click;  2) With Previous;  3) After Previous;  Izohlar: A) Oldingi effekt tugagach boshlash;  B) Oldingi effekt bilan birga boshlash;  C) Bosilganda boshlash;
\textbf{Chap ustun:}\begin{enumerate}[label=\arabic*),leftmargin=9mm]
\item On Click
\item With Previous
\item After Previous
\end{enumerate}
\textbf{O'ng ustun:}\begin{enumerate}[label=\Alph*),leftmargin=9mm]
\item Oldingi effekt tugagach boshlash
\item Oldingi effekt bilan birga boshlash
\item Bosilganda boshlash
\end{enumerate}
\begin{tcolorbox}[enhanced,breakable,colback=teal!4,colframe=teal!55!black,boxrule=0.45pt,arc=1.2mm,left=2.8mm,right=2.8mm,top=2mm,bottom=2mm]
\textbf{\color{teal}To'g'ri javob:} \textbf{1-C, 2-B, 3-A}
\par\smallskip
\textbf{Izoh:} To'g'ri moslik: On Click - Bosilganda boshlash, With Previous - Oldingi effekt bilan birga boshlash, After Previous - Oldingi effekt tugagach boshlash.
\end{tcolorbox}
\vspace{2mm}\hrule height .2pt\vspace{2mm}

% DB 7e65847b-ae26-4669-ac73-79088121a38e status=published
\needspace{6\baselineskip}
\textbf{\color{ink}360.} {ootnotesizecolor{muted}Y1 | published | 7e65847b}
\par
PowerPoint Animation Pane'da ikkita animatsiya belgilandi. Birinchi effekt ishlayotgan vaqtda ikkinchi effektni ham boshlash yoki uning tugashini kutish uchun qaysi vaqt sozlamalari farqlanadi?
\begin{enumerate}[label=\Alph*),leftmargin=9mm]
\item With Previous birinchi effektdan oldin, After Previous esa bosilganda boshlanadi
\item With Previous avvalgisidan keyin, After Previous esa bir vaqtda boshlanadi
\item With Previous birga, After Previous oldingi tugagach boshlanadi
\item With Previous faqat slayd almashadi, After Previous esa obyektni o'chiradi
\end{enumerate}
\begin{tcolorbox}[enhanced,breakable,colback=teal!4,colframe=teal!55!black,boxrule=0.45pt,arc=1.2mm,left=2.8mm,right=2.8mm,top=2mm,bottom=2mm]
\textbf{\color{teal}To'g'ri javob:} \textbf{C) With Previous birga, After Previous oldingi tugagach boshlanadi}
\par\smallskip
\textbf{Izoh:} With Previous avvalgi animatsiya bilan bir vaqtda boshlashga, After Previous esa avvalgi animatsiya tugagach boshlashga xizmat qiladi. On Click bosishga bog'liq alohida sozlamadir.
\end{tcolorbox}
\vspace{2mm}\hrule height .2pt\vspace{2mm}

% DB a88f3eab-5d91-ee34-6121-30231555b187 status=published
\needspace{6\baselineskip}
\textbf{\color{ink}361.} {ootnotesizecolor{muted}Y1 | published | a88f3eab}
\par
Bitta slaydda rasmga Entrance effekti, keyingi slaydga Fade o'tish effekti berildi. Qaysi fikr to'g'ri?
\begin{enumerate}[label=\Alph*),leftmargin=9mm]
\item Entrance slayd ichidagi obyektga, Fade (Transition) esa slaydlar almashishiga tegishli
\item Har ikkisi faqat shriftda ishlaydi
\item Fade slaydni tahrirlashni o'chiradi
\item Entrance faqat slaydlarni almashtiradi
\end{enumerate}
\begin{tcolorbox}[enhanced,breakable,colback=teal!4,colframe=teal!55!black,boxrule=0.45pt,arc=1.2mm,left=2.8mm,right=2.8mm,top=2mm,bottom=2mm]
\textbf{\color{teal}To'g'ri javob:} \textbf{A) Entrance slayd ichidagi obyektga, Fade (Transition) esa slaydlar almashishiga tegishli}
\par\smallskip
\textbf{Izoh:} Animation va Transition turli darajadagi effektlar: obyekt ichida vs slaydlar orasida.
\end{tcolorbox}
\vspace{2mm}\hrule height .2pt\vspace{2mm}

% DB dccd7f88-1821-cdc7-a227-d41af31f6a33 status=published
\needspace{6\baselineskip}
\textbf{\color{ink}362.} {ootnotesizecolor{muted}Y1 | published | dccd7f88}
\par
Duration va Delay farqi?
\begin{enumerate}[label=\Alph*),leftmargin=9mm]
\item Duration effekt davomiyligi, Delay boshlanish oldidan kutish
\item Delay har doim tezlikni ikki marta oshiradi
\item Duration boshlash tugmasi
\item Ikkisi faqat rasm rangini bildiradi
\end{enumerate}
\begin{tcolorbox}[enhanced,breakable,colback=teal!4,colframe=teal!55!black,boxrule=0.45pt,arc=1.2mm,left=2.8mm,right=2.8mm,top=2mm,bottom=2mm]
\textbf{\color{teal}To'g'ri javob:} \textbf{A) Duration effekt davomiyligi, Delay boshlanish oldidan kutish}
\par\smallskip
\textbf{Izoh:} Vaqtning turli bosqichlari.
\end{tcolorbox}
\vspace{2mm}\hrule height .2pt\vspace{2mm}

% DB 18d17a8f-3ac5-4689-bef7-dc42209c1005 status=published
\needspace{6\baselineskip}
\textbf{\color{ink}363.} {ootnotesizecolor{muted}Y1 | published | 18d17a8f}
\par
PowerPoint ish stolida taqdimot ochiq. Slayd namoyishini \textbf{birinchi slayddan} klaviatura orqali boshlash uchun qaysi tugma ishlatiladi?
\begin{enumerate}[label=\Alph*),leftmargin=9mm]
\item F5
\item Shift+F5
\item F7
\item Ctrl+S
\end{enumerate}
\begin{tcolorbox}[enhanced,breakable,colback=teal!4,colframe=teal!55!black,boxrule=0.45pt,arc=1.2mm,left=2.8mm,right=2.8mm,top=2mm,bottom=2mm]
\textbf{\color{teal}To'g'ri javob:} \textbf{A) F5}
\par\smallskip
\textbf{Izoh:} F5 namoyishni boshidan boshlaydi. Shift+F5 joriy slayddan boshlaydi, F7 odatda spelling tekshiruvini, Ctrl+S esa faylni saqlashni chaqiradi. Bu turli vazifalarni farqlash muhim.
\end{tcolorbox}
\vspace{2mm}\hrule height .2pt\vspace{2mm}

% DB 21c997a7-ce28-4f83-b4d4-3925275355aa status=published
\needspace{6\baselineskip}
\textbf{\color{ink}364.} {ootnotesizecolor{muted}Y1 | published | 21c997a7}
\par
Slaydda ko'rinib turgan shakl ma'lum lahzada ekrandan chiqib ketishi kerak. PowerPointdagi qaysi animatsiya guruhi tanlanadi?
\begin{enumerate}[label=\Alph*),leftmargin=9mm]
\item Entrance - obyektni paydo qiladi
\item Emphasis - obyektni ta'kidlaydi
\item Exit - obyektni ekrandan chiqaradi
\item Motion Path - obyektni belgilangan yo'l bo'ylab suradi
\end{enumerate}
\begin{tcolorbox}[enhanced,breakable,colback=teal!4,colframe=teal!55!black,boxrule=0.45pt,arc=1.2mm,left=2.8mm,right=2.8mm,top=2mm,bottom=2mm]
\textbf{\color{teal}To'g'ri javob:} \textbf{C) Exit - obyektni ekrandan chiqaradi}
\par\smallskip
\textbf{Izoh:} Exit animatsiyalari slaydda ko'rinib turgan obyektning chiqib ketishi yoki yashirilishi uchun xizmat qiladi. Entrance kirishni, Emphasis urg'u berishni, Motion Path esa obyekt harakat trayektoriyasini boshqaradi.
\end{tcolorbox}
\vspace{2mm}\hrule height .2pt\vspace{2mm}

% DB 232f7cb7-a9c8-b7ce-c6a1-6b31caf9e8e9 status=published
\needspace{6\baselineskip}
\textbf{\color{ink}365.} {ootnotesizecolor{muted}Y1 | published | 232f7cb7}
\par
Obyektni bosganda paydo bo'lsin. Start parametri?
\begin{enumerate}[label=\Alph*),leftmargin=9mm]
\item After 10 seconds (transition)
\item After Previous
\item On Click
\item With Previous
\end{enumerate}
\begin{tcolorbox}[enhanced,breakable,colback=teal!4,colframe=teal!55!black,boxrule=0.45pt,arc=1.2mm,left=2.8mm,right=2.8mm,top=2mm,bottom=2mm]
\textbf{\color{teal}To'g'ri javob:} \textbf{C) On Click}
\par\smallskip
\textbf{Izoh:} On Click bosishga bog'liq.
\end{tcolorbox}
\vspace{2mm}\hrule height .2pt\vspace{2mm}

% DB 2d385a36-979f-d543-3b2f-06476b14d012 status=published
\needspace{6\baselineskip}
\textbf{\color{ink}366.} {ootnotesizecolor{muted}Y1 | published | 2d385a36}
\par
Rasm slaydga uchib kirsin. Qaysi bo'limdan boshlanadi?
\begin{enumerate}[label=\Alph*),leftmargin=9mm]
\item Animations \ensuremath{\rightarrow} Entrance
\item Transitions \ensuremath{\rightarrow} Apply To All
\item Review \ensuremath{\rightarrow} Spelling
\item Design \ensuremath{\rightarrow} Themes faqat
\end{enumerate}
\begin{tcolorbox}[enhanced,breakable,colback=teal!4,colframe=teal!55!black,boxrule=0.45pt,arc=1.2mm,left=2.8mm,right=2.8mm,top=2mm,bottom=2mm]
\textbf{\color{teal}To'g'ri javob:} \textbf{A) Animations \ensuremath{\rightarrow} Entrance}
\par\smallskip
\textbf{Izoh:} Obyekt animatsiyasi Animationsda.
\end{tcolorbox}
\vspace{2mm}\hrule height .2pt\vspace{2mm}

% DB 3ee2f192-d244-c2a9-c452-6a5df4b2220e status=published
\needspace{6\baselineskip}
\textbf{\color{ink}367.} {ootnotesizecolor{muted}Y2 | published | 3ee2f192}
\par
Amallarni izohi bilan moslang.  1) Duration;  2) Delay;  3) Animation Pane;  Izohlar: A) Effekt qancha davom etishi;  B) Effektlar tartibini boshqarish;  C) Effekt oldidan kutish;
\textbf{Chap ustun:}\begin{enumerate}[label=\arabic*),leftmargin=9mm]
\item Duration
\item Delay
\item Animation Pane
\end{enumerate}
\textbf{O'ng ustun:}\begin{enumerate}[label=\Alph*),leftmargin=9mm]
\item Effekt qancha davom etishi
\item Effektlar tartibini boshqarish
\item Effekt oldidan kutish
\end{enumerate}
\begin{tcolorbox}[enhanced,breakable,colback=teal!4,colframe=teal!55!black,boxrule=0.45pt,arc=1.2mm,left=2.8mm,right=2.8mm,top=2mm,bottom=2mm]
\textbf{\color{teal}To'g'ri javob:} \textbf{1-A, 2-C, 3-B}
\par\smallskip
\textbf{Izoh:} To'g'ri moslik: Duration - Effekt qancha davom etishi, Delay - Effekt oldidan kutish, Animation Pane - Effektlar tartibini boshqarish.
\end{tcolorbox}
\vspace{2mm}\hrule height .2pt\vspace{2mm}

% DB 409777cb-6670-4915-8c34-eff4073f5382 status=published
\needspace{6\baselineskip}
\textbf{\color{ink}368.} {ootnotesizecolor{muted}Y1 | published | 409777cb}
\par
PowerPoint taqdimotida bir slayd ikkinchisiga Fade bilan almashadi, yangi slayddagi rasm esa Fly In bilan paydo bo'ladi. Qaysi izoh ikkala effektni to'g'ri farqlaydi?
\begin{enumerate}[label=\Alph*),leftmargin=9mm]
\item Fade obyektning Exit animatsiyasi, Fly In esa slaydning Transition effekti
\item Fade slaydlararo Transition, Fly In esa slayd ichidagi obyekt Animation effekti
\item Fade faqat ovoz parametrini boshqaradi, Fly In esa slaydning vaqtini belgilaydi
\item Fade va Fly In bir xil turdagi slaydlararo o'tish animatsiyalaridir
\end{enumerate}
\begin{tcolorbox}[enhanced,breakable,colback=teal!4,colframe=teal!55!black,boxrule=0.45pt,arc=1.2mm,left=2.8mm,right=2.8mm,top=2mm,bottom=2mm]
\textbf{\color{teal}To'g'ri javob:} \textbf{B) Fade slaydlararo Transition, Fly In esa slayd ichidagi obyekt Animation effekti}
\par\smallskip
\textbf{Izoh:} PowerPointda Transition (masalan, slaydlararo Fade) to'liq slaydni almashtirishga qo'llanadi. Animation (masalan, Fly In) slayd ichidagi alohida obyektning paydo bo'lishini belgilaydi. Ikkalasini bir xil menyu darajasi deb adashtirmaslik kerak.
\end{tcolorbox}
\vspace{2mm}\hrule height .2pt\vspace{2mm}

% DB 4e43f06a-a5e1-449a-8a68-13f2fd077679 status=published
\needspace{6\baselineskip}
\textbf{\color{ink}369.} {ootnotesizecolor{muted}Y1 | published | 4e43f06a}
\par
PowerPoint slaydida matn, diagramma va rasm mavjud. O'qituvchi ularning har birini alohida kirish animatsiyasi bilan ko'rsatmoqchi. Qaysi mulohaza to'g'ri?
\begin{enumerate}[label=\Alph*),leftmargin=9mm]
\item Faqat matn maydoniga Entrance animatsiya berilishi mumkin
\item Faqat rasm va shakllarga animatsiya berish mumkin, diagrammaga emas
\item Bunday holatda faqat Transition ishlaydi, obyekt Animation ishlamaydi
\item Matn, rasm va diagrammalarga alohida animatsiya qo'llash mumkin
\end{enumerate}
\begin{tcolorbox}[enhanced,breakable,colback=teal!4,colframe=teal!55!black,boxrule=0.45pt,arc=1.2mm,left=2.8mm,right=2.8mm,top=2mm,bottom=2mm]
\textbf{\color{teal}To'g'ri javob:} \textbf{D) Matn, rasm va diagrammalarga alohida animatsiya qo'llash mumkin}
\par\smallskip
\textbf{Izoh:} PowerPoint matn maydonlari, rasmlar, shakllar va diagramma kabi turli obyektlarga Animation effekti bera oladi. Transition slaydlar orasidagi o'tish uchun, obyektlarning individual paydo bo'lishi uchun emas.
\end{tcolorbox}
\vspace{2mm}\hrule height .2pt\vspace{2mm}

% DB 5c1f7b11-a017-141a-a855-10b2d8029467 status=published
\needspace{6\baselineskip}
\textbf{\color{ink}370.} {ootnotesizecolor{muted}Y2 | published | 5c1f7b11}
\par
Animatsiya boshlanish turlari va bajarilish shartlarini moslang.
\textbf{Chap ustun:}\begin{enumerate}[label=\arabic*),leftmargin=9mm]
\item On Click
\item With Previous
\item After Previous
\item Duration
\end{enumerate}
\textbf{O'ng ustun:}\begin{enumerate}[label=\Alph*),leftmargin=9mm]
\item Oldingi effekt tugagach ishga tushish
\item Bajarilish qancha vaqt davom etishini belgilash
\item Sichqoncha bosilganda boshlash
\item Oldingi effekt bilan birga boshlash
\end{enumerate}
\begin{tcolorbox}[enhanced,breakable,colback=teal!4,colframe=teal!55!black,boxrule=0.45pt,arc=1.2mm,left=2.8mm,right=2.8mm,top=2mm,bottom=2mm]
\textbf{\color{teal}To'g'ri javob:} \textbf{1-C, 2-D, 3-A, 4-B}
\par\smallskip
\textbf{Izoh:} 1-C, 2-D, 3-A, 4-B; Delay alohida maydon va bu uch start turini almashtirmaydi.
\end{tcolorbox}
\vspace{2mm}\hrule height .2pt\vspace{2mm}

% DB d3ec77df-1ae5-474e-89a1-438a0f3e1f0c status=published
\needspace{6\baselineskip}
\textbf{\color{ink}371.} {ootnotesizecolor{muted}Y1 | published | d3ec77df}
\par
Slaydda 5 ta animatsiya qo'llangan. Ularning ishga tushish navbati va har bir effektning boshlanish vaqti boshqariladigan oynani tanlang.
\begin{enumerate}[label=\Alph*),leftmargin=9mm]
\item Slide Sorter - slaydlar umumiy tartibini ko'rsatadi
\item Selection Pane - obyektlarni tanlash va qatlamlarda ko'rishga yordam beradi
\item Animation Pane - effektlarning navbati va vaqtini boshqaradi
\item Transitions - slaydlararo almashish effektlarini o'rnatadi
\end{enumerate}
\begin{tcolorbox}[enhanced,breakable,colback=teal!4,colframe=teal!55!black,boxrule=0.45pt,arc=1.2mm,left=2.8mm,right=2.8mm,top=2mm,bottom=2mm]
\textbf{\color{teal}To'g'ri javob:} \textbf{C) Animation Pane - effektlarning navbati va vaqtini boshqaradi}
\par\smallskip
\textbf{Izoh:} Animation Pane slayd ichidagi Animation effektlarini ko'rish, qayta tartiblash, boshlanish shartlari va vaqt sozlamalarini tekshirish imkonini beradi. Slide Sorter slaydlarning o'z tartibini boshqaradi; Transitions esa slaydlararo effektlar uchun.
\end{tcolorbox}
\vspace{2mm}\hrule height .2pt\vspace{2mm}

% DB ddf5a937-7301-8a23-76ce-6f5c9ed77540 status=published
\needspace{6\baselineskip}
\textbf{\color{ink}372.} {ootnotesizecolor{muted}Y1 | published | ddf5a937}
\par
Keyingi slayd Fade bilan almashsin. Qaysi bo'lim?
\begin{enumerate}[label=\Alph*),leftmargin=9mm]
\item Home \ensuremath{\rightarrow} Replace
\item Transitions
\item Insert \ensuremath{\rightarrow} Table
\item Animations \ensuremath{\rightarrow} Emphasis
\end{enumerate}
\begin{tcolorbox}[enhanced,breakable,colback=teal!4,colframe=teal!55!black,boxrule=0.45pt,arc=1.2mm,left=2.8mm,right=2.8mm,top=2mm,bottom=2mm]
\textbf{\color{teal}To'g'ri javob:} \textbf{B) Transitions}
\par\smallskip
\textbf{Izoh:} Fade slayd o'tishi sifatida tanlanishi mumkin.
\end{tcolorbox}
\vspace{2mm}\hrule height .2pt\vspace{2mm}

% DB f933cfc5-846c-42c8-a7b6-a160bc334008 status=published
\needspace{6\baselineskip}
\textbf{\color{ink}373.} {ootnotesizecolor{muted}Y1 | published | f933cfc5}
\par
PowerPointdagi bitta obyekt animatsiyasi uchun \textbf{Delay} va \textbf{Duration} belgilandi. Qaysi izoh to'g'ri?
\begin{enumerate}[label=\Alph*),leftmargin=9mm]
\item Delay animatsiya qancha davom etishini, Duration esa boshlanishigacha kutishni bildiradi
\item Delay ishga tushgandan keyingi kutish vaqtini, Duration effekt bajarilish vaqtini bildiradi
\item Delay effekt boshlanishidan oldingi kutish, Duration effektning davom etish vaqtini bildiradi
\item Delay va Duration ikkalasi ham slayd almashishining avtomatik vaqtini beradi
\end{enumerate}
\begin{tcolorbox}[enhanced,breakable,colback=teal!4,colframe=teal!55!black,boxrule=0.45pt,arc=1.2mm,left=2.8mm,right=2.8mm,top=2mm,bottom=2mm]
\textbf{\color{teal}To'g'ri javob:} \textbf{C) Delay effekt boshlanishidan oldingi kutish, Duration effektning davom etish vaqtini bildiradi}
\par\smallskip
\textbf{Izoh:} Delay - animatsiya ishga tushishi oldidan kechikish; Duration - effektning o'zi bajariladigan davomiylik. Ular alohida boshqariladi; slaydlararo Transitionning After parametri bilan ham aralashtirilmaydi.
\end{tcolorbox}
\vspace{2mm}\hrule height .2pt\vspace{2mm}

% DB ff5ae762-ba37-4c66-86e7-5147a950792d status=published
\needspace{6\baselineskip}
\textbf{\color{ink}374.} {ootnotesizecolor{muted}Y3 | published | ff5ae762}
\par
PowerPoint taqdimotidagi ikkita obyekt animatsiyasining boshlanish usuli va ketma-ketligini Animation Pane yordamida sozlab, namoyishda tekshirish tartibini belgilang.
\begin{enumerate}[label=\Alph*),leftmargin=9mm]
\item Namoyishda ketma-ketlikni tekshirish
\item With Previous/After Previousni sozlash
\item Animation Pane oynasini ochish
\item Ikki obyektga animatsiya qo'shish
\end{enumerate}
\begin{tcolorbox}[enhanced,breakable,colback=teal!4,colframe=teal!55!black,boxrule=0.45pt,arc=1.2mm,left=2.8mm,right=2.8mm,top=2mm,bottom=2mm]
\textbf{\color{teal}To'g'ri javob:} \textbf{D - C - B - A}
\par\smallskip
\textbf{Izoh:} Bosqichlar: Ikki obyektga animatsiya qo'shish \ensuremath{\rightarrow} Animation Pane oynasini ochish \ensuremath{\rightarrow} With Previous/After Previousni sozlash \ensuremath{\rightarrow} Namoyishda ketma-ketlikni tekshirish.
\end{tcolorbox}
\vspace{2mm}\hrule height .2pt\vspace{2mm}

% DB 0e3ed27c-a1dc-3399-0da8-d6b058a5998d status=published
\needspace{6\baselineskip}
\textbf{\color{ink}375.} {ootnotesizecolor{muted}Y2 | published | 0e3ed27c}
\par
Ushbu vaziyatdagi vositalar va vazifalarni moslashtiring.  1) Transition;  2) Motion Path;  3) Apply To All;  Izohlar: A) Tanlangan o'tishni hamma slaydga qo'llash;  B) Slaydlar orasidagi o'tish;  C) Slayddagi obyekt harakat yo'li;
\textbf{Chap ustun:}\begin{enumerate}[label=\arabic*),leftmargin=9mm]
\item Transition
\item Motion Path
\item Apply To All
\end{enumerate}
\textbf{O'ng ustun:}\begin{enumerate}[label=\Alph*),leftmargin=9mm]
\item Tanlangan o'tishni hamma slaydga qo'llash
\item Slaydlar orasidagi o'tish
\item Slayddagi obyekt harakat yo'li
\end{enumerate}
\begin{tcolorbox}[enhanced,breakable,colback=teal!4,colframe=teal!55!black,boxrule=0.45pt,arc=1.2mm,left=2.8mm,right=2.8mm,top=2mm,bottom=2mm]
\textbf{\color{teal}To'g'ri javob:} \textbf{1-B, 2-C, 3-A}
\par\smallskip
\textbf{Izoh:} To'g'ri moslik: Transition - Slaydlar orasidagi o'tish, Motion Path - Slayddagi obyekt harakat yo'li, Apply To All - Tanlangan o'tishni hamma slaydga qo'llash.
\end{tcolorbox}
\vspace{2mm}\hrule height .2pt\vspace{2mm}

% DB 405b138f-f1ca-fb32-b3de-9d4c71a727bd status=published
\needspace{6\baselineskip}
\textbf{\color{ink}376.} {ootnotesizecolor{muted}Y1 | published | 405b138f}
\par
Slayd 10 soniyadan keyin avtomatik o'tsin. Qaysi sozlama?
\begin{enumerate}[label=\Alph*),leftmargin=9mm]
\item Design \ensuremath{\rightarrow} Variants
\item Insert \ensuremath{\rightarrow} Equation
\item Animations \ensuremath{\rightarrow} Entrance 10 s
\item Transitions \ensuremath{\rightarrow} Advance Slide \ensuremath{\rightarrow} After 10 s
\end{enumerate}
\begin{tcolorbox}[enhanced,breakable,colback=teal!4,colframe=teal!55!black,boxrule=0.45pt,arc=1.2mm,left=2.8mm,right=2.8mm,top=2mm,bottom=2mm]
\textbf{\color{teal}To'g'ri javob:} \textbf{D) Transitions \ensuremath{\rightarrow} Advance Slide \ensuremath{\rightarrow} After 10 s}
\par\smallskip
\textbf{Izoh:} Slayd o'tish vaqti Transitionsda.
\end{tcolorbox}
\vspace{2mm}\hrule height .2pt\vspace{2mm}

% DB 89293875-27ae-5cea-31ac-0d53038d7e6b status=published
\needspace{6\baselineskip}
\textbf{\color{ink}377.} {ootnotesizecolor{muted}Y1 | published | 89293875}
\par
Effekt 2 soniya kechikib, 4 soniya davom etsin. Boshlanishidan tugashigacha vaqt?
\begin{enumerate}[label=\Alph*),leftmargin=9mm]
\item 8 soniya
\item 4 soniya
\item 6 soniya
\item 2 soniya
\end{enumerate}
\begin{tcolorbox}[enhanced,breakable,colback=teal!4,colframe=teal!55!black,boxrule=0.45pt,arc=1.2mm,left=2.8mm,right=2.8mm,top=2mm,bottom=2mm]
\textbf{\color{teal}To'g'ri javob:} \textbf{C) 6 soniya}
\par\smallskip
\textbf{Izoh:} Delay+Duration=2+4=6.
\end{tcolorbox}
\vspace{2mm}\hrule height .2pt\vspace{2mm}

% DB 910a5aeb-1ad9-20d7-ea5d-53f452c57d14 status=published
\needspace{6\baselineskip}
\textbf{\color{ink}378.} {ootnotesizecolor{muted}Y1 | published | 910a5aeb}
\par
Slaydda matn qizilga o'tib e'tibor tortsin, lekin yo'qolmasin. Qaysi animatsiya turi?
\begin{enumerate}[label=\Alph*),leftmargin=9mm]
\item Transition faqat
\item Page Break
\item Exit
\item Emphasis
\end{enumerate}
\begin{tcolorbox}[enhanced,breakable,colback=teal!4,colframe=teal!55!black,boxrule=0.45pt,arc=1.2mm,left=2.8mm,right=2.8mm,top=2mm,bottom=2mm]
\textbf{\color{teal}To'g'ri javob:} \textbf{D) Emphasis}
\par\smallskip
\textbf{Izoh:} Emphasis mavjud obyektga urg'u beradi.
\end{tcolorbox}
\vspace{2mm}\hrule height .2pt\vspace{2mm}

% DB b9e56aad-282f-4c2b-b178-7765909cffe3 status=published
\needspace{6\baselineskip}
\textbf{\color{ink}379.} {ootnotesizecolor{muted}Y3 | published | b9e56aad}
\par
PowerPoint taqdimotidagi slaydlar belgilangan vaqt oralig'ida avtomatik almashishi kerak. Shu vaqtni barcha slaydlarga tatbiq qilib, Slide Showda tekshirish amallarini tartiblang.
\begin{enumerate}[label=\Alph*),leftmargin=9mm]
\item Avtomatik almashishni kuzatish
\item Slaydlar almashinuvi uchun vaqt qo'yish
\item Apply To All berish
\item Slide Show rejimini ishga tushirish
\end{enumerate}
\begin{tcolorbox}[enhanced,breakable,colback=teal!4,colframe=teal!55!black,boxrule=0.45pt,arc=1.2mm,left=2.8mm,right=2.8mm,top=2mm,bottom=2mm]
\textbf{\color{teal}To'g'ri javob:} \textbf{B - C - D - A}
\par\smallskip
\textbf{Izoh:} Bosqichlar: Slaydlar almashinuvi uchun vaqt qo'yish \ensuremath{\rightarrow} Apply To All berish \ensuremath{\rightarrow} Slide Show rejimini ishga tushirish \ensuremath{\rightarrow} Avtomatik almashishni kuzatish.
\end{tcolorbox}
\vspace{2mm}\hrule height .2pt\vspace{2mm}

% DB da0b4f3d-f3d8-174d-e815-b0d95e1c519a status=published
\needspace{6\baselineskip}
\textbf{\color{ink}380.} {ootnotesizecolor{muted}Y1 | published | da0b4f3d}
\par
3 ta animatsiya tartibini almashtirish uchun?
\begin{enumerate}[label=\Alph*),leftmargin=9mm]
\item Animation Pane ro'yxatida qayta tartiblash
\item Design \ensuremath{\rightarrow} Themes
\item File \ensuremath{\rightarrow} Save As
\item Word \ensuremath{\rightarrow} Footnote
\end{enumerate}
\begin{tcolorbox}[enhanced,breakable,colback=teal!4,colframe=teal!55!black,boxrule=0.45pt,arc=1.2mm,left=2.8mm,right=2.8mm,top=2mm,bottom=2mm]
\textbf{\color{teal}To'g'ri javob:} \textbf{A) Animation Pane ro'yxatida qayta tartiblash}
\par\smallskip
\textbf{Izoh:} Pane animatsiya tartibini ko'rsatadi.
\end{tcolorbox}
\vspace{2mm}\hrule height .2pt\vspace{2mm}

% DB 0b0dfdda-f294-e28e-c132-0e1ca11ebaf4 status=published
\needspace{6\baselineskip}
\textbf{\color{ink}381.} {ootnotesizecolor{muted}Y1 | published | 0b0dfdda}
\par
Birinchi animatsiya On Click, ikkinchisi After Previous. Bosish hali amalga oshmagan. Ikkinchi animatsiyaning boshlanish vaqti haqida qanday xulosa aniq?
\begin{enumerate}[label=\Alph*),leftmargin=9mm]
\item Ikkinchisi birinchidan oldin boshlanadi
\item Birinchi bosish qachon bo'lishi noma'lumligi sababli mutlaq boshlanish soniyasini topib bo'lmaydi
\item After Previous sharti o'zi birinchi effektni ishga tushiradi
\item Ikkinchisi doim 0-sekundda boshlanadi
\end{enumerate}
\begin{tcolorbox}[enhanced,breakable,colback=teal!4,colframe=teal!55!black,boxrule=0.45pt,arc=1.2mm,left=2.8mm,right=2.8mm,top=2mm,bottom=2mm]
\textbf{\color{teal}To'g'ri javob:} \textbf{B) Birinchi bosish qachon bo'lishi noma'lumligi sababli mutlaq boshlanish soniyasini topib bo'lmaydi}
\par\smallskip
\textbf{Izoh:} Bosish vaqti berilmagani uchun aniq vaqt mavjud emas; After Previous oldingi effekt yakunini kutadi.
\end{tcolorbox}
\vspace{2mm}\hrule height .2pt\vspace{2mm}

% DB 48b79de5-9826-60aa-4cc1-bddd0069b97a status=published
\needspace{6\baselineskip}
\textbf{\color{ink}382.} {ootnotesizecolor{muted}Y1 | published | 48b79de5}
\par
A animatsiyasi 0-sekundda boshlanib 2 soniya davom etadi. B After Previous rejimida 1 soniya Delay bilan ishga tushib 3 soniya davom etadi. B qachon tugaydi?
\begin{enumerate}[label=\Alph*),leftmargin=9mm]
\item 6-sekundda
\item 4-sekundda
\item 3-sekundda
\item 5-sekundda
\end{enumerate}
\begin{tcolorbox}[enhanced,breakable,colback=teal!4,colframe=teal!55!black,boxrule=0.45pt,arc=1.2mm,left=2.8mm,right=2.8mm,top=2mm,bottom=2mm]
\textbf{\color{teal}To'g'ri javob:} \textbf{A) 6-sekundda}
\par\smallskip
\textbf{Izoh:} A 2-sekundda tugaydi; B 3-sekundda boshlanadi va 6-sekundda tugaydi.
\end{tcolorbox}
\vspace{2mm}\hrule height .2pt\vspace{2mm}

% DB 5a54d899-d1f7-6e4e-6923-33c135073a0c status=published
\needspace{6\baselineskip}
\textbf{\color{ink}383.} {ootnotesizecolor{muted}Y3 | published | 5a54d899}
\par
PowerPointdagi yangi slayd hali bo'sh, unda uchta obyekt ham, animatsiya ham yo'q. Uchinchi effekt ikkinchisidan so'ng, ikkinchisi birinchisi bilan birga boshlanishi kerak. Amallarni tartiblang.
\begin{enumerate}[label=\Alph*),leftmargin=9mm]
\item Slide Show'da uch effektning haqiqiy boshlanishini sinab ko'rish
\item Uchta obyektning har biriga Entrance animatsiyasini qo'shish
\item Slaydga uchta obyektni joylashtirish
\item Animation Pane'da effektlarni 1-2-3 tartibga keltirib, \texttt{On Click}, \texttt{With Previous}, \texttt{After Previous} vaqtlarini sozlash
\end{enumerate}
\begin{tcolorbox}[enhanced,breakable,colback=teal!4,colframe=teal!55!black,boxrule=0.45pt,arc=1.2mm,left=2.8mm,right=2.8mm,top=2mm,bottom=2mm]
\textbf{\color{teal}To'g'ri javob:} \textbf{C - B - D - A}
\par\smallskip
\textbf{Izoh:} Avval obyektlar yaratiladi, keyin ularga animatsiya qo'shiladi. Shundan keyingina Animation Pane'dagi tartib va Start: 1=On Click, 2=With Previous, 3=After Previous sozlanadi; oxirida namoyishda tekshiriladi. Bu tartib obyekt/effekt mavjud bo'lishiga bog'liq.
\end{tcolorbox}
\vspace{2mm}\hrule height .2pt\vspace{2mm}

% DB 72637c21-3d71-9612-b35e-43ca7f8995b7 status=published
\needspace{6\baselineskip}
\textbf{\color{ink}384.} {ootnotesizecolor{muted}Y2 | published | 72637c21}
\par
Moslashtiring.  1) Animation;  2) Transition;  3) Duration;  Ta'riflar: A) Effekt davomiyligi;  B) Slaydlar almashinish effekti;  C) Slayd ichidagi obyekt effekti;
\textbf{Chap ustun:}\begin{enumerate}[label=\arabic*),leftmargin=9mm]
\item Animation
\item Transition
\item Duration
\end{enumerate}
\textbf{O'ng ustun:}\begin{enumerate}[label=\Alph*),leftmargin=9mm]
\item Effekt davomiyligi
\item Slaydlar almashinish effekti
\item Slayd ichidagi obyekt effekti
\end{enumerate}
\begin{tcolorbox}[enhanced,breakable,colback=teal!4,colframe=teal!55!black,boxrule=0.45pt,arc=1.2mm,left=2.8mm,right=2.8mm,top=2mm,bottom=2mm]
\textbf{\color{teal}To'g'ri javob:} \textbf{1-C, 2-B, 3-A}
\par\smallskip
\textbf{Izoh:} Mos keladigan ta'riflar: 1-C, 2-B, 3-A.
\end{tcolorbox}
\vspace{2mm}\hrule height .2pt\vspace{2mm}

% DB 83e7c9e4-c3e1-4258-8372-b8e5dbcd50a8 status=published
\needspace{6\baselineskip}
\textbf{\color{ink}385.} {ootnotesizecolor{muted}Y3 | published | 83e7c9e4}
\par
PowerPoint slaydiga video joylashtirildi. Video faqat kerakli qismi ijro etilishi uchun Playback menyusidagi Trim Video oynasida boshi va oxiri qisqartirilishi, so'ng Slide Showda tekshirilishi kerak. Amallarni tartiblang.
\begin{enumerate}[label=\Alph*),leftmargin=9mm]
\item Namoyishda kerakli qismni tekshirish
\item Playback bo'limidan Trim Video oynasini ochish
\item Videoni slaydga joylashtirish
\item Videoning keraksiz qismlarini Trim Video oynasida qisqartirish
\end{enumerate}
\begin{tcolorbox}[enhanced,breakable,colback=teal!4,colframe=teal!55!black,boxrule=0.45pt,arc=1.2mm,left=2.8mm,right=2.8mm,top=2mm,bottom=2mm]
\textbf{\color{teal}To'g'ri javob:} \textbf{C - B - D - A}
\par\smallskip
\textbf{Izoh:} Avval video slaydga joylashtiriladi, so'ng Playback bo'limidan Trim Video oynasi ochiladi. Unda ijro etiladigan segment chegaralari belgilanib, namoyishda tekshiriladi. Tartib: 3 \ensuremath{\rightarrow} 2 \ensuremath{\rightarrow} 4 \ensuremath{\rightarrow} 1.
\end{tcolorbox}
\vspace{2mm}\hrule height .2pt\vspace{2mm}

% DB dcc5d85c-c926-49d2-82ae-7402da50914d status=published
\needspace{6\baselineskip}
\textbf{\color{ink}386.} {ootnotesizecolor{muted}Y3 | published | dcc5d85c}
\par
PowerPoint slaydidagi bitta obyekt uchun animatsiya effektini belgilash, Start/Duration/Delay parametrlarini sozlash va Animation Pane orqali tekshirish kerak. Amallarni tartiblang.
\begin{enumerate}[label=\Alph*),leftmargin=9mm]
\item Animation Pane va namoyishda tekshirish
\item Animatsiya beriladigan obyektni tanlash
\item Start/Duration/Delay parametrlarini sozlash
\item Animations menyusida effekt qo'shish
\end{enumerate}
\begin{tcolorbox}[enhanced,breakable,colback=teal!4,colframe=teal!55!black,boxrule=0.45pt,arc=1.2mm,left=2.8mm,right=2.8mm,top=2mm,bottom=2mm]
\textbf{\color{teal}To'g'ri javob:} \textbf{B - D - C - A}
\par\smallskip
\textbf{Izoh:} Avval obyekt belgilanadi, Animations orqali effekt qo'shiladi, uning vaqt parametrlari sozlanadi va Animation Pane hamda namoyishda tekshiriladi. Tartib: 2 \ensuremath{\rightarrow} 4 \ensuremath{\rightarrow} 3 \ensuremath{\rightarrow} 1.
\end{tcolorbox}
\vspace{2mm}\hrule height .2pt\vspace{2mm}

% DB f5d93870-0d73-7baf-da7b-0c419da59ebc status=published
\needspace{6\baselineskip}
\textbf{\color{ink}387.} {ootnotesizecolor{muted}Y2 | published | f5d93870}
\par
Quyidagi tushunchalarni mazmuniga qarab bog'lang.  1) Advance Slide After;  2) Loop until Esc;  3) Rehearse Timings;  Izohlar: A) Esc bosilguncha namoyishni takrorlash;  B) Slaydni belgilangan vaqtdan so'ng almashtirish;  C) Namoyish vaqtlarini mashq qilish;
\textbf{Chap ustun:}\begin{enumerate}[label=\arabic*),leftmargin=9mm]
\item Advance Slide After
\item Loop until Esc
\item Rehearse Timings
\end{enumerate}
\textbf{O'ng ustun:}\begin{enumerate}[label=\Alph*),leftmargin=9mm]
\item Esc bosilguncha namoyishni takrorlash
\item Slaydni belgilangan vaqtdan so'ng almashtirish
\item Namoyish vaqtlarini mashq qilish
\end{enumerate}
\begin{tcolorbox}[enhanced,breakable,colback=teal!4,colframe=teal!55!black,boxrule=0.45pt,arc=1.2mm,left=2.8mm,right=2.8mm,top=2mm,bottom=2mm]
\textbf{\color{teal}To'g'ri javob:} \textbf{1-B, 2-A, 3-C}
\par\smallskip
\textbf{Izoh:} To'g'ri moslik: Advance Slide After - Slaydni belgilangan vaqtdan so'ng almashtirish, Loop until Esc - Esc bosilguncha namoyishni takrorlash, Rehearse Timings - Namoyish vaqtlarini mashq qilish.
\end{tcolorbox}
\vspace{2mm}\hrule height .2pt\vspace{2mm}

% DB b51442c1-6c3b-47ec-bea7-435678565f44 status=review
\needspace{6\baselineskip}
\textbf{\color{ink}388.} {ootnotesizecolor{muted}Y1 | review | b51442c1}
\par
O'qituvchi 12 slaydli taqdimotni har safar 15 soniyadan keyin avtomatik almashtirmoqchi. Qaysi sozlama mos?
\begin{enumerate}[label=\Alph*),leftmargin=9mm]
\item Transitions \ensuremath{\rightarrow} Advance Slide \ensuremath{\rightarrow} After: 00:15 (zarur slaydlarda)
\item Animations \ensuremath{\rightarrow} Start After Previous \ensuremath{\rightarrow} Duration: 15 soniya (barcha obyektlarga)
\item Slide Show \ensuremath{\rightarrow} Rehearse Timings \ensuremath{\rightarrow} 15 soniya (faqat birinchi slaydda)
\item Transitions \ensuremath{\rightarrow} Duration: 15 soniya, Advance Slide faqat Mouse Click
\end{enumerate}
\begin{tcolorbox}[enhanced,breakable,colback=teal!4,colframe=teal!55!black,boxrule=0.45pt,arc=1.2mm,left=2.8mm,right=2.8mm,top=2mm,bottom=2mm]
\textbf{\color{teal}To'g'ri javob:} \textbf{A) Transitions \ensuremath{\rightarrow} Advance Slide \ensuremath{\rightarrow} After: 00:15 (zarur slaydlarda)}
\par\smallskip
\textbf{Izoh:} PowerPointda slaydni avtomatik almashtirish vaqti Transitions oynasidagi Advance Slide - After maydonida beriladi; 15 soniyalik sozlama kerakli slaydlarga yoki barchasiga qo'llanadi. Transition Duration o'tish animatsiyasi davomiyligini belgilaydi, obyekt Animations esa slayd almashtirish vaqti emas.
\end{tcolorbox}
\vspace{2mm}\hrule height .2pt\vspace{2mm}

% DB c54468f7-ccf1-4188-bcb3-9f530d20e2f9 status=review
\needspace{6\baselineskip}
\textbf{\color{ink}389.} {ootnotesizecolor{muted}Y1 | review | c54468f7}
\par
Bitta slaydda rasmga Entrance effekti, keyingi slaydga Fade o'tish effekti berildi. Qaysi fikr to'g'ri?
\begin{enumerate}[label=\Alph*),leftmargin=9mm]
\item Entrance slayd ichidagi obyektga, Fade (Transition) esa slaydlar almashishiga tegishli
\item Entrance slaydlar almashishiga, Fade esa slayddagi rasm paydo bo'lishiga tegishli
\item Entrance bilan Fade ikkalasi ham faqat obyektlarga biriktiriladigan animatsiya effektlaridir
\item Entrance slayd ichidagi rasmga, Fade esa matn shriftini o'zgartirishga xizmat qiladi
\end{enumerate}
\begin{tcolorbox}[enhanced,breakable,colback=teal!4,colframe=teal!55!black,boxrule=0.45pt,arc=1.2mm,left=2.8mm,right=2.8mm,top=2mm,bottom=2mm]
\textbf{\color{teal}To'g'ri javob:} \textbf{A) Entrance slayd ichidagi obyektga, Fade (Transition) esa slaydlar almashishiga tegishli}
\par\smallskip
\textbf{Izoh:} PowerPointda Entrance animatsiyasi slayd ichidagi obyektning paydo bo'lishini boshqaradi, Transition turidagi Fade esa slaydlar orasidagi o'tishni belgilaydi. Bular obyekt va butun slayd darajasidagi turli mexanizmlardir; obyekt animatsiyasini slayd almashish effekti bilan aralashtirmaslik kerak.
\end{tcolorbox}
\vspace{2mm}\hrule height .2pt\vspace{2mm}
\clearpage\section*{Nazorat va manba haqida}
\addcontentsline{toc}{section}{Nazorat va manba haqida}
389 ta amaldagi savol: 381 ta published, 8 ta review. 148 ta archived va 7 ta draft yozuv chiqarib tashlangan.
Manba bankidagi javoblar, moslashtirish kalitlari va izohlar har bir savolga identifikatori bo'yicha bog'langan.
Rasmlarga havolalar internetdan ochiladi; rasm tasvirlari bu nashrga ichki fayl sifatida biriktirilmagan.
Review savollar mustaqil mutaxassis tomonidan tasdiqlangan hisoblanmaydi.
\end{document}